\documentclass[11pt,a4paper]{article}
% Shared typesetting and notation for the two associated documents.
\usepackage[T1]{fontenc}
\usepackage[utf8]{inputenc}
\usepackage[english]{babel}
\usepackage{lmodern,microtype}
\usepackage[a4paper,margin=24mm,headheight=14pt]{geometry}
\usepackage{amsmath,amssymb,amsthm,mathtools,mathrsfs}
\usepackage{booktabs,longtable,array,tabularx}
\usepackage{xcolor,enumitem,fancyhdr,needspace,etoolbox,placeins}
\usepackage{xr-hyper}
\usepackage[colorlinks=true,linkcolor=blue!40!black,citecolor=blue!40!black,urlcolor=blue!40!black]{hyperref}
\usepackage{bookmark,xurl}
\newcommand{\E}{\mathbb E}
\renewcommand{\P}{\mathbb P}
\newcommand{\N}{\mathbb N}
\newcommand{\Z}{\mathbb Z}
\newcommand{\F}{\mathbb F}
\newcommand{\Law}{\mathcal L}
\newcommand{\Pois}{\operatorname{Pois}}
\newcommand{\TV}{d_{\mathrm{TV}}}
\newcommand{\ind}{\mathbf 1}
\newcommand{\eps}{\varepsilon}
\DeclareMathOperator{\Var}{Var}
\DeclareMathOperator{\Cov}{Cov}
\DeclareMathOperator{\rank}{rank}
\DeclareMathOperator{\codim}{codim}
\renewcommand{\sf}{\operatorname{sf}}
\newtheorem{theorem}{Theorem}[section]
\newtheorem{proposition}[theorem]{Proposition}
\newtheorem{lemma}[theorem]{Lemma}
\newtheorem{corollary}[theorem]{Corollary}
\theoremstyle{definition}
\newtheorem{definition}[theorem]{Definition}
\theoremstyle{remark}
\newtheorem{remark}[theorem]{Remark}
\newtheorem{example}[theorem]{Example}
\numberwithin{equation}{section}
\setlength{\parindent}{1.2em}
\setlength{\parskip}{3pt}
\setlength{\emergencystretch}{2em}
\setlist{itemsep=3pt,topsep=5pt,leftmargin=20pt}
\renewcommand{\arraystretch}{1.12}
\pagestyle{fancy}
\fancyhf{}
\fancyhead[L]{\small Long multiplicative runs}
\fancyhead[R]{\small Brice Pouly}
\fancyfoot[C]{\thepage}
\renewcommand{\headrulewidth}{0.3pt}
\setcounter{tocdepth}{1}
\hypersetup{
 pdftitle={Sharp Poisson approximation and second-order correction for long multiplicative runs},
 pdfauthor={Brice Pouly},
 pdfsubject={Random multiplicative functions, Poisson approximation, hereditary stopping, arithmetic dependencies},
 pdfkeywords={Random multiplicative functions; long runs; optimal Poisson approximation; factorial cumulants; working draft 1.8, 7 November 2026},
 bookmarksdepth=2}
\allowdisplaybreaks[1]
\widowpenalty=10000
\clubpenalty=10000

% Presentation-only identifiers for the separate condensed preview.
\fancyfoot[L]{\scriptsize\color{black!60}Simplified editorial preview}
\hypersetup{pdftitle={Long multiplicative runs: simplified editorial preview},pdfkeywords={Simplified editorial preview; not an adopted manuscript revision}}

\newcommand{\PoneTechnicalCompanion}{}
\externaldocument[][nocite]{build/main}[P1_Simplified_Editorial_Preview_Main_2026-11-07_EN.pdf]
\makeatletter
\patchcmd{\l@section}{1.5em}{2.3em}{}{\PackageError{p1}{Could not widen companion contents numbers}{}}
\makeatother
\renewcommand{\thesection}{T\arabic{section}}
\renewcommand{\theHsection}{T.\arabic{section}}
\fancyhead[L]{\small Long multiplicative runs: technical companion}
\hypersetup{pdftitle={Sharp Poisson approximation for long multiplicative runs: Technical companion (simplified editorial preview)},pdfsubject={Technical proofs, arithmetic estimates, tails and positional comparison}}
\title{\bfseries Sharp Poisson approximation\\and second-order correction\\for long multiplicative runs\\[9pt]\large Technical companion}
\author{Brice Pouly\\[5pt]
 \small With mathematical and editorial assistance from ChatGPT (GPT-6 Astra)}
\date{\small Simplified editorial preview of working draft 1.8 --- 7 November 2026\\[3pt]Technical companion}
\begin{document}
\maketitle
\thispagestyle{plain}
\begin{abstract}
This technical companion accompanies the article \emph{Sharp Poisson
approximation and second-order correction for long multiplicative runs}.
Together the two documents give the proof of the optimal quadratic
Poisson--Charlier error power \(4/3\), the third factorial cumulant
power \(41/30\), and the hierarchy of explicit rational-channel
approximations for critical runs of a completely multiplicative
random sign function.

The companion contains the detailed arithmetic counts, fixed-source
estimates, short-selection moment bounds, algebraic constructions,
and conditional estimates used by the article. It also retains the
regular-set reconstruction needed for unsigned positional comparison,
and the further results on offsets, local blocks, effective tails,
Palm identities, and deletion. Fixed-order and growing-range
estimates have separate hypotheses throughout. The article states
the technical interfaces and retains the main reconstruction and
optimality arguments; the complete proof uses both documents.
\end{abstract}
\clearpage
\begingroup
\small
\setlength{\parskip}{0pt}
\tableofcontents
\endgroup
\clearpage
\section*{Scope and reading guide}
\addcontentsline{toc}{section}{Scope and reading guide}
\label{p14:comp:guide}

This companion and the associated article, \emph{Sharp Poisson
approximation and second-order correction for long multiplicative runs},
form a single mathematical work. They share the model, notation,
version number, and date. The article follows the exact-marginal
proof of the optimal quadratic approximation. This companion gives
the detailed proofs of the technical inputs and retains the further
results on arithmetic geometry, tails, and deletion.

References beginning with T belong to this companion. Other
numbered references belong to the article. The two PDFs have
cross-document links when kept together under their supplied
filenames. Each cited statement also has a stable printed number,
so the proof can be checked without those links.

\paragraph{Common model and conventions.}
The completely multiplicative function \(f\) has independent fair
signs at primes. Both signs of a constant run are counted. With
\[
 I_N=[N,2N)\cap\Z,\quad |L-\log_2N|\le C,\quad
 B=L+1,\quad p=2^{-L},\quad \lambda_N=Np,
\]
the event at \(x\in I_N\) and its count are
\[
 J_x=\ind_{\{f(x-1)=-f(x),\ f(x)=\cdots=f(x+L-1)\}},
 \qquad Z_N=\sum_{x\in I_N}J_x.
\]
We write \(p_x=\E J_x\), \(\mu_N=\E Z_N\),
\(\alpha_N=\Var Z_N-\mu_N\), \(\Pi_\mu=\Pois(\mu)\),
and \(\Delta=\delta_1-\delta_0\). Products of measures
are convolutions, \(\|\nu\|_1=\sum_k|\nu\{k\}|\),
and total variation for probabilities is half this norm.
All unmarked logarithms are natural. Constants may depend on
the fixed critical-band width \(C\) and on parameters explicitly
declared fixed.

The space \(S_x\) is the image, under prime valuations modulo two,
of the even selections in \(V_x=\{x-1,\ldots,x+L-1\}\).
The exact formulation uses all \(x\) with \(p_x>0\), dimensions
\(d_x=\dim S_x\), and the relative nullity
\(\sum_{x\in I}d_x-\dim\sum_{x\in I}S_x\).
The regular-set argument restricts to \(G_N\), where \(d_x=L\)
and \(p_x=p\). Its homogeneous notation and its restricted count
are always identified when that argument is used. The raw arithmetic
nullity \(\rho(I)=|I|L-\dim\sum_{x\in I}S_x\) includes local
kernels; it agrees with the effective nullity on \(G_N\).
Section~\ref{p13:transfer:section} explicitly uses \(\rho\)
as a local abbreviation for the effective nullity.

\paragraph{Technical inputs for the optimal scalar theorem.}
The arithmetic estimates in Sections~\ref{p1:arith:appendix}
and~\ref{p1:arith:supplement} supply the split-product bounds,
variable-coefficient counts, and channel budgets. The proofs
associated with Sections~\ref{p1:fixed:section},
\ref{p13:short:section}, and~\ref{p13:half:section} establish
fixed-source propagation and conditional half counts. The
rigidity and triple-counting proofs then supply the inputs to
Sections~\ref{p13:rig:section} and~\ref{p13:third:section}.
The article states the required hypotheses at each interface
and retains the reconstruction, its separately uniform tail,
and the quartic lower-bound argument.

The short proofs of Proposition~\ref{p1:arith:fixed-support} and
Proposition~\ref{p1:fixed:compressed} return to
Theorem~\ref{p13:half:conditional}. The latter also uses
Lemma~\ref{p13:transfer:dense-targets}, whereas
Proposition~\ref{p1:fixed:primitive} uses
Proposition~\ref{p13:transfer:primitive}. These geometric inputs
are proved independently of the moment theorem and either law reconstruction.

\paragraph{Entering a proof and returning to the article.}
The following local routes lead to the complete proof subsection,
then to the next use in the article. A subsection and a lemma can
have the same printed number; the middle column names the
subsection and links to its heading.

\begin{center}
\small
\begin{tabularx}{\textwidth}{@{}>{\raggedright\arraybackslash}p{0.23\textwidth}>{\raggedright\arraybackslash}p{0.40\textwidth}>{\raggedright\arraybackslash}X@{}}
\toprule
Entry in the article & Proof in this companion & Return to the article\\
\midrule
Proposition~\ref{p13:short:moments}
& Subsection~\ref{p14:comp:short:moments}: uniform moments, using
the prescribed-divisor estimates above it
& Corollary~\ref{p13:short:old-space}: short projections into the old space\\[4pt]
Theorem~\ref{p13:half:conditional}
& Subsection~\ref{p14:comp:half:iteration}: threshold iteration,
using the certificate count of Lemma~\ref{p13:half:saturation}
& Corollary~\ref{p13:half:connected}: connected hosts at fixed order\\[4pt]
Theorem~\ref{p13:rig:common-curve}
& Subsection~\ref{p14:comp:rigidity:common-curve}: one curve for
the whole relation space
& Proposition~\ref{p13:rig:grid-fibres}: integer fibres, followed
by the degree-uniform point count\\[4pt]
Proposition~\ref{p13:third:aligned-family}
& Subsection~\ref{p14:comp:third:aligned-family}: all triples above
the fixed channels, including exceptional pairs
& Corollary~\ref{p13:third:aligned}: strong aligned pairs\\[4pt]
Lemma~\ref{p13:transfer:dense-targets}
& Subsection~\ref{p14:comp:transfer-dense-targets}: two dense targets
over prescribed sources
& Lemma~\ref{p13:transfer:compressed-target}, then the primitive-mass estimates\\
\bottomrule
\end{tabularx}
\end{center}

\paragraph{The regular-set reconstruction and positional results.}
The complete earlier reconstruction is retained in
Sections~\ref{p1:core:section}--\ref{p1:close:section} and
the associated stopping proofs. Its weighted primitive and
compressed-pair estimates support the unsigned positional
comparison and the deletion arguments in
Section~\ref{p1:limits:positions}. This branch is needed for
those results; the scalar optimal proof is assembled directly
from its stated exact-marginal inputs. Any proof below that
invokes Theorem~\ref{p1:thm:main} is a consequence of that
theorem, not an input to it.

\paragraph{The selected version 1.8 strengthenings.}
Section~\ref{p14:comp:confinement} now gives the fixed-arity anchored
small-kernel witness, constructs its hypotheses for primitive families,
and proves the dense-host and mass estimates used in
Proposition~\ref{p13:transfer:primitive}. It also completes the
quasi-linear bounds of Proposition~\ref{p1:arith:hosts}. The earlier
elementary pair-profile calculation still uses its independent
$6/5$ estimate; the later strengthening is not fed back into that proof.
The graph generator uses disjoint stars and nonnegative weights.
The power-threshold tail and the ordinary smooth-window census remain
in Section~\ref{p1:limits:section}, where they serve the supplementary
tail and forcing conclusions. Neither is an input to the optimal
scalar reconstruction. The exact identities of the reference arrivals
are placed with the positional proof and are distinguished from the
arithmetic indicators.

\paragraph{Further consequences and limitations.}
Sections~\ref{p1:offset:section}, \ref{p1:consequences:section},
and~\ref{p1:uniform:appendix} give refinements of the offset
geometry, interval dependence, and uniform counting estimates.
Section~\ref{p1:limits:section} treats exact local blocks,
effective tails, modified counts, Palm identities, and the cost
of positive forcing. Section~\ref{p1:moment-obstruction}
explains why fixed moments alone cannot imply a polynomial
rate for a complete law. These statements retain their own
quantifiers; none enlarges a fixed-order estimate to a growing
range without a separate proof.

\paragraph{The role of the final sections.}
The proof map in Section~\ref{p1:intro:dependencies} records
the inputs to the scalar theorem.
The positional argument in Section~\ref{p1:limits:positions}
proves Lemma~\ref{p1:limits:vector-reference} and uses the
regular-set weighted reconstruction in Sections~\ref{p1:core:section}--\ref{p1:close:section}.
The remaining late sections give further consequences with their own
assumptions; they do not enter the optimal scalar reconstruction.
This is a statement about logical dependence, not separate papers.
In particular, Proposition~\ref{p13:optimal:nominal} derives its
quantized mean identity directly from the exact marginals,
without the supplementary local-block asymptotics.

The preserved version~1.6 source archive provides
\texttt{python3 check\_all.py} for its historical builds and finite
diagnostics. The version~1.8 working draft is rebuilt from the repository with
\texttt{python3 scripts/build\_manuscript\_v18.py}.
Source integrity and statement correspondence are checked by
\texttt{scripts/check\_revised\_sources\_v18.py}; finite diagnostics
do not test the uniform asymptotic estimates.

\section{Uniform arithmetic estimates}
\label{p1:arith:appendix}

This section supplies the uniformity needed when the window length grows.
The degree parameters in the determinant arguments are fixed before $N$
tends to infinity.  Dependence on such fixed parameters is permitted;
dependence on a squarefree coefficient of size $N^{O(1)}$ is not.

\subsection{Split products, defects, and the regular set}

Products in short intervals have a long history; see
Erd\H{o}s and Turk~\cite{erdosturk1984}.  The polynomial square-root
argument for a fixed even number of distinct factors appears in the
proof of Munsch, Shparlinski, Sun, and Xiao~\cite[Lemma~3.3]{munsch2026}.
We record an explicit bound uniform in the degree, as required when the
number of selected factors grows.  No constant depending implicitly on
that degree is used in the following lemma.

\begin{lemma}[A split-product bound uniform in the degree]\label{p1:arith:runge}
Let $d=2k\geq2$ and let $\gamma_1,\ldots,\gamma_d$ be distinct integers
with $|\gamma_i|\leq R$, $R\geq1$.  If $U\geq2R$ is an integer and
$\prod_i(U+\gamma_i)$ is a square, then
\[
 U\leq(64R)^d.
\]
\end{lemma}

\begin{proof}
Write $G(T)=\prod_i(T+\gamma_i)$ and
\[
 \phi(z)=\prod_i(1+\gamma_i z)^{1/2}=\sum_{j\geq0}c_jz^j,
 \qquad P(T)=\sum_{j=0}^k c_jT^{k-j}.
\]
The branches equal one at the origin.  Cauchy's bound on $|z|=(2R)^{-1}$
gives $|c_j|\leq(3/2)^{d/2}(2R)^j$ and, for $j\leq k$,
$|c_j|\leq2^dR^j$.  Comparing coefficients in $\phi^2$ inductively gives
$2^{2j-1}c_j\in\mathbb Z$ for $j\geq1$, so $2^dP\in\mathbb Z[T]$.
For $U>4R$ the tail satisfies
\[
 |\sqrt{G(U)}-P(U)|\leq4\,3^{d/2}R^{k+1}/U.
\]
If this difference is nonzero, it is at least $2^{-d}$, whence
$U\leq(8R)^d$.  If it vanishes, the nonzero integer polynomial
$Q=2^{2d}(P^2-G)$ has root $U$.  It is nonzero because $G$ has simple roots.
Keeping the homogeneous powers of $R$ in the coefficient products gives
\[
 |[T^{d-j}](P^2-G)|\leq(d+2)4^dR^j\leq8^dR^j,
 \qquad H(Q)\leq(32R)^d.
\]
The integral-root bound gives $U\leq1+H(Q)\leq(64R)^d$.
The range $2R\leq U\leq4R$ already satisfies the conclusion.
\end{proof}

The following coefficient-uniform count is classical; see
Cilleruelo and Garaev~\cite[Section~3, Lemma~4 and Proposition~1]{cilleruelogaraev2011}
and the directly comparable formulation of
Tao~\cite[Lemma~2.10]{tao2026anatomy}.  When $ab$ is not a square,
the substitution $(au)^2-abv^2=ah$ reduces it to the former lemma;
the square case is a divisor count.  We include a proof to keep the
polynomial coefficient range explicit.

\begin{lemma}[Uniform quadratic equation count]\label{p1:arith:pell}
Fix $A_0>0$.  Uniformly for positive integers $a,b$ and nonzero integers $h$
with $a,b,|h|\leq N^{A_0}$, the number of positive solutions to
\[
 au^2-bv^2=h,\qquad u,v\leq N^{A_0},
\]
is $N^{o_{A_0}(1)}$.
\end{lemma}

\begin{proof}
If $ab$ is a square, multiply the equation by $a$ and factor the resulting
difference of squares; the divisor bound suffices.  Otherwise write
$ab=Ds^2$ with $D$ squarefree and nonsquare.  Each solution gives the
algebraic integer
$\xi=au+sv\sqrt D$ of norm $ah$ in the maximal order of
$\mathbb Q(\sqrt D)$.  Its principal ideal has norm $|ah|$.
The number of ideals of this norm is at most $\tau(|ah|)$: at a split prime
there are at most $e+1$ ideals of norm $\ell^e$, and at an inert or ramified
prime there are no more.  Elements generating the same ideal differ by a
unit.  A real quadratic unit greater than one is at least
$(1+\sqrt5)/2$, as follows from its integral trace and norm $\pm1$.
The positive values of $\xi$ arising here lie in $[1,N^{O(A_0)}]$;
there are therefore $O_{A_0}(\log N)$ possibilities in each ideal class
of generators.  Finally $\tau(|ah|)=N^{o_{A_0}(1)}$.  This argument makes
the coefficient uniformity explicit and requires no bound for a fundamental
unit or a class number.
\end{proof}

\begin{lemma}[Local and global defect counts]\label{p1:arith:defects}
Suppose $B\asymp\log N$ and $H\asymp B$, with fixed comparison constants.
Any interval of diameter $O(B)$ in $[cN,CN]$ contains
$O(B/\log B)$ $H$-defective integers.  Moreover
\begin{align}
 \#\{n\leq CN:K_H(n)=1\}&\leq N^{1/2+o(1)},\label{p1:arith:onedfect}\\
 \#\{x\in I_N:V_x\text{ contains two }H\text{-defective integers}\}
   &\leq N^{o(1)}.\label{p1:arith:twodefect}
\end{align}
All statements remain valid after fixed dilations and translations of the
intervals.
\end{lemma}

\begin{proof}
For $D$ defective integers in a short interval, form the binary matrix with
their small-prime valuation columns and an additional parity row.  Its
kernel has redundancy at most $\pi(H)+1$.  A nonzero word selects an even
square product, so Lemma~\ref{p1:arith:runge} gives minimum distance
$\delta\gg\log N/\log B$.  Choose
$t\asymp\log N/\log B$ with $2t<\delta$.  The Hamming bound gives
\[
 \sum_{j=0}^t\binom Dj\leq2^{\pi(H)+1}.
\]
If $D\geq2t$, use $\binom Dt\geq(D/t)^t$.  Since
$(\pi(H)+1)/t=O(1)$, this proves $D=O(B/\log B)$; the other case is
immediate.

Write a defective integer uniquely as $n=a^2c$, where $c$ is squarefree
and $H$-smooth.  Then
\[
 \#\{n\leq CN:K_H(n)=1\}
 \leq\sqrt{CN}\prod_{\ell\leq H}(1+\ell^{-1/2})
 \leq\sqrt{CN}\exp(O(\sqrt H)).
\]
For two distinct defective values in one window, write
$x+i=c_1a^2$, $x+j=c_2b^2$.  The offsets and the squarefree smooth parts
have $B^2 4^{\pi(H)}=N^{o(1)}$ choices.  All their heights are polynomial
in $N$, and $c_1a^2-c_2b^2=i-j\ne0$.
Lemma~\ref{p1:arith:pell} supplies $N^{o(1)}$ possibilities.  This proves
\eqref{p1:arith:twodefect} and the stated uniform versions.
\end{proof}

\subsection{Two small kernels with their smooth coefficients retained}

The following form of the approximate determinant method is needed because
the small-prime parts are fixed separately and may be polynomially large.
It is an extension of the proof of \cite[Sections~2.1--2.3]{reuss2014},
rather than an application of its fixed-coefficient statement.

\begin{lemma}[Uniform determinant estimate]\label{p1:arith:reuss}
Fix $A_0,\delta,\varepsilon>0$.  Suppose $1\leq c_1,c_2\leq3N$,
$0<|h|\leq(\log N)^{A_0}$, and
\[
 c_1D^2U\asymp c_2E^2V\asymp N,\quad
 D,E\ll\sqrt N,\quad c_1c_2(DE)^2\geq N,
 \quad\log(DE),\log(UV)\geq\delta\log N.
\]
The comparison constants are fixed.  The number of positive solutions
$d\asymp D,e\asymp E,u\asymp U,v\asymp V$ to
$c_2e^2v-c_1d^2u=h$ is
\begin{equation}\label{p1:arith:reussbound}
 \ll_{A_0,\delta,\varepsilon}N^\varepsilon
 \min\{(DE\mathcal M)^{1/2}+D+E,
          (UV\mathcal M)^{1/2}+U+V\},
 \qquad
 \log\mathcal M=\frac98\frac{\log(DE)\log(UV)}{\log N}.
\end{equation}
Fixed or polylogarithmic factors in the lower bound for $c_1c_2(DE)^2$
can be absorbed by decreasing the error margin.
\end{lemma}

\paragraph{Parameters in the determinant argument.}
Here and in the proof, height of a polynomial means the maximum absolute
value of its integer coefficients.  The ledger below distinguishes the
fixed auxiliary degrees from the coefficients that vary with $N$.
The symbols $A,B_0,\kappa,L_1,T_1$ are specified in Steps~2--3.
\begin{center}
\small
\begin{tabularx}{\textwidth}{@{}>{\raggedright\arraybackslash}p{.23\textwidth}>{\raggedright\arraybackslash}p{.19\textwidth}>{\raggedright\arraybackslash}p{.22\textwidth}>{\raggedright\arraybackslash}X@{}}
\toprule
Object and dimension & Degree & Point or box scale & Coefficient height\\
\midrule
Input hypersurface in $\mathbb A^4$, dimension $3$
& Total degree $3$
& $(D,E,U,V)$, each $O(N)$
& $c_1,c_2\leq3N$; $|h|\leq(\log N)^{A_0}$\\
Auxiliary curve in $\mathbb P^1\!\times\!\mathbb P^1$, dimension $1$
& Bidegree $(a,b)\leq(A,B_0)$, fixed
& Primitive pairs of heights $L_1,T_1\leq N^C$
& $\|G\|\leq N^\kappa$ after the lattice changes\\
Eliminant curve in $\mathbb A^2$, when $a=1$ or $b=1$
& Bounded in terms of $A,B_0$
& The same primitive boxes
& Resultants and nonzero right-hand sides $N^{O(\kappa)}$\\
\bottomrule
\end{tabularx}
\end{center}
All displayed height exponents, including $C$, are fixed, uniformly over
the input boxes.
The relative chart has errors $\alpha=O(M^{-1})$ and
$\beta=O(|h|/N)$, with $M=N^\zeta\mathcal M$; no coefficient ratio
is charged to the latter error.  Choose the final $\varepsilon$ first,
then a sufficiently small $\zeta$, then $A,B_0$ and $\kappa$, and finally
the curve-count exponent $\epsilon'$ so that
$(C+\kappa)\epsilon'\leq\zeta$.  Only after these choices does $N$
tend to infinity.  In particular, the factor $\|G\|^{\epsilon'}$
in Heath-Brown's bound is explicitly included in this budget.
The coefficient uniformity in this ledger is established below;
Reuss's stated theorem allows its constant to depend on the fixed shift
$h$ and does not itself assert this extension.

\begin{proof}
The proof has four steps.  They identify separately the determinant
scale, the uniform count on each auxiliary curve, and the two lattice
summations.  Only the last step sums over the short intervals.

\medskip\noindent
\textbf{Step 1: primitive pairs and relative normalization.}
If $g=\gcd(d,e)$ and $q=\gcd(u,v)$, then $g^2q\mid h$.
Divide by $g^2q$ and replace the four dyadic scales accordingly.
Both pairs are now primitive.  The effective size $N/(g^2q)$ and the
hypotheses on the scales change only by polylogarithmic factors;
in particular the lower bound for $c_1c_2(DE)^2$ has the permitted
polylogarithmic loss.  There are only polylogarithmically many reductions,
which will be summed at the end.  We keep the notation $N,D,E,U,V,h$
for a reduced box.

Set $s=d/e$, $t=v/u$.  The exact curve approximation is
\[
 t=(c_1/c_2)s^2+h/(c_2e^2u).
\]
On a relative interval of length $M^{-1}$ write
\[
 s=s_0(1+\alpha),\qquad
 t=t_0(1+2\alpha+\alpha^2+\beta),\qquad
 t_0=(c_1/c_2)s_0^2\asymp V/U.
\]
Then $|\alpha|\ll M^{-1}$ and $|\beta|\ll|h|/N$, uniformly in
$c_1,c_2$.  The relevant vertical error is relative to $t_0$;
its bound therefore contains no additional factor $c_1/c_2$.

\medskip\noindent
\textbf{Step 2: the determinant with variable coefficients.}
For the auxiliary monomials $s^at^b$, $0\leq a\leq A$,
$0\leq b\leq B_0$, the row denominator is $e^Au^{B_0}$.
Extracting $s_0^at_0^b$ from the columns leaves the same logarithmic
denominator cost as in the normalized calculation:
\[
 A\log E+B_0\log U+\frac A2\log(D/E)+\frac{B_0}2\log(V/U)
 =\frac A2\log(DE)+\frac{B_0}2\log(UV).
\]
Take $M=N^\zeta\mathcal M$, with a small fixed $\zeta>0$, and use
$N/|h|$ for the vertical approximation scale.  The strict determinant
inequality in \cite[Section~2.1, Lemma~7]{reuss2014} then holds, with
degrees fixed in terms of $\zeta,\delta$.  The resulting integer
polynomial has coefficient height $N^{\kappa(\zeta,\delta)}$.
The exponent is uniform in $c_1,c_2$: their sizes, and the rational
column scalings, are polynomial in $N$, while the degrees are fixed.

To check the allowed scale, put $z=\log(DE)/\log N$ and
$w=1-\log(c_1c_2)/(2\log N)$.  Then
\[
 \frac{\log(UV)}{\log N}=2w-2z+O(1/\log N),\qquad
 \frac{\log\mathcal M}{\log N}=\frac94z(w-z)+o(1)
                    \leq\frac9{16}+o(1).
\]
Thus $M\leq N/|h|$ for sufficiently small $\zeta$.
The two logarithmic size ratios stay in a compact subset of the positive
quadrant, so the degree choices are uniform over all boxes.
The parameter order is: choose $\zeta$ from the final error margin,
then the auxiliary degrees and the height exponent $\kappa$, and only
then the internal exponent in the curve count below.

\medskip\noindent
\textbf{Step 3: fibres and the count on one auxiliary curve.}
Apply the two unimodular lattice changes of
\cite[Section~2.2]{reuss2014} to $(d,e)$ and $(v,u)$.
Write the resulting primitive pairs as $\lambda=(\lambda_1,\lambda_2)$
and $\tau=(\tau_1,\tau_2)$, in boxes of side lengths
$L_1\geq L_2$ and $T_1\geq T_2$, respectively.  The changes of
coordinates have polynomial height; enlarging $\kappa$ absorbs their
contribution to the heights of the original and auxiliary equations.
The original equation still has bidegree $(2,1)$.

An auxiliary factor of degree zero in $(d,e)$ fixes only boundedly
many primitive $(u,v)$; Lemma~\ref{p1:arith:pell} counts the remaining
quadratic equation uniformly.  For degree zero in $(u,v)$, fix $(d,e)$
and set $g_f=\gcd(c_1d^2,c_2e^2)$.  Solvability requires $g_f\mid h$.
The step of the resulting one-dimensional lattice in $u$ is
$c_2e^2/g_f$; the step in $v$ is $c_1d^2/g_f$.  Its intersection
with the box has at most
\[
 O\left(1+\frac{Ng_f}{c_1c_2(DE)^2}\right)
 \leq O\left(1+\frac{N|h|}{c_1c_2(DE)^2}\right)
\]
points.  This is polylogarithmic by the stated lower bound.

For an absolutely irreducible factor of bidegree $(a,b)$,
$a,b\geq1$, \cite[Lemma~2]{heathbrown2012} bounds primitive solutions by
$O_{a,b,\epsilon'}(X^{2/b+\epsilon'}\|G\|^{\epsilon'})$ when the
first pair has height at most $X$; the second pair has no height
restriction.  Thus $b\geq2$ gives the required bound in $L_1$, and
$a\geq2$ gives the bound in $T_1$ after exchanging the pairs.
Choose $\epsilon'$ after $\kappa$ so that the coefficient-height
factor costs at most the chosen power $N^{O(\zeta)}$.

The missing bound when one bidegree is one follows directly by
elimination.  In bidegree $(1,b)$ write
$\lambda_1F_1(\tau)+\lambda_2F_2(\tau)=0$.
Primitivity gives
$q\lambda_1=F_2(\tau)$, $q\lambda_2=-F_1(\tau)$ for a nonzero integer
$q$.  The two dehomogenizations of the coprime forms $F_1,F_2$ give
nonzero resultants $R,R'$ of height $N^{O(\kappa)}$ and
\[
 q\mid R\tau_1^m,\qquad q\mid R'\tau_2^{m'}.
\]
Since $(\tau_1,\tau_2)=1$, prime by prime $q$ divides $RR'$.
There are $N^{o(1)}$ possibilities, uniformly in the coefficients.
Substitution into the original equation gives
$F_3(\tau_1,\tau_2)=hq^2$, where $F_3$ is homogeneous of fixed degree
$2b+1$.  If $F_3$ is zero there is no solution.  Otherwise
$F_3-hq^2$ is a nonzero polynomial.  Fixing the first coordinate leaves
boundedly many second coordinates, except for boundedly many full
fibres, each of length $O(T_2)$.  As $T_2\ll T_1$, the count is
$O(T_1)$.  This includes powers of a linear form and needs no uniform
height-free bound for a Thue equation.  Exchanging the pairs handles
bidegree $(a,1)$ with count $O(L_1)$ and right-hand side $hq$.

These cases prove, with arbitrarily small internal losses, the local
bound of \cite[Lemma~8]{reuss2014} in both orientations.  Reducible
auxiliary polynomials have boundedly many factors.  Rational points on
an irreducible factor that is not geometrically integral lie in its
bounded-degree conjugate intersections.  The resulting local count is
therefore uniform in all the variable coefficients.

\medskip\noindent
\textbf{Step 4: the two short-vector summations.}
We give the scale check behind \cite[Section~2.3]{reuss2014}.
For the first lattice, use intervals with
$s_0=(D/E)x_3/M$, $x_3\asymp M$.  A shortest primitive vector
$(x_1,x_2)$ in a box with $L<L_1\leq2L$ satisfies
\[
 x_2x_3=(ME/D)x_1+O(E/L),\qquad
 |x_1|\ll D/L,\quad |x_2|\ll E/L,
 \quad L\gg(DE/M)^{1/2}.
\]
If $x_1x_2\ne0$, the integer product $x_2x_3$ has $O(E/L)$
choices for each of $O(D/L)$ choices of $x_1$.  The divisor bound gives
$N^{o(1)}DE/L^2$ intervals.  Their local contributions sum to
$N^{O(\zeta)}DE/L$; summing the dyadic $L$ gives
$N^{O(\zeta)}(DE M)^{1/2}$.  The primitive axis vectors have one
coordinate equal to $\pm1$ and give the additional costs $D+E$.

For the other lattice choose $U=N/(c_1D^2)$ and $V=N/(c_2E^2)$
within the dyadic comparison constants.  Then
\[
 t_0=(c_1/c_2)(D/E)^2(x_3/M)^2
     =(V/U)(x_3/M)^2.
\]
Replace this by the grid point $(V/U)x_4/M$, where
$x_4=\lfloor x_3^2/M\rfloor$.  The change is within the allowed
interval width.  Since $x_3\asymp M$, this discretization has bounded
multiplicity and $x_4\asymp M$.  In particular no factor $c_1/c_2$
remains.  The same divisor summation gives
$N^{O(\zeta)}\{(UV M)^{1/2}+U+V\}$.
Finally take $\zeta$ and the internal margins sufficiently small,
replace $M$ by $N^\zeta\mathcal M$, and sum the primitive reductions.
This proves \eqref{p1:arith:reussbound}.
\end{proof}

\begin{lemma}[Two small odd kernels]\label{p1:arith:twokernel}
Uniformly for $B\asymp\log N$, fixed $0<\beta<1$, and $\varepsilon>0$,
\[
 G_2(N,N^\beta):=
 \#\{x\in I_N:\#\{n\in V_x:K_B(n)\leq N^\beta\}\geq2\}
 \ll_{\beta,C,\varepsilon}N^{\max(\beta,7/12)+\varepsilon}.
\]
For $7/12\leq\beta<1$ this exponent is sharp in the sense
$G_2(N,N^\beta)=N^{\beta+o(1)}$.
\end{lemma}

\begin{proof}
Two selected integers have the unique forms
$n=c_1d^2u$ and $n+h=c_2e^2v$, where $1\leq h<B$,
$u,v\leq T=N^\beta$, and $c_1,c_2$ are squarefree and $B$-smooth.
There are $N^{o(1)}$ choices for $c_1,c_2,h$.
Fix them and dyadically subdivide $d,e$.  With $D,E$ fixed,
$U\asymp N/(c_1D^2)$ and $V\asymp N/(c_2E^2)$.
The elementary congruence count, for fixed $d,e$, gives in this box
\begin{equation}\label{p1:arith:elementarybox}
 \ll\frac{N|h|}{c_1c_2DE}+DE
 \ll (\log N)^{O(1)}T+DE.
\end{equation}
Indeed the common divisor of the two moduli divides $h$, and
$DE\gg N/(\sqrt{c_1c_2}T)$.
If $c_1c_2(DE)^2<N$, then $DE<\sqrt N$ and
\eqref{p1:arith:elementarybox} suffices.  It also suffices for
$DE\leq N^\delta$, with a small fixed $\delta$.
For $UV\leq N^\delta$ fix $u,v$ and apply
Lemma~\ref{p1:arith:pell}.  These edge regimes cost
$N^{\max(\beta,1/2)+O(\delta)}$.

In the other boxes apply Lemma~\ref{p1:arith:reuss}.  In its notation the
main exponents are
\[
 f_1(z,w)=z/2+\tfrac98z(w-z),\qquad
 f_2(z,w)=w-z+\tfrac98z(w-z).
\]
They meet at $z=2w/3$; $f_1$ is increasing to the left of that point and
$f_2$ decreasing to its right.  Since $0\leq w\leq1+o(1)$,
\[
 \min(f_1,f_2)\leq w/3+w^2/4\leq7/12+o(1).
\]
The remaining terms satisfy $D+E\ll N^{1/2}$ and $U+V\ll T$.
Sum the polylogarithmic boxes and reduce the internal margins.

For the lower bound let $D\asymp\sqrt{N/T}$ with a fixed sufficiently
large comparison constant, and take coprime $d,e\in[D,2D]$.
The conditions $n\equiv0\pmod{d^2}$, $n+1\equiv0\pmod{e^2}$ have
$\gg N/D^4$ solutions in a fixed interior subinterval of $[N,2N)$,
because $\beta>1/2$.  Both remaining cofactors are at most $T$.
The total number of incidences is $\gg D^2N/D^4\asymp T$.
A fixed $n$ is counted at most $\tau(n)\tau(n+1)=N^{o(1)}$ times.
The two neighboring integers lie in a common window, proving the claim.
\end{proof}

\begin{lemma}[A large selected kernel]\label{p1:arith:anchor}
For any fixed $0<\tau<1/2$ and $1\leq T\leq3N$, the number of separated
pairs possessing a mixed square relation with a selected vertex whose
$B$-kernel exceeds $T$ is at most
\[
 N^{3/2+o(1)}T^{-1/2+\tau}.
\]
Also the number of windows containing one vertex with kernel at most $T$
is at most $N^{1/2+o(1)}T^{1/2}$.
\end{lemma}

\begin{proof}
Write $n=a^2cr$, with $c$ squarefree and $B$-smooth and $r=K_B(n)$.
The one-kernel assertion follows from
\[
 \sum_{\substack{n\leq3N\\K_B(n)\leq T}}1
 \leq\sqrt{3N}\prod_{\ell\leq B}(1+\ell^{-1/2})
                       \sum_{r\leq T}r^{-1/2}
 \leq N^{1/2+o(1)}T^{1/2},
\]
with a factor $B$ for its placements.
For the first assertion anchor a selected value $n=x+i$ with $K_B(n)>T$.
Every prime of its kernel must occur at one of the $B$ offsets of the other
window.  There are at most $B^{\omega(r)}$ assignments; each specifies a
residue class modulo $r$.  Since $r\leq3N$, there are $O(N/r)$ possible
other starts, including $r>N$.  Summation over anchors is bounded by
\[
 B N^{3/2}\prod_{\ell\leq B}(1+\ell^{-1/2})
       \sum_{\substack{r>T\ {\rm squarefree}\\\ell\mid r\Rightarrow\ell>B}}
           B^{\omega(r)}r^{-3/2}.
\]
Rankin's inequality bounds the last sum by
\[
 T^{-1/2+\tau}\prod_{\ell>B}(1+B\ell^{-1-\tau})
 =T^{-1/2+\tau}\exp(O_\tau(B^{1-\tau})).
\]
Both products and the placement factor are $N^{o(1)}$.
\end{proof}

\subsection{Multiquadratic surfaces and the curves that can survive}

We use two uniform determinant estimates.  A geometrically integral affine
surface of degree $d$ in $\mathbb A^3$ has, outside
$O_{d,\varepsilon}(H^{2/\sqrt d+\varepsilon})$ integral points of height
at most $H$, a covering by
$O_{d,\varepsilon}(H^{1/\sqrt d+\varepsilon})$ curves of bounded degree
\cite[Lemma~2]{salberger2023}.  An integral curve of degree $e$ in fixed
affine space has $O_{e,n,\varepsilon}(H^{1/e+\varepsilon})$ integral
points in the same box \cite[Theorem~A]{pila1995}.  These constants do
not depend on the defining coefficients.  For a curve irreducible over
$\mathbb Q$ but not geometrically integral, its rational points lie in the
intersection of distinct conjugate components and have bounded cardinality.

Consider equations
\begin{equation}\label{p1:arith:surfaceequations}
 w_\nu^2=d_\nu P_\nu(X)Q_\nu(Y),\quad
 P_\nu(X)=\prod_{i\in I_\nu}(X+i),\quad
 Q_\nu(Y)=\prod_{j\in J_\nu}(Y+j),\quad1\leq\nu\leq k,
\end{equation}
where $I_\nu,J_\nu$ are nonempty, disjoint within each list, and
$|I_\nu|+|J_\nu|\leq K$.  The offsets are $O(B)$ and the nonzero
coefficients $d_\nu$ are squarefree and $O(B)$-smooth.  Constants $k,K$
are fixed.

\begin{lemma}[The uniform geometric counts]\label{p1:arith:geometric}
Fix $A_0,K,\varepsilon>0$.  For sufficiently large fixed $k$, a system
of the form \eqref{p1:arith:surfaceequations}, with arbitrary nonzero integer
coefficients $|d_\nu|\leq N^{A_0}$, has
$O_{A_0,K,\varepsilon}(N^{1/2+\varepsilon})$ integral hosts
$X,Y\asymp N$ at which the even-selection square-class maps on the
two actual sets of selected values are injective and their images are
in direct sum.  If every $Q_\nu$ has degree one, the bound is
$O_{A_0,K,\varepsilon}(N^\varepsilon)$.  If the degrees of the $P_\nu$
vary, select sufficiently many equations of one of their finitely many
possible degrees before applying the assertion.
These estimates are uniform in the coefficients and do not require their
smoothness; any enumeration of such coefficients is a separate cost.
They remain valid after other fixed centers have been specialized into
the $d_\nu$.

For one variable, a bounded-degree system of split-product square equations
with polynomial-height coefficients has $N^{o(1)}$ possible integer values
as soon as its disjoint root sets include at least two distinct selected
offsets.  This includes either one equation of degree at least two or
two equations with disjoint nonempty root sets.

The following graph consequences hold, with the number of component
equations chosen sufficiently large after the stated fixed parameters.
\begin{enumerate}
\item For any $\eta,\varepsilon>0$, pairs of separated nonaligned
logarithmic windows with at least $\eta B$ nontrivial unpinned components
number $O_{\eta,C,\varepsilon}(N^{1/2+\varepsilon})$.
\item For every fixed integer $a\geq1$ and $\varepsilon>0$, there is
$k_a$ such that $k_a$ components having $a$ vertices in one window and
one in the other have $O_{a,C,\varepsilon}(N^\varepsilon)$ dyadic hosts.
For $a=1$ aligned hosts are excluded; for $a\geq2$ no such exclusion
is required.
\item After one start is fixed, a component with at least two vertices
in the other window has $N^{o(1)}$ possible partners, uniformly for a
fixed bound on the component size and polynomial coefficient heights.
\end{enumerate}
\end{lemma}

\paragraph{Parameters in the geometric counts.}
The external estimates are applied in the following embeddings.  Write
$m=\binom{D+2}{2}+k-1$ for the ambient dimension of the affine Veronese chart,
and $q=(a+1+A_0)/2$ in the degree-$(a,1)$ case.
\begin{center}
\small
\begin{tabularx}{\textwidth}{@{}>{\raggedright\arraybackslash}p{.23\textwidth}>{\raggedright\arraybackslash}p{.23\textwidth}>{\raggedright\arraybackslash}p{.15\textwidth}>{\raggedright\arraybackslash}X@{}}
\toprule
External count & Dimension and degree & Point height & Coefficient height\\
\midrule
Surface covering, Steps~1--4
& Surface in $\mathbb A^3$; $2^k\leq d_{\rm im}$,
$d_{\rm im}=O_{k,K}(1)$
& $H\ll N^D$
& Polynomial in $N$; no coefficient-height dependence in the covering constant\\
Lifted curves, Steps~2--4
& Curves in $\mathbb A^m$; $2D\leq e=O_{D,k,K}(1)$
& $H\ll N^D$
& No bound on the coefficients of the covering curves is required\\
Degree-$(a,1)$ surface covering, Steps~5--6
& Surface in $\mathbb A^3$; $2^k\leq d_{\rm im}$,
$d_{\rm im}=O_{k,a}(1)$
& $H\ll N^q$
& Polynomial in $N$; the same coefficient-uniform covering estimate\\
Degree-$(a,1)$ lifted curves, Steps~5--6
& Curves in $\mathbb A^{k+2}$; $E+1\leq e=O_{k,a}(1)$
& $H\ll N^q$
& No coefficient-height bound is required by the curve estimate\\
\bottomrule
\end{tabularx}
\end{center}
The curve rows concern components that can contain the hosts being counted;
Steps~3 and~5--6 exclude the lower degrees.  The input coefficients satisfy
$|d_\nu|\leq N^{A_0}$, and the offsets are $O(\log N)$.
Thus the defining equations and the projected surface equations have
polynomial height with fixed exponents, while the external constants depend
only on the fixed dimensions, degrees and internal error margins.
For the first bound choose $A_0,K,\varepsilon$, then
$D\geq(K+A_0)/2$, then $k$, and finally the internal margin in the
surface and curve counts.  For the degree-$(a,1)$ bound choose $q$, then
$E$ with $q/(E+1)<\varepsilon/8$, then $k$ satisfying the two conditions
in Step~6, and finally that internal margin.  All these choices precede
$N\to\infty$.  The count for one system is uniform over arbitrary
coefficients in the stated height range; the $N^{o(1)}$ enumeration in
the graph application uses smoothness separately.

\begin{proof}
We first count one fixed system, keeping the coefficient height visible;
only afterwards do we sum the patterns arising from the graph.
We will retain a sufficiently large fixed number \(k\) of equations,
chosen after the height parameter \(D\) below.  All degrees and
ambient dimensions below depend on that number and on the fixed parameters,
not on the values of the coefficients.

\medskip\noindent
\textbf{Step 1: geometric integrality and the finite affine projection.}
Over an algebraically closed field of characteristic zero, the square classes on the
right of \eqref{p1:arith:surfaceequations} are independent: the valuation
along a factor $X+i$ detects exactly its own equation.  The coordinate
ring is free of rank $2^k$ over the polynomial ring in $X,Y$, and its
generic localization is a field.  It is therefore integral.

The affine projection $(X,Y,w)\mapsto(X,Y,\sum_\nu w_\nu)$ is finite
and birational onto its image.  Finiteness follows from integrality over
$\mathbb Q[X,Y]$.  Birationality follows because no nontrivial sign change
fixes the sum: the individual radicals are independent basis elements of
the multiquadratic extension.  The image surface in $\mathbb A^3$ has
degree at least $2^k$ and degree bounded in terms of $k,K$.
The projection need not extend across the projective boundary; only its
affine finiteness is used when curves are lifted.

\medskip\noindent
\textbf{Step 2: projective degrees and the height of integral points.}
Choose an integer $D\geq(K+A_0)/2$ before choosing $k$, and homogenize
in the weighted ring
\[
 \mathbb Q[X,Y,Z,W_1,\ldots,W_k]/
 \left(W_\nu^2-d_\nu Z^{2D-|I_\nu|-|J_\nu|}
       \prod_{i\in I_\nu}(X+iZ)\prod_{j\in J_\nu}(Y+jZ)\right),
\]
with $\deg X=\deg Y=\deg Z=1$, $\deg W_\nu=D$.
Its $D$th Veronese ring is generated in degree one by the degree-$D$
monomials in $X,Y,Z$ and the $W_\nu$.  This embeds an integral projective
surface finite of degree $2^k$ over $\mathbb P^2$.  The hyperplane bundle
is the pullback of $\mathcal O_{\mathbb P^2}(D)$, so its degree is
$2^kD^2$ and every integral curve $C$ satisfies
\begin{equation}\label{p1:arith:curvedegree}
 \deg C=D\deg\sigma(C)\,[\mathbb Q(C):\mathbb Q(\sigma(C))].
\end{equation}

Here $\sigma$ is the finite map to $\mathbb P^2$.
At an integral point in the dyadic region,
$|w_\nu|\ll N^{(K+A_0)/2}$.  In the chart $Z=1$, the Veronese
coordinates are integers of height $H\ll N^D$, including the $W_\nu$.
The same bound holds for the affine projection in Step~1.  Thus both
applications of the determinant estimates use one common height bound.

\medskip\noindent
\textbf{Step 3: excluding the curves of smallest possible degree.}
A curve of degree $D$ maps with degree one to a line.  A vertical or
horizontal line meeting the positive region is impossible: a nonconstant
product of distinct linear factors in the other variable would have to be
a rational square.  On an oblique line, square products force every root
on one side to match exactly one root on the other.  Two component
equations give two distinct exact units, hence a positive rational
alignment of height $O(B)$.  Thus a nonaligned point is on no curve of
degree $D$; by \eqref{p1:arith:curvedegree}, its curves have degree at least
$2D$.

In the direct-sum formulation the
exceptional degree-$D$ curves are excluded for a more intrinsic reason:
their two exact units give nonempty even selections, two values on each
side, with equal square classes.  Injectivity makes both classes nonzero,
contradicting the direct sum.  Thus the argument does not require absence
of geometric alignment when the actual spaces are already in direct sum.

\medskip\noindent
\textbf{Step 4: the surface covering and the first graph consequence.}
Apply the surface covering, lift its curves by the finite affine map, and
take their closures in this embedding.  Their numbers of components and
degrees are bounded in terms of the fixed parameters by B\'ezout.
Pila's estimate and the degree lower bound $2D$ give, per pattern,
\[
 O\big(N^{2D2^{-k/2}+\varepsilon'}
        +N^{1/2+D2^{-k/2}+\varepsilon'}\big).
\]
Choose $k$ large enough after $D,\varepsilon$ and reduce the internal
margin $\varepsilon'$.  This proves the per-system bound.

For the graph consequence, smoothness gives $|d_\nu|\leq N^{A_0}$ for a
fixed $A_0$.  At least half the components in a population $c\geq\eta B$ have size
at most $K=\lceil4/\eta\rceil$.  Select $k$ of them.
The offsets and smooth coefficients have
$B^{O_K(k)}2^{k\pi(O(B))}=N^{o(1)}$ patterns.  Choosing $k$ large after
$D,\varepsilon$, and summing these patterns, proves the first assertion.

\medskip\noindent
\textbf{Step 5: reduced fibres for equations of type $(a,1)$.}
For the stronger degree-one assertion we exclude all curves of a fixed
degree at most $E$, before making $k$ large.  For this part use the
ordinary affine embedding with coordinates $X,Y,w_1,\ldots,w_k$.
Write the equations as
$w_\nu^2=d_\nu P_\nu(X)(Y+j_\nu)$ with $\deg P_\nu=a$.
On the normalization of an integral curve of degree $e\leq E$,
the coordinate functions $X,Y$ have degrees $m_X,m_Y\leq e$ and the
genus is at most $(e-1)(e-2)/2$.  Neither function can be constant:
otherwise the many distinct roots in the other variable would all be
branch values, contradicting Riemann--Hurwitz.
Discard equations having a branch value among their selected roots, or a
finite fibre meeting the pole divisor of the other coordinate.
At most
\[
 (2g-2+2m_X)+(2g-2+2m_Y)+m_X+m_Y\leq2e^2
\]
equations are lost.  For each remaining equation all finite fibres are
reduced and avoid the poles.  Parity of divisors gives
\[
 F_Y(-j_\nu)=\sum_{i\in I_\nu}F_X(-i),\qquad m_Y=am_X.
\]
In fact
\[
 \operatorname{div}\!\left(\frac{P_\nu(X)}{Y+j_\nu}\right)
       =F_Y(\infty)-aF_X(\infty),
\]
independently of $\nu$.  For two good equations the quotients therefore
differ by a nonzero constant.  Eliminating this common quotient gives
$Y=R(X)$ with $\deg R\leq a$.  Since
$m_Y=(\deg R)m_X=am_X$, the degree is exactly $a$.

\medskip\noindent
\textbf{Step 6: interpolation on the offsets and exclusion of low-degree curves.}  If
$k\geq4(E+1)^2$, there are at least three good equations.
Their $3a$ distinct grid points $(-i,-j_\nu)$ lie on the graph of $R$:
the reduced-fibre identity shows that $R$ is finite and equals $-j_\nu$
at each selected $-i$.  These grid points determine $R$.
Choose $2a+1$ of them.  Rational interpolation, using maximal minors of
the integer evaluation matrix, gives coprime integer polynomials
$P,Q$ of degrees at most $a$, coefficients $B^{O_a(1)}$, and $R=P/Q$.
The kernel of this interpolation system is one-dimensional: two solutions
would have cross-product of degree at most $2a$ vanishing at $2a+1$
distinct abscissae.  Since the reduced rational function has degree exactly
$a$, no nonconstant polynomial multiple of its coprime numerator and
denominator still has both degrees at most $a$.
This makes the coefficient uniformity explicit: the interpolation data
are the offsets, so their height bound is independent of the $d_\nu$.

If $Q$ is nonconstant, the nonzero integer resultant
$\operatorname{Res}(P,Q)$ has size $B^{O_a(1)}$ and $Q(X)$ divides it
at an integral point.  Thus $|X|\leq B^{O_a(1)}$.  If $Q$ is constant
and $a\geq2$, then $|Y|\gg |X|^a/B^{O_a(1)}$, impossible for
$X,Y\asymp N$.  If $a=1$ and $Q$ is constant, the graph is a rational
line; its grid points give an exact positive alignment.  Thus every
remaining dyadic nonaligned point lies on a curve of degree greater than
$E$.

In the ordinary affine embedding the height is $H\ll N^q$ for a fixed
$q=(a+1+A_0)/2$.  Choose $E$ so that $q/(E+1)<\varepsilon/8$, then
$k\geq4(E+1)^2$ so that $2q2^{-k/2}<\varepsilon/8$.
The same finite projection, surface covering and curve estimate now give
$N^\varepsilon$ per fixed coefficient pattern.  For $a=1$, its exceptional
line is also excluded under the direct-sum hypothesis: two grid matches
give the forbidden even selections on the two actual blocks.  Hence the
degree-one per-system estimate is proved with arbitrary polynomial-height
coefficients.  In the graph application one sums the $N^{o(1)}$ smooth
patterns, proving the second graph assertion, including $a=1$.

\medskip\noindent
\textbf{Step 7: fixing one center and reducing to a quadratic equation.}
For the last assertion fix the first start and the component data.  The
other start satisfies $\prod_{j\in J}(Y+j)=ev^2$, with fixed squarefree
$e\leq N^{O(1)}$ and $|J|\geq2$.  Outside primes dividing
$e\prod_{i\ne j}(i-j)$, every selected factor has even valuation.
Fixing the squarefree parts of two factors costs
$2^{O(\omega(e\prod_{i\ne j}(i-j)))}=N^{o(1)}$ for fixed $|J|$.
Their difference is a nonzero fixed offset, so
Lemma~\ref{p1:arith:pell} gives $N^{o(1)}$ starts.  This is a
bounded-height argument; no height-free theorem for all integral points
is being invoked.
The one-variable assertion with two separate equations follows in exactly
the same localization: take one selected factor from each equation and
put the prime divisors of both coefficients and both within-equation
discriminants into the exceptional prime set.  Outside it each chosen
factor has even valuation.  The two distinct offsets again give one
quadratic equation counted by Lemma~\ref{p1:arith:pell}.
\end{proof}

\begin{lemma}[The dense weighted sector]\label{p1:arith:dense}
For every fixed $\eta,\varepsilon>0$,
\[
 \sum_{\substack{x,y\in I_N,\ |x-y|>L\\
        (x,y)\ {\rm nonaligned},\ c(x,y)\geq\eta B}}
       (2^{\rho(x,y)}-1)
 \ll_{C,\eta,\varepsilon}N^{1+\varepsilon}.
\]
The same bound holds for the full value kernel, without block-parity rows.
\end{lemma}

\begin{proof}
Let $t_2,t_3$ count nontrivial unpinned components with two and three
vertices.  If either exceeds a sufficiently large fixed threshold,
Lemma~\ref{p1:arith:geometric}(2), in both orientations for size three,
gives $N^\varepsilon$ hosts.  Their weight is at most
$2^{B+o(B)}=N^{1+o(1)}$, since each other component needs two vertices
and the number of defective vertices is $o(B)$.
Otherwise $t_2,t_3=O(1)$ and the vertex budget gives
$4c-2t_2-t_3\leq2B$, hence $c\leq B/2+O(1)$.
Lemma~\ref{p1:arith:geometric}(1) gives $N^{1/2+\varepsilon}$ hosts,
each of weight at most $2^{B/2+o(B)}=N^{1/2+o(1)}$.
Reduce the internal margins in both estimates.
\end{proof}

\begin{lemma}[Uniform residual rank on an aligned pair]\label{p1:arith:alignedrank}
For separated aligned windows with reduced height
$b=\max(a,b)\leq B^3$, let $u$ be the number of exact units.  Then
\[
 \rho(x,y)\leq(u-1)_++o(B)\leq L/b+o(B),
\]
uniformly over the channels. For fixed channel height $b$ and fixed
$D\ge6b$, an enlarged prime cutoff $6bB\le H\le DB$ leaves
residual $O_{C,b,D}(B/\log B)$.
\end{lemma}

\begin{proof}
Normalize the two windows to $bV_x$ and $aV_y$, and merge equal values.
All remaining distinct integers lie in an interval of diameter $O(B^4)$
centered at an integer at least $N$.  Remove prime rows $\ell\leq B$
and rows with $\ell\mid ab$.  Every other prime row meets at most two
distinct normalized vertices, since it occurred at most once in each
original window.  Its equations again leave one variable per unpinned
component.  Let $c_K$ count these components of size at most a fixed $K$.

For these components, the removed prime rows and one total-parity row
form a binary matrix of redundancy at most
$s=\pi(B)+\omega(ab)+1=O(B/\log B)$.  A nonzero kernel word gives
an even square product of distinct normalized integers.  By
Lemma~\ref{p1:arith:runge}, it uses at least
$\delta\gg\log N/\log B$ vertices, and hence at least $\delta/K$
components.  The Hamming argument in Lemma~\ref{p1:arith:defects} yields
$c_K=O_{C,K}(B/\log B)$.  Components larger than $K$ number at most
$2B/K$.  Since $K$ can be chosen arbitrarily large before $N$, the
dimension of the value kernel on the distinct normalized set is $o(B)$.

The kernel of the merging map on the original relative selections is
exactly the even-selection space of the $u$ exact units, of dimension
$(u-1)_+$.  Its image lies in the value kernel just bounded, proving
the first assertion.  For fixed height the normalized interval has
diameter $O(B)$.  Private primes above a sufficiently large fixed multiple
of $B$ leave only $O(B/\log B)$ defective normalized vertices, directly
by Lemma~\ref{p1:arith:defects}; this gives the last assertion.
\end{proof}

\subsection{Asymmetric extensions and homogeneous seeds}

\begin{lemma}[An asymmetric pair profile]\label{p1:arith:asymmetric}
Fix integers $D,Q\geq1$ and $D_{\rm cut}\geq2D+1$.
Let $B=L+1$ and $|L-\log N/\log2|\leq C$.
For each $N\leq X<N+QN$, let $U_X$ be any subset of
$X+[-DB,DB]$, and let $V_z$ be a new window with $z\in I_N$.
Choose an integer cutoff
$(2D+1)B\leq H\leq D_{\rm cut}B$.
For this asymmetric comparison call $(X,z)$ aligned at radius $D$ if
there are coprime positive integers $r,s\leq B^3$ such that
$|rz-sX|\leq D(r+s)B$; use its complement in the sum below.
Let $\delta^*(X,z)$ be the dimension of the even selections
of $V_z$ whose prime valuations above $H$ belong to the value space of
$U_X$.  Then
\[
 \sum_{\substack{N\leq X<N+QN,\ z\in I_N\\
                 (X,z)\ {\rm nonaligned\ at\ radius}\ D}}
       (2^{\delta^*(X,z)}-1)
 \ll_{C,D,D_{\rm cut},Q,\varepsilon}N^{31/24+\varepsilon}.
\]
The choice $Q=1$ gives the two-dyadic-window version.
Adding prime coordinates at most $H$ to the old data leaves
$\delta^*$ unchanged.
\end{lemma}

\begin{proof}
All alignments in this proof use the radius $DB$ of the old enclosing
interval.  The cutoff lower bound is larger than both interval diameters,
so a prime above $H$ occurs at most once in each block.
The fixed expanded old-centre range changes only constants in the
biform support estimates and the bounded-degree graph counts.
For the graph above $H$, let $D_z$ be the number of new defective vertices
and let $c$ count nontrivial unpinned components.  Projection to the new
block gives $\delta^*\leq\min(L,D_z+c)$, with $D_z=o(B)$.
A nonzero even selection uses at least two new vertices.  If it uses no
old vertex, the new window has two defects and contributes $N^{1+o(1)}$
hosts after the old parameter is chosen.  Otherwise it uses at least one
old and two new vertices.  Set $T=N^{5/12}$.
If all their kernels are at most $T$, Lemmas~\ref{p1:arith:twokernel}
and~\ref{p1:arith:anchor} give
\[
 N^{7/12+\varepsilon'}N^{1/2+5/24+o(1)}=N^{31/24+\varepsilon'+o(1)}.
\]
If one kernel exceeds $T$, the same CRT anchor proof gives
$N^{3/2+o(1)}T^{-1/2+\tau}=N^{31/24+(5/12)\tau+o(1)}$.
All these arguments concern the large-prime valuations only, so they do
not impose parity on the old selection.  They are unchanged at the
fixed enlarged threshold $H$.

For $c<\eta B$, the weight is at most $2^{\eta B+o(B)}$, and the support
bound suffices.  For $c\geq\eta B$, the proof of
Lemma~\ref{p1:arith:geometric}(1) applies to the unequal fixed widths and
gives $N^{1/2+\varepsilon'}$ hosts.  Fix a large integer $K$.  Components
with one new vertex and more than $K$ old vertices number $O_D(B/K)$.
A sufficiently large fixed population of any other type $(a,1)$,
$1\leq a\leq K$, has only $N^{\varepsilon'}$ hosts by
Lemma~\ref{p1:arith:geometric}(2).  On this exceptional set use
$2^{\delta^*}\leq2^L\ll N$.  Else the number $c_1$ of one-new-vertex
components is $O_D(B/K)+O_K(1)$, and $2c-c_1\leq B$.  Thus
$\delta^*\leq B/2+O_D(B/K)+o(B)$ and the remaining dense mass is
$N^{1+O_D(1/K)+\varepsilon'}$.
Choose $K$, then $\eta,\tau,\varepsilon'$ after the required margin.
Finally shrinking the old set only decreases $\delta^*$, and prime bits
at most $H$ have already been forgotten.
\end{proof}

\paragraph{Alignment convention for asymmetric applications.}
The enlarged radius in this lemma follows the actual old interval.
Original pairs of equal windows retain the alignment convention in
the main article.  Whenever a source is compressed or an enclosing
interval is enlarged below, its aligned sector is enlarged at the
same time; that sector still has $NB^{O(1)}$ hosts and is treated by
the separate aligned-rank estimates.  The radius-one assertion with
arbitrary $D$ in version~1.6 is false, as witnessed by the checked
family below.

\paragraph{Necessity of the enlarged alignment radius.}
In Lemma~\ref{p1:arith:asymmetric}, the alignment radius follows the
old enclosing interval $X+[-DB,DB]$: the exclusion is
$|rz-sX|\le D(r+s)B$. The centres satisfy
$N\le X<N+QN$, $z\in I_N$, with fixed $Q\ge1$, and the cutoff satisfies
$(2D+1)B\le H\le D_{\rm cut}B$. The weight $2^{\delta^*}-1$ and the
power $31/24$ are unchanged. A shifted $2:3$ family with
$L=3k$, $B=3k+1$, $X=2t$, $z=3t+3B$ and
$U_X=[X,X+3B]\cap\Z$ has $\delta^*\ge k-1$. It is nonaligned at
radius one for $t>5B^4$ but aligned at radius two; its dyadic mass
of order $N^{4/3}$ refutes the original radius-one formulation.
Compressed-source applications use the same enlarged radius on both
sides of the alignment partition.

\begin{lemma}[Aligned loss above a rational pair]\label{p1:arith:palmrank}
Fix a geometrically good pair in channel $a:b=2:3$ or $3:4$, regardless
of its affine compatibility.  Among targets $z$ with
$|z-x|,|z-y|\geq L$ and distinct from the two sources, there are
$B^{O(1)}$ aligned targets per source.  Except for $B^{O(1)}$ source
pairs at dyadic endpoints, their additional relative rank loss modulo
the old value space is at most
\[
 L/3+o(L)\quad(a:b=2:3),\qquad
 L/2+o(L)\quad(a:b=3:4).
\]
The result allows the prime bits of $a,b$ in the old information.
\end{lemma}

\begin{proof}
Represent the source by $x=at-i_0$, $y=bt-j_0$ and put $X=abt$.
The true old normalized values form a subset of
$X+[-bB,bB]$.  Applying the residual argument of
Lemma~\ref{p1:arith:alignedrank} to a fixed-width enclosing interval and
the new interval leaves $o(B)$ dimensions after exact units are removed.
If an exact unit of the enlarged graph uses an old value absent from the
true set, its new nondefective vertex is pinned in the true graph.
Consequently the remaining loss is at most the number of new vertices
matched to actual old values, plus $o(B)$.

Write an aligned target as $z=(ct+h)/d$, with $c,d$ positive coprime,
$c,d\leq B^{O(1)}$ and $h/d=O(B)$.  A match to the first old family
requires
$h+dk\equiv0\pmod{c/(c,a)}$; the second family requires the same
condition modulo $c/(c,b)$.  Their intersection is the class modulo $c$,
because $(a,b)=1$.  Thus the union has at most
\begin{equation}\label{p1:arith:matchdensity}
 B\frac{(c,a)+(c,b)-1}{c}+O(1)
\end{equation}
new offsets.  Conditions on the true intervals can only reduce this count.

For $2:3$, the source range is
$N/2-O(B)\leq t\leq2N/3+O(B)$, so $3/2\leq c/d\leq4$ apart
from endpoint errors of order $B/N$.  A distinct rational of polynomial
height cannot cross either endpoint by such a small amount.
For $c\geq7$ the density in \eqref{p1:arith:matchdensity} is at most
$1/3$ (check $7\leq c\leq12$, then bound the numerator by four).
For $c=5$ it is $1/5$.  The cases $c/d=2,3$ physically reproduce one
old family.  Their supports intersect the new support in at most one
integer, because the starts are at distance at least $L$.  The other
family supplies at most $B/3+O(1)$ units, counted in its actual window
of width $B$.  The only remaining possibilities with $c\leq6$ are
$c/d=4,3/2$; these force $t=N/2+O(B)$ or $t=2N/3+O(B)$.
They account for only $B^{O(1)}$ source pairs.

For $3:4$, one has $N/3-O(B)\leq t\leq N/2+O(B)$ and
$2\leq c/d\leq6$.  The density is at most $1/2$ for $c\geq9$
(check $9\leq c\leq12$, then use numerator at most six), and also
for $c=5,7,8$.  The physical cases $c/d=3,4$ leave at most one unit
in one old family and $B/4+O(1)$ in the other.  The remaining fractions
$2,6$ force $t=N/2+O(B)$ or $N/3+O(B)$ and are endpoint sources.
Finally the canonical slopes and displacements have polynomially many
choices in $B$, proving the target count.
\end{proof}

\begin{proposition}[Homogeneous seeds for fixed extensions]
\label{p1:arith:fixed-homogeneous-seeds}
Let $\mathcal A_{\rm raw}$ be the geometric good-edge family obtained by
omitting the compatibility condition in the definition of $\mathcal E_N$.
For a finite set $I\subset G_N$ put
\[
 W(I)=p^{|I|}2^{\rho(I)}.
\]
All configurations in the following sums have distinct, pairwise separated
starts.  For every $\varepsilon>0$:
\begin{enumerate}
\item The sum of $W(\{x,y\})$ over dependent pairs outside
$\mathcal A_{\rm raw}$ is $O_{C,\varepsilon}(N^{-4/5+\varepsilon})$.
Moreover $\sum_{e\in\mathcal A_{\rm raw}}W(e)\ll_C LN^{-2/3}$.
\item For a source $e$ in channel $a:b$, sum over targets $z$ for which
$S_z\cap\sum_{x\in e}S_x\ne\{0\}$.  Then
\[
 \sum_{e,z}W(e\cup\{z\})
 \ll_{C,\varepsilon}
 \begin{cases}N^{-4/3+\varepsilon},&a:b=2:3,\\
               N^{-5/4+\varepsilon},&a:b=3:4.
 \end{cases}
\]
\item For two distinct geometric good edges, whose union has three or
four starts,
\[
 \sum_{e\ne f}W(e\cup f)
 \ll_{C,\varepsilon}N^{-5/4+\varepsilon}+(LN^{-2/3})^2.
\]
\end{enumerate}
These are rank bounds without any affine compatibility restriction.
\end{proposition}

\begin{proof}
The homogeneous indicators
$H_x=\ind_{\{f(x-1)=\cdots=f(x+L-1)\}}$ have zero affine target and
therefore $\E\prod_{x\in I}H_x=W(I)$ exactly.
The first assertion follows from the positive profile and the same
exceptional-channel estimates used in Theorem~\ref{p1:arith:pairs}.
On a dependent pair $2^\rho\leq2(2^\rho-1)$, so the support contribution
is already controlled by that profile.

For a dependent extension let $\delta=\dim(S_z\cap\sum_{x\in e}S_x)>0$.
Then $W(e\cup\{z\})=W(e)p2^\delta$, with
$2^\delta\leq2(2^\delta-1)$.
For the compressed old interval, test alignment at its fixed enclosing
radius $DB$ and use the same radius in the aligned-target estimate.
Lemma~\ref{p1:arith:asymmetric} gives
the nonaligned mass
$N^{-3+1/b+31/24+\varepsilon}\leq N^{-11/8+\varepsilon}$.
Lemma~\ref{p1:arith:palmrank} and $NB^{O(1)}$ aligned hosts give
$N^{-4/3+\varepsilon}$ or $N^{-5/4+\varepsilon}$.
Endpoint sources cost at most $N^{-5/3+o(1)}$.  No contact was dropped
on a compatibility argument: all starts in this proposition are separated.

Two edges sharing a vertex have $NB^{O(1)}$ hosts.  The same three-window
rank bound gives their mass $N^{-5/4+\varepsilon}$.
For disjoint edges first treat nonaligned compressed blocks, using
a fixed radius enclosing both true value sets in the alignment test.
Its complementary enlarged aligned sector is retained below.  Write
$r_f=|U_f|-1\geq5L/3-O(1)$ and let $\delta^*$ be the loss of the
even-selection rank on $U_f$ after the old normalized values and small
primes have been exposed.  The excess above $W(e)W(f)$ is at most
$W(e)W(f)(2^{\delta^*}-1)$.
The weak-component support proof of Lemma~\ref{p1:arith:asymmetric}
still applies to two subsets of fixed-width logarithmic intervals:
it needs two selected new values and one selected old value.
It gives $N^{31/24+\varepsilon}$ hosts.  In this sector
$W(e)W(f)\leq N^{-10/3+o(1)}$, so its mass is
$N^{-49/24+\varepsilon}$.

In the dense sector, the multiquadratic argument gives
$N^{1/2+\varepsilon}$ hosts.  A fixed large population of components
with one new vertex and boundedly many old vertices has only
$N^\varepsilon$ hosts, by Lemma~\ref{p1:arith:geometric}(2).
On that exceptional population use $W(f)2^{\delta^*}\leq1$ and
$W(e)\leq N^{-5/3+o(1)}$.  Otherwise all but
$O(B/K)+O_K(1)$ components use two new vertices.  Hence
\[
 \delta^*\leq|U_f|/2+O(B/K)+o(B),\qquad
 W(f)2^{\delta^*}\leq N^{-5/6+O(1/K)+o(1)}.
\]
The remaining dense mass is $N^{-2+\varepsilon}$.
The entire nonaligned excess is therefore $N^{-5/3+\varepsilon}$.

Aligned compressed pairs, including overlapping envelopes, have
$NB^{O(1)}$ hosts.  Away from endpoint sources their first three original
windows have nullity at most $3L/4+o(L)$, by
Lemma~\ref{p1:arith:palmrank}.  Adding the fourth loses at most $L$
further ranks.  Their mass is at most
$N\,N^{-4+7/4+o(1)}=N^{-5/4+o(1)}$.
The endpoint source contribution is again $N^{-5/3+o(1)}$.
Finally the sum of all baseline products is bounded by
$(\sum_{e\in\mathcal A_{\rm raw}}W(e))^2$.
All estimates were made with homogeneous ranks, proving the assertion
independently of the original phases.
\end{proof}

\subsection{Support and branching at a fixed number of blocks}

\begin{lemma}[Small-kernel counts for fixed-width blocks]
\label{p1:arith:small-kernels}
Let $U_X\subset X+[-DB,DB]$ with $D$ fixed and $X\asymp N$, and let
$H\asymp B$ exceed the diameter.  Uniformly for $1\leq T\leq CN$,
the number of translations containing one value with $K_H(n)\leq T$
is at most $N^{1/2+o(1)}T^{1/2}$.  For fixed $0<\theta<1$, the number
containing two such distinct values with $T=N^\theta$ is
$O_{C,D,\theta,\varepsilon}(N^{\max(\theta,7/12)+\varepsilon})$.
In particular it is $O(N^{\theta+\varepsilon})$ for $\theta\geq7/12$.
\end{lemma}

\begin{proof}
Enlarge the selected set to its interval.  The proofs of
Lemmas~\ref{p1:arith:anchor} and~\ref{p1:arith:twokernel} use only a
fixed logarithmic width and a prime threshold of that order.  In particular
$2^{\pi(H)}=N^{o(1)}$, the smooth parts have fixed polynomial height,
and all changes of offset count are powers of $B$.
\end{proof}

Here a projection is an element of the effective even-selection image
in the square-class space, rather than a formal coordinate vector.

\begin{proposition}[A relative relation meeting several blocks]
\label{p1:arith:fixed-support}
Fix integers $r\geq s\geq2$.  For $r$ separated logarithmic windows,
the number of tuples having a relative square relation whose projection
is nonzero in at least $s$ blocks is
\[
 O_{C,r,\varepsilon}\left(N^{r-2s/(2s+1)+\varepsilon}\right).
\]
The same assertion holds for injective even-selection spaces on subsets
of fixed-width logarithmic intervals, including the compressed good edges,
with only $B^{O_r(1)}$ choices of fixed offsets and fixed dilation factors.
\end{proposition}

\begin{proof}
Choose $s$ blocks on which a witnessing relation has nonzero effective
projection, and fix the other $r-s$ centres. Their values generate the
old space in Theorem~\ref{p13:half:conditional}. The selected $s$
nonzero projections sum into that old space; they need not remain
nonzero modulo it. The theorem therefore gives
$O(N^{s/2+\varepsilon'})$ choices for the selected centres.
Summing the $O_r(1)$ choices of blocks and $O(N^{r-s})$ other centres
gives the stronger bound $O(N^{r-s/2+\varepsilon'})$.
For $s\ge2$, $s/2\ge2s/(2s+1)$, so this implies the displayed bound.

The same theorem permits the effective images of all fixed-width blocks
in its stated height range. Fixed dilations preserve that range, and
the $B^{O_r(1)}$ offset and motif choices are absorbed by reducing
$\varepsilon'$. The injectivity in the statement identifies the
nonzero projections with nonempty even selections; no count for
formal vectors with zero effective image is used.
\end{proof}

\begin{proposition}[A fixed number of exposed source windows]
\label{p1:arith:fixed-branching}
Fix $j\geq1$ and source starts $x_1,\ldots,x_j\in I_N$.
Outside a deterministic set of $N^{1/2+o_{C,j}(1)}$ possible starts,
the new relative vector is uniform and independent of every value in
the old windows.  An explicit bound for the exceptional set is
\[
 j(2L+1)+B\sqrt{3N}\exp(2\sqrt{m_{j,N}}),\qquad
 m_{j,N}=\left\lceil\pi(B)+\frac{jB\log(3N)}{\log B}\right\rceil.
\]
Fixed dilations, subsets, and fixed enlarged logarithmic windows satisfy
the same $N^{1/2+o(1)}$ conclusion.
\end{proposition}

\begin{proof}
Let $\mathcal P$ contain all primes at most $B$ and all primes dividing
an old value.  It has at most $m_{j,N}$ members.
The number of integers at most $3N$ whose squarefree part is supported
on $\mathcal P$ is at most
\[
 \sqrt{3N}\prod_{\ell\in\mathcal P}(1+\ell^{-1/2})
 \leq\sqrt{3N}\exp(2\sqrt{|\mathcal P|}),
\]
because the $i$th distinct prime is at least $i$.
Discard overlapping starts and windows containing two such values;
the stated estimate follows even by charging each such window to one
of its values and multiplying by $B$.
If a nonzero new relative combination belonged to the old value space,
its even selection would contain at least two integers.  Every large
odd prime of each selected integer is private inside the new window,
so it would have to divide an old value.  Both selected integers would
have squarefree part supported on $\mathcal P$, contrary to the retained
set.  The new rank therefore increases by exactly $L$ over the old
value space, which proves the asserted conditional uniformity.
For fixed $j$, $m_{j,N}=O_{C,j}((\log N)^2/\log\log N)$.
\end{proof}

\subsection{Products grouped by their ternary support}

\begin{proposition}[A ternary group budget]\label{p1:arith:ternary-budget}
Let $U_i\subset X_i+[-DB,DB]$, $X_i\asymp N$, be three sets with
injective even-selection square-class maps whose images are pairwise
in direct sum.  The width constant and the number of offset patterns
are fixed.  Choose $H$ greater than the three diameters.  For an offset
triple $u=(u_1,u_2,u_3)$ let
\[
 G_u=\prod_{\substack{\ell>H\ {\rm prime}\\
       v_\ell(X_i+u_i)\ {\rm odd}\ (i=1,2,3)}}\ell.
\]
Retain the groups greater than one and write their decreasing values
$G_{(1)}>G_{(2)}>\cdots$.  Put
$\mathcal G_k=\prod_{i=1}^kG_{(i)}$, or zero if fewer than $k$ exist.
For every fixed $k\geq3$, $\varepsilon>0$, and $1\leq T\leq N$,
\begin{equation}\label{p1:arith:ternarycount}
 \#\{\text{tuples}:\mathcal G_k>T\}
 \ll_{C,D,k,\varepsilon}N^\varepsilon
      (N^3/T^2+N^{6/5}).
\end{equation}
\end{proposition}

\begin{proof}
A prime larger than $H$ has only one occurrence in each block, so the
groups are pairwise coprime and
$G_u\mid\gcd(X_1+u_1,X_2+u_2,X_3+u_3)$.
Two groups satisfy $G_uG_v\ll NB$.
If they share an occurrence this follows from divisibility of that
integer.  Otherwise choose two blocks in which the offsets differ.
Their product divides a determinant $nm'-n'm$ of size $O(NB)$.
Its vanishing would give a nonzero relative relation on those two
blocks, contrary to their direct sum.

Three groups with product greater than a sufficiently large constant
times $NB^2$ force their three value vectors to have zero determinant.
The corresponding offset triples cannot be collinear.  Indeed, writing
$u_j=u_1+t_jd$, with $d$ a primitive integer vector and $t_j$ integers,
in the collinear case, all three coprime groups divide
$d\times(X+u_1)$.  If this vector is nonzero, its size $O(NB)$ gives a
contradiction.  If it is zero, two positive value vectors are proportional
with distinct coordinates in every block, again giving a forbidden pair
relation.  Thus the centers lie on a genuine plane
\[
 ((u_2-u_1)\times(u_3-u_1))\cdot X
 =-((u_2-u_1)\times(u_3-u_1))\cdot u_1.
\]
There are only $B^{O(1)}$ such planes, determined by the offsets alone.

For fixed $k$ group first by the squarefree product $Q=\mathcal G_k$.
The ordered factorizations cost $d_k(Q)\ll_{k,\varepsilon'}Q^{\varepsilon'}$,
and the supports cost $B^{3k}$.  CRT fixes one residue class modulo $Q$
for each center.  Summing $O((N/Q+1)^3)$ gives, for
$1\leq T\leq N\leq R\leq N^{O(1)}$,
\begin{equation}\label{p1:arith:groupcrt}
 \#\{T<\mathcal G_k\leq R\}
 \ll_{k,\varepsilon}N^\varepsilon(N^3/T^2+R).
\end{equation}
For $\mathcal G_3>R\gg NB^2$, use its forced plane.  The two largest
groups have product $P>R^{2/3}$ and $P\ll NB$.  Two free coordinates
are fixed by CRT modulo $P$; the plane determines the third.
Summing over the two positive group sizes and their offsets gives
\begin{equation}\label{p1:arith:planecrt}
 \#\{\mathcal G_3>R\}
 \ll_\varepsilon N^\varepsilon(N^2R^{-2/3}+N).
\end{equation}
Here one uses
$\sum_{ab>A}(ab)^{-2}\ll\log(2A)/A$,
$\sum_{ab\leq U}(ab)^{-1}\ll\log^2(2U)$, and
$\#\{ab\leq U\}\ll U\log(2U)$.
The third numerical group is not enumerated again: it certifies the
plane and the lower bound for $P$.

For $k=3$, take $R=N^{6/5}$ in
\eqref{p1:arith:groupcrt}--\eqref{p1:arith:planecrt}.
For induction on fixed $k$, the region
$T<\mathcal G_k\leq N^{6/5}$ is already covered by
\eqref{p1:arith:groupcrt}.  Beyond it,
$\mathcal G_{k-1}>N^{6(k-1)/(5k)}$.
At $k=4$ the preceding $N^3/T^2$ term is $N^{6/5}$; at $k=5$
it is smaller; for $k\geq6$ the threshold is at least $N$ and the
preceding tail at $N$ suffices.  Only a fixed number of induction steps
is taken.  This proves \eqref{p1:arith:ternarycount}.
\end{proof}

\subsection{The probability of two active good edges}

\begin{lemma}[A positive second moment for good edges]
\label{p1:arith:twogoodpositive}
Let $\Xi_N=\sum_{e\in\mathcal E_N}\prod_{x\in e}J_x$.
Then
\[
 \E(\Xi_N)_2=\beta_N^2+o_C(\beta_N^2)
 \ll_C L^2N^{-4/3}.
\]
In particular
$0\leq\beta_N-\P(\Xi_N\geq1)\leq\tfrac12\E(\Xi_N)_2$.
\end{lemma}

\begin{proof}
We only need a strict gap below $2L/3$ for the nullity of three aligned
windows containing a good source.  Here is a direct bound that avoids
a complete classification of the optimal three-window configuration.
Keep the target parametrization and congruences from
Lemma~\ref{p1:arith:palmrank}.  If its fraction $c/d$ is physical,
two of the three parameter intervals have the same coefficient and
disjoint interiors, up to one endpoint.  Their intersections with the
third interval have total length at most that third interval.  Thus
the full three-window nullity is at most $L/2+o(L)$ for $2:3$, and
$L/3+o(L)$ for $3:4$.  This bound includes the original source nullity.
The residual on this fixed-width normalized union is $o(L)$ by the
private-prime defect estimate, just as in the fixed-height case of
Lemma~\ref{p1:arith:alignedrank}.

For a nonphysical target above $2:3$, the additional rank loss is at
most $3L/10+o(L)$.  Indeed the density
$((c,2)+(c,3)-1)/c$ already suffices except for $c=9,12$ and the
physical or endpoint fractions.  For $c\geq14$ it is at most $2/7$;
$c=5,7,8,10,11,13$ also satisfy the desired bound directly.
For $c=9$, the allowed reduced denominators are $4,5$.
Count the new class modulo nine everywhere, and the remainder of
the class modulo three inside the true three-family parameter
interval of length $3B/d$.  This gives
\[
 B/9+(3B/d)(1/3-1/9)+O(1)
 =B/9+2B/(3d)+O(1)\leq5B/18+O(1)<3B/10+O(1).
\]
Counting only the three-family units would miss possible matches
outside that actual interval; the union count retains them.
For $c=12$ the possible denominators are $5,7$.
At $d=7$ the two true old intervals supply at most $2B/7+O(1)$ units.
At $d=5$ count the entire new class modulo six and the part of the
class modulo four outside it.  The latter lies in a true parameter
interval of length $4B/5$, giving
\[
 B/6+(4B/5)(1/4-1/12)+O(1)=3B/10+O(1).
\]
The full nullity is therefore at most $(1/3+3/10)L+o(L)=19L/30+o(L)$.

Above $3:4$, the cruder additional bound $2L/5+o(L)$ is enough.
The density test proves it for $c\geq15$ and for
$c=5,7,9,10,11,13,14$.  The only other nonphysical interior
fractions are $8/3$ and $12/5$.
For $8/3$, count the class modulo eight everywhere and the remainder
of the class modulo two inside the true four-family parameter
interval of length $2B/3$.  The bound is
\[
 B/8+(2B/3)(1/2-1/8)+O(1)=3B/8+O(1).
\]
For $12/5$, each true old family supplies at most $B/5+O(1)$ units.
Thus the full nullity is at most $(1/4+2/5)L+o(L)=13L/20+o(L)$.
The finitely many endpoint fractions concern only $B^{O(1)}$ sources,
as before.  Since $19/30<13/20<2/3$, the bound $13L/20+o(L)$
works for both channels outside this exception.

For two distinct good edges, configurations sharing a vertex and
aligned configurations with four vertices have $NB^{O(1)}$ hosts.
In the four-vertex case adding the last window loses at most $L$
ranks.  The joint probabilities in both cases therefore total at most
$N^{-27/20+o(1)}$.  Endpoint sources cost $N^{-5/3+o(1)}$.
Nonaligned compressed edges have total absolute covariance
$O_{C,\varepsilon}(N^{-5/3+\varepsilon})$ by the two-block argument in
Proposition~\ref{p1:arith:fixed-homogeneous-seeds}; that argument also
applies to compatible affine sources.
Cross-contacts at distance less than $L$ have zero joint probability.
Their baseline products, those of adjacent edges, and the diagonal
$\sum_e q_e^2$ have total $N^{-7/3+o(1)}$, since the geometric edge
degree is $O(B)$.  Contacts at distance exactly $L$ were included in
the physical bounds above.  Consequently
\[
 \E(\Xi_N)_2=\beta_N^2+O_{C,\varepsilon}(N^{-27/20+\varepsilon}).
\]
Take $\varepsilon<1/60$ and use
$\beta_N\asymp LN^{-2/3}$.  The final inequality is the elementary
pointwise bound $(m-1)_+\leq m(m-1)/2$ for integers $m\geq0$.
\end{proof}


\section{Coefficient budgets, coordinate restrictions, and fixed channels}
\label{p1:arith:supplement}

The first two tools below do not depend on the critical-window scaling.
They control, respectively, the small-prime coefficients of selected
values and the repetition of coordinate forms in a binary relation space.
We then give the extension estimates for any fixed rational channel.
In that part, the channel and every finite family of channels are fixed
before \(N\) tends to infinity.

\subsection{A deterministic budget for the small-prime part}

For a positive integer \(n\), write
\[
 u_B(n)=\frac{\sf(n)}{K_B(n)}
       =\prod_{\substack{\ell\leq B\ {\rm prime}\\v_\ell(n)\ {\rm odd}}}\ell.
\]

\begin{lemma}[A small-prime coefficient budget]
\label{p1:arith:smallpart-budget}
Let \(V\) be any interval of \(B\geq2\) consecutive positive integers.
Then
\[
 \sum_{n\in V}\log u_B(n)\leq B\log B+O(B),
\]
with an absolute constant independent of the position of \(V\).
For every fixed \(A>0\), at most
\[
 \bigl(A^{-1}+O_A(1/\log B)\bigr)B
\]
of its values satisfy \(u_B(n)>B^A\).
\end{lemma}

\begin{proof}
Dropping the parity condition gives
\[
 \sum_{n\in V}\log u_B(n)
 \leq\sum_{\ell\leq B}
       \left(\left\lfloor\frac B\ell\right\rfloor+1\right)\log\ell
 \leq\log(B!)+\sum_{\ell\leq B}\log\ell.
\]
Here the sums are over primes.  Legendre's formula gives the second
inequality, and the elementary Chebyshev bound makes the last sum
\(O(B)\).  Finally \(\log(B!)\leq B\log B\).
Each value counted in the last assertion contributes more than
\(A\log B\) to the nonnegative sum.
\end{proof}

\subsection{Regularising coordinate forms by independent restrictions}

\begin{lemma}[Coordinate regularisation]
\label{p1:arith:coordinate-regularisation}
Let \(R\subset\F_2^V\) have dimension \(r\), put \(n=|V|\), and let
\(\ell_v\in R^*\) be the coordinate forms.  Suppose that \(s\) distinct
nonzero forms occur among them.  For integers \(K,d\geq1\), there is
\(D\subset V\) such that
\[
 R'=\{x\in R:x_v=0\text{ for every }v\in D\},\qquad
 \operatorname{codim}_R R'=|D|
 \leq\frac{n}{K+1}+\frac{s-r}{d+1},
\]
and the restricted coordinate forms have the following properties:
\begin{enumerate}
 \item every nonzero form has multiplicity at most \(K\);
 \item every nonzero form \(\ell\) has at most \(d\) unordered
 decompositions \(\ell=a+b\), where \(a,b\) are distinct nonzero
 coordinate forms present on \(R'\).
\end{enumerate}
Coordinates whose restricted forms vanish may then be removed without
any additional loss of dimension.
\end{lemma}

\begin{proof}
The coordinate forms span \(R^*\), and their restrictions continue to
span the dual of every subspace encountered below.
At an intermediate stage let \(S\) be the set of distinct nonzero forms
and let \(r_0\) be their rank.  Setting a present nonzero form \(\ell\)
equal to zero reduces the dimension by one and replaces the forms by
their images modulo \(\langle\ell\rangle\).
The form \(\ell\) disappears.  All other collisions are exactly the
disjoint unordered pairs \(\{a,a+\ell\}\subset S\).
If there are \(t(\ell)\) such pairs, then
\[
 |S_{\rm new}|=|S|-1-t(\ell),\qquad
 (|S|-r_0)_{\rm new}=|S|-r_0-t(\ell).
\]
Thus the nonnegative potential \(|S|-r_0\) never increases.

If a form has multiplicity greater than \(K\), set it equal to zero.
This permanently kills at least \(K+1\) previously nonzero coordinates.
There are at most \(n/(K+1)\) such steps, since these sets of killed
coordinates are disjoint.
When the multiplicity condition holds but a form has more than \(d\)
decompositions, set that form equal to zero instead.
Such a step decreases the potential by at least \(d+1\), so there are
at most \((s-r)/(d+1)\) of these steps.
Each step decreases the dimension, and the procedure therefore ends
with both required properties.

At every step select one coordinate carrying the form just killed.
Its restriction was nonzero before that step, so its original coordinate
form is independent of the previously selected forms.
The selected coordinates constitute \(D\), and consequently
\(\operatorname{codim}_R R'=|D|\).
\end{proof}

\begin{corollary}[Constraints meeting three blocks]
\label{p1:arith:ternary-degree}
Suppose \(V\) is partitioned into three blocks and each row of a family
of binary constraints annihilating \(R\) has at most one occurrence in
each block.  Apply Lemma~\ref{p1:arith:coordinate-regularisation} and
delete the coordinates whose forms on \(R'\) are zero.
Every remaining row has weight zero, two, or three.
A row of weight two joins equal nonzero forms.
A row of weight three has three distinct nonzero forms \(a,b,c\) with
\(a+b+c=0\).

The connected components formed by the rows of weight two have at most
\(K\) vertices.  After repeated rows with the same support are identified,
the number of ternary rows incident to any vertex is at most \(2dK^2\).
\end{corollary}

\begin{proof}
The sum of the forms on a row is zero, since that row annihilates \(R'\).
A row of weight one is therefore impossible.  At weight two its forms
are equal.  At weight three a repeated form would force the third form
to be zero, which has already been excluded.
Binary components stay inside a class of equal forms, of size at most
\(K\).  At a vertex of form \(\ell\), a ternary row chooses one of at
most \(d\) decompositions \(\ell=a+b\), then the other two coordinates
from classes of size at most \(K\).  Allowing their two assignments to
the remaining blocks gives the asserted \(2dK^2\) bound.
\end{proof}

\begin{remark}[Order of restrictions]
If some coordinates must first be killed for an arithmetic reason, do
this before applying the regularisation lemma.  A further restriction
after regularisation can identify previously distinct forms and destroy
both bounds.  The lemma is an existence statement about a subspace;
its proof does not enumerate possible relation spaces or histories of
restrictions.
\end{remark}

\subsection{A finite rank bound for every fixed channel}

Fix coprime integers \(a,b\) with
\[
 b\geq3,\qquad b/2<a<b.
\]
Let \(\mathcal G_{a:b}\) be the symmetric family of ordered separated
pairs in this channel with at least two exact units and with a private
prime coordinate on every distinct normalized value in
\(bV_x\cup aV_y\), above a fixed threshold \(H_{a,b}=C_{a,b}B\).
Choose the constant large enough for the fixed normalized widths.
For \(\alpha=(x,y)\) put
\[
 q_\alpha=\P(J_x=J_y=1),\qquad
 M_{a:b}=\sum_{\alpha\in\mathcal G_{a:b}}q_\alpha.
\]
Pairs of probability zero remain in the family and are omitted from
conditional expressions.  The private-prime condition makes the
relative kernel exactly the canonical even-selection kernel.

\begin{lemma}[Mass of a fixed channel]
\label{p1:arith:fixed-channel-mass}
With \(S_{a,b}(L)\) as defined before
\eqref{p1:arith:channelidentity}, and
\(w_{a,b}=2(2a-b)/(ab)\), one has
\[
 M_{a:b}=p^2Nw_{a,b}S_{a,b}(L)
          +O_{a,b,C,\varepsilon}(N^{-1+\varepsilon}).
\]
In particular \(M_{a:b}\asymp_{a,b,C}L N^{-1+1/b}\).
\end{lemma}

\begin{proof}
In either orientation the canonical parameter satisfies
\(N/a\leq t<2N/b\), up to fixed multiples of \(B\) at the endpoints.
The two orientations therefore supply \(Nw_{a,b}\) translations per
choice of the within-window placements.
The diameter and sign calculation giving
\eqref{p1:arith:channelidentity} applies for this fixed \(a,b\).

Removing the normalized defective pairs costs
\(N^{-3/2+1/b+o(1)}\): there are \(N^{1/2+o(1)}\) such hosts,
their nullity is at most \(L/b+o(L)\), and their weights include \(p^2\).
Endpoint errors are smaller.
On the retained dependent pairs the signed kernel contribution is
\(q_\alpha-p^2\).  Adding back \(p^2\) costs \(O_{a,b}(p^2NB)\),
because each start has \(O_{a,b}(B)\) possible partners.
This proves the displayed formula.
Its leading term is positive and of order \(L N^{-1+1/b}\), by the
finite identity, uniformly in the critical band.
\end{proof}

Define the finite, explicitly computable number
\begin{equation}\label{p1:arith:gamma-fixed-channel}
 \gamma_{a:b}=
 \max\left(
 \left\{\frac1b\right\}\cup
 \left\{\frac{(c,a)+(c,b)-1}{c}:
 \begin{array}{l}
 1\leq c\leq b(a+b-1),\ d\geq1,\ (c,d)=1,\\
 b/2<c/d<2a,\quad c\nmid a,\quad c\nmid b
 \end{array}\right\}\right).
\end{equation}
All variables in the second set are integers.

\begin{lemma}[Additional rank in a fixed channel]
\label{p1:arith:fixed-channel-rank}
One has \(\gamma_{a:b}\leq2/3\).
Apart from \(B^{O_{a,b}(1)}\) source pairs at the dyadic endpoints,
every aligned target \(z\ne x,y\) compatible with a source
\((x,y)\in\mathcal G_{a:b}\) satisfies
\[
 \delta_{x,y}(z):=
 L-\dim\frac{S_x+S_y+S_z}{S_x+S_y}
 \leq\gamma_{a:b}L+o_{a,b,C}(L).
\]
This definition includes any internal rank deficiency of the new window.
The bound remains true when the prime signs of \(a,b\) are included in
the old information.
\end{lemma}

\begin{proof}
Write \(x=at-i_0\), \(y=bt-j_0\), choosing the lower orientation.
As in Lemma~\ref{p1:arith:palmrank}, the normalized old values lie in
a fixed-width interval about \(abt\).  An aligned target has the form
\(z=(ct+h)/d\), with \((c,d)=1\), numerator, denominator and offsets
of polynomial size in \(B\).
The dyadic positions require
\(b/2\leq c/d\leq2a\), up to \(B^{O(1)}/N\).
Separation of rational fractions excludes strict violations of this
interval once \(N\) is large.
Either equality forces \(at=N+B^{O(1)}\) or \(bt=2N+B^{O(1)}\);
these give the stated exceptional sources.

For the interior fractions, a new offset \(k\) matched to a true old
value satisfies one of
\[
 h+dk\equiv0\pmod{c/(c,a)},\qquad
 h+dk\equiv0\pmod{c/(c,b)}.
\]
Their intersection has modulus \(c\).
The union contains at most
\[
 B\,\frac{(c,a)+(c,b)-1}{c}+O(1)
\]
new positions.  The same residual argument as in
Lemma~\ref{p1:arith:palmrank} leaves only \(o(L)\) further dimensions.
Matching to a value in an enclosing interval but absent from the true
old set contributes no new freedom.

Suppose \(c\nmid a\) and \(c\nmid b\), and put \(g=(c,a)\), \(h_1=(c,b)\).
These numbers are coprime and \(gh_1\mid c\).
If one equals one, the other is a proper divisor of \(c\), giving
density at most \(1/2\).
Otherwise \(g,h_1\geq2\) and
\[
 \frac{g+h_1-1}{c}
 \leq\frac1g+\frac1{h_1}-\frac1{gh_1}\leq\frac23.
\]
The last maximum occurs at the smallest coprime pair \(\{2,3\}\).
When \(c>b(a+b-1)\), the density is less than \(1/b\).
All other nonphysical fractions are included in
\eqref{p1:arith:gamma-fixed-channel}.

If \(c\mid a\), a proper divisor has \(c\leq a/2<b/2\), which is
impossible here; thus \(c=a,d=1\).
Similarly \(c\mid b\) leaves only \(c=b,d=1\).
These are the physical fractions reproducing one old family.
Compatibility gives distance at least \(L\) from that source, so the
two true windows share at most one value.
Matches to the other old family have steps \(a,b\) in intervals of
width \(B\), and number at most \(B/b+O(1)\).
This proves the rank bound in these cases as well.
\end{proof}

\begin{proposition}[Palm mean for a fixed channel]
\label{p1:arith:fixed-channel-palm}
Put
\[
 \eta_{a:b}=\min\{17/24,1-\gamma_{a:b}\}\geq1/3.
\]
Then, with the sum over \(z\in I_N\setminus\{x,y\}\),
\[
 R_{a:b}:=
 \frac1{M_{a:b}}\sum_{\alpha=(x,y)\in\mathcal G_{a:b}}q_\alpha
 \sum_z\bigl(\P(J_z=1\mid J_x=J_y=1)-p\bigr)_+
 \ll_{a,b,C,\varepsilon}N^{-\eta_{a:b}+\varepsilon}.
\]
Moreover
\[
 \frac1{M_{a:b}}\sum_{\alpha=(x,y)\in\mathcal G_{a:b}}q_\alpha
 \E[Z_N-2\mid J_x=J_y=1]
 =\lambda+
 O_{a,b,C,\varepsilon}
 \bigl(N^{-\min(\eta_{a:b},1/2)+\varepsilon}\bigr).
\]
These bounds are not asserted to be optimal, or uniform in growing
channel height.
\end{proposition}

\begin{proof}
The source probability is \(O_{a,b,C}(p^2 2^{L/b})\).
For the compressed old interval choose a fixed enclosing radius \(DB\).
Use this same radius in the alignment test and the complementary
aligned-target sector; the large-prime cutoff is chosen accordingly.
The nonaligned asymmetric estimate
in Lemma~\ref{p1:arith:asymmetric}, with the fixed larger widths and
\(B^{O_{a,b}(1)}\) patterns, bounds the total positive excess by
\[
 p^3 2^{L/b}N^{31/24+\varepsilon'}.
\]
Division by \(M_{a:b}\asymp NLp^2 2^{L/b}\) gives
\(N^{-17/24+\varepsilon}\).

Each source has \(B^{O_{a,b}(1)}\) aligned targets.
Lemma~\ref{p1:arith:fixed-channel-rank} bounds their positive conditional
excess by \(pN^{\gamma_{a:b}+o(1)}\) per target.
The exceptional endpoint sources have total weight
\(B^{O(1)}p^2 2^{L/b}\).
Even bounding every conditional probability by one, their aligned
contribution after division by \(M_{a:b}\) is \(N^{-1+o(1)}\).
This proves the first assertion, including all new windows.

For the mean, Proposition~\ref{p1:arith:fixed-branching} with two
exposed source windows leaves at most \(N^{1/2+o(1)}\) exceptional
targets for each source.
Outside them the new relative vector is uniform and independent of
the old values, so its conditional probability is exactly \(p\).
The total negative deviation is therefore at most \(pN^{1/2+o(1)}\).
Combining it with the positive estimate and the removal of the two
sources, which costs \(2p\), proves the assertion.
\end{proof}

\subsection{The ternary mass generated by a finite family of pairs}

\begin{corollary}[Triples containing a good pair]
\label{p1:arith:finite-family-triples}
Let \(\mathcal F\) be a nonempty finite fixed family of admissible
channels, and set
\[
 M_{\mathcal F}=\sum_{a:b\in\mathcal F}M_{a:b},\qquad
 \nu_{\mathcal F}=
 \min\{1/2,\min_{a:b\in\mathcal F}\eta_{a:b}\}.
\]
The sum of \(q(\{x,y,z\})\) over all ordered triples of distinct starts
containing at least one good pair from this family is
\[
 3\lambda M_{\mathcal F}
 +O_{\mathcal F,C,\varepsilon}
    (M_{\mathcal F}N^{-\nu_{\mathcal F}+\varepsilon}).
\]
In particular, for fixed \(Q\geq5\), the triples containing a good
pair of height \(5\leq b\leq Q\) have total mass
\(O_{Q,C,\varepsilon}(N^{-4/5+\varepsilon})\).
\end{corollary}

\begin{proof}
Choose which two positions of an ordered triple form the designated
pair.  For each of the three choices,
Proposition~\ref{p1:arith:fixed-channel-palm} gives
\(\lambda M_{\mathcal F}
 +O(M_{\mathcal F}N^{-\nu_{\mathcal F}+\varepsilon})\).
The pair orientations are already included in \(M_{\mathcal F}\).

This designated-pair sum counts a triple once for each good pair that
it contains.  If the third window is independent of all old values
of the designated pair, it cannot form a dependent good pair with
either source.
Every overcount therefore has its third start in the exceptional set
of Proposition~\ref{p1:arith:fixed-branching}.
For each designated source, the conditional mass of this set is at most
\[
 pN^{1/2+o(1)}
 +\sum_z\bigl(\P(J_z=1\mid J_x=J_y=1)-p\bigr)_+.
\]
Its weighted average is
\(O_{\mathcal F,C,\varepsilon}(N^{-\nu_{\mathcal F}+\varepsilon})\).
This bounds the entire overcount, since a triple has at most three pairs.
The finite channel families are disjoint for large \(N\), by the separation
of distinct reduced slopes.
Finally, Lemma~\ref{p1:arith:fixed-channel-mass} gives
\(M_{\mathcal F}\ll_{Q,C}L N^{-4/5}\) for heights between five and \(Q\).
\end{proof}

The last corollary concerns the part of the third moment containing a
good pair.  It makes no assertion about the primitive remainder or about
a family whose heights grow with \(N\).


\section{Pair profiles and the regular-set estimates}
\label{p14:comp:pair-profiles}

This section supplies the estimates used in Section~\ref{p1:arith:section}.
The regular-set transfer is used for the explicit channel formulas and for
the positional comparison. The scalar reconstruction with exact marginals
does not discard these exceptional windows.

\subsection{Regular-set transfer}
\label{p14:comp:regular-transfer}

\begin{proof}[Proof of Proposition~\ref{p1:arith:core}]
The one-window estimates proved in Lemma~\ref{p1:arith:defects} give
$m_B(x)=O_C(B/\log B)$ uniformly, and only $N^{o(1)}$ windows contain two
$B$-defective integers.  A nonzero relative relation is an even nonempty
selection.  Private prime coordinates therefore force it to use at least two
defects.  Its dimension is at most $m_B(x)$.  Consequently
\[
 \sum_{x\in I_N}(2^{L-\dim S_x}-1)=N^{o(1)}.
\]
Fourier inversion proves \eqref{p1:arith:marginalsum}.  The sum of the
probabilities of active discarded starts is at most $pN^{o(1)}$; the coupling
using the original prime signs proves the total-variation assertion.  On $G_N$
the marginal is exactly $p$, giving both assertions about the mean.

For the last statement, let $D=\{Z_N\ne Z_N^\circ\}$.  Since
$0\leq Z_N^\circ\leq Z_N$, H\"older's inequality gives, for any fixed $s>1$,
\[
 \E\big[Z_N^j\ind_D\big]
 \leq(\E Z_N^{js})^{1/s}\P(D)^{1-1/s}.
\]
Choose a sufficiently large integer $s$, and then use the fixed-moment bound.
This last step is used only after that bound has been established; a small
coupling error alone would not justify a variance transfer.
\end{proof}

\subsection{Sparse hosts and weighted residual mass}
\label{p14:comp:sparse-hosts}

\begin{proof}[The elementary bounds preceding Proposition~\ref{p1:arith:hosts}]
A relative relation that is internal to one window contributes at most
$N^{1+o(1)}$ hosts, by Proposition~\ref{p1:arith:core}.  Otherwise it selects
at least two vertices in each window.  Put $T=N^{3/5}$.  If all selected
kernels are at most $T$, each window contains two small kernels.  The estimate
\begin{equation}\label{p1:arith:twokernelused}
 \#\{x\in I_N:\#\{n\in V_x:K_B(n)\leq N^\beta\}\geq2\}
   \ll_{C,\beta,\varepsilon}N^{\max(\beta,7/12)+\varepsilon}
\end{equation}
proved in Lemma~\ref{p1:arith:twokernel} gives $N^{6/5+\varepsilon}$
choices.  If one selected kernel exceeds $T$, Lemma~\ref{p1:arith:anchor}
gives
$N^{3/2+o(1)}T^{-1/2+\tau}$, for every fixed $\tau>0$.
Choose $\tau$ sufficiently small.

For the weighted assertion, first restrict to $c(x,y)<\eta B$.
Equation~\eqref{p1:arith:graphbound} bounds the weight by
$2^{\eta B+o(B)}$; choose $\eta$ after the desired error exponent.
The complementary sector has total weight $O_{C,\eta,\varepsilon}(N^{1+\varepsilon})$
by Lemma~\ref{p1:arith:dense}.  The latter estimate uses the component
equations themselves, rather than just their support: many small components
put the pair on a multiquadratic surface, and its low-degree curves are either
aligned or absent from the dyadic region.
This proves both bounds with exponent $6/5+\varepsilon$, which is all
that the pair-profile proof below uses. The exponent $1+\varepsilon$
in the revised statement is proved later in
Section~\ref{p14:comp:confinement}; no late confinement estimate is
being inferred from this elementary calculation.
\end{proof}

\subsection{The complete pair profile}
\label{p14:comp:complete-profile}

\begin{proof}[Proof of Theorem~\ref{p1:arith:pairs}]
For the channel $2:3$, either orientation permits
$N/2\leq t<2N/3$, up to $O(B)$ endpoint changes.  For $3:4$ the corresponding
interval is $N/3\leq t<N/2$.  Both intervals have length $N/6+O(B)$.
The two orientations therefore give $N/3$ times the finite sums above.
Translations near a boundary contribute $O(B^3 2^{L/b})$.

For either fixed channel choose $H$ a sufficiently large fixed multiple of
$B$ and normalize the values to $bV_x\cup aV_y$.  Its diameter is $O(B)$.
If its distinct values are all nondefective at level $H$, each has a private
prime and the only relations before normalization are the even selections
of exact units.  Lemma~\ref{p1:arith:defects} shows that the exceptional
channel pairs number $N^{1/2+o(1)}$.  Lemma~\ref{p1:arith:alignedrank} bounds
their nullity by $L/b+o(L)$.  Replacing the rational kernel by the full
arithmetic kernel consequently costs at most
$N^{1/2+1/b+o(1)}$, for either the positive or the signed sum.

There are $N B^{O(1)}$ aligned hosts altogether.  The slope of height two
forces both starts within $O(B)$ of the dyadic endpoints; its mass is
$N^{1/2+o(1)}$.  Height one has no exact unit in separated windows.
The only reduced slopes of heights three and four that occur in the interior
are $2:3$ and $3:4$.  Every remaining channel has height at least five and
nullity at most $L/5+o(L)$, again by Lemma~\ref{p1:arith:alignedrank}; its
total contribution is $N^{6/5+o(1)}$. The elementary version just
proved controls the nonaligned sector at the same sufficient scale.
The inequality following
\eqref{p1:arith:affine} permits exactly the same error estimates for the signed
sum.

At distances $1,\ldots,L-1$, two run starts are incompatible.  At distance
$L$, their edge systems form a tree on an interval of diameter $2L$;
Lemma~\ref{p1:arith:defects}, at threshold $2L+1$, bounds its nullity by
$o(L)$.  Hence the entire contact band changes the second factorial moment
by $O(N^{-1+o(1)})$.  Multiplying \eqref{p1:arith:signedprofile} by $p^2$,
and using $\mu_N-\lambda=O(N^{-1+o(1)})$, proves
\eqref{p1:arith:cumulantexpansion}.  Finally
$2^{(L-r)/b}=2^{-r/b}(N/\lambda)^{1/b}$ in
\eqref{p1:arith:S23}--\eqref{p1:arith:S34} gives the two equivalents.
\end{proof}

\subsection{The good-edge mass}
\label{p14:comp:good-edge-mass}

\begin{proof}[Proof of Corollary~\ref{p1:arith:edgefamily}]
Each start has $O(B)$ possible partners in a fixed channel.  The rational
kernel has dimension one less than the number of exact units when that
number is positive, giving the rank bound.  All pairs excluded from the good
family have precisely the error mass estimated in the proof of
Theorem~\ref{p1:arith:pairs}, together with the contact band.
For good compatible pairs the centered contribution is $q_e-p^2$.
Since $p^2|\mathcal E_N|=O(B/N)$, replacing it by $q_e$ has the stated cost.
The asymptotic follows from \eqref{p1:arith:leadingphase}.
\end{proof}

\subsection{A further window after a good edge}
\label{p14:comp:pair-palm}

The following channel-specific estimate is used in the first reconstruction
and in the positional comparison. Its good-edge notation is that of
Section~\ref{p1:arith:section}.

\begin{proposition}[Positive Palm cost of a good edge]\label{p1:arith:palm}
For $a:b\in\{2:3,3:4\}$ let $\mathcal E_{a,b}$ be its good edges and put
$\beta_{a,b}=\sum_{e\in\mathcal E_{a,b}}q_e$.  Then
\begin{align}
 \frac1{\beta_{2,3}}\sum_{\{x,y\}\in\mathcal E_{2,3}}q_{xy}
    \sum_{z\in G_N\setminus\{x,y\}}
       (\P(J_z=1\mid J_x=J_y=1)-p)_+
   &\ll_{C,\varepsilon}N^{-2/3+\varepsilon},\label{p1:arith:palm23}\\
 \frac1{\beta_{3,4}}\sum_{\{x,y\}\in\mathcal E_{3,4}}q_{xy}
    \sum_{z\in G_N\setminus\{x,y\}}
       (\P(J_z=1\mid J_x=J_y=1)-p)_+
   &\ll_{C,\varepsilon}N^{-1/2+\varepsilon}.\label{p1:arith:palm34}
\end{align}
Consequently the unnormalised positive two-block masses are
$O_{C,\varepsilon}(N^{-4/3+\varepsilon})$ and
$O_{C,\varepsilon}(N^{-5/4+\varepsilon})$, respectively.
\end{proposition}

\begin{proof}
We give the two geometric estimates needed for this bound in
Lemma~\ref{p1:arith:asymmetric} and Lemma~\ref{p1:arith:palmrank}.
Compress the source pair using Lemma~\ref{p1:arith:compression}.
Choose a fixed $D$ enclosing the true compressed old values and use
radius $DB$ both for the asymmetric nonaligned sector and for its
complementary aligned sector.  The expanded old-centre version of
Lemma~\ref{p1:arith:asymmetric} applies with the new centre in $I_N$.
Away from alignment with this source, the asymmetric profile gives
$N^{31/24+\varepsilon}$ for the sum of $2^{\delta}-1$, where $\delta$ is
the loss of new rank.  There are $B^{O(1)}$ source offset patterns and
$q_{xy}\ll p^2 2^{L/b}$.  Multiplication by the additional factor $p$
and division by $\beta_{a,b}\asymp LN^{-1+1/b}$ give
$N^{-17/24+\varepsilon}$.

For aligned targets there are $B^{O(1)}$ possibilities per source.  Outside
$B^{O(1)}$ source pairs at the dyadic endpoints, Lemma~\ref{p1:arith:palmrank}
bounds the new rank loss by $L/3+o(L)$ for $2:3$ and $L/2+o(L)$ for $3:4$.
Thus their positive conditional cost is at most
$pB^{O(1)}2^{L/3+o(L)}$ or $pB^{O(1)}2^{L/2+o(L)}$.
The exceptional source mass is $N^{-2+1/b+o(1)}$ and each has only
$B^{O(1)}$ aligned targets, so after normalization it contributes
$N^{-1+o(1)}$.  This proves both estimates.  The source mass itself is
$N^{-1+1/b+o(1)}$, which gives the final assertions.
\end{proof}

\section{Fixed-source geometry and the independent moment proof}
\label{p14:comp:fixed}

The statements of the fixed-source estimates and the fixed-moment theorem
are in Section~\ref{p1:fixed:section}. This section supplies the geometric
count and the positive classification used in their proofs. The original
window estimates below use the regular-set convention $p=2^{-L}$;
Section~\ref{p13:exact:section} gives their extension to effective spaces.
The moment proof is independent of either reconstruction of the law.

\subsection{Uniformly few targets of large rank loss}

We first record the one-parameter geometric estimate in the
precise form needed here.

\begin{lemma}[Several disjoint square equations]
\label{p1:fixed:curve}
Fix \(K,A,D_0>0\) and \(\varepsilon>0\).
There is a fixed integer \(j\) such that the following holds.
Let \(P_1,\ldots,P_j\) be nonconstant products of distinct
linear factors \(Z+u\), of degrees at most \(K\), with
\(|u|\le D_0B\). Suppose their root lists are pairwise disjoint,
and let \(0<|D_\nu|\le N^A\) be integers. Then
\[
 \#\{z\asymp N:
       w_\nu^2=D_\nu P_\nu(z)\ (1\le\nu\le j)
       \text{ for some integers }w_\nu\}
       \ll_{K,A,D_0,\varepsilon}N^\varepsilon .
\]
The constants are uniform in the displayed coefficients and
offsets.
\end{lemma}

\begin{proof}
The square classes of the functions \(D_\nu P_\nu(Z)\) are
geometrically independent: a root of \(P_\nu\) has odd
valuation in that class and even valuation in every other
class. Thus the resulting curve is geometrically integral
and has degree \(2^j\) over the \(Z\)-line.

Choose an integer \(D\ge(K+A)/2\). Give \(T,S\) degree one
and each \(W_\nu\) degree \(D\), and use the equations
\[
 W_\nu^2
   =D_\nu S^{2D-\deg P_\nu}
       \prod_{u\in P_\nu}(T+uS).
\]
The graded ring is free of rank \(2^j\) over \(\mathbb Q[T,S]\);
its generic localization is a field. Its \(D\)-th Veronese
ring is generated in degree one by the monomials in \(T,S\)
of degree \(D\) and the \(W_\nu\). The corresponding
projective curve has degree \(D2^j\), since its hyperplane
bundle is the pullback of \(\mathcal O_{\mathbb P^1}(D)\).
On \(S=1\), the coordinates are integral and have height
\(O_{K,D}(N^D)\).

Pila's uniform curve estimate
\cite[Theorem~A]{pila1995} bounds their number by
\[
 O_{K,A,j,\varepsilon'}
 \bigl((N^D)^{1/(D2^j)+\varepsilon'}\bigr)
       =O(N^{2^{-j}+D\varepsilon'}).
\]
Choose \(j\) and then \(\varepsilon'\) to put this exponent
below \(\varepsilon\). Both choices precede \(N\).
\end{proof}

\begin{lemma}[A product certificate with fixed sources]
\label{p1:fixed:CRT}
Fix \(s,k\ge1\). Uniformly over sources
\(x_1,\ldots,x_s\in I_N\), the number of centres \(z\in I_N\)
for which there are pairwise coprime \(G_1,\ldots,G_k>1\)
satisfying
\[
 G_i\mid x_{a_i}+u_i,\qquad G_i\mid z+w_i,\qquad
 \prod_iG_i>N,
 \quad u_i,w_i\in\{-1,\ldots,L-1\},
\]
is \(O_{s,k,C,\varepsilon}(N^\varepsilon)\) for every
\(\varepsilon>0\).
\end{lemma}

\begin{proof}
There are \((sB)^kB^k=N^{o(1)}\) choices of positions.
For fixed positions, the product \(Q=\prod_iG_i\) divides
\(\prod_i(x_{a_i}+u_i)\le(3N)^k\). Its divisors and their
ordered factorizations into \(k\) factors have
\(N^{o(1)}\) possibilities. Coprimality fixes one residue
class for \(z\) modulo \(Q>N\), hence at most one centre
in \(I_N\). This counts the certificate without requiring
the groups to be fixed in advance.
\end{proof}

\subsection{Proof of Proposition~\ref{p1:fixed:targets}}
\label{p14:comp:fixed-targets}

\begin{proof}
Take \(H=D_0B\), with fixed \(D_0>1\). A prime exceeding
\(H\) has at most one occurrence of odd valuation in each
window. Apply Lemma~\ref{p1:fixed:recovery} to
\(R(\boldsymbol x,z)\), with \(d=s\), choosing \(t\) so that
the loss of dimension is at most \(\eta L/4\).
The surviving target projection has dimension at least
\(3\eta L/4\). Nonzero labels have multiplicity at most
\(K=t-1\).

For each prime \(q>H\) incident to an active target coordinate
whose label has odd multiplicity in the \(q\)-row, choose a
minimal circuit of distinct labels containing this label.
Such a circuit exists: a minimal zero-sum subset containing
the label cannot contain a proper zero-sum subset, because
it could be removed. Choose incident coordinates for its
other labels. They all belong to sources, because \(q\)
has only one target occurrence. Group primes according to
the selected physical support. These groups are pairwise
coprime. The number of supports through a fixed target
coordinate is at most
\[
 D=\max\left(1,\sum_{j=2}^{s}j!t^{j-1}K^j\right).
\]

Choose the integer \(j\) in Lemma~\ref{p1:fixed:curve}
for degree at most \(K\) and coefficient height
\(N^{A_0+2K+2}\), where
\(\prod_{q\le H}q\le N^{A_0}\).
Choose \(k\) large enough that
\(\theta=KD/k\le1\) and \(j\theta<\varepsilon/4\).
If the product of the \(k\) largest nontrivial groups
exceeds \(N\), Lemma~\ref{p1:fixed:CRT} applies. It remains
to count the complementary case.

Discard from selection the label classes meeting the target
positions of these \(k\) groups. At most \(k\) classes are
discarded; no further restriction of the relation space is
made. Every group incident to a remaining target class
has size at most \(N^{1/k}\).
For such a class \(A\), write
\[
 M_A=\prod_{v\in A\ {\rm at\ sources}} n_v,\qquad
 P_A(Z)=\prod_{v\in A\ {\rm at\ target}}(Z+w_v).
\]
Its squarefree product \(d_A=\sf(M_AP_A(z))\) can be
factored as
\[
 d_A=d_{0,A}d_{1,A}e_A,\qquad
 d_{0,A}\mid\prod_{q\le H}q,\quad
 d_{1,A}\mid M_A,\quad 1\le e_A\le N^\theta .
\]
Indeed, every remaining prime with odd multiplicity and
not dividing \(M_A\) lies in one of at most \(KD\) incident
groups. The choices of \(d_{0,A}\) and \(d_{1,A}\) are
\(N^{o(1)}\), uniformly for fixed source positions.

There are at least \(3\eta L/4-k\) available classes
meeting the target. Select \(j\) of them. Their target
polynomials have simple disjoint roots. After choosing
their positions and coefficients, the equations
\[
 w_\nu^2=d_{A_\nu}M_{A_\nu}P_{A_\nu}(z)
\]
are counted by Lemma~\ref{p1:fixed:curve}.
The positions cost \(B^{O(1)}\), and the coefficients cost
\(N^{j\theta+o(1)}\). The stated choices, with the error
margin split between these factors, give
\eqref{p1:fixed:targets-bound}.
Only the selected positions and coefficients have been
enumerated, not the possible relation spaces or their
regularizations.
\end{proof}
\subsection{Primitive families at a fixed order}

A family is \emph{primitive} if it has a nonzero relative
relation and every proper subfamily has none. Its relation
space projects injectively to each block. For original
windows its nullity is consequently at most \(L\).
For the homogeneous prescription define
\[
 W(I)=p^{|I|}2^{\rho(I)}.
\]
This is an actual joint probability, regardless of the
compatibility of the run-start prescriptions. Those have
joint probability either \(W(I)\) or zero.

\begin{proposition}[Fixed primitive mass]
\label{p1:fixed:primitive}
For every fixed \(r\ge3\) and \(\varepsilon>0\),
\begin{equation}
 \sum_{\substack{(x_1,\ldots,x_r)\ {\rm primitive}\\
                  {\rm distinct,\ separated}}}
          p^r2^{\rho(x_1,\ldots,x_r)}
        \ll_{r,C,\varepsilon}N^{-4/5+\varepsilon}.
 \label{p1:fixed:primitive-bound}
\end{equation}
In particular the same estimate holds for the affine
joint probabilities of the original events.
\end{proposition}

\begin{proof}
On the regular set used in this section, $p_x=p$, so
Proposition~\ref{p13:transfer:primitive} bounds the unordered homogeneous
mass by $O_{r,C,\varepsilon}(N^{-r/2+\varepsilon})$.
Passing to ordered tuples costs at most $r!$, and separation only
restricts the sum. Since $r\ge3$, this implies
\eqref{p1:fixed:primitive-bound}.
This uses the fixed-arity geometric proof of that proposition,
before either reconstruction of the law or the moment theorem.
\end{proof}

The proof is homogeneous throughout. It therefore bounds
incompatible configurations as well as compatible ones.
This feature will be needed for centered Fourier coefficients.

\begin{proposition}[Fixed primitive mass after pair compression]
\label{p1:fixed:compressed}
Fix a channel \(a:b\in\{2:3,3:4\}\), and put
\(\kappa_b=2-1/b\). For every fixed \(r\ge3\), the total
homogeneous mass of primitive families consisting of
one exactly compressed pair in this channel and
\(r-1\) further original windows is
\begin{equation}
 O_{r,C,\varepsilon}
       (N^{-\kappa_b+1/5+\varepsilon}).
 \label{p1:fixed:compressed-bound}
\end{equation}
For a compressed pair and just one further original
window the corresponding masses are, respectively,
\(O_\varepsilon(N^{-4/3+\varepsilon})\) and
\(O_\varepsilon(N^{-5/4+\varepsilon})\).
The affine positive masses satisfy the same bounds.
\end{proposition}

\begin{proof}
Put $t=r-1\ge2$ and fix a compression motif. The compressed block
is the even-selection image of a genuine set of $O(B)$ values of
width $O(B)$; write $r_e$ for its rank.
Use its normalized centre $bx$; it lies above the original target
range and determines the source pair once the motif is fixed.
There are $B^{O(1)}$ motifs and $NB^{O(1)}$ source pairs.
For its homogeneous mass, Lemma~\ref{p1:arith:alignedrank} gives
$q_0(e)=2^{-r_e}\ll N^{-\kappa_b+\varepsilon'}$, including
affine-incompatible sources. Primitivity gives $\rho\le L$ and
the exact homogeneous weight $q_0(e)p^t2^\rho$.

Every primitive family has a nonzero effective projection in each
of its $t+1$ blocks. Theorem~\ref{p13:half:conditional}, with no old
sources, counts at most $N^{(t+1)/2+\varepsilon'}$ hosts,
including the compressed centre. For $t\ge3$, their total mass is
\[
 N^{(t+1)/2+\varepsilon'}N^{-\kappa_b-t+1+\varepsilon'}
   =N^{-\kappa_b+(3-t)/2+2\varepsilon'}
   \le N^{-\kappa_b+2\varepsilon'}.
\]
For $t=2$, the same count on $\rho\le L/4$ gives
$N^{-\kappa_b-1/4+2\varepsilon'}$.
On $\rho\ge L/4$, fix the compressed source image $W$.
Every proper subfamily is direct, so each target is disjoint from
$W$ and the two targets are direct. Lemma~\ref{p13:transfer:dense-targets}
therefore gives conditional homogeneous target mass
$O(N^{-1+\varepsilon'})$; the original targets are regular, so their
homogeneous marginals are exactly $p$.
The source sum is $\sum_e q_0(e)\ll N^{1-\kappa_b+2\varepsilon'}$,
giving total $O(N^{-\kappa_b+3\varepsilon'})$.
Reduce the margins and absorb the motif choices. This proves a
stronger bound than \eqref{p1:fixed:compressed-bound}.

The two-block case is the positive extension estimate
in Proposition~\ref{p1:arith:fixed-homogeneous-seeds}.
For affine compatible extensions it also follows
directly from the positive Palm bounds:
if the additional rank loss is an integer \(d\ge1\),
\[
 p2^d\le2(p2^d-p)_+.
\]
The ordered pair mass times its mean positive Palm
excess is \(O(N^{-4/3+\varepsilon})\) in \(2:3\)
and \(O(N^{-5/4+\varepsilon})\) in \(3:4\).
\end{proof}

\subsection{Proof of Theorem~\ref{p1:fixed:moments}}
\label{p14:comp:fixed-moments}

\begin{proof}
Expand \(\E(Z_N)_r\) over ordered distinct centres.
We describe all configurations removed before evaluating
the two principal classes.

Local defects and contacts have mass \(N^{-1+\varepsilon}\)
after propagation. A dependent separated pair outside
the two retained channels has mass
\(N^{-4/5+\varepsilon}\), also after propagation.
Tuples containing a primitive original subfamily of
size at least three have that same bound, by
Proposition~\ref{p1:fixed:primitive}.
The positive mass of two good pairs is
\(O(N^{-5/4+\varepsilon}+a_N^2)\), including pairs sharing
a vertex, by
Proposition~\ref{p1:arith:fixed-homogeneous-seeds};
propagation again preserves the required precision.

There remains the case of one good pair with an
additional relation. Compress this pair exactly.
A minimal dependent subfamily of the compressed
system has at least two blocks and contains the
compressed block, unless an original primitive family
has already been removed. Apply
Proposition~\ref{p1:fixed:compressed} to this seed and
then propagate. Each of its exponents exceeds \(4/5\).
This exhausts the dependent residual configurations:
after their removal the spaces are either all direct,
or one good pair is direct from all remaining
singletons.

The support estimate
of Proposition~\ref{p1:arith:fixed-support}, together with
the \(O_r(N^{r-1}B)\) contact tuples, shows that the
number of full-rank tuples is
\[
 N^r+O_{r,C,\varepsilon}(N^{r-4/5+\varepsilon}).
\]
Their joint probability is exactly \(p^r\), and their
contribution is \(\lambda_N^r+O(N^{-4/5+\varepsilon})\).

Let \(A_N\) be the ordered positive mass of good pairs.
The pair calculation gives
\[
 A_N=a_N+O_{C,\varepsilon}(N^{-4/5+\varepsilon}).
\]
Fix one such pair. After any fixed number of already
chosen windows, at most \(N^{1/2+o(1)}\) target centres
are dependent on their relative space, by
Proposition~\ref{p1:arith:fixed-branching}. Thus the number of
ordered direct extensions by \(r-2\) windows is
\[
 N^{r-2}\bigl(1+O_{r,C}(N^{-1/2+o(1)})\bigr),
\]
uniformly in the pair. Their joint weight is
\(q(e)p^{r-2}\). Placing the ordered pair in two of
the \(r\) positions gives contribution
\[
 \binom r2\lambda_N^{r-2}A_N
       +O_{r,C}(A_NN^{-1/2+o(1)}).
\]
Since \(A_N=O_C(LN^{-2/3})\), the last error is
smaller than the stated remainder. The two principal
classes give the formula with \((\lambda_N,a_N)\).
The mean transfer
\(\mu_N-\lambda_N=O(N^{-1+\varepsilon})\) and
\(\alpha_N-a_N=O(N^{-1+\varepsilon})\) give
\eqref{p1:fixed:moment-expansion}.

The Stirling-number identity
\(Z^r=\sum_{j=0}^r\left\{\begin{smallmatrix}r\\j\end{smallmatrix}\right\}(Z)_j\)
proves \eqref{p1:fixed:ordinary} and boundedness.
Finally, to any fixed formal order, the factorial
generating function equals
\[
 e^{\mu_N t}\left(1+\frac{\alpha_Nt^2}{2}\right)
       +O_{\rm coeff}(N^{-4/5+\varepsilon}).
\]
Taking its formal logarithm is legitimate at that
fixed order. Terms quadratic in \(\alpha_N\) are
\(O(N^{-4/3+o(1)})\), which proves the cumulant claim.
\end{proof}

\section{Short split products: kernels and moment enumeration}
\label{p14:comp:short}

This section proves the short-support moment input from
Section~\ref{p13:short:section}.  The degree is allowed to grow as a
sufficiently small multiple of $\log\log N$; all estimates below
retain that dependence.

\begin{lemma}[A quadratic kernel with a prescribed divisor]
\label{p13:short:quadratic}
Let \(1\le h\le2DB\), where \(D\) is fixed, and \(N\le n\le2N\).
Suppose \(N^\beta\le q=\sf(n(n+h))<2N^\beta\) and a fixed squarefree
integer \(e\) divides \(q\). Put \(\delta=\log e/\log N\).
Uniformly in these data, the number of \(n\) is at most
\[
 N^{c_2(\beta,\delta)+o(1)},\qquad
 c_2(\beta,\delta)=
 \min\{1,\beta-\delta,\max(\beta/2-\delta,1-\beta/2)\}
 \le \tfrac12+\tfrac\beta4-\tfrac\delta2.
\]
Every \(o(1)\) in this lemma is uniform and may be replaced
by an arbitrary fixed positive exponent.
The uncapped minimum without the entry \(1\) also gives a valid
counting exponent.  Its exact comparison with the affine expression
requires \(\beta-2\delta\le2\); the displayed capped comparison is
valid without this additional range restriction.
\end{lemma}

\begin{proof}
For a fixed \(q\), the equation \(n(n+h)=qv^2\) becomes
\((2n+h)^2-4qv^2=h^2\). Lemma~\ref{p1:arith:pell}
gives \(N^{o(1)}\) solutions, uniformly for the polynomial
heights in question. When \(q=1\), factor the difference of
squares instead. Summing over the \(O(1+N^\beta/e)\)
possible kernels gives the first exponent.

For the other estimate, \(v\asymp N^{1-\beta/2}\) and
\(ev^2\mid n(n+h)\). If
\(\rho_h(M)=\#\{a\bmod M:a(a+h)\equiv0\pmod M\}\), then
\[
 \rho_h(\ell^a)
 \le2\ell^{\min(v_\ell(h),\lfloor a/2\rfloor)},\qquad
 \rho_h(M)\le h\,2^{\omega(M)}.
\]
For \(a\le2v_\ell(h)\), a root is divisible by
\(\ell^{\lceil a/2\rceil}\); for larger \(a\), the two
separated roots give the displayed bound. The Chinese
remainder theorem proves the second inequality. Here
\(M=ev^2\le N^{2+o(1)}\), so summing residue classes gives
\[
 N^{o(1)}
 \sum_{v\asymp N^{1-\beta/2}}\left(\frac{N}{ev^2}+1\right)
 \le N^{o(1)}
       \{N^{\beta/2-\delta}+N^{1-\beta/2}\}.
\]
The trivial \(O(N)\) count supplies the entry \(1\) in the
minimum.  If \(\beta-\delta\ge1\), the maximum is
\(\beta/2-\delta\), and the capped exponent is at most
\[
 \min(1,\beta/2-\delta)
 \le \frac{1+\beta/2-\delta}{2}
 =\frac12+\frac\beta4-\frac\delta2.
\]
Otherwise the maximum is \(1-\beta/2\), and the minimum of this
number and \(\beta-\delta\) is at most their average, with the
same affine value.  This proves the exact comparison, including
admissible cases with \(\beta>2\).
\end{proof}

\paragraph{Necessity of the exponent cap.}
In Lemma~\ref{p13:short:quadratic}, the counting exponent is capped:
\[
 c_2(\beta,\delta)
   =\min\{1,\beta-\delta,\max(\beta/2-\delta,1-\beta/2)\}.
\]
Its affine upper bound is then valid for every real $\beta,\delta$.
The uncapped expression still gives a valid counting estimate, but
its exact affine comparison requires $\beta-2\delta\le2$ when
$\delta\ge0$. This restriction is necessary on actual positive-divisor
data: squarefree consecutive $N,N+1$ with $h=e=1$ give
$\beta=\log(N(N+1))/\log N>2$ and $\delta=0$.
The extra entry $1$ is supplied by the trivial $O(N)$ count.


\begin{lemma}[Kernels of a split product]
\label{p13:short:split}
Let
\[
 P_A(X)=\prod_{a\in A}(X+a),\qquad
 A\subset[-DB,DB]\cap\Z,\qquad
 2\le d=|A|\le c\log\log N,
\]
where \(d\) is even and \(D,c>0\) are fixed.
For a fixed squarefree \(e\), the number of \(n\asymp N\)
with \(q=\sf(P_A(n))\asymp N^\beta\) and \(e\mid q\) is at most
\begin{equation}
 N^{1/2+\beta/(2d)-\delta/d+c\log2+o(1)},
 \qquad \delta=\frac{\log e}{\log N}.
 \label{p13:short:split-count}
\end{equation}
For a fixed kernel \(q\), the fibre has at most
\(N^{2c\log2+o(1)}\) elements. All estimates are uniform in \(A\).
\end{lemma}

\begin{proof}
Pair the \(d\) linear factors into \(m=d/2\) quadratics
\(Q_j\), put \(q_j=\sf(Q_j(n))\), and write
\[
 D_A=\prod_{\substack{a,b\in A\\a<b}}|a-b|.
\]
Its logarithm is \(O_D(d^2\log B)=o(\log N)\).
Outside \(D_A\), the \(q_j\) are pairwise coprime, and
\begin{equation}
 \frac{\prod_jq_j}{q}\mid D_A^2.
 \label{p13:short:resultant}
\end{equation}
Indeed, a prime occurring in \(s\) of the \(q_j\) divides at
least \(s\) distinct linear factors. Its exponent in the
quotient is \(2\lfloor s/2\rfloor\), whereas its exponent
in \(D_A^2\) is at least \(2\binom{s}{2}\).

Remove the primes of \(D_A\) from \(e\), obtaining
\(e'=e/\gcd(e,\operatorname{rad}D_A)\).
There is a factorization \(e'=\prod_j e_j\), with
\(e_j\mid q_j\), such that the factors are pairwise coprime.
Set \(\beta_j=\log q_j/\log N\) and
\(\delta_j=\log e_j/\log N\). Then
\[
 \sum_j\beta_j=\beta+o(1),\qquad
 \sum_j\delta_j=\delta+o(1).
\]
For at least one \(j\),
\(\beta_j-2\delta_j\le(\beta-2\delta)/m+o(1)\).
Lemma~\ref{p13:short:quadratic} bounds its quadratic count
by \(N^{1/2+\beta/(2d)-\delta/d+o(1)}\).
We enumerate only that \(j\), the divisor \(e_j\mid e'\),
and its dyadic size. This costs at most
\(m\tau(e')O(\log N)\), not the number of all factorizations
of \(e'\). Since \(e'\le(C_DN)^d\), the product of its distinct
prime factors is at least \(\omega(e')!\). Stirling's estimate
therefore gives
\[
 \omega(e')\le(c+o(1))\log N,\qquad
 \tau(e')\le N^{c\log2+o(1)}.
\]
This proves \eqref{p13:short:split-count}.

For a fixed \(q\), choose two distinct offsets \(a,b\).
Both \(\sf(n+a)\) and \(\sf(n+b)\) divide
\(\operatorname{rad}(qD_A)\), so there are at most
\(4^{\omega(qD_A)}\le N^{2c\log2+o(1)}\) choices.
The effective coefficients \(u,v\) satisfy \(u,v\ll_D N\).
For each such choice,
\(us^2-vt^2=a-b\ne0\) has \(N^{o(1)}\) solutions by
Lemma~\ref{p1:arith:pell}, including its factorizable case.
Thus Pell is applied to coefficients of height \(O(N)\),
even though \(q\) itself may be as large as \(N^{O(d)}\).
\end{proof}

\subsection{Proof of Proposition~\ref{p13:short:moments}}
\label{p14:comp:short:moments}

\begin{proof}
The moment counts the tuples for which
\(\prod_iP_A(n_i)\) is a square. Write \(q_i=\sf(P_A(n_i))\).
There is a unique representation
\[
 q_i=\prod_{\substack{S\subseteq\{1,\ldots,v\}\\
                  |S|\ge2\text{ even},\ i\in S}}d_S
\]
with pairwise coprime squarefree \(d_S\). Fix dyadic size
classes for these numbers, and put
\(\alpha_S=\log d_S/\log N\) and
\(\beta_i=\sum_{S\ni i}\alpha_S\).
There are \(N^{o(1)}\) size profiles, since \(v\) is fixed.

Expose \(v-1\) of the integers in a chosen order. Before
exposing \(i\), the product of the \(d_S\) whose sets meet
an earlier index is a prescribed divisor \(e_i\mid q_i\).
If its exponent is \(\delta_i\), Lemma~\ref{p13:short:split}
gives the cost
\[
 N^{1/2+\beta_i/(2d)-\delta_i/d+c\log2+o(1)}.
\]
After choosing \(n_i\), assign the new prime factors of
\(q_i/e_i\) to the at most \(2^{v-1}\) possible sets \(S\).
Their number is at most
\(N^{c(v-1)\log2+o(1)}\); inconsistent assignments are
discarded. Every \(S\) has at least two members, so the last
kernel is now fixed and its fibre costs
\(N^{2c\log2+o(1)}\). The enumeration costs in total at most
\(N^{c[v(v-1)+2]\log2+o(1)}\).

For the remaining exponent, average over orders. A set \(S\)
of size \(a\) has coefficient
\[
 \frac1d\left(1-\frac a2+
               \frac12\ind_{\{\text{last index}\in S\}}\right)
\]
in \(\sum_{i\ne\mathrm{last}}(\beta_i/(2d)-\delta_i/d)\).
Its average is \(d^{-1}(1-a/2+a/(2v))\), which is
\(1/(dv)\) for \(a=2\) and nonpositive for \(a\ge4\).
As \(2\sum_{|S|=2}\alpha_S\le\sum_i\beta_i\le vd+o(1)\),
some order has total exponent at most \(v/2+o(1)\).
The order can depend on the size profile; its \(v!\) choices
do not affect the estimate. This proves
\eqref{p13:short:moment-bound}.
\end{proof}

\section{Conditional half-counts and the certificate iteration}
\label{p14:comp:half}

We prove Theorem~\ref{p13:half:conditional} with its stated uniformity
in arbitrary fixed source blocks.  Blocks and effective spaces have the
meaning specified in Section~\ref{p13:half:section}; in particular,
the nonzero projections are required in the individual new spaces,
not in their quotients by the old space.

\subsection{Small kernels and smooth moduli}

Choose \(H=C_DB\) larger than every block diameter, and set
\[
 K_H(n)=\prod_{\substack{\ell>H\\v_\ell(n)\ {\rm odd}}}\ell.
\]
For \(q>0\), let \(\mathcal X_q\) be the set of centres
containing distinct values \(n_1,n_2\) with
\(K_H(n_1)K_H(n_2)\le N^q\).

\begin{lemma}[Two small kernels]
\label{p13:half:small-kernels}
For every fixed \(q>0\) and \(\varepsilon>0\),
\(\#\mathcal X_q\ll_{q,D,C,\varepsilon}N^{q+\varepsilon}\).
\end{lemma}

\begin{proof}
Write \(n_i=c_i k_i a_i^2\), where \(k_i=K_H(n_i)\)
and \(c_i\) is squarefree and supported on primes at most \(H\).
The \(c_i\) have \(2^{2\pi(H)}=N^{o(1)}\) choices, while
the pairs \(k_1k_2\le N^q\) have \(N^{q+o(1)}\) choices.
Fix the two distinct offsets. Their nonzero difference \(h\)
gives
\[
 c_1k_1a_1^2-c_2k_2a_2^2=h.
\]
The effective coefficients are \(O_D(N)\).
Lemma~\ref{p1:arith:pell} gives \(N^{o(1)}\) solutions,
including the factorizable case, and each solution determines
the centre. The \(O_D(B^2)\) offsets are absorbed.
\end{proof}

For \(Y=(\log N)^A\), with fixed \(A\ge2\), write
\[
 n=s_Y(n)a_Y(n)^2K_Y(n).
\]
Here \(s_Y(n)\) is the entire part supported on primes at most
\(Y\); \(a_Y(n)\) contains the halves of the valuations at
larger primes; and \(K_Y(n)\) is their odd squarefree part.

\begin{lemma}[A large smooth modulus]
\label{p13:half:smooth}
Fix \(t\), \(c>0\), and \(A\ge2\).
There is a deterministic set \(\mathcal M_Y\) of centres,
of cardinality \(O(N^{1-1/A+\varepsilon})\), with the
following property. If a selection has more than
\(c\log\log N\) values, its centre lies outside
\(\mathcal X_{1/2}\), and it contains fewer than \(t\)
values with \(K_Y(n)>N^{1/5}\), then its centre belongs
to \(\mathcal M_Y\).
\end{lemma}

\begin{proof}
Outside \(\mathcal X_{1/2}\), at most one value satisfies
\(a_Y(n)^2>N^{3/4+1/100}\): two such values would give
\[
 K_H(n_1)K_H(n_2)
 \le\frac{O_D(N^2)}{a_Y(n_1)^2a_Y(n_2)^2}<N^{1/2}.
\]
Remove this possible value and the at most \(t-1\)
large kernels. For sufficiently large \(N\), select
26 remaining values. Their smooth parts satisfy
\[
 s_i\gg_D N^{1-3/4-1/100-1/5}=N^{1/25}.
\]
Since \(\gcd(s_i,s_j)\le |n_i-n_j|=O_D(B)\),
\[
 \operatorname{lcm}(s_1,\ldots,s_{26})
 \ge\frac{\prod_i s_i}{\prod_{i<j}\gcd(s_i,s_j)}
 \ge N^{26/25-o(1)}>C_1N,
\]
where \(C_1N\) exceeds the length of the centre interval.
Multiply its prime factors, one at a time and with
multiplicity, until the first passage above \(C_1N\).
The resulting smooth modulus satisfies
\[
 C_1N<D_0\le C_1NY,\qquad
 D_0\mid\operatorname{lcm}(s_1,\ldots,s_{26}).
\]
For fixed \(1-1/A<\sigma<1\), Rankin's elementary bound
counts such moduli by
\[
 (C_1NY)^\sigma
       \prod_{\ell\le Y}(1-\ell^{-\sigma})^{-1}
 =N^{\sigma+o(1)}.
\]
Indeed the logarithm of the product is
\(O_\sigma(Y^{1-\sigma})=o(\log N)\).
Assign each maximal prime power of \(D_0\) to one of the
26 offsets on which it occurs. This costs
\(26^{\omega(D_0)}=N^{o(1)}\).
The Chinese remainder theorem determines at most one centre
in its interval, since \(D_0>C_1N\). The offsets themselves
have \(B^{26}=N^{o(1)}\) choices.
Choose \(\sigma\) sufficiently close to \(1-1/A\).
\end{proof}

\subsection{Several anchors force a congruence certificate}

\begin{lemma}[Incoming and outgoing certificate counts]
\label{p13:half:saturation}
Suppose a relation in \eqref{p13:half:condition} selects
\(t\) values in one new block whose kernels exceed \(N^b\),
where \(b=1/5\) and \(tb>r+s-1\).
Some other old or new block receives a congruence certificate
whose modulus exceeds a fixed multiple of \(NB\).
For fixed \(r,s,t,A\), the directed graph of these
source--receiver certificates has both incoming and outgoing
degrees at most \(N^{1/A+o(1)}\).
\end{lemma}

\begin{proof}
Represent the old vector by a selection of actual old values;
parity in the old blocks is not required. Together with the
new selections this is a square product.
Primes above \(Y\) cannot occur in two distinct values
of one block. The \(t\) anchor kernels are therefore
pairwise coprime, and their product \(Q\) satisfies
\(Q>N^{tb}\), \(Q\ll_D N^t\).
Assign each prime of \(Q\) to one of its other odd occurrences.
The \(r+s-1\) receiver moduli have product \(Q\).
For sufficiently large \(N\), at least one exceeds \(C_2NB\).

In that modulus keep the primes only up to its first passage
above \(C_2NB\). The retained list has at most
\[
 k=\left\lfloor\frac{\log(C_2NB)}{\log Y}\right\rfloor+1
\]
members. For a fixed source, the \(t\) offsets and a subset
of the \(O_t(\log N/\log Y)\) available primes cost
\(N^{o(1)}\). Assigning their destinations among \(O_D(B)\)
receiver positions costs at most
\(B^k=N^{1/A+o(1)}\). The congruences determine the receiver.

For a fixed receiver, its values contain
\(R=O_D(B\log N/\log Y)\) primes above \(Y\).
Choose at most \(k\) of them and assign each to an anchor:
the number of codes is at most
\[
 B^t\sum_{j\le k}\binom Rj t^j=N^{1/A+o(1)}.
\]
When \(R\ge2k\), use \((k+1)(etR/k)^k\);
when \(R<2k\), use \((1+t)^R=N^{o(1)}\).
The retained modulus again exceeds the source interval.
Forgetting all other primes enlarges the count and does not
remove the original full-support relation from the remaining
blocks.
\end{proof}

\subsection{The iteration}
\label{p14:comp:half:iteration}

\begin{proof}[Proof of Theorem~\ref{p13:half:conditional}]
We induct strongly on \(r\), simultaneously for every fixed
number \(s\) of sources. Put \(T_0=1\). The case \(r=1,s=0\)
is empty, since its sole projection is required to be nonzero.
For all other cases, choose
\[
 t=5(r+s-1)+1.
\]
First choose a small fixed step \(\delta>0\) within the
final error margin. Let
\[
 x_0=0,\qquad x_j=\min\{\tfrac12,x_{j-1}+\delta\},
 \qquad A_j=x_j^{-1},\qquad Y_j=(\log N)^{A_j}.
\]
There are \(J=\lceil1/(2\delta)\rceil\) stages.
All these numbers are fixed before \(N\); choose the
internal margins after \(J\).

Corollary~\ref{p13:short:old-space} counts lists possessing
a witness with all new selections short.
For every remaining list choose one witness, without
enumerating it. It has at least one long new selection.
At the start of stage \(j\), every surviving long new
source outside \(\mathcal X_{1/2}\) belongs to a set
of size \(N^{1-x_{j-1}+\varepsilon'}\).
This is trivial for \(j=1\).

If such a source selects \(t\) large \(K_{Y_j}\) kernels,
Lemma~\ref{p13:half:saturation} supplies a receiver.
When the receiver is new, count the source--receiver pair
by \(N^{1-x_{j-1}+x_j+O(\varepsilon')}\).
The other \(r-2\) projections are still nonzero in their
own spaces; their sum belongs to the old space enlarged
by these two fixed blocks. The induction hypothesis counts
them by \(N^{(r-2)/2+\varepsilon'}\).
The total is
\[
 N^{r/2+x_j-x_{j-1}+O(\varepsilon')}
 \le N^{r/2+\delta+O(\varepsilon')}.
\]
If the receiver is old, it is already fixed. The incoming
degree counts the new source by \(N^{x_j+O(\varepsilon')}\).
Adjoin that source to the old space and apply induction
to the \(r-1\) remaining targets. The cost is at most
\[
 N^{x_j+(r-1)/2+O(\varepsilon')}
 \le N^{r/2+O(\varepsilon')}.
\]
Only a fixed number of choices of roles is involved.

Remove all these saturated cases. Every surviving long
selection outside \(\mathcal X_{1/2}\) has fewer than \(t\)
large kernels, and Lemma~\ref{p13:half:smooth} places its
centre in a set of size \(N^{1-x_j+\varepsilon'}\).
This proves the stage invariant.
At the final stage every long source lies in
\(\mathcal X_{1/2}\cup\mathcal M_{Y_J}\), whose size is
\(N^{1/2+\varepsilon'}\).
Fix one such centre and use the induction hypothesis for
the other \(r-1\) targets. This costs
\(N^{1/2+\varepsilon'}N^{(r-1)/2+\varepsilon'}\).

Every application of induction uses fewer targets,
although the fixed number of old sources can increase.
Thus the simultaneous induction is legitimate. Choosing
\(\delta\) and then all internal margins within
\(\varepsilon\) proves the theorem.
\end{proof}

\section{Pair partitions and the cost of one dependent target}
\label{p14:comp:activity}

The signed triple identities in Section~\ref{p13:activity:section}
use one quantitative input: the independent target deficit following
a pair.  We first prove the finite partition needed for that estimate.

\begin{lemma}[A finite partition of the pair activity]
\label{p13:activity:pair-partition}
For each fixed \(\eta>0\), dependent pairs can be partitioned
into good pairs from a finite fixed set of rational channels,
a weak class, and an exceptional class, with
\[
 A_{\rm w}\ll_{C,\eta,\varepsilon}N^{-1+\eta+\varepsilon},
 \qquad A_{\rm s}\ll_{C,\eta,\varepsilon}N^{-1+\varepsilon}.
\]
A good pair in a channel \(a:b\), \(b\ge3\), has
\[
 |c(e)|\ll_{a,b,C}N^{-(2-1/b)+o(1)}
\]
and an exact compression into a block of fixed logarithmic
width. Moreover \(A_2\ll_{C,\varepsilon}N^{-2/3+\varepsilon}\).
\end{lemma}

\begin{proof}
Choose \(Q\) with \(1/(Q+1)<\eta/2\), and retain the
channels \(3\le b\le Q\), \(b/2<a<b\), \((a,b)=1\).
Lemma~\ref{p1:arith:alignedrank} bounds the rank in all
larger aligned channels by \(L/(Q+1)+o(L)\).
Put these pairs, nonaligned pairs of rank at most a
sufficiently small fixed fraction of \(\eta L\), and
affinely incompatible pairs in the weak class.
Their coefficients have magnitude at most
\(N^{-2+\eta+o(1)}\).
Corollary~\ref{p13:half:pairs} counts their dependent hosts
by \(N^{1+\varepsilon'}\), proving the weak bound.

Nonaligned pairs of macroscopic rank have total activity
\(O(N^{-1+\varepsilon'})\) by Lemma~\ref{p1:arith:dense}.
The local defects contribute only bounded factors under
\eqref{p13:exact:marginals}.
Pairs meeting the irregular set have the same bound:
its marginal mass is \(N^{-1+o(1)}\), by
Proposition~\ref{p1:arith:core}, and fixed-source
propagation \eqref{p13:exact:propagation} bounds their
joint mass. The independent products have the same cost.
At each retained fixed height, normalized defective pairs
cost \(N^{-3/2+1/b+o(1)}\), as in the proof of
Lemma~\ref{p1:arith:fixed-channel-mass}.
Endpoint height-two pairs and non-strict height-one contacts
cost \(N^{-1+o(1)}\), by the same rank and endpoint estimates.
These are the exceptional class.
The remaining pairs have the private coordinates defining
the good fixed-channel compression. Its rank is at most
\(L/b+O(1)\), giving the stated individual weight.

The finite channel sums are \(O(N^{-2/3+\varepsilon'})\),
by Lemma~\ref{p1:arith:fixed-channel-mass}. Together with
the other two classes, for a sufficiently small fixed
\(\eta\), this also proves the assertion about \(A_2\).
\end{proof}

\subsection{Proof of Proposition~\ref{p13:activity:deficit}}
\label{p14:comp:activity:deficit}

\begin{proof}
After any fixed pair, Theorem~\ref{p13:half:conditional}
with one target counts at most \(N^{1/2+\varepsilon'}\)
centres. Their total independent mass is at most
\(N^{-1/2+\varepsilon'}\).
The weak and exceptional pair classes thus cost
\(N^{-3/2+\eta+O(\varepsilon')}\).

For a good pair in a fixed channel, multiply its two
windows by the channel coefficients and merge identical
values. An even selection in each original window is
multiplied by a square, so its square class is unchanged.
Merging cancels repeated occurrences and preserves total
even parity. Consequently every nonzero class in
\((S_x+S_y)\cap S_z\) is represented by a nonempty even
selection in the normalized block and in the target.
The normalized block has fixed logarithmic width.
Corollary~\ref{p13:half:pairs} counts these two centres
by \(N^{1+\varepsilon'}\).
The offsets and their inverse normalization have
\(B^{O_{a,b}(1)}\) multiplicity. Each host carries weight
at most \(N^{-(2-1/b)-1+o(1)}\), so its contribution is
\[
 N^{1+\varepsilon'}N^{-(2-1/b)-1+o(1)}
 \ll N^{-5/3+2\varepsilon'}.
\]
Sum over the finite channels, then choose \(\eta\) and
\(\varepsilon'\) within the final margin.
\end{proof}

\section{Polynomial confinement and dense arithmetic profiles}
\label{p14:comp:confinement}

We use the block and alignment conventions of
Section~\ref{p13:conf:section}. The group argument treats nonaligned
pairs and primitive families at every fixed arity, then proves the
two triple interfaces stated there.

\begin{lemma}[Polynomial confinement of common divisors]
\label{p13:conf:groups}
Fix \(\tau>0\), \(0<\delta<1/2\), and \(C_0>0\).
Let \(x,y\asymp N\), and let pairwise coprime integers
\(g_a>N^\tau\) satisfy
\[
 g_a\mid x+i_a,\quad g_a\mid y+j_a,\qquad
 |i_a|,|j_a|\le C_0B.
\]
Suppose no line through \((-x,-y)\) contains two
distinct points among the \((i_a,j_a)\).
Then their number, with multiplicity, is
\(O_{\tau,\delta,C_0}(B^{1/2+\delta})=o(B)\).
The same estimate holds after points on exact alignments
have been removed.
\end{lemma}

\begin{proof}
Choose a fixed degree \(d\) such that, for
\(M=(d+1)(d+2)/2\), \((M-1)\tau>d\).
Let \(v_d(s,t)\) list the \(M\) monomials of total degree
at most \(d\), including \(1\).
If \(M\) of the vectors \(v_d(i_a,j_a)\) form a basis,
replace each row of that basis in turn by \(v_d(-x,-y)\).
Subtracting the new row from the other rows shows that
the resulting determinant is divisible by the product
of their \(M-1\) pairwise coprime divisors.
Its absolute value is at most
\[
 C_dN^dB^{d(M-1)}<N^{(M-1)\tau}.
\]
All the determinants vanish.
The vector \(v_d(-x,-y)\) would then belong to the
intersection of the \(M\) coordinate hyperplanes of
this basis, which is \(\{0\}\). Its constant coordinate
is one, a contradiction.

Thus a nonzero polynomial \(F(S,T)\) of degree at most
\(d\) vanishes on all the offset points.
For a geometrically integral component of degree at least
two, Pila's uniform estimate
\cite[Theorem~A]{pila1995} gives
\(O_d(B^{1/2+\delta})\) integer points of height \(O(B)\).
A rationally irreducible component that is not geometrically
integral has only \(O_d(1)\) rational points, by intersecting
two conjugate components.

A line through at least two distinct grid points has a
primitive equation \(aS+bT=c\) with
\(|a|,|b|\ll B\), \(|c|\ll B^2\).
Every group on it divides \(ax+by+c\).
This integer is nonzero by the hypothesis, and has size
\(O(NB)\). Coprimality bounds the number of these groups
by \(O_\tau(1)\). A line through only one grid point
needs no separate point count.
Finally, all groups at one point divide its true value
\(x+i_a=O(N)\); their multiplicity is \(O_\tau(1)\).
There are at most \(d\) components. These observations
prove the bound.
\end{proof}

The groups in this lemma can be composite. Only divisibility
and coprimality are used, and the auxiliary curve is not
enumerated. Its degree is chosen after \(\tau\) and before
\(N\).

\paragraph{An anchored small-kernel witness at fixed arity.}
\label{p18:conf:anchored}
Fix integers $r\ge2$, $t\ge2$ and real numbers $\eta,\tau>0$.
Let $M$ be a binary space of physical square relations on $r$ true
windows with centres in $I_N$: its words select an even set in each
window and their total product has even prime valuations.
Write $\ell_v\in M^*$ for the coordinate forms and call a position
active when $\ell_v\ne0$. Designate one window as the anchor.
Assume that:
\begin{itemize}
\item $\dim M\ge\eta L$ and the physical projection into the anchor
is injective;
\item no nonzero form on $M$ has $t$ representations by pairwise
disjoint coordinate sets of size between $1$ and $r-1$;
\item between the anchor and each other block, no two distinct pairs
of active offsets have the same ratio of true values.
\end{itemize}
For every cutoff $H\ge B=L+1$, the anchor then contains two distinct
values $n_1,n_2$ with $K_H(n_i)\le N^{D\tau}$, where
\[
 D=\sum_{j=1}^{r-1}j!(t-1)^{j-1}2^{j(t-1)}.
\]
One threshold, chosen after $r,t,C,\eta,\tau$, precedes the centres,
$M$ and $H$. No injectivity or nonalignment between the other
blocks is assumed.

\begin{proof}
The recovery-exclusion hypothesis gives multiplicity at most $t-1$
for every nonzero coordinate label, by considering singleton recovery
sets. The final counting argument of Lemma~\ref{p1:fixed:recovery}
also bounds the independent $j$-label representations of a nonzero
label by $j!(t-1)^{j-1}$: a maximal disjoint family of their supports
has fewer than $t$ members; its union, of size at most $j(t-1)$,
meets every representation, and induction removes a label. A zero
remaining sum is excluded by independence. Permitting every subset
of the $j$ label fibres gives at most $2^{j(t-1)}$ physical lifts.
Thus at most $D$ physical supports of size at most $r-1$ can
represent a fixed active anchor label.

A prime $\ell>H$ has at most one odd occurrence in each true window.
If it occurs at an active anchor position $v$, its parity row on
$M$ expresses $\ell_v$ as a sum of at most $r-1$ other labels.
Remove zero labels and choose a minimal subrepresentation.
Its labels are distinct, nonzero and independent: otherwise a
nonempty zero-sum subset could be removed without changing the
represented nonzero form. Choose their actual positions in the
prime row. The resulting support $A$ consists of $v$ and positions
in other blocks and uses only active coordinates.

Assign every prime with an active anchor occurrence to one such
support, and let $g_A$ be the product assigned to $A$. Distinct
nontrivial groups are coprime and divide every true value selected
by their support. Every anchor position meets at most $D$ groups,
and its entire odd kernel above $H$ divides their product.
Assign each group to one other block it meets.
For each of the $r-1$ choices, Lemma~\ref{p13:conf:groups} gives
$o(B)$ groups larger than $N^\tau$, with their full multiplicities.
Its no-line hypothesis is exactly the stated active-offset
condition. Equal projected points are covered by the multiplicity
bound in that lemma.

Impose zero only on the anchor positions touched by these large
groups. Each group has a unique anchor position, so the dimension
cost is at most the number of such groups, namely $o(B)$.
A nonzero space remains for large $N$ and its anchor projection
is still injective. A nonzero even projected word uses at least
two positions. Each of those positions was active on the original
space $M$ and met only groups at most $N^\tau$, so its odd kernel
is at most $N^{D\tau}$. The incidence bound is used on $M$ before
erasure; no recovery condition or prime-row classification on the
erased space is assumed.
\end{proof}

\paragraph{Dense primitive hosts and their exact weights.}
\label{p18:conf:fixed-primitive}
For each fixed $r\ge3$, $\eta>0$ and $\varepsilon>0$, primitive
families of $r$ true windows with
$\rho_{\mathrm{rel}}\ge\eta L$ have $O(N^\varepsilon)$ hosts.
Indeed, every effective relation-space projection is injective. Choose
linear sections of the local even-selection maps to lift that space
without increasing its dimension or adding local square kernels.
Lemma~\ref{p1:fixed:recovery}, with $d=r-1$ and $t$ chosen after
$r,\eta$, loses at most $2rB\,t^{-1/(r-1)}\le\eta L/4$ dimensions.
The surviving physical space has dimension at least $3\eta L/4$,
retains all injective projections, and has the required recovery
exclusion. Two exact
units between any pair would give equal nonzero even square classes
in those two blocks: the selected two-value products are nonsquares
by Lemma~\ref{p1:arith:runge}. This contradicts primitivity, since
the pair is a proper subfamily. Thus the active-offset condition
holds for each choice of anchor. Apply the witness with $H=B$ and
density parameter $\eta/2$ at each of the $r$ windows, with the same
fixed constant $D$.
Every centre lies in $\mathcal X_{2D\tau}$.
Lemma~\ref{p13:half:small-kernels} therefore counts at most
$N^{2rD\tau+r\varepsilon'}$ unordered families. Choose $\tau$ and
$\varepsilon'$ after $r,\eta,\varepsilon$ so that this exponent is
at most $\varepsilon$. No auxiliary support or chosen subspace is
enumerated in this last count.

Projection injectivity also gives
$\rho_{\mathrm{rel}}(I)\le\min_{x\in I}d_x\le L$.
The actual marginals satisfy $p_x\le2p\ll_C N^{-1}$, so each exact
weight $2^{\rho_{\mathrm{rel}}(I)}\prod_{x\in I}p_x$ is
$O_{r,C}(N^{1-r})$. Its dense sum is
$O_{r,C,\eta,\varepsilon}(N^{1-r+\varepsilon})$.
The same bound holds for homogeneous prescriptions; their one-window
masses are $2^{-d_x}\le2p$, including locally deficient windows.
For the complete fixed-size estimate, split at
$\rho_{\mathrm{rel}}=\eta L$. The independent-mass bound in
Corollary~\ref{p13:half:connected} makes the low sector
$O(N^{-r/2+\eta+\varepsilon'})$, while the dense sector has the bound
just proved. Choosing $\eta,\varepsilon'$ within the final margin
gives $O(N^{-r/2+\varepsilon})$. Each of these choices is fixed
before $N$; none gives a growing-$r$ estimate.

\begin{theorem}[Dense primitive triples]
\label{p13:conf:dense-primitive}
For fixed \(\eta,\varepsilon>0\), primitive triples of
true windows with \(\rho_{\mathrm{rel}}(T)\ge\eta B\)
have \(O_{\eta,C,\varepsilon}(N^\varepsilon)\) hosts.
Their homogeneous mass is \(O(N^{-2+\varepsilon})\).
The total absolute centered activity of all primitive
triples is \(O(N^{-3/2+\varepsilon})\).
\end{theorem}

\begin{proof}
Use the preceding fixed-arity argument with $r=3$.
Since $\eta B\ge\eta L$, it gives the dense host count.
For each such triple the homogeneous probability is
$2^{\rho_{\rm rel}}\prod_x2^{-d_x}\le(2p)^3\,2^L\ll_C N^{-2}$.
This gives the dense homogeneous mass.
The complete fixed-size weight bound is $O(N^{-3/2+\varepsilon})$.
For a primitive family, $c(I)=q(I)-\prod_Ip_x$ and
$|c(I)|\le2^{\rho_{\rm rel}(I)}\prod_Ip_x$, even when $q(I)=0$.
Summing proves the final assertion.
\end{proof}

Here homogeneous mass is the sum of the probabilities of the
homogeneous events, with zero right-hand side in place of the affine
start prescription.
It must be distinguished from a weight using the original marginals
$p_x$, which can vanish. For instance, a start at $x=3$ of length
$L=7$ is impossible: the square $4$ forces the common bit on
$3,\ldots,9$ to be zero, and $6=2\cdot3$ then forces the bit at $2$
to be zero, contradicting the left-boundary equation. The corresponding
homogeneous event is compatible and has positive probability.

\begin{corollary}[Dense nonaligned pairs]
\label{p13:conf:dense-pairs}
For fixed \(\eta,\varepsilon>0\), nonaligned pairs
with \(d_{xy}\ge\eta B\) have \(O(N^\varepsilon)\)
hosts.
\end{corollary}

\begin{proof}
Realize their effective relation space with injective
projections. Regularize with
Lemma~\ref{p1:fixed:recovery}, now with \(d=1\),
losing at most \(\eta B/4\).
Every active prime row has two equal labels.
A coordinate belongs to at most \(K=t-1\) active supports.
Confinement bounds the groups above \(N^\tau\) by \(o(B)\).
Kill both endpoints of every such group.
The remaining positive-dimensional space has two active
positions in each block, with kernels at most \(N^{K\tau}\).
The centres lie in \(\mathcal X_{2K\tau}\).
Their count is at most \(N^{4K\tau+2\varepsilon'}\);
choose \(\tau,\varepsilon'\) sufficiently small.
\end{proof}


\paragraph{Completion of the quasi-linear pair profile.}
\label{p18:conf:pair-profile}
We now prove the stronger bounds in Proposition~\ref{p1:arith:hosts},
retaining its full arithmetic nullity
$\rho(x,y)=2L-\dim(S_x+S_y)$.
If both windows are regular, positive $\rho$ is exactly a nonzero
effective intersection; Corollary~\ref{p13:half:pairs} counts these
pairs by $N^{1+\varepsilon'}$. A locally deficient coordinate gives
at most $2N\,|I_N\setminus G_N|=N^{1+o(1)}$ further ordered pairs.
This proves the support bound, even before imposing separation.

For the weighted nonaligned sum, split at
$d_{xy}=\eta L$ with fixed $\eta>0$. Local rigidity gives, uniformly,
\[
 \rho(x,y)=d_{xy}+(L-d_x)+(L-d_y)\le d_{xy}+2,
 \qquad \rho(x,y)\le L+2.
\]
In the low sector, the just proved support bound and
$2^{\eta L+2}\ll_C N^\eta$ give
$N^{1+\eta+\varepsilon'}$. In the high sector,
Corollary~\ref{p13:conf:dense-pairs}, with a smaller fixed density
parameter if necessary to replace $L$ by $B$, gives
$N^{\varepsilon'}$ hosts, each of weight $O_C(N)$.
Choose $\eta,\varepsilon'$ after the final $\varepsilon$.
This proves both assertions with $N^{1+\varepsilon}$ and with a
common threshold before all pairs. Only conditional half-counts,
local rigidity and confinement were used; the pair-profile theorem
and its elementary $6/5$ proof are not inputs to this completion.

\subsection{Proof of Proposition~\ref{p13:conf:near-primitive}}
\label{p14:comp:confinement:near}

\begin{proof}
The kernel of each projection of the relation space
has dimension at most \(\theta B\).
Intersect three complements to these kernels.
The resulting subspace loses at most \(3\theta B\)
dimensions and has all projections injective.

In a given pair, two exact units impose a reduced rational
alignment of height \(O(B)\). Two distinct such alignments
would bound the centres by \(O(B^3)\), which is impossible
at the present scale. The alignment is therefore unique.
If it has \(u\) units, their even-selection space has
dimension \(u-1\); its image loses at most one more
dimension by Lemma~\ref{p1:limits:local-rigidity}.
Thus \(u\le d_{xy}+2\le\theta B+2\).
Kill all the positions of these units in the three pairs,
at a cost of at most \(6\theta B+12\) dimensions.
The remaining space still has macroscopic dimension
and injective projections.

Regularize with a further loss at most \(\eta B/8\).
No pair of surviving offsets lies on an exact alignment,
so the confinement argument in
Theorem~\ref{p13:conf:dense-primitive} applies.
After all positions incident to large groups have been
killed, the centres again belong to
\(\mathcal X_{2C_t\tau}\).
The same count gives \(N^\varepsilon\) hosts.
We never enumerate the chosen subspaces or deletion sets.

Finally,
\(\rho_{\mathrm{rel}}(T)\le d_z+d_{xy}\le L+\theta B\).
The homogeneous weight is at most
\(N^{-2+\theta+o(1)}\), proving the mass estimate.
\end{proof}

\subsection{Proof of Proposition~\ref{p13:conf:strong-profile}}
\label{p14:comp:confinement:strong}

\begin{proof}
Choose the strong pair \(e=\{x,y\}\). Put
\(\delta_z=\dim(S_z\cap(S_x+S_y))
 =\rho_{\mathrm{rel}}(T)-d_{xy}\).
Connectedness gives \(\delta_z>0\), and compatibility gives
\(q(T)=q(e)p_z2^{\delta_z}\).

For \(0<\delta_z<\eta'L\), the conditional count
in Theorem~\ref{p13:half:conditional} gives
\(N^{1/2+\varepsilon'}\) targets of weight at most
\(N^{-1+\eta'+o(1)}\).
Lemma~\ref{p1:arith:dense}, together with
Corollary~\ref{p13:half:pairs} for the independent
part, bounds the total joint mass of strong nonaligned
pairs by \(N^{-1+\varepsilon'}\).
This branch costs \(N^{-3/2+\eta'+O(\varepsilon')}\).

For \(\delta_z\ge\eta'L\), the effective-space form of
Proposition~\ref{p1:fixed:targets}, explained after
\eqref{p13:exact:kappa3}, gives \(N^{\varepsilon'}\)
targets after the fixed pair.
Corollary~\ref{p13:conf:dense-pairs} also gives
\(N^{\varepsilon'}\) source pairs.
Each triple has \(q(T)\ll N^{-3+\sigma+o(1)}\).
Choose \(\eta'\) and the internal margins after
\(\varepsilon\).
\end{proof}

\section{The short-basis and common-curve constructions}
\label{p14:comp:rigidity}

The notation is that of Section~\ref{p13:rig:section}.  This section
supplies the complete local estimates, basis extraction, and algebraic
construction used by the main rank argument.  The final
degree-uniform point count and the nonaligned contradiction are proved
in that section of the article.

\subsection{Local square relations and defective positions}

For primes larger than the window length, a local square relation
can use only positions having no odd valuation at any such prime.
The next elementary argument bounds those positions without an
admissibility assumption.

\begin{lemma}[A local square threshold]
\label{p13:rig:local-square}
Let $B\ge2$, let $I\subset[-B,B]\cap\Z$ have positive even
cardinality $m$, and let $x$ be a positive integer.  If
$x>(B+1)^{12m}$, then $\prod_{a\in I}(x+a)$ is not an integer
square.
\end{lemma}

\begin{proof}
The assumed threshold implies $x\ge2B$.
If the product were a square, Lemma~\ref{p1:arith:runge},
with $R=B$ and $d=m$, would give $x\le(64B)^m$.
For $B\ge2$, we have $64B\le(B+1)^5$.
Thus $x\le(B+1)^{5m}\le(B+1)^{12m}$,
contrary to the hypothesis.
\end{proof}

\begin{lemma}[Few defective positions in every critical window]
\label{p13:rig:defective}
The number of $n\in V_x$ for which
\[
             \prod_{\substack{p>10B\\v_p(n)\ {\rm odd}}}p=1
\]
is $O_C(B/\log B)$, uniformly over all critical windows.
\end{lemma}

\begin{proof}
Apply the local assertion of Lemma~\ref{p1:arith:defects}
with $H=10B$ to $V_x$.
In the critical band its diameter is $O(B)$ and its values
lie in $[N/2,3N]$ for all sufficiently large $N$.
The comparison constants are uniform in $x$.
\end{proof}

We use the elementary bound $\pi(t)\ll t/\log t$.  For
completeness, the product of the primes in $(u,2u]$ divides
$\binom{2u}{u}\le4^u$ for integer $u$; dyadic summation gives
$\vartheta(t)=\sum_{p\le t}\log p\ll t$.  Splitting at
$\sqrt t$ gives the stated bound for $\pi(t)$.

We will also need the elementary prime harmonic estimate
\begin{equation}
 \sum_{B^a<p\le10B}\frac1p
       =\log(1/a)+O_a(1/\log B),\qquad 0<a<1\text{ fixed}.
 \label{p13:rig:mertens}
\end{equation}
Here is a derivation recording its error term.  Chebyshev's bound
above implies $\psi(t)=\sum_{p^j\le t}\log p\ll t$.
The identity
$\log(n!)=\sum_{d\le n}\Lambda(d)\lfloor n/d\rfloor$
then gives $\sum_{d\le t}\Lambda(d)/d=\log t+O(1)$.
The contribution of prime powers with exponent at least two is
bounded, so $A(t)=\sum_{p\le t}(\log p)/p=\log t+O(1)$.
Partial summation yields
\[
 \sum_{p\le t}\frac1p
 =\frac{A(t)}{\log t}+\int_2^t\frac{A(u)}{u(\log u)^2}\,du
 =\log\log t+c_0+O(1/\log t).
\]
The difference at $10B$ and $B^a$ proves
\eqref{p13:rig:mertens}.

\subsection{Large primes turn intersections into components}

Construct a graph on the disjoint formal copies of $V_x,V_y$.
A prime $p>10B$ has at most one occurrence of odd valuation
in either block.  If it has two occurrences, join their vertices;
if it has only one, mark that vertex as forced to zero.
The corresponding binary equation says that the endpoints of
an edge have equal coordinates, or that a marked coordinate is
zero.  A component is free when it contains no marked vertex.

\begin{lemma}[Free components]
\label{p13:rig:components}
Let $c$ be the number of free components containing an edge.
Then
\[
                         c=d_{xy}+O_C(B/\log B).
\]
Each such component meets both blocks and is selected in its
entirety or omitted in any solution to the large-prime equations.
\end{lemma}

\begin{proof}
The isolated unmarked vertices are precisely the defective
positions.  Their total number $D_0$ is $O_C(B/\log B)$ by
Lemma~\ref{p13:rig:defective}.  The kernel of the large-prime
equations has dimension $c+D_0$.  Imposing the remaining
prime equations and the two block parities costs at most
$\pi(10B)+2$ dimensions.  If $\mathcal R$ is the space of
complete relations and $k_x=\dim\ker\phi_x$, then the map
taking a relation to its common class gives an exact sequence
with kernel $\ker\phi_x\oplus\ker\phi_y$.  Hence
\[
 c+D_0-\pi(10B)-2\le\dim\mathcal R=k_x+k_y+d_{xy}
                                      \le c+D_0.
\]
A local relation can use only defective positions, so
$k_x+k_y\le D_0$.  These inequalities prove the assertion.
\end{proof}


\subsection{Proof of Theorem~\ref{p13:rig:short-basis}}
\label{p14:comp:rigidity:short-basis}

\begin{proof}
\emph{A graph estimate.}
A multigraph with $n$ vertices and $e/n\to\infty$ has a cycle
of length $O(1+\log(n+1)/\log(e/n))$.  Loops and parallel
edges give cycles of length one and two.  Otherwise delete
vertices of degree less than $t=\lfloor e/(2n)\rfloor$.
Not all vertices can disappear, since fewer than $nt\le e/2$
edges would then be deleted.  The remaining graph has minimum
degree at least $t$.  If its girth exceeds $2r+1$, a breadth
first search contains at least
$1+t\sum_{j<r}(t-1)^j$ vertices.  This bounds $r$ and proves
the estimate.

Consequently, when $n\ll B/\log B$ and $e\le B$, the binary
cycle space contains at least $e-O(B/\sqrt{\log B})$
independent cycles of length $O(\log B/\log\log B)$.
Indeed, while at least $B/\sqrt{\log B}$ edges remain,
record a short cycle and delete just one of its edges.
The recorded cycles are independent: the first deleted edge
belongs to no later cycle, and the same argument applies
successively to all deleted edges.

\emph{Medium-prime columns.}
Start with the free nontrivial components of
Lemma~\ref{p13:rig:components}, and discard those with more
than $M=\lfloor\log B\rfloor$ vertices.  This loses at most
$2B/(M+1)$ components.  The remaining number is
$d_{xy}+O_C(B/\log B)$.  Put $z=B^a$.  To each retained
component attach its parity column on the primes $z<p\le10B$,
and let $w_j$ be that column's Hamming weight.  Each nonzero
entry consumes a divisible position, whence
\[
 \sum_j w_j
 \le\sum_{z<p\le10B}\#\{n\in V_x\sqcup V_y:p\mid n\}
 \le2B\sum_{z<p\le10B}\frac1p+2\pi(10B)
 \le2B\log(1/a)+O_a(B/\log B).
\]
Discard only columns of weight at least three.  Their number
is at most one third of this sum, so at least
\[
                 e\ge d_{xy}-\frac23B\log(1/a)
                                   -O_{a,C}(B/\log B)
\]
columns remain.  There is no factor $M$ in this count: it
bounds the number of discarded columns, not their physical size.

Use the primes $z<p\le10B$ and one extra vertex $\infty$
as vertices of a multigraph.  A column of weight two joins its
two primes; a column of weight one joins its prime to $\infty$;
a zero column is a loop at $\infty$.  Its binary incidence
matrix has the same kernel as the parity matrix, because the
row at $\infty$ is the sum of the other rows.  The graph has
$n\le\pi(10B)+1\ll B/\log B$ vertices and $e\le B$ edges.
The preceding graph argument therefore supplies a space
$\mathcal U$ of dimension at least
$e-O(B/\sqrt{\log B})$ with a basis of short cycles.
Each basis vector uses
$O(\log^2 B/\log\log B)$ physical positions.

\emph{Small primes and parity.}
All large and medium primes already cancel on $\mathcal U$.
The primes at most $z$ and the two block parities impose at
most $s=\pi(z)+2\ll_a B^a/\log B+2$ further equations on its
cycle coordinates.  Their kernel is $\mathcal K_a$; since
$s=o(B/\sqrt{\log B})$, it has the dimension in
\eqref{p13:rig:short-dimension}.  A fundamental kernel basis
for a matrix with at most $s$ rows uses at most $s+1$ columns
per vector: express each nonpivot column in terms of pivot
columns.  Thus each resulting vector uses at most $s+1$
short cycles, giving \eqref{p13:rig:short-support}.

Finally, the original free nontrivial components are disjoint
and meet both blocks.  Every nonzero selection of them contains
a component carrying a large prime which occurs once in each
block and in no other component.  Its common square class is
therefore nonzero.  This also proves injectivity of both physical
projections and of the common-class map.  The imposed parity
equations complete the proof.
\end{proof}

\subsection{Two relations produce a polynomial of controlled height}

The short basis will have support $B^{a+o(1)}$.  We therefore
need an auxiliary polynomial whose degree and logarithmic
height grow linearly with the support, and a simultaneous
construction whose threshold is quadratic in that support.
The following elementary factor bound will keep the height
under control when common polynomial factors are removed.

\begin{lemma}[Height of a bivariate factor]
\label{p13:rig:factor-height}
Let $P\in\Z[X,Y]\setminus\{0\}$, with partial degrees
$d_X,d_Y$.  Every primitive integer factor $Q$ of $P$, and
the corresponding integer quotient, satisfy
\[
 H(Q)\le2^{d_X+d_Y}\|P\|_2
       \le2^{d_X+d_Y}\sqrt{(d_X+1)(d_Y+1)}\,H(P).
\]
\end{lemma}

\begin{proof}
The multivariate Mahler measure $M(P)$ is multiplicative,
at most $\|P\|_2$, and at least one for a nonzero integer
polynomial.  Jensen's formula in one variable, followed by
the leading coefficient in that variable, proves the last
assertion by iteration.  The one-variable coefficient bound
also iterates: integrating first in $X$ and then in $Y$ gives
\[
 |[X^iY^j]P|\le\binom{d_X}{i}\binom{d_Y}{j}M(P).
\]
Thus $H(Q)\le2^{\deg_XQ+\deg_YQ}M(Q)$; multiplicativity
and $M(P/Q)\ge1$ prove the result.  The quotient has the same
bound, with the integer content handled separately.
\end{proof}

\begin{proposition}[Two complete witnesses]
\label{p13:rig:two-witnesses}
There is an absolute constant $c>0$ with the following property.
Let $u,v$ be independent formal words, even in each block,
with offsets in $[-B,B]$.  For $w=u,v,u+v$, write
$|I_w|=2r_w$, $|J_w|=2s_w$, and assume
$\max(r_w,s_w)\le q$, where $1\le q\le B$.
Suppose that $F_u(x,y)$ and $F_v(x,y)$ are integer squares,
with $x,y\in[N,2N)$.  If $N>(B+1)^{cq}$, there is a
nonzero $R\in\Z[X,Y]$ such that
\begin{equation}
 R(x,y)=0,\qquad
 \deg_XR,\deg_YR\le12q,\qquad \deg R\le24q,
 \qquad H(R)\le(B+1)^{cq}.
 \label{p13:rig:two-polynomial}
\end{equation}
Moreover, $R$ is the norm of a biquadratic auxiliary function
whose four polynomial coefficients have no common nonconstant
factor.
\end{proposition}

\begin{proof}
Set $F_1=F_u$, $F_2=F_v$, $G=\gcd(F_1,F_2)$, with monic
gcd, and $F_3=F_{u+v}=F_1F_2/G^2$.
Formal independence makes
\[
 K=\mathbb Q(X,Y)(W_1,W_2),\qquad W_i^2=F_i,
\]
a degree-four extension.  With $W_0=1$ and
$W_3=W_1W_2/G$, the four $W_\chi$ form a basis.  At the
given integer point use positive square roots.  Then all
$W_\chi(x,y)$ are integers: $F_3(x,y)$ is an integer which
is a rational square, hence an integer square, and
$W_3(x,y)=\sqrt{F_1(x,y)}\sqrt{F_2(x,y)}/G(x,y)$.

\emph{The linear system at infinity.}
Put $r_0=s_0=0$, $k=3q$, and consider
\[
 \mathcal A=\sum_{\chi=0}^3A_\chi(X,Y)W_\chi,
 \qquad \deg_XA_\chi\le k-r_\chi,
 \quad \deg_YA_\chi\le k-s_\chi.
\]
There are
\[
 n=\sum_{\chi=0}^3(k-r_\chi+1)(k-s_\chi+1)
   \ge21q^2+18q+4
\]
unknown coefficients, and $n\le4(3q+1)^2$.
Writing $\xi=1/X$, $\eta=1/Y$, the positive formal branches
are
\[
 W_\chi=X^{r_\chi}Y^{s_\chi}
 \prod_{u\in I_\chi}(1+u\xi)^{1/2}
 \prod_{v\in J_\chi}(1+v\eta)^{1/2}.
\]
Thus $f(\xi,\eta)=\xi^k\eta^k\mathcal A$ is a formal
power series.  Require its coefficients of total degree at
most $2k$ to vanish.  This gives at most
\[
 t=\frac{(2k+1)(2k+2)}2=18q^2+9q+1,
 \qquad n-t\ge3q^2+9q+3,\qquad \frac t{n-t}<6
\]
homogeneous linear equations.

All coefficients required from the binomial series have
denominator dividing $4^{2k}$.  Cauchy's estimate on circles
of radius $(2B)^{-1}$ bounds their numerators, after this
common denominator is cleared, by $(B+1)^{c_0q}$ for an
absolute $c_0$.  An elementary pigeonhole form of Siegel's
lemma yields a nonzero integer solution of height
$(B+1)^{c_1q}$.  Indeed, if the integer matrix has height $A$,
the images of $\{0,\ldots,M\}^n$ take at most
$(2nAM+1)^t$ values.  Choosing
$M=\lceil(3nA)^7\rceil$ makes two images equal, because
$t/(n-t)<6$.  Their difference is a nonzero kernel vector;
$\log n=O(\log B)$ is absorbed in the stated height.

\emph{Primitive coefficients.}
Let $D$ be the primitive polynomial gcd of the four
$A_\chi$, omitting zero coefficients, and write
$A_\chi=DA'_\chi$.  The factor bound
of Lemma~\ref{p13:rig:factor-height} keeps
$H(A'_\chi)\le(B+1)^{O(q)}$.  Write
$a=\deg_XD$, $b=\deg_YD$, $e=\deg D$.  Then
\[
 f=\widetilde D f',\qquad
 \widetilde D=\xi^a\eta^bD(\xi^{-1},\eta^{-1}),\qquad
 f'=\xi^{k-a}\eta^{k-b}\sum_\chi A'_\chi W_\chi.
\]
These are power series, and
$\operatorname{ord}\widetilde D=a+b-e$ for total order.
Since order is additive in the power-series domain,
\begin{equation}
 \operatorname{ord}f'\ge2k+1-a-b+e
                          =(k-a)+(k-b)+e+1.
 \label{p13:rig:primitive-order}
\end{equation}
Thus removing $D$ preserves vanishing beyond the remaining
pole weight.  This is a formal division, never a division
by the possibly zero value $D(x,y)$.  Replace the auxiliary
by its primitive version from now on.

\emph{Integer separation and the norm.}
On the bidisc $|\xi|,|\eta|\le(4B)^{-1}$, the analytic
numerator is bounded by $(B+1)^{O(q)}$.  Cauchy's estimate
and the order of vanishing give, in the comparable range
$x,y\in[N,2N)$,
\begin{equation}
                  |\mathcal A(x,y)|\le
                                   \frac{(B+1)^{c_2q}}N.
 \label{p13:rig:integer-gap}
\end{equation}
To see the power of $N$ explicitly before removing $D$, the
tail of the numerator is bounded by its coefficient bound
times
\[
 \sum_{\ell\ge2k+1}(\ell+1)(4B/N)^\ell
       \ll(k+1)(4B/N)^{2k+1}.
\]
Multiplication by $x^ky^k\le(2N)^{2k}$ gives
\eqref{p13:rig:integer-gap}.  After removing $D$, both pole
weights decrease, and \eqref{p13:rig:primitive-order} gives
the same estimate with the new height bound.

The specialization $\mathcal A(x,y)$ is an integer, so it
vanishes for the stated threshold.  The auxiliary is nonzero
in $K$ by independence of the four characters.  Its norm is
therefore nonzero.  An explicit polynomial expression is
\[
 \begin{split}
 U&=A_0^2+A_1^2F_1-A_2^2F_2-A_3^2F_3,\\
 V&=A_0A_1-A_2A_3(F_2/G),\\
 R&=U^2-4F_1V^2.
 \end{split}
\]
Here $F_2/G$ is a polynomial.  The partial degrees of $U$
are at most $2k$, and those of $V$ are at most
$2k-r_1$, $2k-s_1$.  For the second term in $V$, use
$r_3=r_1+r_2-\deg_XG$ and its $Y$ analogue.
It follows that the partial degrees of $R$ are at most
$4k=12q$.  The coefficient $\ell^1$ norms of all factors
in this formula are $(B+1)^{O(q)}$, which proves the height
bound and completes \eqref{p13:rig:two-polynomial}.
\end{proof}


\subsection{Proof of Theorem~\ref{p13:rig:common-curve}}
\label{p14:comp:rigidity:common-curve}

\begin{proof}
For $i<j$, the half-cardinalities of $w_i,w_j,w_i+w_j$
are at most $q=\min(2q_0,B)$, since an even subset of
$[-B,B]\cap\Z$ has size at most $2B$.
Proposition~\ref{p13:rig:two-witnesses} gives a nonzero
polynomial $R_{ij}$ of degree at most $48q_0$ and height
$(B+1)^{O(q_0)}$, vanishing at $(x,y)$ and arising from
an auxiliary with primitive coefficients.
Choose an irreducible primitive factor $C_0$ of $R_{12}$
which vanishes at this point.  The factor-height lemma
preserves its height bound.  If $C_0$ depended on only one
variable, its integer root of size at least $N$ would
contradict Cauchy's root bound.  Thus it depends on both.

If $C_0$ failed to divide some $R_{ij}$, the resultant
$\operatorname{Res}_Y(C_0,R_{ij})$ would be a nonzero integer
polynomial vanishing at $x$.  Its height is
$(B+1)^{O(q_0^2)}$: its Sylvester matrix has size $O(q_0)$,
entries of coefficient norm $(B+1)^{O(q_0)}$, and
$O(q_0)!$ determinant terms.  This again contradicts the
root bound under the stated threshold.  Hence $C_0$ divides
every $R_{ij}$.

No factor $X+u$ or $Y+v$ vanishes identically on this curve,
since $C_0(x,y)=0$ with both coordinates larger than $B$.
The displayed square classes are therefore defined in
$K_0$.  If the classes of $F_{w_i},F_{w_j}$ were independent
there, they would still generate a degree-four field
extension.  Specializing the polynomial norm formula gives
zero, since $C_0\mid R_{ij}$; in a field this forces the
auxiliary to vanish.  Independence of its four characters
would force all four coefficient polynomials to be divisible
by $C_0$, contrary to their primitivity.  Thus the images
of every pair of basis vectors are dependent.  Over $\F_2$
this means that all nonzero images are equal, so the image
has dimension at most one.
\end{proof}

\subsection{Proof of Proposition~\ref{p13:rig:aligned}}
\label{p14:comp:rigidity:aligned}

\begin{proof}
Fix $\tau<1$.  Theorem~\ref{p13:rig:short-basis} gives a
space $V=\mathcal K_\tau$ of dimension
\[
 t\ge d_{xy}-\frac23B\log(1/\tau)
                     -O_{\tau,C}(B/\sqrt{\log B}),
\]
with a basis of half-cardinalities at most
$q_0=B^{\tau+o(1)}$.  The case $t<2$ is covered by
the displayed error term.  For $t\ge2$, apply
Proposition~\ref{p13:rig:two-witnesses} to each pair of
basis vectors.  Its threshold holds because
$q_0\log(B+1)=o(B)=o(\log N)$.

Put $\ell(X,Y)=bX-aY-h$.  Every resulting norm polynomial
$R_{ij}$ is divisible by $\ell$.  Otherwise, with
$d_Y=\deg_YR_{ij}$, the nonzero integer polynomial
\[
           a^{d_Y}R_{ij}\left(X,\frac{bX-h}{a}\right)
\]
would vanish at $x$.  Its height is $(B+1)^{O(q_0)}$
since $\max(a,b,|h|)\le2B^4$, contradicting the
integer root bound for $x\ge N$.
The primitive-norm argument in
Theorem~\ref{p13:rig:common-curve} now applies to this
line and shows that the image of $V$ in
$\mathbb Q(\ell=0)^\times/\mathbb Q(\ell=0)^{\times2}$
has dimension at most one.

Let $V_0$ be its kernel, of dimension at least $t-1$.
On the line, $X+u$ and $Y+v$ have the same root precisely when
\begin{equation}
                            h+bu=av.
 \label{p13:rig:line-matching}
\end{equation}
Every word of $V_0$ must pair all its first-block roots
with second-block roots.  The integer solutions of
\eqref{p13:rig:line-matching} advance by $(a,b)$ in two
intervals of $B$ consecutive offsets, so their number $r$
satisfies
\[
                        r\le1+\frac{B-1}{\max(a,b)}.
\]
Injectivity of the physical projection gives
$t-1\le\dim V_0\le r$, proving
\eqref{p13:rig:aligned-quantitative} after absorbing the
bounded terms.

For compatible distinct centres with $\max(a,b)\ge2$,
this gives $d_{xy}\le B/2+(2/3)B\log(1/\tau)+o_{\tau,C}(B)$.
Choose $\tau$ sufficiently close to one after fixing a
desired error margin, then let $N$ tend to infinity.
This order of choices proves the uniform $B/2+o_C(B)$ bound.

For $a=b=1$, compatibility implies $|x-y|\ge L$:
if $x<y<x+L$, the first run makes $f(y-1)=f(y)$,
whereas the second requires their signs to be opposite.
At distance exactly $L=B-1$, the true windows share one
position; at larger distance they are disjoint.  Thus
\eqref{p13:rig:line-matching} has at most one solution.
A nonzero word in $V_0$ would need at least two pairings,
because its selections are nonempty and even.  Consequently
$V_0=0$, $t\le1$, and the dimension estimate followed by
$\tau\uparrow1$ gives $d_{xy}=o_C(B)$.
\end{proof}

\section{All triples above finitely many rational pair channels}
\label{p14:comp:third}

The canonical overlap inequality and the positive $10:12:15$ witness
are proved in Section~\ref{p13:third:section}.  The remaining technical
input is uniformity of the upper bound across every triple above a
fixed finite set of aligned channels, including arithmetically
exceptional source pairs.  We prove that input here.

\subsection{Proof of Proposition~\ref{p13:third:aligned-family}}
\label{p14:comp:third:aligned-family}

\begin{proof}
Fix one channel $a:b$ and orient the source pair $e=\{x,y\}$.
Choosing $B^{O_{a,b}(1)}$ offsets, write
$x=at-i_0$, $y=bt-j_0$ and $X=abt$.
Merge repeated values in $bV_x\cup aV_y$ to obtain
$U_X\subset X+[-D B,D B]$, where $D$ depends only on the fixed channel.
Use radius $DB$ for alignment with this compressed source, and choose
the large-prime cutoff above $(2D+1)B$.  The old centre lies in a fixed
expanded interval starting at $N$, while the new centre remains in $I_N$,
as required by the expanded version of Lemma~\ref{p1:arith:asymmetric}.
Multiplying an even selection by a constant does not change its square
class, so
\begin{equation}
 S_x+S_y\subseteq
       \operatorname{span}\{v(n):n\in U_X\}.
 \label{p13:third:true-compression}
\end{equation}
This inclusion holds regardless of injectivity on $U_X$.
The uniform aligned rank bound of
Lemma~\ref{p1:arith:alignedrank} gives
\begin{equation}
 q(e)\leq N^{-\kappa_b+o(1)},\qquad
                   \kappa_b=2-\frac1b\geq\frac53.
 \label{p13:third:source-weight}
\end{equation}
For fixed height the residual defect is uniform, so the same bound
includes admissible local rank defects.

\paragraph{Targets not aligned with the compressed source.}
For $z\notin e$, put
$\delta=\dim(S_z\cap(S_x+S_y))$.
Connectedness of $T=e\cup\{z\}$ implies $\delta>0$.
The inclusion \eqref{p13:third:true-compression} bounds $\delta$ by
the asymmetric rank loss $\delta^*(X,z)$ of
Lemma~\ref{p1:arith:asymmetric}.  That lemma permits an arbitrary true
set $U_X$ of this fixed logarithmic width; it does not require injectivity.
It gives
\[
 \sum_{X,z\ {\rm nonaligned}}(2^{\delta^*(X,z)}-1)
                         \ll N^{31/24+\varepsilon'}.
\]
Since $q(T)=q(e)p_z2^\delta$, $p_z\leq2p$, and
$2^\delta\leq2(2^{\delta^*}-1)$, these targets have total mass at most
\begin{equation}
 N^{-\kappa_b-1+31/24+\varepsilon'+o(1)}
                     \ll N^{-11/8+\varepsilon}.
 \label{p13:third:unaligned-target-cost}
\end{equation}
The number of source-offset patterns is polynomial in $B$ and is
absorbed into the exponent margin.

\paragraph{The residual kernel for an aligned target.}
This sector includes every target aligned at the same enlarged radius
$DB$.  Enlarging that fixed radius changes only the constants in the
polynomial offset and height bounds below.
Write an aligned target as
$z=(ct+h)/d$, with $(c,d)=1$ and
$c,d,|h|\leq B^{O_{a,b}(1)}$, as in the proof of
Lemma~\ref{p1:arith:fixed-channel-rank}.
Before this final normalization, $U_X$ and $V_z$ each have fixed
logarithmic width.  Choose $H>C_D B$ larger than both diameters and
remove the prime rows at most $H$ and those dividing the normalization
multipliers.  Each remaining prime row meets at most one distinct old
value and at most one new value.  This uses the fact that repeated
old values were merged before the target was added.

After the common normalization, merge equal values once more.
There are at most $3B$ distinct integers, all at height at least a
constant multiple of $N$ and in an interval of diameter
$B^{O_{a,b}(1)}$.  The remaining prime rows have weight at most two,
and leave one variable per unpinned component.  We give the uniform
residual estimate, since the coefficients $c,d$ need not be fixed.

For a fixed integer $K$, restrict to components of at most $K$ vertices.
The removed prime rows and a total-parity row form a code on these
components of redundancy
\[
 s\leq\pi(H)+O_{a,b}(\log B)+1=O_{a,b}(B/\log B).
\]
A nonzero word in its kernel gives an even square product of distinct
normalized integers.  Lemma~\ref{p1:arith:runge}, applied with the
polynomial-in-$B$ diameter just specified, forces it to contain
$\gg_{a,b}\log N/\log B$ physical vertices and hence
$\gg_{a,b,K}\log N/\log B$ components.
If $c_K$ is the number of these components, the Hamming inequality with
$t\asymp_{a,b,K}\log N/\log B$ gives
\[
 \sum_{j=0}^{t}\binom{c_K}{j}\leq2^s,
 \qquad c_K=O_{a,b,K}(B/\log B).
\]
Here $t$ is chosen below half the component minimum distance.
Components larger than $K$ number at most $3B/K$.
Choosing $K$ arbitrarily large before $N$ therefore bounds the entire
value kernel on the distinct normalized set by $o_{a,b}(L)$,
uniformly over $c,d,h$.

The relevant map factors as
\[
 \bigoplus_{v\in T}E_v
   \longrightarrow\{\text{merged normalized coordinates}\}
   \longrightarrow\{\text{square classes}\}.
\]
Its nullity is at most the nullity of merging plus the residual $o(L)$
just proved.  Internal defects of the old windows are included in this
composition and do not require any assumption of good normalization.

\paragraph{Counting the exact coincidences.}
In the common parameter the physical coefficients are $ad,bd,c$.
Suppose first that they are distinct.  Divide by their common gcd
and order them as $u<v<w$.  Unless the parameter belongs to a
$B^{O_{a,b}(1)}$ endpoint population, dyadic geometry gives $w<2u$.
Indeed a rational ratio of polynomial height cannot cross two by
$B^{O(1)}/N$, and equality forces the extreme windows to the two
dyadic endpoints.
The lattice calculation \eqref{p13:third:density} and
Lemma~\ref{p13:third:rational-rank} show that merging loses at most
$19L/30+O(1)$ dimensions.  The residual estimate therefore gives
\[
 \dim(S_x+S_y+S_z)\geq\frac{71}{30}L-o(L),\qquad
 q(T)\leq N^{-71/30+o(1)}.
\]
There are $NB^{O_{a,b}(1)}$ aligned triples of this form, including the
choices of offsets and rational coefficients.  Their total mass is
$O(N^{-41/30+\varepsilon})$.

If two physical coefficients are equal, compatibility makes the
interiors of the corresponding parameter intervals disjoint, since
distinct compatible starts have separation at least $L$.
Their combined coincidences with the third lattice are bounded by the
length of that third interval: at most $L/a+O(1)$ or $L/b+O(1)$ units.
As $a\geq2$, the total merging nullity is at most $L/2+O(1)$.
Including the residual, this gives mass $O(N^{-3/2+\varepsilon})$.
Endpoint populations have only $B^{O_{a,b}(1)}$ source and target
choices; the bound $q(T)\leq q(e)$ from
\eqref{p13:third:source-weight} gives $O(N^{-5/3+\varepsilon})$.
These estimates include exact contacts.  Strict contacts are incompatible
and do not enter the sum.

The least of the exponents obtained is $41/30$.  Sum over the fixed
finite family $\mathcal F$ to conclude.
\end{proof}

\section{Conditional dense targets and small-kernel counts}
\label{p14:comp:transfer}

These are the two geometric proofs used in
Section~\ref{p13:transfer:section}. Their source counts, widths, and rank
thresholds are fixed before $N$ tends to infinity. They do not use the
finite-transfer theorem or the law approximation.

\subsection{Proof of Lemma~\ref{p13:transfer:dense-targets}}
\label{p14:comp:transfer-dense-targets}

\begin{proof}
We first count target pairs, keeping the number of binary components in
\emph{each} target explicit. All component counts below include isolated
active coordinates.

\paragraph{The two injective projections.}
Put
\[
 R=\ker\bigl(W\oplus S_y\oplus S_z\longrightarrow W+S_y+S_z\bigr),
 \qquad \dim R=d.
\]
If an element of $R$ has zero $y$-projection, its other two entries
belong to $W\cap S_z=0$. Hence the projection $R\to S_y$ is
injective; the other projection is injective by $W\cap S_y=0$.
In particular $d\leq\min(\dim S_y,\dim S_z)\leq L$.
Represent $W$ by a linear section in the physical source coordinates,
and each target space by a section in its formal even-selection space.
These sections preserve both injective target projections and all actual
prime-valuation equations, including on an irregular window.

Fix $0<\delta<1/2$, to be chosen inside the final exponent margin,
and choose fixed positive $\xi,\theta$ with
\[
 \xi\leq\tfrac14\min(\eta,\delta),\qquad
 \theta\leq\tfrac18\min(\eta,\delta).
\]
Apply Lemma~\ref{p1:fixed:recovery} to obtain $R'\subseteq R$ with
$d':=\dim R'\geq d-\xi L$. Choose its recovery parameter large
enough to give this loss and to control circuits in the $s+2$ physical
blocks. Every nonzero coordinate label on $R'$ then has a fixed
bounded multiplicity $K$.

For each prime $\ell>H\asymp B$, pair equal active labels in its
row. All edges of the resulting binary graph join equal nonzero
labels, so every component lies within one label class and has at
most $K$ vertices. The coordinate labels in either target span
$(R')^*$, since its projection is injective. A component contributes
only one label to this span, even if it meets that target several times.
Consequently each target meets at least $d'$ binary components.
This is a separate lower bound for each target, not just for their union.

\paragraph{Enumerating coefficients from the fixed sources.}
After the equal labels have been paired, the remaining distinct labels
in a prime row split into minimal circuits. Each such circuit uses at
least three physical blocks: $\ell>H$ has at most one occurrence in
any block. Since only two blocks are free, every circuit contains a
source occurrence. The product of primes assigned to a circuit support
therefore divides a fixed source value.
The recovery bound gives a fixed bound on the number of circuit supports
incident to any selected coordinate. For a fixed number of selected
binary components, enumerate their positions, these circuit supports,
one source occurrence on each support, and divisors of the corresponding
source values. The cost is
$B^{O(1)}N^{o(1)}=N^{o(1)}$. Small-prime parities cost
$2^{O(\pi(H))}=N^{o(1)}$ as well. Thus the squarefree coefficients
of their component equations have $N^{o(1)}$ choices and height
$N^{A_0}$ for a fixed $A_0$. The height bound per component is fixed
once $s$, the recovery parameters and the width multiplier are fixed;
it does not increase when more disjoint component equations are selected.
The number of equations required by Lemma~\ref{p1:arith:geometric}
is chosen afterwards. No subspace or full deletion pattern is enumerated.

\paragraph{Removing the small aligned curves.}
Two exact units between the targets determine a reduced alignment of
height $O(B)$. Two distinct alignments of this height would force
$y,z=O(B^3)$, impossible in the critical range. If there are fewer than
two exact units, no removal is needed. Otherwise let $u\geq2$ be the
number of units on the unique alignment. Their even selections have dimension
$u-1$ before passage to square classes, and a target has at most one
local rank defect. Their common image therefore gives
\[
 u-2\leq\dim(S_y\cap S_z)\leq\theta L
\]
in the weak aligned case. Keep at most one of these units and exclude
from selection every binary component meeting an endpoint of another
unit. At most $2(u-1)\leq2\theta L+2$ components are excluded.
This deletes components from the selectable collection; it does not
restrict $R'$ again. The surviving positions contain no two exact units
on an alignment, which excludes the small-curve alternatives in the
proof of Lemma~\ref{p1:arith:geometric}. No removal is needed for a
nonaligned pair.

For $t=y,z$, write $m_t$ for the number of selectable components
meeting target $t$. The separate projection bounds now give
\begin{equation*}
 m_t\geq d'-2\theta L-2
      \geq d-(\xi+2\theta)L-2.
\end{equation*}
Thus $m_t\geq\eta L/3$ for large $N$ whenever $d\geq\eta L$;
if $d>(1/2+\delta)L$, then
$m_t>(1/2+\delta/3)L$ for both targets.

\paragraph{The first stratum: a positive rank proportion.}
A component type records its number of coordinates in each physical
block. There are finitely many types, because the component size is at
most $K$. If a type meeting both targets occurs sufficiently often,
select the fixed number of its components required by the two-variable
geometric count. With the small curves excluded above and coefficients
of height $N^{A_0}$, the proof of Lemma~\ref{p1:arith:geometric}
gives $N^{1/2+\varepsilon'}$ target pairs, including the
$N^{o(1)}$ pattern and coefficient choices.

Otherwise only a bounded number of components meet both targets.
Each target still has a positive proportion of $L$ components avoiding
the other target. In each target choose a univariate type occurring with
positive frequency, and then the fixed number of its components needed
by the one-variable part of Lemma~\ref{p1:arith:geometric}, taking
at least two. Their component equations form a one-variable system
with fixed coefficients and at least two distinct selected offsets;
different components have disjoint vertices. In the degree-one case,
these are precisely the two linear factors for
Lemma~\ref{p1:arith:pell}. The one-variable count gives
$N^{\varepsilon'}$ target pairs after allocating the margin between
the two centers and the coefficient choices. Every selected set has fixed size, so its even-selection
square-class map is injective for large $N$ by
Lemma~\ref{p1:arith:runge}; irregularity of the entire window is
irrelevant to this local use. We have proved
\[
 \#\{(y,z):d\geq\eta L\}\ll N^{1/2+\varepsilon'}.
\]
Both counts here and below impose the pair restriction in the lemma.

\paragraph{The second stratum: rank greater than half a window.}
Suppose $d>(1/2+\delta)L$. Let $n_{1,t}$ be the number of selectable
components using exactly one coordinate in target $t$. Every other
component meeting that target uses at least two of its $B=L+1$
coordinates. Hence, separately for $t=y,z$,
\[
 B\geq n_{1,t}+2(m_t-n_{1,t}),\qquad
 n_{1,t}\geq2m_t-B>\frac{2\delta}{3}L-1.
\]
If a mixed type with one coordinate in either target occurs sufficiently
often, the degree-one case of the geometric count gives
$N^{\varepsilon'}$ pairs directly. Otherwise only boundedly many of
the $n_{1,t}$ components are mixed. The displayed inequality then
provides two components for each target, each using one coordinate of
that target and no coordinate of the other. The resulting two distinct
linear factors give the uniform two-value Pell count for each center.
Thus
\[
 \#\{(y,z):d>(1/2+\delta)L\}\ll N^{\varepsilon'}.
\]
This explains the threshold $L/2$: above it, the coordinate budget
forces singleton occurrences in each target.

Finally $p_y,p_z\ll N^{-1}$ and $2^L\asymp_C N$.
The stratum $\eta L\leq d\leq(1/2+\delta)L$ has weighted cost
$N^{1/2+\varepsilon'}N^{-2+1/2+\delta}$, and the stratum above it
has cost $N^{\varepsilon'}N^{-2+1}$. Choose $\delta>0$ and then
$\varepsilon'>0$ so that $\delta+\varepsilon'<\varepsilon$;
all preceding choices are fixed before $N$. This proves
\eqref{p13:transfer:dense-target-bound}.
In a primitive family, each target is direct from the prescribed
$r-2$ sources separately, as well as from the other target; the full
nullity is at most $L$. Multiplication by the $r-2$ source marginals
and summation over the prescribed source collection gives
\eqref{p13:transfer:restricted-dense}, after dividing the final exponent
margin between the source count and the conditional estimate.
\end{proof}

\subsection{Proof of Lemma~\ref{p13:transfer:small-kernels}}
\label{p14:comp:transfer-small-kernels}

\begin{proof}
For the first assertion fix the $m$ offsets, at a polylogarithmic cost,
and write $d_i=n_i/K_Y(n_i)=s_i a_i^2$. Since
$d_i\gg N^{1-b}$ and $\gcd(d_i,d_j)\ll B$,
\[
 \operatorname{lcm}(d_1,\ldots,d_m)
 \geq\frac{\prod_i d_i}{\prod_{i<j}\gcd(d_i,d_j)}
 \geq N^{m(1-b)-o(1)}>N.
\]
The $a_i$ are pairwise coprime. If $\prod_i a_i^2>N$, multiply the
square factors in a fixed order until their product first exceeds $N$.
The resulting square $D$ lies between $N$ and $N^{1+u}$, so it has
at most $N^{(1+u)/2}$ possible roots. Assign its complete prime powers
to the $m$ offsets; the number of assignments is
$m^{\omega(D)}=N^{o(1)}$. The CRT leaves at most one center modulo
$D>N$ in the interval.

If $\prod_i a_i^2\leq N$, start with this product and add the small
prime factors from the least common multiple until the first crossing
of $N$. Now $N<D\leq NY$ and $D=sa^2$ with $s$ $Y$-smooth.
For fixed $\sigma>\max(1/2,1-1/A)$, the number of such modules is at most
\[
 (NY)^\sigma\sum_{a\geq1}a^{-2\sigma}
       \prod_{\ell\leq Y}(1-\ell^{-\sigma})^{-1}
       =N^{\sigma+o(1)}.
\]
Choose $\sigma$ just above $1-1/A$. Assigning the complete prime powers
again costs $N^{o(1)}$ and the CRT gives at most one center. This proves
\eqref{p13:transfer:small-kernels-bound}.

For the second assertion, use $H=C_DB$ as in
Lemma~\ref{p13:half:small-kernels}.
That lemma gives the result for $q>0$.
For $q=0$, use $\mathcal X_0\subseteq\mathcal X_{\varepsilon/2}$
and apply the lemma with error margin $\varepsilon/2$.
\end{proof}

\section{The first retained class on the regular set}
\label{p14:comp:stopped}

This section specializes the general stopping identity of
Section~\ref{p1:stop:section} to the two leading pair channels. It
supplies the boundary estimates used by the first reconstruction and
the positional comparison. The optimal scalar reconstruction uses the
larger class $\mathscr D_2$ in Section~\ref{p13:transfer:section}.
Here and in the regular-set reconstruction below, write
\[
 G=G_N,\quad n=|G|,\quad Z_G=Z_N^\circ,\quad
 \mu=np,\quad E_N=\mathcal E_N,
 \qquad a_N=\E(Z_N)_2-\lambda_N^2.
\]
The pair calculation and mean estimate give
$a_N-\alpha_N=O_{C,\varepsilon}(N^{-1+\varepsilon})$ and
$a_N\asymp_C LN^{-2/3}>0$ for large $N$.
Every space on $G$ has dimension $L$, and every marginal is $p$.

\subsection{Retaining a direct family and one compressed pair}

We now use the regular set $G$ of arithmetic windows.
For $I\subseteq G$, recall
\[
 \rho(I)=|I|L-\dim\sum_{x\in I}S_x.
\]
Every $S_x$ has dimension $L$, and a family is called direct
when $\rho(I)=0$.  Let $E_N$ be the family of good pairs
from Corollary~\ref{p1:arith:edgefamily}.  Its pairs are
compatible and dependent, and
\begin{equation}
 \rho(e)\le L/3+O(1)\qquad(e\in E_N).
 \label{p1:stop:pair-nullity}
\end{equation}
For $e=\{x,y\}$ put $V_e=S_x+S_y$ and $q_e=q(e)$.  Define
\begin{equation}
 \Gamma_e(h)=\left\{
 U\subseteq G\setminus e:|U|=h,\quad
 V_e,\ (S_z)_{z\in U}\ \text{are in direct sum}
                    \right\}.
 \label{p1:stop:gamma}
\end{equation}
The direct-sum condition concerns the single compressed block
$V_e$ and the other blocks.  It does not assert independence
of the two spaces inside $e$.  The affine probability formula
gives
\begin{equation}
 q(e\cup U)=q_ep^{|U|}
              \qquad\bigl(U\in\Gamma_e(|U|)\bigr).
 \label{p1:stop:direct-pair-probability}
\end{equation}

Let $\mathscr D$ consist of all direct families together with
all families $e\cup U$ occurring in
\eqref{p1:stop:gamma}.  Thus a family in $\mathscr D$ is allowed
one good pair, treated as a single block, but no further dependence.

\begin{lemma}[Heredity and boundary nullity]
\label{p1:stop:class}
The family $\mathscr D$ is hereditary.  Each non-direct member
contains a unique dependent pair $e\in E_N$, and its nullity
equals $\rho(e)$.  Hence, for all sufficiently large $N$,
\begin{equation}
 \partial\mathscr D
 \subseteq\{I\notin\mathscr D:
                 0<\rho(I)\le 4L/3+O(1)<3L/2\}.
 \label{p1:stop:boundary-nullity}
\end{equation}
\end{lemma}

\begin{proof}
Every subfamily of a direct family is direct.
Consider $e\cup U$ in the second part of the definition.
A subfamily which retains both endpoints of $e$ still consists
of $V_e$ and some of the independent singleton blocks.  A subfamily
which removes an endpoint is direct, since the remaining endpoint
space is contained in $V_e$.  This proves heredity.

The direct-sum condition excludes a dependent pair inside $U$
and excludes a dependent pair joining $U$ to either endpoint of
$e$.  Thus $e$ is the unique dependent pair.  Additivity of the
ranks of the compressed blocks then gives
$\rho(e\cup U)=\rho(e)$.

Adding one $L$-dimensional space increases nullity by at most
$L$.  The proper prefix of a boundary member has nullity at most
$L/3+O(1)$, so the member itself has nullity at most
$4L/3+O(1)$.  Finally, nullity zero would make it a direct
family and therefore a member of $\mathscr D$.
\end{proof}

The retained dependence has the size of the principal correction.
The strict upper bound on the boundary nullity is what makes
the stopped remainder easier to estimate than the unrestricted
absolute expansion.

\subsection{The factorial tail of the boundary}

Write $n=|G|$ and $\mu=np$, and choose a fixed upper bound
$\Lambda$ for $\mu$ throughout the critical band.  We need a
cutoff large enough to absorb the loss $2^{3L/2}$ in
\eqref{p1:stop:boundary-nullity}, yet small enough for the
critical-size component estimate.  The following explicit
choice meets both requirements:
\begin{equation}
 K_N=\min\left\{K\in\mathbb N:
 K\ge\max(4,\lceil4e\Lambda\rceil),\quad
 K\log\frac{K}{2e\Lambda}\ge3\log N\right\}.
 \label{p1:stop:cutoff}
\end{equation}
It satisfies
\begin{equation}
 K_N=(3+o(1))\frac{\log N}{\log\log N}=o(L);
 \qquad K_N\le L/4
 \quad\text{for all sufficiently large }N.
 \label{p1:stop:cutoff-range}
\end{equation}
The coefficient $3$ is a convenient sufficient choice.

\begin{lemma}[Large boundary families]
\label{p1:stop:tail}
Uniformly in the fixed critical band,
\begin{equation}
 \sum_{\substack{I\in\partial\mathscr D\\|I|>K_N}}
           2^{|I|}q(I)=O_C(N^{-3/2}).
 \label{p1:stop:tail-bound}
\end{equation}
\end{lemma}

\begin{proof}
The affine law and \eqref{p1:stop:boundary-nullity} imply
\[
 q(I)\le 2^{3L/2}p^{|I|}
          \ll_C N^{3/2}p^{|I|}.
\]
This remains valid for incompatible families, whose probability
is zero.  Therefore
\[
 \sum_{\substack{I\in\partial\mathscr D\\|I|>K_N}}
       2^{|I|}q(I)
 \ll_C N^{3/2}
       \sum_{r>K_N}\binom nr(2p)^r
 \le C_C N^{3/2}
       \sum_{r>K_N}\frac{(2\mu)^r}{r!}.
\]
By $K!\ge(K/e)^K$ and \eqref{p1:stop:cutoff},
$(2\mu)^{K_N}/K_N!\le N^{-3}$.
The ratios of subsequent terms are at most
$2\Lambda/(K_N+1)<1/2$, so the last tail is at most
$2N^{-3}$.  This proves the assertion.
\end{proof}

Only families of size at most $K_N$ remain.  They have nullity
at most $3L/2$, a bound fixed relative to $L$.  The component
estimate in Theorem~\ref{p1:core:bounded-nullity} supplies the
arithmetic input for this range.  Section~\ref{p1:close:section}
will convert that centered estimate into the positive boundary
mass required by \eqref{p1:stop:norm}.

\section{Proofs of the uniform certificate tools}
\label{p14:comp:certificate-tools}

All definitions and parameter conditions are those of
Section~\ref{p14:cert:section}. The counts are geometric, and all
constants are uniform in the size of the family being summed.

\subsection{Proof of Lemma~\ref{p1:core:rankin}}
\label{p14:comp:certificate-rankin}

\begin{proof}
Write uniquely $v=k s a^2$, where $k=K_Y(v)$ and $s$ is the
squarefree product of the primes at most $Y$ occurring to odd order.
The inequality $1\le (C_1N/(ksa^2))^\sigma$ on the range being
counted gives an upper bound
\[
 (C_1N)^\sigma
 \left(\sum_{k\le T}k^{-\sigma}\right)
 \left(\sum_{a\ge1}a^{-2\sigma}\right)
 \prod_{\ell\le Y}(1+\ell^{-\sigma}).
\]
Enlarging the sum over $k$ to all positive integers only increases it.
The three factors are respectively $O_\sigma(T^{1-\sigma})$,
$\zeta(2\sigma)$, and at most
$\exp(\sum_{j\le Y}j^{-\sigma})$. This proves the bound.
\end{proof}

\subsection{Proof of Lemma~\ref{p1:core:moderate}}
\label{p14:comp:certificate-moderate}

\begin{proof}
There are at most $rNB^2$ choices for the anchor and its two offsets.
The product $Q$ has at most $m$ prime factors, and each has at most
$rB$ destinations. Once these assignments are fixed, the Chinese
remainder theorem leaves at most $1+N/D_j$ possibilities for centre
$j$. At most two of the $D_j$ exceed $N$, since $Q\le9N^2$.
For those two modules, replacing $1+N/D_j$ by $2N/D_j$ costs at most
$(6B)^2$. Therefore
\[
 \prod_{j\ne\mathrm{anchor}}(1+N/D_j)
 \le 2^{r-1}(6B)^2N^{r-1}T^{-2}.
\]
This proves \eqref{p1:core:moderate-count}; allowing collisions and
incompatible assignments only enlarges the count.

After multiplication by $(up)^r/r!$, the remaining series is, up to
a fixed factor,
\[
 N^{-2\beta}B^{m+4}\sum_{r\ge0}\frac{z^r r^{m+1}}{r!},
 \qquad z=2uNp=O_u(1).
\]
Since $r^{m+1}\le(m+1)!e^r$, this series is at most
$N^{-2\beta}B^{m+4}(m+1)!\exp(ez)$. Finally
\[
 m=\left(\frac2{A_0}+o(1)\right)\frac{H}{\log H},\qquad
 \log(B^m(m+1)!)=\left(\frac4{A_0}+o(1)\right)H.
\]
This yields \eqref{p1:core:moderate-baseline}. The factorial here
divides the orderings of a set of distinct windows; it is not an
assumption on the growth of a fixed-order arithmetic constant.
\end{proof}

\subsection{Proof of Lemma~\ref{p1:core:degrees}}
\label{p14:comp:certificate-degrees}

\begin{proof}
For an outgoing edge, choose two anchor offsets in $B^2$ ways. Each of
the at most $m$ primes is either unused or assigned to one of $B$
offsets in the receiver. A qualifying code has modulus $D>N$ and
therefore determines at most one receiving centre in $I_N$. This gives
$B^2(B+1)^m=N^{2/A_0+o(1)}$ possibilities.

For the incoming degree, fix the receiving window $y$ and let
$\Omega_y$ be the primes greater than $Y$ dividing any of its values.
Writing $R=\lceil B\log(3N)/\log Y\rceil$, we have $|\Omega_y|\le R$.
A code chooses two anchor offsets, at most $m$ primes from $\Omega_y$,
and one of the two offsets for each chosen prime. For a qualifying
code their product exceeds $N$, so its congruences determine at most
one anchor centre. Thus the number of possible incoming neighbours is
at most
\[
 B^2\sum_{j=0}^m\binom Rj2^j
 \le B^2(m+1)(2eR/m)^m
 =N^{2/A_0+o(1)}.
\]
The displayed binomial bound holds for large $N$, since $m\le R/2$.
Nonqualifying codes may be included in the number of codes; the
one-centre conclusion is applied only to qualifying ones. No receiving
offset needs to be chosen in this incoming count: the receiver has
already been fixed. Symmetrisation adds the two degree bounds.
\end{proof}

\subsection{Proof of Lemma~\ref{p1:core:generator}}
\label{p14:comp:certificate-generator}

\begin{proof}
Fix an admissible set $D$. Its isolated vertices are marked singletons.
On the other vertices choose an inclusion-minimal edge cover of the
induced graph. Every edge of this cover has an endpoint of degree one
in the cover, since otherwise that edge could be removed. Its
components are therefore nontrivial stars. Thus $D$ has a partition
into marked singletons and stars, each with one centre and a nonempty
subset of its neighbours as leaves.

Write $w(A)=\prod_{v\in A}w_v$. Choose one such partition for each $D$.
The product over all possible blocks of $1+w(A)$ contains a distinct
term of weight $w(D)$ for this chosen partition: its union recovers
$D$. Dropping disjointness only adds nonnegative terms. Since
$1+u\le e^u$ for $u\ge0$, the weighted sum is at most the exponential
of the total block weight. The singleton contribution is
$\sum_{v\in\mathcal M}w_v$. For a fixed star centre $v$, its leaf sum
is exactly
\[
 w_v\sum_{\varnothing\ne A\subseteq N(v)}\prod_{u\in A}w_u
   =w_v\left(\prod_{u\sim v}(1+w_u)-1\right).
\]
Counting the same two-vertex star from both centres is harmless.
This proves the heterogeneous assertion. Set $w_v=z$ and use
$\deg(v)\le\Delta$ to obtain \eqref{p1:core:generator-bound}.
There is no small-activity hypothesis. When $\Delta z<2$, the
elementary bound $e^u-1\le2u/(2-u)$ also gives the convenient
majorant $\exp\{Mz+2n\Delta z^2/(2-\Delta z)\}$.
\end{proof}
\section{From fixed orders to all primitive families}
\label{p1:core:section}

The estimates at each fixed order do not control the growing order needed
to reconstruct a distribution. We now obtain that uniformity without
estimating the growth of the fixed-order constants. The argument has two
parts. A large squarefree kernel produces congruences which count the
possible centres. When too much of this kernel is received by one other
window, it produces an edge in a graph of small degree. A generating
function then controls all large families at once. Only a finite set of
small cardinalities is left to the fixed-order estimates.

Throughout this section, a family of distinct windows is \emph{primitive}
if its spaces are not in direct sum but every proper subfamily is in
direct sum. Connectedness refers to the canonical decomposition into direct-sum
components in Lemma~\ref{p1:alg:components}.
All families in this section consist of windows in the regular set $G_N$.
The quantities $q(I)$, $c(I)$, $\rho(I)$ and $\rho_*(I)$ retain their
definitions from Section~\ref{p1:alg:section}. All weights and all auxiliary
parameters in this section are fixed before $N$ tends to infinity.

\subsection{The centred estimate at each fixed order}

The small-cardinality part of the argument needs an estimate for absolute
centred coefficients, including incompatible affine systems. Positive
probabilities alone do not give this estimate: a joint event can be
impossible while its centred coefficient is nonzero. For separated windows
we remove this difficulty by using a homogeneous majorant. Contacts require
a different argument, since their homogeneous probabilities can be large.

For $x\in G_N$, let $H_x$ be the indicator of the homogeneous event that
every character in $S_x$ takes the value $1$. Equivalently, the random
multiplicative function is constant on the $L+1$ values in window $x$. Set
\begin{equation}
 W(I)=\E\prod_{x\in I}H_x
      =p^{|I|}2^{\rho(I)}.
 \label{p1:core:homogeneous-weight}
\end{equation}
All homogeneous systems are compatible. Moreover, the full-support Fourier
formula of Lemma~\ref{p1:alg:fourier}, after taking absolute values, gives
\begin{equation}
 |c(I)|\le W(I).
 \label{p1:core:homogeneous-majorant}
\end{equation}
Here the geometric good-pair class includes pairs whose prescribed affine
run-start events are incompatible; denote this enlarged class by
$\mathcal A_{\rm raw}$.

\begin{lemma}[Absolute centred mass at a fixed order]
\label{p1:core:fixed-centred}
For every fixed integer $r\ge3$ and $\varepsilon>0$,
\begin{equation}
 \sum_{\substack{C\subseteq G_N,\ |C|=r\\C\ \mathrm{connected}}}
       |c(C)|\ll_{r,\varepsilon}N^{-4/5+\varepsilon}.
 \label{p1:core:fixed-centred-bound}
\end{equation}
The sum over all $r$-element families containing a contact, with no
connectedness restriction, is $O_{r,\varepsilon}(N^{-1+\varepsilon})$.
\end{lemma}
\begin{proof}
First suppose that all windows are separated. Call a family residual if
its relation space is neither zero nor exactly the relation space of one
geometric good pair, extended by zero on the other blocks. A connected
family of at least three windows is residual. We prove the stronger bound
obtained by summing $W(C)$ over all residual separated families of size $r$.

Proposition~\ref{p1:fixed:propagation} applies to the homogeneous targets
as well as to the prescribed affine targets. Its rank increment
$\dim(S_x\cap\sum_{y\in I}S_y)$ is geometric, and the conditional
probability is $p$ times the corresponding power of $2$ in either case.
Consequently a seed may be completed to any fixed total cardinality at
a cost $N^{\varepsilon}$, after reducing the error exponent used for
the seed. The following four cases exhaust residual families.

If a family contains a dependent pair outside $\mathcal A_{\rm raw}$,
the first seed estimate of Proposition~\ref{p1:arith:fixed-homogeneous-seeds}
gives total weight $O_\varepsilon(N^{-4/5+\varepsilon})$. Propagation
adds the other $r-2$ windows at the stated cost. If it contains an
original primitive subfamily of size $h\ge3$, with $h\le r$,
Proposition~\ref{p1:fixed:primitive} gives the same bound for its
homogeneous weight. There are only finitely many such values of $h$,
so propagation again gives the required estimate.

Next suppose that the family contains two distinct geometric good pairs.
Their union has three or four windows. By
Proposition~\ref{p1:arith:fixed-homogeneous-seeds}, the total homogeneous
weight of these unions is
\[
 O_\varepsilon(N^{-5/4+\varepsilon}+L^2N^{-4/3})
       =O_\varepsilon(N^{-5/4+\varepsilon}).
\]
Completing the union to $r$ windows costs only $N^{\varepsilon}$.
This case includes a pair sharing one endpoint with the other pair.

The remaining family has a unique geometric good pair $e$ and an
additional relation. Compress $e$ exactly, and choose a minimal dependent
subfamily of the resulting blocks. If it avoids $e$, it is one of the
original primitive seeds already considered. Otherwise it consists of
the compressed pair and $h-1$ original windows. For $h=2$, the
pair-and-target seed estimate of
Proposition~\ref{p1:arith:fixed-homogeneous-seeds} gives
$O_\varepsilon(N^{-4/3+\varepsilon})$ for channel $2:3$ and
$O_\varepsilon(N^{-5/4+\varepsilon})$ for channel $3:4$.
For $h\ge3$, Proposition~\ref{p1:fixed:compressed} gives
\[
 O_{h,\varepsilon}(N^{-\kappa_b+1/5+\varepsilon}),
 \qquad \kappa_b=2-1/b,\quad b\in\{3,4\}.
\]
The largest of these bounds is $O(N^{-22/15+\varepsilon})$.
Propagation is applied after replacing the compressed seed by its original
windows. The number of seeds and placements in a family of fixed size is
bounded in terms of $r$. These cases prove the separated estimate by
\eqref{p1:core:homogeneous-majorant}. No compatibility restriction has
been introduced in this argument.

For contacts, return to the actual affine events. Expanding the centred
product gives
\begin{equation}
 |c(C)|\le\sum_{U\subseteq C}p^{r-|U|}q(U).
 \label{p1:core:contact-expansion}
\end{equation}
Fix $h=|U|$ and write $u=r-h$. If $U$ already contains a contact,
the positive contact seed followed by
Proposition~\ref{p1:fixed:propagation} gives
\[
 \sum_{\substack{|U|=h\\U\ \mathrm{contains\ a\ contact}}}q(U)
       \ll_{h,\varepsilon}N^{-1+\varepsilon}.
\]
The remaining $u$ centres and their factors $p^u$ cost at most
$(Np)^u=O_r(1)$. If $U$ is separated but $C$ contains a contact,
at least one endpoint of a contact is among the $u$ new centres.
For fixed $U$ there are at most $O_r(BN^{u-1})$ such completions:
either a new centre is within $O(B)$ of $U$, or two new centres are
within $O(B)$ of each other. The cases where no completion is possible
contribute zero. Since
\[
 \sum_{|U|=h}q(U)
       =\E\binom{Z_N^\circ}{h}=O_h(1)
\]
by the fixed moments in Theorem~\ref{p1:fixed:moments}, these
completions contribute $O_r(B/N)$. Summing over the finitely many
values of $h$ proves the contact assertion. The moments used here are
proved from the positive affine classification before this lemma; they
do not depend on any all-cardinality centred estimate.
\end{proof}

\subsection{Primitive families at all cardinalities}

\begin{theorem}[Primitive families at all cardinalities]
\label{p1:core:primitive}
For every fixed $w>0$ and $\varepsilon>0$,
\begin{equation}
 \sum_{\substack{C\subseteq G_N\ \mathrm{primitive}\\|C|\ge3}}
 w^{|C|}\bigl(q(C)+|c(C)|\bigr)
 \ll_{w,\varepsilon}N^{-4/5+\varepsilon}.
 \label{p1:core:primitive-bound}
\end{equation}
The sum includes overlapping windows and incompatible affine systems.
\end{theorem}

\subsection{One rank loss for a primitive family}

If $C$ is primitive, the projection of its relation space onto any one
block is injective. Indeed, an element in the kernel of that projection
would be a relation supported on a proper subfamily. Since a regular
window has rank $L$, a primitive family of $r\ge3$ windows satisfies
\begin{equation}
 1\le\rho(C)\le L,\qquad
 q(C)\le p^{r-1},\qquad c(C)=q(C)-p^r.
 \label{p1:core:primitive-rank}
\end{equation}
The last identity follows by expanding the centred product: every proper
subfamily has joint probability equal to the product of its marginals.
Consequently $|c(C)|\le p^{r-1}$, including when $q(C)=0$.

A nonzero relation of a primitive family uses every block. Its projection
on a regular window corresponds to a nonempty even selection of the
$B=L+1$ values in that window. It therefore selects at least two distinct
values there. This elementary fact is what permits a two-value counting
argument, rather than a one-value friability estimate.

The kernel, moderate-certificate, saturation graph and generating-function
estimates used next are stated in Section~\ref{p14:cert:section} and
proved in Section~\ref{p14:comp:certificate-tools}. Their $p$-weighted
forms are used here on the regular set.

\subsection{The finite-small-order, uniform-large-order split}

Take
\begin{equation}
 A_0=40,\quad\beta=\frac{24}{25},\quad\sigma=\frac{49}{50},
 \quad\delta=\frac1{1250},\quad\theta=\frac2{A_0}=\frac1{20}.
 \label{p1:core:primitive-parameters}
\end{equation}
These values leave a strict margin in all subsequent inequalities.
For a fixed $w>0$, substitute $z=wpN^{\delta/2}$ into
Lemma~\ref{p1:core:generator}. The three quantities controlling its
exponent satisfy
\[
 Mz\le N^{-\delta/2+o(1)},\qquad
 \Delta z\le N^{-1+\theta+\delta/2+o(1)},\qquad
 n\Delta z^2\le N^{-1+\theta+\delta+o(1)}.
\]
They tend to zero, and
$nz((1+z)^\Delta-1)\le nz(e^{\Delta z}-1)=O(n\Delta z^2)$.
Consequently, for every fixed integer $s\ge3$,
\begin{equation}
 \sum_{\substack{C\ \mathrm{primitive},\ |C|\ge s\\
                         C\ \mathrm{without\ a\ moderate\ certificate}}}
 w^{|C|}\bigl(q(C)+|c(C)|\bigr)
 \ll_w N^{1-s\delta/2}.
 \label{p1:core:primitive-tail}
\end{equation}
Indeed, removing the tilt contributes $N^{-s\delta/2}$ and the
primitive rank bound contributes at most $2p^{-1}=O(N)$.
For the moderate families, Lemma~\ref{p1:core:moderate} and the same
rank bound give $N^{1-2\beta+4/A_0+o(1)}=N^{-41/50+o(1)}$.

\begin{proof}[Proof of Theorem~\ref{p1:core:primitive}]
Use $s=5000$, so the nonmoderate tail in
\eqref{p1:core:primitive-tail} is $O_w(N^{-1})$.
For each of the finitely many sizes $3\le r<s$, the fixed-order
primitive estimate in Proposition~\ref{p1:fixed:primitive} gives
$O_{r,\varepsilon}(N^{-4/5+\varepsilon})$
for the positive mass of separated families. Compatible contacts and
their fixed-order extensions contribute $O_{r,\varepsilon}(N^{-1+\varepsilon})$.
Lemma~\ref{p1:core:fixed-centred} controls $|c(C)|$,
including contacts and incompatibility. Alternatively, for a primitive
one can use $|c(C)|\le q(C)+p^r$ and the fixed-order support count.
Replacing its fixed weight by $w^r$ costs only a constant on this
finite set of sizes. Add these estimates to the two large-size bounds.
No fixed-order estimate has been applied with $r=r(N)$.
\end{proof}

\section{One distinguished pair after compression}
\label{p1:core:compression-section}

The approximation will keep the contribution of one good pair.
To control additional relations above it, the pair must first be
treated as a single affine block. Centring this new block is different
from centring its two original indicators; we retain the distinction
until the final identity of this section.

For an affine-compatible good pair $e=\{x,y\}$ of channel
$(a,b)=(2,3)$ or $(3,4)$,
let $J_e=J_xJ_y$, $q_e=\E J_e$ and $V_e=S_x+S_y$.
Its exact compression has centre $X_e=bx$, offset
$h_e=ay-bx=O(B)$, and values
\[
 U_e=b\mathcal V_x\cup a\mathcal V_y,
 \qquad \mathcal V_x=\{x-1,x,\ldots,x+L-1\}.
\]
Thus $U_e$ is a set of integer values, while $V_e$ is a vector space.
Lemma~\ref{p1:arith:compression} gives an
injective even-selection representation on the distinct values of
$U_e$, width $O(B)$, at most $2B$ values, and
\begin{equation}
 r_e=\dim V_e\ge\kappa_bL-O(1),\qquad
 q_e=2^{-r_e}\ll N^{-\kappa_b},\qquad \kappa_b=2-1/b.
 \label{p1:core:compressed-rank}
\end{equation}
There are $B^{O(1)}$ possible motifs; for a fixed motif its centre
determines the pair subject to the integral congruences. The affine
event itself is preserved, without an additional random sign from the
multipliers.

For a set $A$ of original windows disjoint from $e$, put
\[
 q(e,A)=\E\left[J_e\prod_{z\in A}J_z\right],\qquad
 \widehat c(e,A)=\E\left[(J_e-q_e)\prod_{z\in A}(J_z-p)\right].
\]
The configuration is primitive after compression when the spaces
$V_e,(S_z)_{z\in A}$ are primitive. Its number of compressed blocks
is $r=|A|+1$ and its number of original windows is $r+1$.

\begin{theorem}[A distinguished compressed pair at all cardinalities]
\label{p1:core:compressed}
For fixed $w>0$, let
\[
 \mathcal P_b(w)=\sum_{\substack{(e,A)\ \mathrm{primitive\ after\ compression}\\
                  e\ \mathrm{in\ channel}\ a:b,\ |A|\ge1}}
 w^{|A|+2}\{q(e,A)+|\widehat c(e,A)|+|c(e\cup A)|\}.
\]
Then, for each $\varepsilon>0$,
\begin{align}
 \mathcal P_3^{\mathrm{sep}}(w)&\ll_{w,\varepsilon}N^{-4/3+\varepsilon},&
 \mathcal P_4^{\mathrm{sep}}(w)&\ll_{w,\varepsilon}N^{-5/4+\varepsilon},
 \label{p1:core:compressed-separated}\\
 \mathcal P_b(w)&\ll_{w,\varepsilon}N^{-1+\varepsilon},&&b=3,4.
 \label{p1:core:compressed-all}
\end{align}
The superscript restricts all original windows to be separated. The
pair remains distinguished in each sum.
\end{theorem}

\begin{proof}
Projection onto any original singleton block is injective on the
primitive relation space. Thus its nullity is at most $L$, and
\begin{equation}
 \widehat c(e,A)=q(e,A)-q_ep^{r-1},\qquad
 q(e,A),\ |\widehat c(e,A)|\le q_ep^{r-2}
       \ll N^{2-\kappa_b}p^r.
 \label{p1:core:compressed-cost}
\end{equation}
For an incompatible configuration, the centred value is exactly
$-q_ep^{r-1}$. These estimates retain the saving from the pair rank.

Use the parameters \eqref{p1:core:primitive-parameters}. The two-anchor
certificate works on the compressed block as well: all its distinct
values lie in an interval of width $O(B)$ and are $O(N)$, so their
large kernels are coprime. Replace the moderate threshold by $C_2NB$
for a sufficiently large fixed constant, and put
$m_*=\lceil2\log(C_5N)/\log Y\rceil$, where $C_5$ bounds all
compressed values divided by $N$. With $P_N=B^{O(1)}$ possible
pair motifs, the ordered moderate host count is at most
\[
 P_N C_3^r r^2B^4N^rT^{-2}(C_4rB)^{m_*}.
\]
There is only one factor $P_N$, since only one compressed block is
present. Dividing by $r!$, using \eqref{p1:core:compressed-cost} and
summing the exponential series as before gives
\begin{equation}
 N^{2-\kappa_b-2\beta+4/A_0+o(1)}
       =N^{-\kappa_b+9/50+o(1)}
 \label{p1:core:compressed-moderate}
\end{equation}
for each of the first two masses. The factors $P_N$ and $C_4^{m_*}$
are $N^{o(1)}$.

For the nonmoderate sector use a typed graph whose vertices are all
original windows and all possible compressed good-pair blocks of the
fixed channel. It has $n_*=N^{1+o(1)}$ vertices,
$M_*=N^{1-\delta+o(1)}$ marked vertices, and maximum saturation degree
$\Delta_*=N^{\theta+o(1)}$. In the incoming and outgoing degree
counts a fixed motif adds only $B^{O(1)}$ choices. Every unmarked
vertex of a primitive configuration of at least three blocks has a
saturated neighbour. Lemma~\ref{p1:core:generator}, with
$z=wpN^{\delta/2}$, is still bounded because
$M_*z$, $\Delta_*z$, and $n_*\Delta_*z^2$ tend to zero with strict
power margins. We may enlarge the generating function to all typed
families, then retain \eqref{p1:core:compressed-cost} for the
configurations with exactly one distinguished pair. For each fixed
$s\ge3$, this bounds their tail by
\[
 O_w(N^{2-\kappa_b-s\delta/2}).
\]
At $s=5000$ this is $O_w(N^{-\kappa_b})$.

It remains to treat finitely many $2\le r<5000$. For separated
configurations with fixed $r\ge3$,
Proposition~\ref{p1:fixed:compressed} bounds their homogeneous mass by
$O_{r,\varepsilon}(N^{-\kappa_b+1/5+\varepsilon})$.
This bounds both $q$ and $|\widehat c|$, since the latter is at most
the homogeneous mass even when the full affine configuration is
incompatible. For $r=2$, the pair-and-target estimate of
Proposition~\ref{p1:arith:fixed-homogeneous-seeds} gives respectively
$O_\varepsilon(N^{-4/3+\varepsilon})$ and
$O_\varepsilon(N^{-5/4+\varepsilon})$ for the same homogeneous
majorant. Thus no phase restriction is needed in the separated estimate.

For the contact sector, Proposition~\ref{p1:fixed:propagation} bounds the
mass of the $r+1$ original windows by $O_{r,\varepsilon}(N^{-1+\varepsilon})$.
The number of possible distinguished pairs in such a fixed-size
family is bounded in terms of $r$. The deterministic host count for
these contacts is $N^{r-1+o(1)}$: the good pair costs $NB^{O(1)}$,
and a contact removes one remaining free centre. Multiplying by the
centering term $q_ep^{r-1}$ gives $N^{-\kappa_b+o(1)}$.
The identity in \eqref{p1:core:compressed-cost} therefore bounds
$|\widehat c|$ on contacts, including incompatible configurations.
Together with the two tails, these bounds prove
\eqref{p1:core:compressed-separated}--\eqref{p1:core:compressed-all}
for $q$ and $|\widehat c|$.

Finally, $A$ is a nonempty direct family, and expansion of
$(J_x-p)(J_y-p)$ gives the exact identity
\begin{equation}
 c(e\cup A)=\widehat c(e,A)
                -p\,c(\{x\}\cup A)-p\,c(\{y\}\cup A).
 \label{p1:core:centering-change}
\end{equation}
Each endpoint family is direct or primitive: a proper subset
containing the endpoint is contained in a proper direct family of
compressed blocks. A fixed original primitive family of size $r$
can be extended by a distinguished companion in at most $O(rB)$ ways,
since the deterministic good-pair degree is $O(B)$. Hence the sum of
the two corrections is bounded by a fixed multiple of
\[
 wpB\sum_{r\ge2}r w^r
             \sum_{\substack{D\ \mathrm{primitive}\\|D|=r}}|c(D)|.
\]
For $r\ge3$, absorb $r$ into a slightly larger fixed weight and use
Theorem~\ref{p1:core:primitive}; for $r=2$, use the total pair activity
$O(LN^{-2/3})$. The result is
$O_{w,\varepsilon}(pB(LN^{-2/3}+N^{-4/5+\varepsilon}))
 =O_{w,\varepsilon}(N^{-5/3+\varepsilon})$.
This is smaller than each claimed bound and proves the assertion for
the original centred coefficient as well.
\end{proof}

\subsection{The active structure outside a small event}

\begin{lemma}[A single pair is the only remaining active component]
\label{p1:core:exceptional}
There exists an event $E_*$ with
$\P(E_*)\ll_\varepsilon N^{-4/5+\varepsilon}$ such that, outside
$E_*$, all active spaces are a direct sum of singleton spaces and at
most one compressed good-pair space. In particular every active
subfamily has nullity at most $L/3+O(1)$.
\end{lemma}
\begin{proof}
Exclude an active contact, an active nongood dependent pair, an active
primitive family of at least three original windows, two distinct
active good pairs, and an active primitive family after compression of
a good pair. The pair and contact estimates, the two-pair seed bound in
Proposition~\ref{p1:arith:fixed-homogeneous-seeds}, and
Theorems~\ref{p1:core:primitive} and~\ref{p1:core:compressed}
bound the union by the asserted probability. All sums here are positive.

If no good pair is active, a minimal dependent active subfamily would
be one of the excluded families. If exactly one good pair is active,
compress it. A minimal dependent subfamily of the resulting blocks
either avoids the compressed block, giving an excluded original
primitive, or contains it, giving an excluded compressed primitive.
Thus the remaining compressed blocks are direct. The only rank loss
is that of the good pair, at most $L/3+O(1)$.
\end{proof}

\section{Connected coefficients with bounded compatible nullity}
\label{p1:core:nullity-section}

The preceding theorems concern minimal dependent families. The
boundary of the stopped expansion can contain larger connected
families. We first control all cardinalities when the compatible
nullity is less than two window ranks. We then remove the restriction
on that fixed nullity level over the growing range needed by the
expansion.

\subsection{A two-anchor estimate at every cardinality}

\begin{theorem}[All cardinalities below two ranks]
\label{p1:core:low-nullity}
For fixed $0\le\eta<2$, $w>0$ and $\varepsilon>0$,
\begin{equation}
 \sum_{\substack{C\subseteq G_N\ \mathrm{connected},\ |C|\ge3\\
                  \rho_*(C)\le\eta L}}
 w^{|C|}|c(C)|
 \ll_{\eta,w,\varepsilon}
 N^{-\min\{4/5,\,2-\eta\}+\varepsilon}.
 \label{p1:core:low-nullity-bound}
\end{equation}
\end{theorem}
\begin{proof}
Expansion of the centred product gives
\begin{equation}
 |c(C)|\le 2^{|C|}p^{|C|}2^{\rho_*(C)}
       \ll_\eta N^\eta(2p)^{|C|}.
 \label{p1:core:compatible-majorant}
\end{equation}
The Fourier expansion also shows that $c(C)\ne0$ implies a relation
meeting every block: centring removes the zero Fourier coordinate in
each factor. Such a relation again selects at least two values per
block. Therefore a family with nonzero coefficient and no moderate
certificate has only marked isolated vertices in the saturation graph.
No independence assertion about the edges of that graph is used.

With the parameters in \eqref{p1:core:parameter-conditions}, the
moderate sector, by Lemma~\ref{p1:core:moderate}, costs
$N^{\eta-2\beta+4/A_0+o(1)}$. The graph generator bounds its
nonmoderate tail beyond any fixed size $s$ by
$O_{\eta,w}(N^{\eta-s\delta/2})$.
For $0<\varepsilon<\min\{1,2-\eta\}$ choose
\[
 A_0>\max\{40,32/\varepsilon\},\qquad
 \beta=1-\varepsilon/16,\qquad
 \sigma=1-1/(2A_0).
\]
Then $\delta=\varepsilon/(32A_0)$, $2/A_0+\delta<1$, and
$2-2\beta+4/A_0<\varepsilon/4$. Choose a fixed integer $s$ with
$s\delta/2\ge\eta+1$. The moderate part has the required
$2-\eta$ exponent, the nonmoderate tail is $O(N^{-1})$, and
the finitely many connected sizes $3\le r<s$ are handled by
Lemma~\ref{p1:core:fixed-centred}.
Shrinking the margins there gives \eqref{p1:core:low-nullity-bound}.
Larger values of $\varepsilon$ follow from smaller ones.
\end{proof}

At $\eta=1$, this already gives the $4/5$ exponent at all sizes.
At $\eta=3/2$, however, the two-anchor estimate would give only
the exponent $1/2$. We need more anchors where a block carries many
independent relation coordinates.

\subsection{Projection dimensions and positive cores}

For a family $D$, let $\mathcal K(D)$ be the kernel of the sum map
$\bigoplus_{x\in D}S_x\to\sum_{x\in D}S_x$. The dimension of its
projection onto $S_x$ is
\begin{equation}
 d_x(D)=\rho(D)-\rho(D\setminus\{x\})
       =\dim\left(S_x\cap\sum_{y\in D\setminus\{x\}}S_y\right).
 \label{p1:core:projection-dimension}
\end{equation}
The kernel of that projection is exactly
$\mathcal K(D\setminus\{x\})$, which proves the identity.

If $t$ is even, a subspace of dimension at least $t-1$ in the
even-selection space contains a vector of weight at least $t$.
Indeed, choose $t-1$ pivot coordinates and set all of them equal to
one. The resulting vector has weight at least $t-1$, and its weight
is even. Call $D$ a \emph{$t$-core} when $d_x(D)\ge t-1$ for every
$x\in D$. At each of its blocks a relation therefore selects at least
$t$ distinct values. The selected relation may depend on the block.

\begin{proposition}[Positive mass of high-nullity cores]
\label{p1:core:positive-core}
Fix $\eta>1$, an even integer $t\ge\eta+2$, and $v>0$. Then
\begin{equation}
 \sum_{\substack{D\ \mathrm{compatible},\ D\ t\mathrm{-core}\\
                  L/2<\rho(D)\le\eta L}}
 v^{|D|}q(D)\ll_{\eta,t,v,\varepsilon}N^{-4/5+\varepsilon}
 \label{p1:core:positive-core-bound}
\end{equation}
for every $\varepsilon>0$, with no restriction on the cardinality.
\end{proposition}
\begin{proof}
We give the multiple-anchor count explicitly. Set
\[
\begin{gathered}
 A_0=8t,\qquad Y=(\log N)^{A_0},\qquad
 \beta=1-\frac1{4t},\qquad T=N^\beta,\\
 \sigma=1-\frac1{16t},\qquad\delta=\frac1{64t^2},\qquad
 m_t=\left\lceil\frac{t\log(3N)}{\log Y}\right\rceil.
\end{gathered}
\]
The number of marked windows remains $N^{1-\delta+o(1)}$.
At an unmarked anchor choose $t$ selected values. Their kernel
product $Q$ is squarefree, with $T^t<Q\le(3N)^t$ and at most $m_t$
prime factors. A moderate allocation has $D_j\le6NB$ for each
receiver. At most $t$ received modules exceed $N$, so its CRT cost is
\[
 \prod_{j\ne\mathrm{anchor}}(1+N/D_j)
 \le2^{r-1}(6B)^tN^{r-1}T^{-t}.
\]
The corresponding ordered host count is at most
$6^t rB^{2t}2^{r-1}N^rT^{-t}(rB)^{m_t}$.
Multiplication by $(up)^r/r!$ and the same exponential-series
estimate as before give, for fixed $u>0$,
\[
 \sum_r\frac{(up)^r}{r!}|\mathcal U_{r,t}|
 \le N^{-t\beta+2t/A_0+o(1)}.
\]
Here $\log(B^{m_t+2t}(m_t+1)!)=(2t/A_0+o(1))\log N$.
Since $q(D)\ll_\eta N^\eta p^{|D|}$, the positive moderate mass is
\begin{equation}
 N^{\eta-t\beta+2t/A_0+o(1)}
       =N^{\eta-t+1/2+o(1)}\le N^{-3/2+o(1)}.
 \label{p1:core:multi-anchor-gain}
\end{equation}

For the remaining cores, define saturation using $t$ anchor values.
The outgoing degree is bounded by $B^t(B+1)^{m_t}$. For the incoming
degree, the fixed receiver has at most
$R=\lceil B\log(3N)/\log Y\rceil$ possible primes; the codes are
bounded by
$B^t\sum_{j\le m_t}\binom Rj t^j$.
As in Lemma~\ref{p1:core:degrees}, a qualifying code has modulus
greater than $N$ and determines at most one centre. Both degrees are
$N^{t/A_0+o(1)}=N^{1/8+o(1)}$.
Every unmarked vertex of a nonmoderate $t$-core has a saturated
neighbour. Lemma~\ref{p1:core:generator}, tilted by $N^{\delta/2}$,
therefore bounds the positive mass beyond fixed size $s$ by
$O_{\eta,v}(N^{\eta-s\delta/2})$.
Choose $s\ge128t^2(\eta+1)$, so this is $O(N^{-1})$.

For a fixed size $h<s$, compatibility and $\rho(D)>L/2$ imply that
activation of $D$ occurs inside the event $E_*$ of
Lemma~\ref{p1:core:exceptional}. Indeed, its complement permits
nullity at most $L/3+O(1)$. For every fixed integer $k>1$, H\"older
and the bounded fixed moments of Theorem~\ref{p1:fixed:moments} give
\[
 \sum_{\substack{|D|=h,\ D\ \mathrm{compatible}\\\rho(D)>L/2}}q(D)
 \le\E\left[\binom{Z_N^\circ}{h};E_*\right]
 \le C_{h,k}\P(E_*)^{1-1/k}.
\]
Choose $k$ large enough that $4/(5k)<\varepsilon/3$, then use a
sufficiently small fixed margin in the bound for $\P(E_*)$.
The result is $O_{h,\varepsilon}(N^{-4/5+\varepsilon})$.
Only finitely many sizes $h$ are involved. Adding these to
\eqref{p1:core:multi-anchor-gain} and the saturated tail proves
\eqref{p1:core:positive-core-bound}. The distributional approximation
being proved in this paper is not used in this argument.
\end{proof}

\subsection{Removing blocks with small relation projections}

\begin{theorem}[Bounded compatible nullity on the growing range]
\label{p1:core:bounded-nullity}
For every fixed $\eta\ge0$, $w>0$, $\varepsilon>0$ and every
$K_N=o(L)$,
\begin{equation}
 \sum_{\substack{C\subseteq G_N\ \mathrm{connected}\\
          3\le|C|\le K_N,\ \rho_*(C)\le\eta L}}
 w^{|C|}|c(C)|\ll_{\eta,w,\varepsilon}N^{-4/5+\varepsilon}.
 \label{p1:core:bounded-nullity-bound}
\end{equation}
For $\eta>1$, choose a fixed even $t\ge\eta+2$ and set $d=t-2$.
The estimate holds uniformly whenever $K_N\le L/(2d)$. In particular,
for $\eta=3/2$ one may take $t=4$, $d=2$ and $K_N\le L/4$.
\end{theorem}
\begin{proof}
Theorem~\ref{p1:core:low-nullity} handles $\rho_*(C)\le L$.
For $L<\rho_*(C)\le\eta L$, choose a compatible subfamily
$A\subseteq C$ attaining $\rho(A)=\rho_*(C)$, using a fixed
lexicographic rule if needed. Remove from $A$ any block whose
relation projection has dimension at most $d$, until none remains.
Let $D$ be the resulting family. Equation~\eqref{p1:core:projection-dimension}
shows that
\[
 0\le\rho(A)-\rho(D)\le d(|A|-|D|)\le dK_N\le L/2.
\]
Thus $D$ is a nonempty compatible $t$-core with
$L/2<\rho(D)\le\eta L$.
If $r=|C|$ and $h=|D|$, the centred expansion and the rank bound give
\begin{equation}
 |c(C)|\le2^r p^r2^{\rho_*(C)}
        \le2^r(2^dp)^{r-h}q(D).
 \label{p1:core:peeling-weight}
\end{equation}
The elements of $C\setminus A$ are included in $p^{r-h}$; they need
not have been removed by the algorithm. This also covers incompatible
families $C$.

Sum \eqref{p1:core:peeling-weight} over all admissible cores and all
their supersets. Dropping restrictions on the latter only increases
the sum. With $n=|G_N|$ and $np\le Np=\lambda$, this gives
\begin{align*}
 \sum_{\substack{|C|\le K_N\\L<\rho_*(C)\le\eta L}}
 w^{|C|}|c(C)|
 &\le\sum_{D\ \mathrm{admissible}}(2w)^{|D|}q(D)
        (1+2^{d+1}wp)^{n-|D|}\\
 &\le e^{2^{d+1}w\lambda}
        \sum_{D\ \mathrm{admissible}}(2w)^{|D|}q(D).
\end{align*}
The last sum is bounded by Proposition~\ref{p1:core:positive-core}.
There is no factor for the deletion order: a deterministic choice
assigns one core to each family, and summing over all supersets already
overcounts the original family. The exponential factor is constant
because $d,w$ are fixed and $\lambda$ is bounded. Adding the sector
$\rho_*(C)\le L$ proves the theorem.
\end{proof}

For the required value $\eta=3/2$, the choices in the positive-core
proof are
$A_0=32$, $\beta=15/16$, $\sigma=63/64$, $\delta=1/1024$ and
$s=5120$. Its moderate mass is $N^{-2+o(1)}$, its saturated tail is
$O(N^{-1})$, and only fixed sizes below $5120$ use fixed moments.
This supplies the whole range containing
$K_N\asymp\log N/\log\log N$ without promoting a fixed-order
estimate to a growing-order one.

The same deletion argument gives the positive corollary
\[
 \sum_{\substack{D\ \mathrm{compatible},\ |D|\le K_N\\
                    L<\rho(D)\le\eta L}}
 w^{|D|}q(D)\ll_{\eta,w,\varepsilon}N^{-4/5+\varepsilon}.
\]
The centred form \eqref{p1:core:bounded-nullity-bound}, rather than
this additional corollary, is the estimate used by the stopped
expansion below.

\section{Recovering the law from the retained pair data}
\label{p1:close:section}

The stopped expansion requires a positive sum of joint
probabilities.  The component theorem controls centered
coefficients instead.  The two quantities cannot be interchanged
by taking absolute values of the entire expansion: a single good
pair has mass of the same order as the correction we wish to
retain.  We separate that pair explicitly, estimate its failed
extensions after compression, and use the component bound only
for the remaining terms.

Throughout this section, $\mathscr D$ is the class from
Section~\ref{p14:comp:stopped}, $n=|G|$, $\mu=np$, and
\[
 \beta_N=\sum_{e\in E_N}q_e.
\]
All asymptotic constants may depend on the fixed critical band.

\subsection{Two deficits under independent counting weights}

For $r\ge0$ define
\begin{align}
 F_r&=p^r\#\{I\subseteq G:|I|=r,\ I\text{ is not direct}\},
       \label{p1:close:deficit-F}\\
 G_r&=\sum_{e\in E_N}q_ep^{r-2}
      \left\{\binom{n-2}{r-2}-\#\Gamma_e(r-2)\right\}
                   \qquad(r\ge2),
       \label{p1:close:deficit-G}
\end{align}
and put $G_0=G_1=0$.  We use zero for binomial coefficients
whose lower index lies outside the usual range.  In particular,
both sequences vanish for $r>n$, and $F_0=F_1=0$.
The weighted deficits are
\begin{equation}
 F=\sum_{r\ge0}2^rF_r,\qquad
 \mathfrak G=\sum_{r\ge0}2^rG_r.
 \label{p1:close:total-deficits}
\end{equation}
These are deterministic counting deficits.  The added centers
carry independent weights $p$, even when the corresponding
arithmetic events are dependent.

\begin{lemma}[Bounds for the independent deficits]
\label{p1:close:deficit-bounds}
For every fixed $\varepsilon>0$,
\begin{equation}
 F\ll_\varepsilon N^{-4/5+\varepsilon}+B/N,
 \qquad
 \mathfrak G\ll_\varepsilon N^{-1+\varepsilon}.
 \label{p1:close:deficit-rates}
\end{equation}
\end{lemma}

\begin{proof}
Let $\mathcal P$ denote the minimal non-direct families,
including pairs, and set
\[
 S_0=\sum_{C\in\mathcal P}(2p)^{|C|}.
\]
Every non-direct family contains some $C\in\mathcal P$.
After choosing it, sum all added centers with weight $2p$.
The resulting union bound gives
\begin{equation}
 F\le(1+2p)^n S_0\le e^{2\mu}S_0.
 \label{p1:close:ordinary-union}
\end{equation}
For a minimal non-direct family of size at least three,
every proper subfamily is direct.  Expanding the centered
product therefore gives
$c(C)=q(C)-p^{|C|}$, and hence
\[
 p^{|C|}\le q(C)+|c(C)|.
\]
Theorem~\ref{p1:core:primitive}, with fixed weight $2$, bounds
the contribution of these families by
$O_\varepsilon(N^{-4/5+\varepsilon})$.
For pairs, the deterministic host bounds of
Theorem~\ref{p1:arith:pairs}, including the $O(NB)$ contacts
and good-pair hosts, give
\[
 \sum_{\substack{C\in\mathcal P\\|C|=2}}(2p)^2
    \ll_\varepsilon N^{-4/5+\varepsilon}+B/N.
\]
These host counts include incompatible pairs; their independent
weights are still $p^2$.  This proves the bound for $S_0$ and $F$.

Fix now $e\in E_N$ and a family of singleton blocks which is not
direct after adjoining $V_e$.  A minimal non-direct subfamily
either lies entirely among the singletons or contains $V_e$.
The first possibility contains an original member of $\mathcal P$.
Write $(e,C)$ for a minimal non-direct family of the second
kind, where $\varnothing\ne C\subseteq G\setminus e$.
Summing additional centers with their independent weights gives
\begin{equation}
 \mathfrak G
 \le4e^{2\mu}\left\{
       \beta_N S_0+
       \sum_{(e,C)\ {\rm minimal}}
                 2^{|C|}q_ep^{|C|}
             \right\}.
 \label{p1:close:compressed-union}
\end{equation}
The factor $4$ accounts for the two original endpoints of $e$.

To use the compressed estimate, regard
$J_e=J_xJ_y$ as one indicator of mean $q_e$ and set
\[
 \widehat c(e,C)
   =\E\left[(J_e-q_e)\prod_{z\in C}(J_z-p)\right].
\]
For a minimal non-direct family of these blocks, every proper
block subfamily is direct.  The same centered expansion yields
\[
 \widehat c(e,C)=q(e\cup C)-q_ep^{|C|},
 \qquad
 q_ep^{|C|}\le q(e\cup C)+|\widehat c(e,C)|.
\]
Theorem~\ref{p1:core:compressed}, with the pair distinguished,
bounds the second sum in \eqref{p1:close:compressed-union}
by $O_\varepsilon(N^{-1+\varepsilon})$.
The factor $4\,2^{|C|}=2^{|C|+2}$ is exactly the weight of
the original windows; it introduces no loss depending on
the size.

Finally, $\beta_N=O(a_N)$ and
$a_N\asymp LN^{-2/3}$.  The already obtained bound for $S_0$
therefore gives $\beta_NS_0=O_\varepsilon(N^{-1+\varepsilon})$,
after choosing a smaller fixed error margin in the input bounds
if necessary.  Equation \eqref{p1:close:compressed-union}
proves the second assertion.
\end{proof}

The extensions in this proof are summed under the independent
weights.  In particular, no identity
$q(C\cup U)=q(C)p^{|U|}$ has been used for a family which is
not in direct sum.

\subsection{The exact measure retained by the hereditary class}

Let $b=\delta_0+p\Delta$, the Bernoulli law of mean $p$, and
define
\begin{equation}
 \mathsf B_\beta=b^{*n}
          +\beta_N\Delta^{*2}*b^{*(n-2)}.
 \label{p1:close:bernoulli-reference}
\end{equation}
At this stage it is convenient to work with a signed correction
to the binomial law.

\begin{proposition}[Evaluation of the retained polynomial]
\label{p1:close:retained-reference}
For the class $\mathscr D$,
\begin{equation}
 T_{\mathscr D}
 =\mathsf B_\beta-\sum_{r\ge0}(F_r+G_r)\Delta^{*r},
 \qquad
 \|T_{\mathscr D}-\mathsf B_\beta\|_1\le F+\mathfrak G.
 \label{p1:close:retained-identity}
\end{equation}
\end{proposition}

\begin{proof}
The total joint probability of direct families of size $r$ is
$\binom nrp^r-F_r$.
By Lemma~\ref{p1:stop:class}, every non-direct member of
$\mathscr D$ contains a unique good pair.  Equations
\eqref{p1:stop:gamma} and \eqref{p1:stop:direct-pair-probability}
show that their total contribution is
\[
 \sum_{e\in E_N}q_ep^{r-2}\#\Gamma_e(r-2)
   =\beta_N\binom{n-2}{r-2}p^{r-2}-G_r.
\]
This term is zero for $r<2$.  Summing against
$\Delta^{*r}$ and applying the binomial formula gives the
identity.  The norm bound follows from
$F_r,G_r\ge0$ and $\|\Delta^{*r}\|_1=2^r$.
\end{proof}

The coefficient $\beta_N$ is the sum of the joint probabilities
of good pairs.  It will become the covariance correction after
the independent contribution of those pairs and the binomial
diagonal have been accounted for.

\subsection{From centered components to positive boundary mass}

For $K\le L/4$ put
\[
 \mathscr Q=\{I\subseteq G:\rho(I)\le3L/2\},
 \qquad
 R_K=\sum_{\substack{I\in\mathscr Q\setminus\mathscr D\\|I|\le K}}
          2^{|I|}q(I).
\]
The boundary of size at most $K$ is contained in this sum.
Define the component activities
\begin{align}
 A_{\rm good}&=4\sum_{e\in E_N}c(e),\notag\\
 A_{\rm bad}&=4\sum_{\substack{I\subseteq G,\ |I|=2\\I\notin E_N}}
                                      |c(I)|,\notag\\
 B_K&=\sum_{\substack{C\ {\rm connected},\ 3\le|C|\le K\\
                                  \rho_*(C)\le3L/2}}
                              2^{|C|}|c(C)|.
 \label{p1:close:activities}
\end{align}
For a good pair, compatibility and positive nullity give
$c(e)=q_e-p^2>0$.  In particular,
$0\le A_{\rm good}\le4\beta_N=O(a_N)$.

\begin{proposition}[Positive mass outside the retained class]
\label{p1:close:boundary-mass}
For every $K\le L/4$,
\begin{equation}
 R_K\le F+\mathfrak G+
 e^{2\mu}\bigl\{
          \exp(A_{\rm good}+A_{\rm bad}+B_K)-1-A_{\rm good}
          \bigr\}.
 \label{p1:close:boundary-inequality}
\end{equation}
Consequently, uniformly in that range,
\begin{equation}
 R_K\ll_\varepsilon
        N^{-4/5+\varepsilon}+B/N+a_N^2.
 \label{p1:close:boundary-rate}
\end{equation}
\end{proposition}

\begin{proof}
Expand each indicator about its mean:
\begin{equation}
 q(I)=\sum_{J\subseteq I}p^{|I|-|J|}c(J),
 \qquad c(\varnothing)=1.
 \label{p1:close:centered-expansion}
\end{equation}
If $I\in\mathscr Q$, then every $J\subseteq I$ and every
canonical component $C$ of $J$ satisfy
$\rho_*(C)\le\rho(C)\le\rho(I)\le3L/2$.
The first inequality follows from the definition of compatible
nullity, and the second from monotonicity of nullity.

The term $J=\varnothing$, summed over
$I\in\mathscr Q\setminus\mathscr D$, contributes at most $F$:
all such $I$ are non-direct.

For a term $J=e\in E_N$, write $I=e\cup U$.
Since $I\notin\mathscr D$, the family $U$ does not belong to
$\Gamma_e(|U|)$.  Using $0<c(e)\le q_e$ and then dropping
the restrictions on $I$ gives the upper bound
\[
 4\sum_{e\in E_N}q_e
      \sum_{\substack{U\subseteq G\setminus e\\
                          U\notin\Gamma_e(|U|)}}(2p)^{|U|}
       =\mathfrak G.
\]
This is the only part of the argument where a term of order
$a_N$ has to be removed before taking an absolute majorant.

It remains to sum nonempty $J$ which are not a single good pair.
The coefficient $c(J)$ factors over its canonical direct-sum
components.  A singleton component has zero coefficient.
Every remaining nonzero family is therefore a collection of
components of size at least two, containing either a bad pair,
a component of size at least three, or at least two good pairs.
Its component activities are included in
\eqref{p1:close:activities}.

The sum over these collections is bounded by
\begin{equation}
 \sum_{\text{remaining }J}2^{|J|}|c(J)|
 \le \exp(A_{\rm good}+A_{\rm bad}+B_K)-1-A_{\rm good}.
 \label{p1:close:component-exponential}
\end{equation}
Indeed, a collection of $m$ distinct components has $m!$
orders, cancelling the factorial in the exponential series.
Allowing repeated, intersecting or non-independent components
only adds nonnegative terms to the majorant.  The empty
collection and a single good-pair component are precisely
the two terms subtracted on the right.

For each fixed $J$, summing the additional members of
$I\setminus J$ in \eqref{p1:close:centered-expansion}
costs at most
$(1+2p)^{n-|J|}\le e^{2\mu}$.  This proves
\eqref{p1:close:boundary-inequality}.

Theorem~\ref{p1:arith:pairs} and
Theorem~\ref{p1:core:bounded-nullity}, with fixed
$\eta=3/2$ and $w=2$, give
\[
 A_{\rm bad}+B_K
      \ll_\varepsilon N^{-4/5+\varepsilon}+B/N.
\]
Use a small fixed error margin so that this quantity tends to
zero.  For $u,v\ge0$ tending to zero,
$e^{u+v}-1-u=O(u^2+v)$.
The bounds for $A_{\rm good}$, $F$ and $\mathfrak G$
now give \eqref{p1:close:boundary-rate}.
An estimate for a smaller error margin implies the displayed
estimate for every prescribed $\varepsilon>0$.
\end{proof}

\subsection{Converting the binomial reference to Poisson}

Set
\begin{equation}
 C_2=\sum_{\substack{I\subseteq G\\|I|=2}}c(I),
 \qquad
 \alpha_G=\operatorname{Var}(Z_G)-\mu=2C_2-np^2.
 \label{p1:close:variance}
\end{equation}
The second identity is the variance formula for a sum of
indicators of common mean $p$.  Moreover,
\begin{equation}
 \beta_N-C_2
       =|E_N|p^2-\sum_{\substack{|I|=2\\I\notin E_N}}c(I),
 \qquad
 |\beta_N-C_2|
       \le |E_N|p^2+\frac14A_{\rm bad}.
 \label{p1:close:beta-covariance}
\end{equation}
Thus this difference is
$O_\varepsilon(N^{-4/5+\varepsilon}+B/N)$.

\begin{lemma}[Poisson--Charlier conversion]
\label{p1:close:conversion}
With the parameters of the critical band,
\begin{equation}
 \|\mathsf B_\beta-Q_{\mu,\alpha_G}\|_1
 \le 4|\beta_N-C_2|+
           O_C\bigl(N^{-2}+a_N/N\bigr).
 \label{p1:close:conversion-bound}
\end{equation}
\end{lemma}

\begin{proof}
Use the algebraic form of Lemma~\ref{p13:transfer:conversion}.
For large $N$, the regular-set count gives $n\ge2$.
Take common marginal $p$ and put coefficient $C_2$ on one fixed
pair, all other pair coefficients being zero. The reference is
$b^{*n}+C_2\Delta^{*2}*b^{*(n-2)}$, its absolute pair mass is
$|C_2|$, and its quadratic parameter is $2C_2-np^2=\alpha_G$.
For zero pair coefficients, the same estimate gives
\begin{equation}
 b^{*n}
 =\Pi_\mu-\frac{np^2}{2}\Pi_\mu*\Delta^{*2}
           +O_{\ell^1}\bigl(np^3+(np^2)^2\bigr).
 \label{p1:close:binomial-expansion}
\end{equation}
For the chosen coefficient, it gives an error
$O_C(np^3+(np^2)^2+|C_2|(np^2+p))$.
Replacing $C_2$ by $\beta_N$ costs at most $4|\beta_N-C_2|$,
since $b^{*(n-2)}$ has norm one. Now $n=O(N)$,
$p=O_C(N^{-1})$ and $C_2=O_C(a_N)$ give
\eqref{p1:close:conversion-bound}.
\end{proof}

\subsection{The first approximation and its uniform boundary}

\begin{proposition}[The first quadratic approximation]
\label{p13:base-approximation}
Uniformly in the critical band, for every \(\varepsilon>0\),
\[
 \|\Law Z_N-Q_{\mu_N,\alpha_N}\|_1
       \ll_{C,\varepsilon}N^{-4/5+\varepsilon}.
\]
The same bound holds with \(\mu_N\) replaced by \(\lambda_N\).
\end{proposition}

This estimate identifies the leading correction and proves the
boundary estimates used below. The optimal scalar reconstruction in
Section~\ref{p13:transfer:section} keeps exact marginals
on the whole set of admissible windows.

\begin{proof}
Choose $\mathscr D$ and $K_N$ as in
Section~\ref{p14:comp:stopped}.  For all sufficiently large $N$,
$K_N\le L/4$.  The stopped identity, the boundary-nullity bound,
Lemma~\ref{p1:stop:tail} and
Proposition~\ref{p1:close:boundary-mass} give
\[
 \|\Law Z_G-T_{\mathscr D}\|_1
 \le R_{K_N}+O_C(N^{-3/2})
 \ll_\varepsilon N^{-4/5+\varepsilon}
                      +B/N+a_N^2+N^{-3/2}.
\]
Proposition~\ref{p1:close:retained-reference} adds the two
deficits $F+\mathfrak G$.
Lemma~\ref{p1:close:conversion}, together with
\eqref{p1:close:beta-covariance}, then gives
\begin{align}
 \|\Law Z_G-Q_{\mu,\alpha_G}\|_1
 &\ll_\varepsilon N^{-4/5+\varepsilon}
             +B/N+a_N^2+N^{-3/2}+N^{-2}+a_N/N \notag\\
 &\ll_\varepsilon N^{-4/5+\varepsilon},
 \label{p1:close:regular-law}
\end{align}
since $B\asymp\log N$ and
$a_N\asymp LN^{-2/3}$.

It remains to restore the excluded windows and the desired
parameters.  The independent reference bounds
\eqref{p1:close:poisson-lipschitz} and
\eqref{p1:close:reference-lipschitz} will transfer the parameters.

Proposition~\ref{p1:arith:core} gives
\[
 \|\Law Z_N-\Law Z_G\|_1
   +|\mu_N-\mu|+|\alpha_N-\alpha_G|
   +|\mu_N-\lambda_N|
      \ll_\varepsilon N^{-1+\varepsilon}.
\]
Apply \eqref{p1:close:reference-lipschitz} first to transfer
$(\mu,\alpha_G)$ to $(\mu_N,\alpha_N)$, and then to replace
$\mu_N$ by $\lambda_N$ while retaining $\alpha_N$.
Both costs are absorbed in \eqref{p1:close:regular-law}.
This proves the two asserted forms of the approximation,
uniformly in the fixed critical band.

Finally, if $0<\varepsilon<2/15$, then
\[
 \frac{N^{-4/5+\varepsilon}}{a_N}
       =O_C\left(\frac{N^{-2/15+\varepsilon}}{\log N}\right)
       \longrightarrow0.
\]
Thus the corrected error is smaller than the retained pair term.
\end{proof}


\subsection{The total-variation coefficient}
\label{p14:comp:tv}

\begin{proof}[Proof of Corollary~\ref{p1:cor:tv}]
Theorem~\ref{p1:thm:main} and the reverse triangle
inequality first give
\[
 \TV(\Law Z_N,\Pi_{\mu_N})
 =\alpha_N\mathcal C(\mu_N)
            +O_{C,\varepsilon}(N^{-4/3+\varepsilon}).
\]
Use \(|\mu_N-\lambda_N|\ll N^{-1+\varepsilon}\) and
\eqref{p1:close:reference-lipschitz} to obtain
\[
 \Law Z_N-\Pi_{\lambda_N}
   =\frac{\alpha_N}{2}\Pi_{\lambda_N}*\Delta^{*2}
          +O_{\ell^1,C,\varepsilon}(N^{-1+\varepsilon}).
\]
The reverse triangle inequality, followed by the factor
$1/2$ converting the $\ell^1$ norm to total variation,
yields
\begin{equation}
 \TV(\Law Z_N,\Pi_{\lambda_N})
      =\alpha_N\mathcal C(\lambda_N)
          +O_{C,\varepsilon}(N^{-1+\varepsilon}).
 \label{p1:close:tv-exact-parameter}
\end{equation}
The pair asymptotic from
Theorem~\ref{p1:arith:pairs} now gives
\begin{equation}
 \TV(\Law Z_N,\Pi_{\lambda_N})
 \sim \mathcal C(\lambda_N)
       \frac{r_N+2}{18\,2^{r_N/3}}\,
       \lambda_N^{5/3}L N^{-2/3},
 \qquad r_N=L\bmod3.
 \label{p1:close:tv-asymptotic}
\end{equation}
The lower bound in \eqref{p1:close:constant-bounds} is uniform
on the fixed critical band, and the remainder is
$o(LN^{-2/3})$.  This proves the sharp order asserted in the
corollary.  The equivalence is relative to the displayed
oscillating expression; no convergence of its phase or
intensity factors is required.
\end{proof}

Replacing $\alpha_N$ by its leading equivalent proves
\eqref{p1:close:tv-asymptotic}.  It does not by itself preserve
the error $O(N^{-1+\varepsilon})$ in
\eqref{p1:close:tv-exact-parameter}, which uses the exact
second factorial cumulant.

Proposition~\ref{p1:discussion:charlier-norm} gives
$\|\Pi_1*\Delta^{*2}\|_1=3/e$, hence $\mathcal C(1)=3/(4e)$.
Along $N=2^L$ with $3\mid L$, the pair coefficient therefore gives
\begin{equation}
 \TV(\Law Z_N,\Pi_1)
       \sim \frac1{12e}\,L N^{-2/3}.
 \label{p1:close:dyadic-example}
\end{equation}
This is the coefficient in the example announced in the introduction.

The total-variation asymptotic uses the exact-marginal reconstruction
and the explicit pair coefficient. The independent positivity and
reference-continuity estimates do not use the reconstruction.

\section{Norms of Poisson differences}
\label{p14:comp:poisson-differences}

The following calculation proves Proposition~\ref{p1:discussion:charlier-norm}.
It fixes both the signed-difference convention and the factor of two between
$\ell^1$ distance and total variation.

\begin{proof}[Proof of Proposition~\ref{p1:discussion:charlier-norm}]
Write $v_k=(\Pi_\mu*\Delta^{*2})\{k\}$.  Direct calculation gives
\[
 v_k=\pi_\mu(k)-2\pi_\mu(k-1)+\pi_\mu(k-2)
     =\frac{\pi_\mu(k)}{\mu^2}\bigl((k-\mu)^2-k\bigr).
\]
The quadratic has roots $u_-,u_+$ and is strictly negative
exactly at the integers $a,\ldots,b$.  These form a nonempty
interval; integer roots contribute zero and are excluded by
the definitions of $a,b$.  Since $\sum_kv_k=0$,
\[
 \begin{aligned}
 \sum_k|v_k|
   &=-2\sum_{k=a}^b v_k\\
   &=2\bigl\{\pi_\mu(a-1)-\pi_\mu(a-2)
                 +\pi_\mu(b-1)-\pi_\mu(b)\bigr\},
 \end{aligned}
\]
by telescoping.  This proves \eqref{p1:discussion:finite-norm}.
For $\mu=2$, the quadratic is $(k-1)(k-4)$: the negative
masses are $v_2=-e^{-2}$ and $v_3=-2e^{-2}/3$.
Their sum, multiplied by $-2$, gives
\eqref{p1:discussion:lambda-two}.  The same calculation at
$\mu=1$ gives the value $3/e$ used in
\eqref{p1:close:dyadic-example}.

For the general bound, define the normalized Charlier polynomials
\[
 C_r(k;\mu)=\sum_{j=0}^r
       \binom rj(-1)^{r-j}\frac{(k)_j}{\mu^j}.
\]
Then $(\Pi_\mu*\Delta^{*r})\{k\}=\pi_\mu(k)C_r(k;\mu)$.
Their exponential generating function is
\[
 \sum_{r\ge0}C_r(k;\mu)\frac{t^r}{r!}
       =e^{-t}(1+t/\mu)^k.
\]
For $P\sim\Pois(\mu)$, the expectation of the product of
the generating functions with variables $s,t$ is
\[
 \E\!\left[e^{-s-t}(1+s/\mu)^P(1+t/\mu)^P\right]
       =e^{st/\mu}.
\]
Coefficient comparison, justified in a neighborhood of
$(s,t)=(0,0)$ by the Poisson exponential moments, yields
\[
 \E[C_r(P;\mu)C_\ell(P;\mu)]
       =\ind_{\{r=\ell\}}\,\frac{r!}{\mu^r}.
\]
Cauchy--Schwarz now gives the second bound in
\eqref{p1:discussion:general-norm}.  The first follows from
$\|\Pi_\mu\|_1=1$, $\|\Delta\|_1=2$, and submultiplicativity
of the convolution norm.
\end{proof}

\section{Tails, deletion, and limits of positive comparison methods}
\label{p1:limits:section}

The stopped expansion controls a distribution without imposing an
exponential moment bound on the original count.  This distinction is
necessary in the present model.  We first record the exact support and
local structure, then give a nearby count with factorial tails.  A real
arithmetic cylinder shows why these tails cannot simply be transferred
back.  The final parts examine deletion of active pairs and the positive
cost of forcing private prime signs.

\subsection{Exact support and the local rank-one structure}

Write $Z=Z_N$, $p=2^{-L}$ and $R_N=\lceil N/L\rceil$.
For $j\ge1$, $(z)_j=z(z-1)\cdots(z-j+1)$.

\begin{proposition}[Support and factorial lattice]
\label{p1:limits:lattice}
For every $N\ge2$ and $L\ge1$,
\begin{equation}
 0\le Z\le R_N,\qquad
 \E(Z)_j\in j!p^j\mathbb Z_{\ge0}\quad(j\ge1).
 \label{p1:limits:lattice-eq}
\end{equation}
\end{proposition}
\begin{proof}
Two starts at distance less than $L$ cannot both be active: the
predecessor of the later start lies inside the earlier constant run.
Thus the active starts in an interval of $N$ integers are $L$-separated,
which gives the support bound.  For $j$ distinct starts, their $jL$
relative sign prescriptions form an affine binary system.  Its joint
probability is either zero or $2^{-r}$ for an integer $r\le jL$.
It is therefore an integer multiple of $p^j$.  Finally,
$(Z)_j=j!\sum_{|I|=j}\prod_{x\in I}J_x$.
\end{proof}

The local exceptional windows have an exact description. We use the
rank-one conclusion of Lemma~\ref{p1:limits:local-rigidity}; a bound $o(L)$
for the local nullity would not suffice for what follows.



The proof of Lemma~\ref{p1:limits:local-rigidity} also gives raw-value nullity at most two at critical length,
because the even relation code has codimension at most one in the
full relation code.  It also gives exact $k$-wise independence of the star coordinates
$Y_i=f(U)f(U+i)$, $1\le i<B$, whenever
$(k+1)\log(64B)<\log U$: every nonempty product of at most $k$
relative signs selects at most $k+1$ values, and Runge excludes a
square product.  These are sufficient bounds.  They do not establish
that critical windows with a nonzero relation actually exist.

Let $\mathcal R$ be the collection of nonempty even sets $S$ of integers
whose product is a square and which are contained in at least one
$V_x$, $x\in I_N$.  Write
\[
 \begin{gathered}
 a_S=\min S,\qquad b_S=\max S,\qquad d_S=b_S-a_S,\\
 B_S=\{b_S-L+1,\ldots,a_S+1\}\cap I_N,\qquad
 e_S=\ind_{\{a_S+1\in I_N\}}.
 \end{gathered}
\]

\begin{proposition}[Exact local blocks]
\label{p1:limits:local-blocks}
For sufficiently large $N$ in a fixed critical band, the intervals
$B_S$, $S\in\mathcal R$, are nonempty and pairwise disjoint.
They partition the locally deficient starts.  On such a block,
\begin{equation}
 \P(J_x=1)=
 \begin{cases}0,&x=a_S+1,\\2p,&x\in B_S\setminus\{a_S+1\},\end{cases}
 \qquad
 \sum_{x\in B_S}J_x\sim\mathrm{Bernoulli}\bigl(2p(|B_S|-e_S)\bigr).
 \label{p1:limits:block-law}
\end{equation}
Outside these blocks the marginal is $p$.  In particular,
\begin{align}
 \E Z-Np&=p\sum_{S\in\mathcal R}(|B_S|-2e_S),
       \label{p1:limits:block-mean}\\
 \sum_{x\in I_N}|\P(J_x=1)-p|
       &=p\sum_{S\in\mathcal R}|B_S|.
       \label{p1:limits:block-absolute}
\end{align}
A full block contributes $p(L-d_S-1)$ to the mean discrepancy.
There are at most two blocks truncated by the endpoints of $I_N$,
whose total boundary correction is at most $2Lp$ in absolute value.
\end{proposition}
\begin{proof}
The condition $S\subseteq V_x$ is exactly
$b_S-L+1\le x\le a_S+1$.  If two distinct such sets belonged
to the same window, they would give two independent nonzero
relative relations, contrary to Lemma~\ref{p1:limits:local-rigidity}.
Thus their intervals are disjoint and every deficient window lies
in exactly one.  Its unique nonzero relation has affine phase
$(-1)^{\ind_{\{x-1\in S\}}}$, so Fourier inversion gives
$p(1+(-1)^{\ind_{\{x-1\in S\}}})$.  Since $x-1\le a_S$,
the predecessor belongs to $S$ exactly when $x=a_S+1$.
This proves the marginal formulas.  The diameter of $B_S$ is at
most $L-1$; consequently at most one of its starts is active.
Its count is Bernoulli with mean the sum of the displayed
marginals.  Summation proves \eqref{p1:limits:block-mean} and
\eqref{p1:limits:block-absolute}.  A full interval has
$L-d_S+1$ starts.  Disjoint intervals allow at most one truncation
at each endpoint, and each full interval has at most $L$ starts.
\end{proof}

No independence between the Bernoulli block counts is asserted.
For the complete relative sign vector $Y_x\in\{\pm1\}^L$ one
also has the exact identity
\[
 \TV(\Law Y_x,\mathrm{Unif}\{\pm1\}^L)=1-2^{-r_x},
 \qquad r_x=L-\dim S_x.
\]
Indeed $Y_x$ is uniform on a binary image of size $2^{L-r_x}$.
Thus this local distance is either zero or $1/2$ at critical
length.  A nonzero $S$ must have diameter at least
$(1-o(1))\log N$, by the subcritical part of Lemma~\ref{p1:limits:local-rigidity};
its full block has length at most
$(1/\log2-1+o(1))\log N$.

\subsection{A nearby count with factorial tails at every order}

Let $G=G_N$, $n=|G|$, $\mu=np$, and let $E'_N$ be the exceptional
event of Lemma~\ref{p1:core:exceptional}.  It contains the
occurrence of a primitive family of size at least three, an active
contact, an active bad dependent pair, two good active pairs, and
an additional dependence after compressing the unique good pair.
Its probability $\delta_N$ is $O_\varepsilon(N^{-4/5+\varepsilon})$.
Outside it, the active original spaces are a direct sum of single
windows and at most one good pair.  Deleting one endpoint of that
pair, if present, therefore gives a direct family.

\begin{theorem}[A modified count with all factorial tails]
\label{p1:limits:modified}
On the original probability space define
$\widetilde Z_N=Z_N^\circ\ind_{(E'_N)^c}$.
Then
\begin{equation}
 \TV(\Law Z_N,\Law\widetilde Z_N)
       \ll_\varepsilon N^{-4/5+\varepsilon}.
 \label{p1:limits:modified-tv}
\end{equation}
Uniformly for every integer $k\ge1$, every real $s\ge1$, and
every integer $r\ge1$,
\begin{align}
 \P(Z_N^\circ\ge k)&\le\frac{\mu^{k-1}}{(k-1)!}+\delta_N,
       \label{p1:limits:real-tail}\\
 \P(\widetilde Z_N\ge k)&\le\frac{\mu^{k-1}}{(k-1)!},
       \label{p1:limits:modified-tail}\\
 \E s^{\widetilde Z_N}&\le s(1+p(s-1))^n
                         \le s e^{(s-1)\mu},
       \label{p1:limits:modified-exponential}\\
 \E(\widetilde Z_N)_r&\le\mu^r+r\mu^{r-1}.
       \label{p1:limits:modified-moments}
\end{align}
\end{theorem}
\begin{proof}
Let $R$ be the maximum cardinality of a direct active subfamily of
$G$.  Directness makes the prescribed events exactly independent.
A union bound gives
\[
 \P(R\ge k)\le\sum_{\substack{|I|=k\\I\ \text{direct}}}q(I)
       \le\binom nkp^k\le\frac{\mu^k}{k!}.
\]
Choose one maximizing direct family in each realization.  Its
$\binom Rr$ subsets of size $r$ are all direct and active, whence
\[
 \E\binom Rr\le\binom nrp^r\quad(r\ge0).
\]
Outside $E'_N$ we have $Z_N^\circ\le R+1$, and everywhere
$\widetilde Z_N\le R+1$.  The two tail bounds follow.
Expanding $s^R=(1+(s-1))^R$ with nonnegative coefficients gives
$\E s^R\le(1+p(s-1))^n$.  Finally
$\binom{R+1}{r}=\binom Rr+\binom R{r-1}$ proves the factorial
moment bound.  The coupling error is at most $\delta_N$ on $G$;
Proposition~\ref{p1:arith:core} restores the deleted starts.
\end{proof}

These inequalities concern the modified variable.  The event defining
it depends on the whole active configuration, so its indicators no
longer have the original affine marginals or centred factorization.
In particular, \eqref{p1:limits:modified-tv} does not bound
$\E[s^{Z_N};E'_N]$.

\subsection{Effective tails of the original count}

A different argument gives a quantitative tail bound for $Z_N$ itself,
at a larger threshold.  It tracks the moment order explicitly.
Put $H=\log N$ and $h=\log H$.

\begin{lemma}[Branching with an explicit order dependence]
\label{p1:limits:branching}
For any $j$ old starts, let $\mathcal F_{\rm old}$ contain all
their raw sign values.  There is a deterministic exceptional set of
new starts of size at most
\[
 D_{j,N}=j(2L+1)+B\sqrt{3N}
       \exp\left(C\sqrt{m_{j,N}/\log(m_{j,N}+2)}\right),
 \qquad
 m_{j,N}=\left\lceil\pi(B)+\frac{jB\log(3N)}{\log B}\right\rceil,
\]
outside which the $L$ new relative signs are uniform and independent
of $\mathcal F_{\rm old}$.  In particular, uniformly for $1\le j\le H$,
\begin{equation}
 \log(\lambda_N+D_{j,N})
       \le\frac H2+C_C\left(h+\frac{\sqrt j\,H}{h}\right).
 \label{p1:limits:branching-uniform}
\end{equation}
If $F_k=\E(Z_N)_k$, then
\begin{equation}
 F_k\le C_C\exp\left\{
   \frac{k-2}{2}H+C_C\left(kh+\frac{k^{3/2}H}{h}\right)
                     \right\}\quad(2\le k\le H).
 \label{p1:limits:growing-moments}
\end{equation}
\end{lemma}
\begin{proof}
Let $\mathcal P$ consist of the primes at most $B$ and the primes
larger than $B$ dividing one of the old values.  It has at most
$m_{j,N}$ members.  The number of integers at most $3N$ whose
squarefree part is supported on $\mathcal P$ is at most
\[
 \sqrt{3N}\prod_{\ell\in\mathcal P}(1+\ell^{-1/2}).
\]
For a set of $m$ primes, the estimate
$\ell_i\gg i\log(i+1)$ for its $i$th member implies
\[
 \log\prod_{\ell\in\mathcal P}(1+\ell^{-1/2})
       \le C\sqrt{m/\log(m+2)}.
\]
Delete starts whose new window overlaps an old one or contains
at least two integers with squarefree part supported on
$\mathcal P$.  Counting each such integer in at most $B$ windows
gives the displayed bound for $D_{j,N}$.

If a nontrivial combination of the new relative rows belonged
to the old raw-value span, it would select an even nonempty set
of new values.  A prime larger than $B$ in such a value is private
in the new window, so must also divide an old value.  Every
selected value would consequently have squarefree part supported
on $\mathcal P$.  At least two are selected, contradicting the
deletion rule.  Thus the rank increases by exactly $L$ even after
conditioning on all old raw values, which proves independence.

For $j\le H$, $m_{j,N}=O_C(jH^2/h)$ and
$\log(m_{j,N}+2)\gg h$.  This proves
\eqref{p1:limits:branching-uniform}.  Summing the conditional
probabilities of the next start, using the bound one on the
exceptional set, yields the finite recurrence
$F_{j+1}\le(\lambda_N+D_{j,N})F_j$.
The pair estimate gives $F_2=O_C(1)$.  Multiplying from $j=2$
to $k-1$ and summing $\sqrt j$ proves
\eqref{p1:limits:growing-moments}.
\end{proof}

\begin{proposition}[A quantified effective support]
\label{p1:limits:effective-tail}
For each fixed $D>0$ and integer $r\ge0$, there is a constant
$K=K(C,D,r)$ such that
\begin{equation}
 \E[Z_N^r;Z_N>A_N]\le N^{-D},\qquad
 A_N=\left\lceil\sqrt N\exp\left(\frac{K\log N}
                                    {(\log\log N)^{2/3}}\right)\right\rceil
 \label{p1:limits:effective-threshold}
\end{equation}
for sufficiently large $N$.  Hence, for every fixed
$\eta,D>0$ and integer $r\ge0$,
$\E[Z_N^r;Z_N>N^{1/2+\eta}]=O(N^{-D})$.
\end{proposition}
\begin{proof}
Take $k=\lceil h^{2/3}\rceil$.  Eventually $k>r$ and
$A_N\ge2k$.  On $Z_N>A_N$,
$Z_N^r\le2^kA_N^{r-k}(Z_N)_k$.
The logarithm of the resulting bound from
\eqref{p1:limits:growing-moments} is at most
\[
 \frac{r-2}{2}H+C_C\left(kh+\frac{k^{3/2}H}{h}\right)
        -(k-r)\frac{KH}{h^{2/3}}+k\log2+O_C(1).
\]
Here $kh=o(H)$, $k^{3/2}/h=1+o(1)$ and
$(k-r)/h^{2/3}=1+o(1)$.  A sufficiently large fixed $K$
makes this at most $-DH$.  The last assertion follows by
comparison of thresholds.
\end{proof}

The all-fixed-order bounds of Theorem~\ref{p1:fixed:moments}
also imply $\E[Z_N^r;Z_N>N^\eta]=O(N^{-D})$ for fixed
$\eta,D,r$, by choosing one sufficiently large fixed moment.
The value of \eqref{p1:limits:effective-threshold} is its explicit
control of a growing order, without extrapolating fixed-order
constants.

A larger threshold admits a much smaller exceptional probability.
The distinction is useful when the moment order itself increases
as fast as $\log N$.

\paragraph{Branching through a quadratic logarithmic order.}
\label{p18:limits:quadratic-order}
For every fixed $K\ge0$ and $\xi>0$, one threshold, depending on
$C,K,\xi$, precedes all critical lengths and all integers
$0\le k\le K(\log N)^2$ for which
\[
                  F_k\le\bigl(2N^{3/4+\xi}\bigr)^k.
\]
This estimate is obtained from $F_0=1$, independently of any
fixed-order moment bound or law approximation.

To prove it, let $\mathcal P$ contain the primes at most $B$ and
all prime divisors of the raw values in at most $j\le K(\log N)^2$
old windows. Their logarithmic budget satisfies
\[
 \sum_{\ell\in\mathcal P}\log\ell
 \le \sum_{\ell\le B}\log\ell+jB\log(3N)
 \le A_{C,K}(\log N)^4=:Y .
\]
For any finite prime set with this budget, splitting at $Y$ and
using the elementary prime upper bound gives
\[
 \sum_{\ell\in\mathcal P}\ell^{-3/4}
 \ll\frac{Y^{1/4}}{\log Y}.
\]
Indeed the primes at most $Y$ have this bound by partial summation;
for the others use
$\ell^{-3/4}\le Y^{-3/4}\log\ell/\log Y$ and sum the budget.
Every positive integer with squarefree part supported on $\mathcal P$
has a unique form $a^2m$, with $m$ squarefree on that set.
Rankin's inequality at $3/4$ therefore bounds their number up to $3N$ by
\[
 (3N)^{3/4}\zeta(3/2)
       \prod_{\ell\in\mathcal P}(1+\ell^{-3/4})
 \le N^{3/4}\exp\{O_{C,K}(\log N/\log\log N)\}.
\]
This is uniform in the actual old windows.

Delete only new windows containing at least two such supported
values. The number deleted is at most $B$ times the displayed
population, hence at most $N^{3/4+\xi}$ after reducing the input
margin. Outside this set the same parity-row argument as in
Lemma~\ref{p1:limits:branching} makes the whole new relative vector
fair and independent of every old raw value: a nonzero relative
relation would select at least two supported values.
This argument also excludes an old start itself, so the recurrence
for distinct ordered tuples is valid. It yields
$F_{j+1}\le(\lambda_N+N^{3/4+\xi})F_j$.
Since $\lambda_N\le2^C$, increasing the common threshold and
iterating from $F_0=1$ proves the claim.

\begin{proposition}[A very small tail from a power threshold]
\label{p1:limits:colored-tail}
For all sufficiently large $N$ in a fixed critical band, uniformly
for every real $b\ge N^{9/10}$,
\begin{equation}
 \P(Z_N>b)\le e^{-2(\log N)^3},\qquad
 \E[Z_N^r;Z_N>b]\le e^{-(\log N)^3/2}
       \quad(r\in\mathbb N,\ 0\le r\le\log N).
 \label{p1:limits:colored-tail-bound}
\end{equation}
In particular, with $S=\sqrt{\log N\log\log N}$, the same conclusions
hold at $A_{N,c}=\lceil R_Ne^{-cS}\rceil$ for every fixed $c>0$.
The threshold may depend on $c$; no assertion is made for varying
$c=c(N)$.
\end{proposition}
\begin{proof}
Use the preceding bound with $K=129$, $\xi=1/20$, and
$k=\lceil128(\log N)^2\rceil$. Eventually
$k\le129(\log N)^2$ and $2k\le N^{9/10}$.
On $Z_N>b$ the falling factorial satisfies
$(Z_N)_k\ge(b/2)^k$. Thus
\[
 \P(Z_N>b)
 \le\left(\frac{4N^{3/4+1/20}}b\right)^k
 \le\bigl(4N^{-1/10}\bigr)^k
 \le\exp\{-k\log N/20\}
 \le e^{-2(\log N)^3}.
\]
Here $4\le N^{1/20}$ is included in the common cutoff.
Since $Z_N\le N$, multiplying this probability by $N^r$ gives
the stated truncated-moment bound simultaneously for $r\le\log N$.
Finally, in a fixed critical band,
\[
 \log(R_Ne^{-cS})
   =\log N-O_C(\log\log N)-c\sqrt{\log N\log\log N}
   \ge\tfrac9{10}\log N
\]
eventually for each fixed $c$. This proves the last assertion.
Neither a two-smooth-pivot estimate nor a coloring argument is needed.
The result concerns the original count, not the modified count, and
does not give a uniform bound for its real exponential moments.
\end{proof}

\subsection{An arithmetic cylinder with large exponential moments}

\begin{proposition}[Divergence of the real exponential moment]
\label{p1:limits:cylinder}
Along odd primes $q\to\infty$, put $N_q=2^{q-1}$ and
$L_q=q-1$.  For every fixed $t>2$,
\begin{equation}
 \liminf_{q\to\infty}\frac{\log N_q}{N_q}
        \log\E e^{tZ_{N_q,L_q}}\ge(t-2)\log2>0.
 \label{p1:limits:exponential-divergence}
\end{equation}
The same bound holds for the count of positive runs alone.
\end{proposition}
\begin{proof}
Let $M=2N_q+q$ and prescribe the prime-sign cylinder
\[
 E_q=\{f(q)=-1,\ f(\ell)=1
             \text{ for every prime }\ell\le M,\ \ell\ne q\}.
\]
It has probability $2^{-\pi(M)}$.  On this event,
$f(n)=(-1)^{v_q(n)}$ for all integers used by the windows.
Every start $x=uq+1\in I_{N_q}$ with $q\nmid u(u+1)$ has a
negative predecessor, then exactly $q-1$ positive values, and
a negative next value.  Their number is
\[
 K_q=(1-2/q)N_q/q+O(1)=(1-o(1))N_q/q.
\]
Therefore
$\log\E e^{tZ}\ge tK_q-\pi(2N_q+q)\log2$.
Since $q=\log N_q/\log2+1$, the prime number theorem gives
\eqref{p1:limits:exponential-divergence}.
\end{proof}

\paragraph{A stronger finite prime-cylinder bound.}
The construction in Proposition~\ref{p1:limits:cylinder} also gives a
stronger finite estimate.  Let $q\leq N$ be prime, put $L=q-1$ and
$M=2q\lfloor N/q\rfloor\leq2N$, and prescribe $f(q)=-1$ and all other prime signs up to
$M$ to be positive.  Retain the starts $x=qu+1$ with
\[
 \lfloor N/q\rfloor+1\leq u<2\lfloor N/q\rfloor,
 \qquad q\nmid u.
\]
Each has a negative predecessor and at least $q-1$ positive values,
all at integers at most $M$.  The signs above $M$ remain free, so this
cylinder has probability $2^{-\pi(M)}$.
The extra condition $q\nmid(u+1)$ would force a negative right endpoint;
it is unnecessary for a start of length at least $q-1$.
Counting the single excluded residue class, including the floor and
endpoint errors, gives at least $N/q-N/q^2-3$ retained starts.
Their restricted count is bounded by the full positive count;
integrability and monotonicity give, for every real $t\geq0$,
\[
 \E e^{tZ_{N,q-1}}\ \geq\ \E e^{tZ^+_{N,q-1}}
 \ \geq\
 \exp\!\left\{t\left(\frac Nq-\frac N{q^2}-3\right)
                    -\pi(M)\log2\right\},
\]
where $Z^+$ counts positive starts only.  The odd-prime-scale consequence remains the one stated in
Proposition~\ref{p1:limits:cylinder}.
At the original scale $q=7$, $N=64$, the new cutoff $126$ charges
30 primes, whereas the earlier sufficient cutoff $134$ charges 32.
With the same retained-count contribution, the finite exponential lower
bound therefore gains a factor four; this exact example is also checked.

\paragraph{Exact valuations and finite layers.}
Within the displayed multiplier interval, a positive start occurs
precisely when $\nu_q(u)$ is even: its predecessor has valuation
$1+\nu_q(u)$, and its following $q-1$ values are positive.
The coprime multipliers form the layer $\nu_q(u)=0$.
For example, $N=18$, $q=3$ also admits $u=9$, adding the start $28$.
All these additional starts use the same supported cylinder.

For every integer depth $k\geq0$, retaining the disjoint layers
$\nu_q(u)=2j$, $0\leq j<k$, gives the explicit lower count
\[
 B_k(N,q)=\frac{N}{q+1}\bigl(1-q^{-2k}\bigr)-3k.
\]
Exact counts of multiples in translated intervals give the endpoint
error $3k$; the geometric sum is evaluated exactly.  For every $t\geq0$,
both full exponential moments are therefore at least
$\exp\{tB_k(N,q)-\pi(M)\log2\}$.
The choice $k=1$ recovers $N/q-N/q^2-3$.
At the original critical scale $q=17$, $N=2^{16}$, the budget for
$k=2$ is strictly greater, giving a strictly stronger exponential
lower bound for every $t>0$ at the same cost.
Larger depth need not improve this explicit estimate because its
endpoint error grows.  Here $k$ is the number of retained valuation layers.

\paragraph{Effective density of the retained starts.}
Let $C(N,q)$ denote the complete retained start count above, and put
$\kappa(N,q)=\lceil\tfrac12\lceil\log_q N\rceil\rceil$.
Every retained multiplier satisfies $0<u<N\leq q^{2\kappa(N,q)}$.
Its even valuation is therefore $2j$ for exactly one
$0\leq j<\kappa(N,q)$, so these finitely many disjoint layers cover
the complete retained set.  Their exact multiple counts give
\[
 \frac{N-1}{q+1}-3\kappa(N,q)\leq C(N,q)
 \leq\frac{N}{q+1}+2\kappa(N,q).
\]
Both full exponential moments consequently have the lower bound
$\exp\{t[(N-1)/(q+1)-3\kappa(N,q)]-\pi(M)\log2\}$ for $t\geq0$.
This depth gives a complete cover; it need not maximize $B_k$.
For each fixed prime $q$, the ratio $C(N,q)/N$ tends to $1/(q+1)$.
Moreover, simultaneously for all primes $q\leq N$,
\[
 \left|\frac{C(N,q)}N-\frac1{q+1}\right|
 \leq\frac{3\kappa(N,2)+1}{N}\longrightarrow0.
\]
The approximation threshold depends only on the requested absolute
normalized error.  This is the density of the deterministic forced
subset.

\paragraph{Charge the rank of the forcing event.}
Let $G$ be the complete retained start set.  On the actual prime-bit
space up to $M$, write $f(n)=(-1)^{Y_n}$ and impose the linear equations
$Y_{x-1}=1$ and $Y_{x+j}=0$ for $x\in G$, $0\leq j<q-1$.
The one-negative-prime sample proves that this affine system is
compatible.  If its effective rank is $r$, its exact mass is $2^{-r}$,
and rank--nullity gives $r\leq\pi(M)$.  Both full exponential moments
are therefore at least $\exp\{tC(N,q)-r\log2\}$ for $t\geq0$.
Either earlier lower count budget may replace $C(N,q)$ in this bound.

A prime dividing none of the observed integers supplies a nonzero
homogeneous kernel vector by flipping just its bit, so $r<\pi(M)$.
For $N=18$, $q=3$, the retained multipliers are exactly $7,8,9,10,11$,
giving five starts.  Their observed integers lie between $21$ and $35$.
The supported cutoff $36$ contains eleven primes, but the prime $19$
is unused. In the homogeneous system the observed values
$27,32,30,28,22,26,34,23,29,31$ successively determine all ten other
prime bits. Hence the rank is exactly $r=10$, and the compatible
forcing event consists of exactly two prime-sign samples, with
probability $2^{-10}$. Both full moments are therefore at least
$\exp(5t)/1024$, compared with $\exp(5t)/2048$ from the same-count
single-atom bound.  This gives at least a factor-two gain, with no
assertion of a new asymptotic threshold or an exact general rank formula.

The number $2$ is a sufficient threshold in this construction;
its optimality is not asserted.  For every fixed $r$, the same
cylinder contributes at most $R_{N_q}^r2^{-\pi(M)}=o(1)$ to
the $r$th moment.  It is thus compatible with all fixed-moment
estimates and with the total-variation approximation.

The cylinder also produces connected families that lie far beyond
the bounded-nullity boundary used in the proof.  The next elementary
rank argument extracts one such family with a positive centred
coefficient.

\begin{lemma}[Deleting large private ranks]
\label{p1:limits:peeling}
Let $A$ be a compatible family of $m$ affine events, each with
constraint-space dimension $L$, and put
$R=\dim\sum_{x\in A}S_x$ and $p=2^{-L}$.
For an integer $1\le t\le L$, repeatedly remove a member whose
current private rank is at least $t$.  The final family $A'$ satisfies
\[
 |A'|\ge m-R/t,\qquad
 \dim\sum_{x\in A'}S_x\le R,\qquad
 r_x(A')<t\quad(x\in A'),
\]
where $r_x(I)=\dim\sum_{y\in I}S_y-
\dim\sum_{y\in I\setminus\{x\}}S_y$.
Each direct component $C$ of $A'$ satisfies
\begin{equation}
 c(C)\ge(1-|C|p2^t)q(C)\ge(1-mp2^t)q(C).
 \label{p1:limits:positive-core}
\end{equation}
If $m-R/t>0$ and $D=R/(m-R/t)$, some such component obeys
\begin{equation}
 \frac{\dim\sum_{x\in C}S_x}{|C|}\le D,\qquad
 |C|\ge L/D,\qquad \rho(C)\ge L^2/D-L.
 \label{p1:limits:rank-density}
\end{equation}
\end{lemma}
\begin{proof}
Each deletion reduces the current rank by at least $t$, so at
most $R/t$ deletions occur.  For a compatible family, all Fourier
phases of its relations are positive.  The centred Fourier formula is
\[
 c(C)=p^{|C|}\#\{(s_x)\in\ker(\bigoplus_{x\in C}S_x
                        \longrightarrow\sum_{x\in C}S_x):
                                      s_x\ne0\ (x\in C)\}.
\]
The image of the projection of this kernel on $S_x$ has dimension
$L-r_x(C)\ge L-t+1$.  In the uniform kernel, the probability
that $s_x=0$ is consequently at most $p2^t$.
A union bound and $q(C)=p^{|C|}2^{\rho(C)}$ prove
\eqref{p1:limits:positive-core}.  The size-weighted average of
the rank-to-size ratios of the components is at most $D$.
Choose a component no larger than that average.  Its rank is
at least $L$, proving $|C|\ge L/D$.  If $D<L$, use
$\rho(C)\ge|C|(L-D)$; if $D\ge L$, the stated lower bound
for nullity is nonpositive.
\end{proof}

\begin{proposition}[A real component of quadratic nullity]
\label{p1:limits:large-component}
Along the subsequence of Proposition~\ref{p1:limits:cylinder},
there are compatible connected families $C_q\subset G_{N_q}$
containing no good pair and satisfying
\begin{equation}
 \begin{gathered}
 |C_q|\ge(\kappa-o(1))L_q,\qquad
 \rho_*(C_q)=\rho(C_q)\ge(\kappa-o(1))L_q^2,\\
 c(C_q)\ge(1-O(L_q^{-1/2}))q(C_q)>0,\qquad
 \kappa=\frac{6\log2}{17}.
 \end{gathered}
 \label{p1:limits:quadratic-nullity}
\end{equation}
\end{proposition}
\begin{proof}
Restrict the starts in the cylinder to
\[
 T_q=\{x=uq+1:3N_q/2+4q\le x<2N_q,\quad q\nmid u(u+1)\},
 \qquad A_q=T_q\cap G_{N_q}.
\]
The regular-set estimate gives
$m=|A_q|\sim N_q/(2q)$; the cylinder witnesses compatibility.
For a good edge $x<y$ in channel $a:b$, its exact unit gives
$y=(b/a)x+(b/a)i-j$, with $-1\le i,j\le L_q-1$ and
$b/a\in\{4/3,3/2\}$.  The lower bound
$y\ge(b/a)x-(1+b/a)q$ exceeds $2N_q$ in this interval.
Thus $A_q$ has no good edge.

All prime coordinates used by these windows belong to the union
\[
 [1,2N_q/3+O(q)],\quad
 [3N_q/4-O(q),N_q+O(q)],\quad
 [3N_q/2-O(q),2N_q+O(q)].
\]
Indeed the multiplier of a prime in a used value is respectively
at least three, two, or one.  The prime number theorem bounds
their number by $(17/12+o(1))N_q/\log N_q$.
Hence the total rank $R$ satisfies
$R/m\le17/(6\log2)+o(1)$.
Apply Lemma~\ref{p1:limits:peeling} with
$t=\lfloor\frac12\log_2L_q\rfloor$.  Then $R/t=o(m)$,
$D\le17/(6\log2)+o(1)$ and $mp2^t=O(L_q^{-1/2})$.
Equations \eqref{p1:limits:positive-core}--\eqref{p1:limits:rank-density}
give the required component.  Every subfamily is compatible,
and nullity is nondecreasing under inclusion, so its compatible
nullity equals its full nullity.
\end{proof}

This is a statement about which arithmetic terms exist.  The
probability of the witnessing cylinder is not identified with
$q(C_q)$, and the proposition gives no lower bound for the total
variation error or for a sum of absolute component activities.

\subsection{Keeping positions while justifying pair deletion}
\label{p1:limits:positions}
\label{p14:comp:positional-proof}

Deleting an active pair is a transformation of the whole field.
The scalar approximation alone does not justify this operation.
We retain coordinates in the same finite expansion, using the
boundary and deficit estimates already proved.
Lemma~\ref{p1:limits:vector-reference} compares the unsigned
field \((J_x)_{x\in I_N}\) with independent Poisson single
arrivals and arrivals along good arithmetic pairs, with an
\(O_{C,j,\varepsilon}(N^{-4/5+\varepsilon})\) error for
every fixed polynomial weight of degree \(j\) in the total count.
The field retains integer positions, without run signs or exact
excess lengths; the target is on the same grid, not a diffuse
spatial process. We first work on the regular set \(G=G_N\)
and then restore the excluded starts.

\paragraph{The finite reference law and its exact identities.}
\label{p18:limits:reference-identities}
The following construction also covers heterogeneous rates and packets
of any finite size. Let $F$ be a finite ground, $E$ a finite family of
subsets of $F$, and suppose
\[
 q_e\ge0,\qquad
 p_x^{\mathrm{ref}}\ge\iota_x:=\sum_{e\in E:x\in e}q_e
 \quad(x\in F).
\]
Take mutually independent Poisson variables $U_x$ of parameters
$\nu_x=p_x^{\mathrm{ref}}-\iota_x$ and $V_e$ of parameters $q_e$,
and put $H_x=U_x+\sum_{e\ni x}V_e$. The superscript distinguishes
these reference intensities from the actual arithmetic marginals
$p_x=\E J_x$. A finite product probability space constructs the
variables; every independent realization has the same field law.

Each $H_x$ is exactly $\Pois(p_x^{\mathrm{ref}})$. In particular,
its falling-factorial test is integrable for every $j\in\mathbb N$
and $\E(H_x)_j=(p_x^{\mathrm{ref}})^j$. The coordinates are square
integrable and
\[
 \Cov(H_x,H_y)
  =\ind_{\{x=y\}}\nu_x+\sum_{e\in E:x,y\in e}q_e.
\]
The variance is $p_x^{\mathrm{ref}}$. For pairs, the distinct-site
covariance is the retained pair rate, or zero when that pair is absent.

Two finite blocks of coordinates, even possibly overlapping, are
independent as complete natural-valued vectors if and only if all
their cross covariances vanish. For disjoint blocks and pair packets,
this says exactly that no retained edge of positive rate crosses
between them. Thus the positive-rate graph is a dependency graph for
the full reference field; zero-rate edges create no dependence, and
measurable events on independent blocks have product probabilities.

\begin{proof}
Represent a coordinate by the finite set of cells it observes: its
own $U_x$ and the incident $V_e$. The sum of their parameters is
$\nu_x+\iota_x=p_x^{\mathrm{ref}}$, proving the marginal law.
Poisson factorial moments give the stated integrability and moment
formula; their first two orders also give square integrability.
Bilinearity of covariance leaves precisely the cells observed in
both coordinates, since distinct cells are independent and a Poisson
cell has variance equal to its parameter. This proves the formula,
including its diagonal.

For two blocks, every cell observed on both sides contributes its
nonnegative rate to at least one cross covariance. If all those
covariances vanish, every shared cell has rate zero and vanishes
almost surely. Remove these cells simultaneously from the finite
block vectors. Each remaining vector is a measurable function of a
disjoint subfamily of independent cells, so the whole vectors are
independent. Conversely, independent square-integrable block vectors
have zero cross covariance. If the blocks overlap, this also requires
zero variance at each common site, as the criterion correctly records.
This proves the block assertion and its specialization to crossing
pair rates. Event factorization follows from independence.
\end{proof}

\paragraph{The bivariate mass formula.}
\label{p18:limits:bivariate-reference}
Put $r_{xy}=\Cov(H_x,H_y)$ and write
\[
 a=p_x^{\mathrm{ref}}-r_{xy},\qquad
 b=p_y^{\mathrm{ref}}-r_{xy},\qquad c=r_{xy}.
\]
All three rates are nonnegative:
$0\le r_{xy}\le\min(p_x^{\mathrm{ref}},p_y^{\mathrm{ref}})$.
The couple $(H_x,H_y)$ has the law of $(A+C,B+C)$ with independent
$\Pois(a),\Pois(b),\Pois(c)$ variables. Consequently
\[
 \P(H_x=m,H_y=n)
 =e^{-(a+b+c)}\sum_{k=0}^{\min(m,n)}
       \frac{a^{m-k}b^{n-k}c^k}{(m-k)!(n-k)!k!}
 \qquad(m,n\in\mathbb N).
\]
Indeed partition the observed cells into those seen only at $x$,
only at $y$, and at both sites. Their independent sums have the
three displayed parameters. For the point probability, the common
sum is $k$ and the two other sums are $m-k,n-k$; these are exactly
the finitely many possibilities in the display. The formula includes
coincident sites and zero rates with $0^0=1$. If $x=y$, then
$a=b=0$ and only the diagonal Poisson mass remains.

In the application $F=I_N$, $E=E_N$ and
$p_x^{\mathrm{ref}}=p$ even outside the regular set; $E_N$ has no
incident edge there. Admissibility is the uniform positivity of
$p-\iota_x$ established in Section~\ref{p14:positions:section}.
These reference identities do not identify the arithmetic covariance:
for a good regular pair it is $q_e-p^2$ for $(J_x,J_y)$, whereas it
is $q_e$ for $(H_x,H_y)$. The covariance matrix alone also does not
give an independence criterion for a general arithmetic field.

On $\mathbb N^G$, let $\boldsymbol e_x$ denote the coordinate
unit vector and set
\[
 \Delta_x=\delta_{\boldsymbol e_x}-\delta_{\boldsymbol0},
 \qquad b_x=\delta_{\boldsymbol0}+p\Delta_x,
 \qquad \Delta_I=\mathop{*}_{x\in I}\Delta_x.
\]
For a hereditary class $\mathscr D$, the same partition by first
excluded prefix gives the exact identity
\begin{equation}
 \Law((J_x)_{x\in G})-
       \sum_{I\in\mathscr D}q(I)\Delta_I
 =\sum_{I\in\partial\mathscr D}\Delta_I*\boldsymbol\nu_I,
 \quad
 \boldsymbol\nu_I=
 \E\left[J_I\delta_{\sum_{y>\max I}J_y\boldsymbol e_y}\right].
 \label{p1:limits:vector-stop}
\end{equation}
One can obtain it directly by expanding
$\prod_x(1+u_x)$ with $u_x=J_x\Delta_x$ in the convolution
algebra.  The proof of Lemma~\ref{p1:stop:polynomial} permits
a different algebra element at every site.  Each
$\boldsymbol\nu_I$ is positive of mass $q(I)$, and
$\|\Delta_I\|_1=2^{|I|}$.  Thus the boundary estimate is
unchanged.

Use the field laws $\mathsf H_N$ and $\mathsf U_N$ and the intensities
$q_e,\iota_x,\nu_x$ defined in Section~\ref{p14:positions:section}. In
particular, $\beta=\sum_e q_e$, $\theta=2\beta$, $\lambda=Np$, and
$\iota_x=0$ outside $G$. The total-count laws are given by
\eqref{p1:limits:compound-target}.

\begin{proof}[Proof of Lemma~\ref{p1:limits:vector-reference}]
First work on $G$ in the unweighted norm.  The retained reference
is now
\[
 \mathsf B_\beta=\mathop{*}_{x\in G}b_x+
       \sum_{e=\{x,y\}\in E_N}q_e\Delta_x*\Delta_y*
                                      \mathop{*}_{z\in G\setminus e}b_z.
\]
The exact retained identity subtracts the missing direct
families, with coefficient $p^{|I|}$, and the missing compressed
extensions, with coefficient $q_ep^{|U|}$.  Taking norms of
these subtractions costs at most $F+\mathfrak G$, with exactly
the deficits of \eqref{p1:close:total-deficits}.  There is no
scalar cancellation in that estimate.  The vector version of
the first-boundary remainder is bounded by the same positive
sum as \eqref{p1:stop:norm}.  Hence
\[
 \|\Law((J_x)_{x\in G})-\mathsf B_\beta\|_1
       \ll_\varepsilon N^{-4/5+\varepsilon}.
\]

Put $\mathsf P=\exp_*\{p\sum_{x\in G}\Delta_x\}$.
Replacing the product of Bernoulli laws by $\mathsf P$ costs
$O(np^2)$.  In each edge term, replacing the product outside
$e$ by $\mathsf P$ costs $O(np^2+p)$ before multiplication
by $4q_e$.  Thus
\[
 \mathsf B_\beta=
 \mathsf P*\left(\delta_{\boldsymbol0}
                     +\sum_{e=\{x,y\}}q_e\Delta_x*\Delta_y\right)
              +O_{\ell^1}(np^2+\beta np^2+\beta p).
\]
The exponential remainder has norm at most
$e^{4\beta}-1-4\beta=O(\beta^2)$.  Moreover
\[
 p\sum_x\Delta_x+\sum_{e=\{x,y\}}q_e\Delta_x*\Delta_y
 =\sum_x(p-\iota_x)\Delta_x+
       \sum_{e=\{x,y\}}q_e
                    (\delta_{\boldsymbol e_x+\boldsymbol e_y}
                                      -\delta_{\boldsymbol0}),
\]
so this convolution exponential is exactly the positive law
$\mathsf H_N$ on $G$.  The errors $np^2$ and $\beta^2$ are
$O_\varepsilon(N^{-1+\varepsilon})$.  Restoring the deleted
starts in both laws costs $N^{-1+o(1)}$.

It remains to justify the weights; total variation alone is
insufficient.  Both total counts have every fixed moment
uniformly bounded: for the real count this is
Theorem~\ref{p1:fixed:moments}, and for the compound target
it follows from \eqref{p1:limits:compound-target} with bounded
parameters.  For their signed difference $\sigma$ and any fixed
integer $s>j$, H\"older's inequality gives
\[
 \sum(1+|\boldsymbol z|)^j|\sigma\{\boldsymbol z\}|
 \le\|\sigma\|_1^{1-j/s}
       \left(\sum(1+|\boldsymbol z|)^s
                         |\sigma\{\boldsymbol z\}|\right)^{j/s}.
\]
Choose $s$ sufficiently large after $\varepsilon$, and use the
unweighted estimate with a smaller error margin.  This proves
\eqref{p1:limits:vector-weighted} for each fixed $j$.
\end{proof}

The weights in this lemma have fixed polynomial degree.  They
are not exponential weights and do not contradict
Proposition~\ref{p1:limits:cylinder}.

Let $\Xi=\sum_{e\in E_N}J_e$, $C_N=\{\Xi\ge1\}$ and
$r_N=\P(C_N)$.  On $C_N$, choose one active good edge uniformly
and delete its two endpoints, obtaining a binary field $J^\flat$.
Off $C_N$ keep the field unchanged.  Also let $\widehat J$ be
the field obtained by deleting every vertex incident to an active
good edge.  Their counts are
\[
 W=\sum_xJ_x^\flat=Z-2\ind_{C_N},\qquad
 \widehat W=\sum_x\widehat J_x.
\]
The auxiliary random choice affects the field $J^\flat$ but
not $W$.

\begin{proposition}[Deletion, isolation, and remaining covariances]
\label{p1:limits:erasure}
Uniformly in the critical band, for every $\varepsilon>0$,
\begin{align}
 \TV(\Law J^\flat,\mathsf U_N)
       +\TV(\Law\widehat J,\mathsf U_N)
       &\ll_\varepsilon N^{-4/5+\varepsilon},
       \label{p1:limits:erasure-tv}\\
 \sum_x\big|\E J_x^\flat-(p-\iota_x)\big|
       +\sum_x\big|\E\widehat J_x-(p-\iota_x)\big|
       &\ll_\varepsilon N^{-1+\varepsilon}.
       \label{p1:limits:erasure-intensity}
\end{align}
For a counting field $X$ put
$K_{xy}(X)=\Cov(X_x,X_y)-\ind_{\{x=y\}}\E X_x$.
Then
\begin{equation}
 \sum_{x,y}|K_{xy}(J^\flat)|+
 \sum_{x,y}|K_{xy}(\widehat J)|
       \ll_\varepsilon N^{-4/5+\varepsilon}.
 \label{p1:limits:erasure-covariance}
\end{equation}
In particular, the absolute second factorial cumulants of
$W$ and $\widehat W$ have this bound.  On the natural coupling,
\begin{align}
 0\le\beta-r_N&\le\tfrac12\E(\Xi)_2=O(a_N^2),
       \label{p1:limits:edge-occurrence}\\
 \TV(\Law W,\Law\widehat W)&\le\tfrac12\E(\Xi)_2=O(a_N^2),
       \label{p1:limits:one-all-tv}\\
 W_1(\Law W,\Law\widehat W)&\le\E(\Xi)_2=O(a_N^2).
       \label{p1:limits:one-all-wasserstein}
\end{align}
Here $a_N\asymp L N^{-2/3}$ as in the main theorem.
\end{proposition}
\begin{proof}
Extend the one-edge deletion kernel to $\mathbb N^{I_N}$ by
choosing uniformly an edge whose two endpoint counts are positive
and subtracting one at each endpoint; if no such edge exists,
leave the configuration unchanged.  It agrees with the binary
deletion kernel.  Apply it to the real and compound fields.  The unweighted comparison error cannot increase.
In the target representation, $T=\sum_eV_e\sim\Pois(\beta)$.
If $T=0$, deletion leaves $U$ unchanged unless two occupied
single sites form an edge.  This event has probability at most
$|E_N|p^2=O(B/N)$.  If $T=1$, let $e$ be its edge.
Deleting $e$ recovers $U$ exactly.  A different edge can be
active only if its endpoints are occupied singles, or if it
shares one endpoint with $e$ and its other endpoint is a
single.  The corresponding bounds are $|E_N|p^2$ and
$O(Bp)$, uniformly in $e$.  Finally $\P(T\ge2)=O(\beta^2)$.
This couples the once-deleted target to $U$ with error
$O(B/N+\beta^2)$.  For isolation, set all coordinates incident
to an active edge to zero; the same argument also excludes
singles at the two endpoints of $e$, at cost $2p$ when
$T=1$.  This proves \eqref{p1:limits:erasure-tv}.

Every deleted count is bounded by its original total count.
The weighted argument at the end of
Lemma~\ref{p1:limits:vector-reference} therefore applies also
to the two deleted laws and $\mathsf U_N$.  With weight
$(1+|\boldsymbol z|)^2$ it controls, in the sum of absolute
entries, their first and second factorial moment measures.
The total intensities are bounded, so it controls the tensor
squares of their means as well.  The independent Poisson field
has zero factorial covariance matrix.  This proves
\eqref{p1:limits:erasure-covariance}.

For the finer intensity assertion, let $B=J-J^\flat$.
Nonnegativity gives the exact identity
\[
 \sum_x(\iota_x-\E B_x)=2(\beta-r_N),\qquad
 0\le\E B_x\le \iota_x.
\]
For isolation, put $D_x=\sum_{y:\{x,y\}\in E_N}J_y$.
Then
\[
 \E\widehat J_x=\E J_x-\iota_x+
       \E[J_x(D_x-1)_+],\qquad
 \sum_x\E[J_x(D_x-1)_+]\le2(\beta-r_N).
\]
The last inequality holds realization by realization: if there
are active edges, at least two vertices are incident to them,
whereas the sum of their active degrees is $2\Xi$.
Now use $\sum_x|\E J_x-p|=N^{-1+o(1)}$ and the positive
two-edge bound in Lemma~\ref{p1:arith:twogoodpositive}.

Finally $0\le k-\ind_{\{k\ge1\}}\le(k)_2/2$ for integers
$k\ge0$.  Also $W=\widehat W$ when $\Xi\le1$, and
\[
 0\le W-\widehat W\le2(\Xi-\ind_{C_N})\le(\Xi)_2.
\]
These inequalities prove
\eqref{p1:limits:edge-occurrence}--\eqref{p1:limits:one-all-wasserstein}.
\end{proof}

The same covariance conclusion applies to the positive cubic
count used to approximate isolation.  With the active degrees
$D_x$ from the proof, put
\[
 Y_x^{[3]}=J_x\left(1-D_x+\binom{D_x}{2}\right),\qquad
 Y^{[3]}=\sum_xY_x^{[3]},\qquad
 \Sigma=\sum_xJ_x\binom{D_x}{2}.
\]
For an integer $d\ge1$, $1-d+\binom d2=\binom{d-1}{2}\ge0$,
whereas its value at zero is one.  Thus
$0\le\widehat J_x\le Y_x^{[3]}$, and the fields agree when
$\Xi\le1$.  More precisely, if
$O=\sum_xJ_x(D_x-1)_+$, then
\[
 Y^{[3]}-\widehat W=\Sigma-O\ge0,\qquad
 \Sigma\le\tfrac12(\Xi)_2,\qquad
 \E Y^{[3]}=m-2\beta+\E\Sigma.
\]
The inequality for $\Sigma$ counts pairs of distinct good
edges sharing a distinguished endpoint; in a simple graph
each unordered pair can share at most one endpoint.
It follows that both the total-variation and the
$W_1$ distances from $Y^{[3]}$ to $\widehat W$ are
$O(a_N^2)$.  Since the deterministic degree is $O(B)$,
$Y^{[3]}\ll B^2Z$.  The fixed moments of $Z$, H\"older's
inequality on $\{\Xi\ge2\}$, and
Lemma~\ref{p1:arith:twogoodpositive} give, for every fixed
polynomial weight, a difference
$O_\varepsilon(N^{-4/3+\varepsilon})$ between the two laws.
In particular,
\begin{equation}
 |\Var(Y^{[3]})-\E Y^{[3]}|
       \ll_\varepsilon N^{-4/5+\varepsilon}.
 \label{p1:limits:cubic-cumulant}
\end{equation}
This deduction needs no asymptotic for the rarer configurations
where the cubic and isolated counts differ.

Let $P_0=\Law(Z\mid C_N^c)$, $\mu_0=\E[Z\mid C_N^c]$,
and $m=\E Z$.  The preceding proof also gives, for every fixed
polynomial weight, a comparison of $P_0$ with
$\Pois(\lambda-\theta)$ at order $N^{-4/5+\varepsilon}$.
Indeed absence of an observed edge in the compound field forces
$T=0$; the remaining conditional law is $U$ conditioned on
having no accidental edge.  Removing this conditioning costs
$O(B/N)$, and $\P(C_N^c)=1-o(1)$ permits conditioning the
comparison in Lemma~\ref{p1:limits:vector-reference}.
Consequently
\begin{equation}
 2r_N=\theta+O(a_N^2),\qquad
 \E W=m-2r_N=\lambda-\theta+O_\varepsilon(N^{-1+\varepsilon}),\qquad
 \mu_0=\lambda-\theta+O_\varepsilon(N^{-4/5+\varepsilon}).
 \label{p1:limits:background-means}
\end{equation}
The identity for the mean of $W$ is exact before its last
estimate.  One may replace $\theta$ by the two-channel expression
of Theorem~\ref{p1:arith:pairs}, or by the actual second
factorial cumulant, at cost $O_\varepsilon(N^{-4/5+\varepsilon})$.

For clarity, the corresponding localization statements are
\begin{align}
 &\big|\TV(\Law Z,H_{\lambda,\theta})
       -(1-r_N)\TV(P_0,\Pois(\mu_0))\big|
       \ll_\varepsilon N^{-4/5+\varepsilon},
       \label{p1:limits:background-localization}\\
 &\big|\TV(\Law Z,H_{\lambda,\theta})
       -\TV(\Law W,\Pois(\E W))\big|
       \ll_\varepsilon N^{-4/5+\varepsilon},
       \label{p1:limits:erased-localization}\\
 &\big|\TV(\Law((J_x)_x),\mathsf H_N)
       -\TV(\Law J^\flat,\mathsf U_N)\big|
       \ll_\varepsilon N^{-4/5+\varepsilon}.
       \label{p1:limits:spatial-localization}
\end{align}
They follow by the triangle inequality from the comparisons just
proved; the changes of Poisson means cost at most their absolute
differences.  Thus the full stopped expansion supplies, in
particular, the first-order reduction to the background.
Small second cumulants alone would not supply any of these
distributional comparisons.

\subsection{Signed Palm identities and weighted recurrence errors}

It is useful to state precisely which information a scalar
Stein comparison retains.  For a probability $P$ on
$\mathbb N$ with finite first moment, put $P_k=P\{k\}$ and
$P_k=0$ for $k<0$, and define
\begin{equation}
 r_k(P;\mu,\theta)
       =kP_k-(\mu-\theta)P_{k-1}-\theta P_{k-2},\qquad
 \mathcal R(P;\mu,\theta)=\sum_{k\ge0}|r_k(P;\mu,\theta)|.
 \label{p1:limits:raw-residual}
\end{equation}

\begin{lemma}[A scalar recurrence criterion]
\label{p1:limits:stein}
If $0\le\theta<\min(\mu,1/2)$, then
\begin{equation}
 \TV(P,H_{\mu,\theta})
       \le\frac{\mathcal R(P;\mu,\theta)}{1-2\theta}.
 \label{p1:limits:stein-bound}
\end{equation}
If $m=\E X>0$, $\kappa=\Var X-\E X$ and
$0\le\kappa<\min(m,1/2)$, the size-biased variable $X^s$
satisfies the exact identity
\begin{equation}
 \mathcal R(\Law X;m,\kappa)
 =2m\TV\left(\Law X^s,
       (1-\kappa/m)\Law(X+1)+(\kappa/m)\Law(X+2)\right).
 \label{p1:limits:size-bias}
\end{equation}
\end{lemma}
\begin{proof}
For an indicator test $h$, a Poisson Stein solution satisfies
$\mu f(k+1)-kf(k)=h(k)-\Pi_\mu h$ and may be chosen with
$\|f\|_\infty\le1$, $\|\Delta f\|_\infty\le1$.
One proof uses the immigration--death process at time $t$,
whose law from $k$ is
$\mathrm{Bin}(k,e^{-t})+\Pois(\mu(1-e^{-t}))$.
One extra initial individual survives with probability $e^{-t}$,
and two survive with probability $e^{-2t}$.  Integrating these
couplings in the Poisson equation gives respectively the bounds
$\int_0^\infty e^{-t}dt=1$ and
$\int_0^\infty2e^{-2t}dt=1$; set $f(0)=f(1)$.

The operator
$A_\theta f(k)=\mu f(k+1)+\theta\Delta f(k+1)-kf(k)$
has mean zero under $H_{\mu,\theta}$, by applying the elementary
Poisson identity separately to its single and double counts.
Changing indices gives
$P A_\theta f=-\sum_kr_kf(k)$.
The difference of the expectations of $\Delta f(k+1)$ under
$P$ and $H_{\mu,\theta}$ is at most
$2\TV(P,H_{\mu,\theta})$.  Taking the supremum over $h$
and absorbing this perturbation proves
\eqref{p1:limits:stein-bound}.  The two measures in
\eqref{p1:limits:size-bias} differ at $k$ by
$r_k(\Law X;m,\kappa)/m$; total variation is half their
$\ell^1$ distance.
\end{proof}

The constant in \eqref{p1:limits:stein-bound} is not claimed
optimal.  The criterion is sufficient, and the raw norm is not
bounded by total variation alone.  Parameter changes satisfy
the elementary bound
\begin{equation}
 |\mathcal R(P;\mu,\theta)-\mathcal R(P;\mu',\theta')|
       \le|\mu-\mu'|+2|\theta-\theta'|.
 \label{p1:limits:residual-parameters}
\end{equation}
When the first three moments exist and $\mu=m$, $\theta=\kappa$,
the signed residual has
$\sum_kr_k=\sum_kkr_k=0$ and
$\sum_kk^2r_k=\kappa_3^{(F)}(X)$.
The last identity tests it by an unbounded function and is not
an estimate for its norm.

\begin{proposition}[The signed Palm measure]
\label{p1:limits:signed-palm}
For any finite family of indicators, let $Z=\sum_xJ_x$,
$p_x=\E J_x$, $m=\E Z$, $P=\Law Z$, and $T_x=Z-J_x$.
Omit sources with $p_x=0$, and define
\[
 \mathcal D=\sum_xp_x
       \{\Law(T_x\mid J_x=1)-\Law(T_x)\}.
\]
Then
\begin{equation}
 \big|\mathcal R(P;m,\theta)
       -\|\mathcal D-\theta\Delta*P\|_1\big|
       \le2\sum_xp_x^2.
 \label{p1:limits:signed-palm-bound}
\end{equation}
Writing $q_A=\E\prod_{x\in A}J_x$, the coefficient of
$(z-1)^j$ in its generating polynomial is exactly
\begin{equation}
 d_j=\sum_{|A|=j}\sum_{x\notin A}(q_{A\cup\{x\}}-p_xq_A).
 \label{p1:limits:signed-palm-coefficients}
\end{equation}
If all $n$ marginals equal $p$ and
$B_j=\E\binom Zj$, this becomes
$d_j=(j+1)B_{j+1}-p(n-j)B_j$.
\end{proposition}
\begin{proof}
Let $SP$ translate a measure by one and let $(KP)_k=kP_k$.
The size-bias identity is
$KP=S\sum_xp_x\Law(T_x\mid J_x=1)$.
Consequently the residual in \eqref{p1:limits:raw-residual}
is the translate of
\[
 \mathcal D+\mathcal B-\theta\Delta*P,\qquad
 \mathcal B=\sum_xp_x\{\Law(T_x)-P\}.
\]
In the natural coupling, $T_x$ and $Z$ differ only when
$J_x=1$, so $\|\mathcal B\|_1\le2\sum_xp_x^2$.
Translation preserves $\ell^1$ norm, and the reverse triangle
inequality proves \eqref{p1:limits:signed-palm-bound}.
Expanding $\prod_{y\ne x}(1+(z-1)J_y)$ before summing over
$x$ gives \eqref{p1:limits:signed-palm-coefficients}.
\end{proof}

The order of summation matters: both the sum over intersections
of a fixed size and the sum over sources occur before the norm.
In this model $p_x\le2p$ for large $N$, so the diagonal error is
$O(N^{-1})$.  Taking absolute values separately for each
forced coordinate would remove cancellations present in
$\mathcal D$.

For reference, two sufficient uncorrected criteria follow without
any affine assumption.  If
\[
 v_x=\TV(\Law T_x,\Law(T_x\mid J_x=1)),\qquad
 w_x=W_1(\Law T_x,\Law(T_x\mid J_x=1)),
\]
then
\begin{equation}
 \mathcal R(P;m,0)\le2\sum_xp_x(p_x+v_x),\qquad
 \TV(P,\Pois(m))\le\sum_xp_x(p_x+w_x).
 \label{p1:limits:palm-sufficient}
\end{equation}
The first follows from the preceding proposition.  For the
second, couple the two scalar $T_x$ laws optimally and lift
that coupling to the original joint law using conditional
distributions given $T_x$.  The extra difference from $Z$ to
$T_x$ has mean $p_x$.  The Poisson Stein equation with
$\|\Delta f\|_\infty\le1$ gives the displayed bound.

The polynomially weighted comparison already proved also makes
the scalar recurrence comparisons legitimate after deletion.
In particular, with
$\mathcal R_0(P;\mu)=\mathcal R(P;\mu,0)$,
\begin{align}
 &\big|\mathcal R(\Law Z;\lambda,\theta)
       -(1-r_N)\mathcal R_0(P_0;\mu_0)\big|
       \ll_\varepsilon N^{-4/5+\varepsilon},
       \label{p1:limits:background-residual}\\
 &\big|\mathcal R(\Law Z;\lambda,\theta)
       -\mathcal R_0(\Law W;\E W)\big|
       \ll_\varepsilon N^{-4/5+\varepsilon}.
       \label{p1:limits:erased-residual}
\end{align}
Indeed the operator in \eqref{p1:limits:raw-residual} is bounded
from the norm $\sum_k(1+k)|\sigma_k|$ to $\ell^1$, uniformly
for the present parameters; its value on the compound target
is zero.  The analogous statement holds for the Poisson
background.  Lemma~\ref{p1:limits:vector-reference}, the deletion
comparison, and \eqref{p1:limits:residual-parameters} prove
both formulas.  This argument uses fixed moments as well as
the distributional approximation.

\subsection{Removing the remaining active pair dependencies}

The recurrence norm weighted by $1/k$ has a useful stability
property under small changes of a law.  For any probability
$P$ on $\mathbb N$ and $\mu>0$ define
\[
 \mathcal V(P;\mu)=\sum_{k\ge1}\frac{|kP_k-\mu P_{k-1}|}{k}.
\]
It is always finite, even without a first-moment assumption.

\begin{lemma}[A recurrence norm equivalent to total variation]
\label{p1:limits:weighted-residual}
For every $P$ and $\mu>0$,
\begin{equation}
 e^{-\mu}\TV(P,\Pi_\mu)\le\mathcal V(P;\mu)
       \le2(1+\mu)\TV(P,\Pi_\mu).
 \label{p1:limits:weighted-equivalence}
\end{equation}
Moreover, if $0<\mu,\nu\le M$, then
\begin{equation}
 |\mathcal V(P;\mu)-\mathcal V(Q;\nu)|
       \le2(1+M)\TV(P,Q)+|\mu-\nu|.
 \label{p1:limits:weighted-stability}
\end{equation}
\end{lemma}
\begin{proof}
Write $a_k=P_k-\Pi_\mu\{k\}$ and
$b_k=a_k-(\mu/k)a_{k-1}$ for $k\ge1$.
The upper bound follows by summing absolute values.
The forward recurrence gives
\[
 a_k=\frac{\mu^k}{k!}a_0+
           \sum_{j=1}^k\frac{\mu^{k-j}j!}{k!}b_j.
\]
The sum over $k\ge j$ of the coefficient of $b_j$ is at most
$e^\mu$.  Since $\sum_ka_k=0$, it follows first that
$|a_0|\le\sum_j|b_j|$ and then that
$\sum_k|a_k|\le2e^\mu\sum_j|b_j|$.
All series are absolutely convergent by the same coefficient
bound.  This proves the lower estimate.  Applying the upper
operator bound to $P-Q$ and bounding
$\sum_{k\ge1}Q_{k-1}/k\le1$ proves stability.
\end{proof}

Let $E_N^{\rm all}$ be the graph on $G$ containing contacts
at distance exactly $L$ and all separated pairs with nonzero
joint relative kernel.  It contains $E_N$.  Define
\[
 \Xi^{\rm bad}=\sum_{e\in E_N^{\rm all}\setminus E_N}J_e,\qquad
 \Lambda^{\rm loc}=\sum_{x\notin G}J_x,\qquad
 \widetilde J_x=J_x\ind_{\{x\in G\}}
       \ind_{\{\sum_{y:\{x,y\}\in E_N^{\rm all}}J_y=0\}},\qquad
 \widetilde W=\sum_x\widetilde J_x.
\]
There is no degree bound imposed on $E_N^{\rm all}$.

\begin{proposition}[All active pair dependencies can be removed]
\label{p1:limits:all-pairs}
For every $\varepsilon>0$,
\begin{align}
 \tau_N:=\E\Xi^{\rm bad}+\E\Lambda^{\rm loc}
       &\ll_\varepsilon N^{-4/5+\varepsilon},
       \label{p1:limits:bad-mass}\\
 \TV(\Law\widetilde W,\Law\widehat W)&\le\tau_N,\qquad
 0\le\E\widehat W-\E\widetilde W\le2\tau_N.
       \label{p1:limits:all-pairs-cost}
\end{align}
In every realization, distinct retained original spaces are
pairwise direct and their starts are strictly separated.
The values of $\mathcal V$ at the respective actual means differ
by $O_\varepsilon(N^{-4/5+\varepsilon})$.  Thus, for a fixed
$0<\varepsilon<2/15$, the condition
\[
 \mathcal V(\Law\widetilde W;\E\widetilde W)=o(a_N)
\]
is equivalent to the Poisson approximation of the erased
background, and hence to
$\TV(\Law Z,H_{\lambda,\theta})=o(a_N)$.
\end{proposition}
\begin{proof}
A compatible separated dependent pair in $G$ has
$q_e=p^22^{\rho(e)}\ge2p^2$, whence
$q_e\le2|q_e-p^2|$.  The absolute bad-pair estimate therefore
bounds their positive event mass.  Contacts have mass
$N^{-1+o(1)}$ and locally deficient windows have total marginal
mass $N^{-1+o(1)}$.  This proves \eqref{p1:limits:bad-mass}.
Pointwise,
\[
 0\le\widehat W-\widetilde W
       \le2\Xi^{\rm bad}+\Lambda^{\rm loc},\qquad
 \{\widehat W\ne\widetilde W\}
       \subseteq\{\Xi^{\rm bad}+\Lambda^{\rm loc}>0\}.
\]
Markov's inequality gives \eqref{p1:limits:all-pairs-cost}.
Two retained starts cannot be closer than $L$ by incompatibility,
cannot be at distance $L$ by the filter, and cannot have a
nonzero pair kernel by its definition.  The means stay in a
fixed compact subset of $(0,\infty)$, so
Lemma~\ref{p1:limits:weighted-residual} transfers the criterion.
Use \eqref{p1:limits:one-all-tv} and
\eqref{p1:limits:erased-localization} for the final equivalence.
\end{proof}

Pairwise directness here is a property of the retained original
constraint spaces in a given realization.  The filtered
indicators themselves depend on the absence of all neighboring
events.  It does not make these indicators independent, or
exclude primitive dependencies among three or more original
spaces.  The stopped expansion supplies the further information
needed for their distribution.

\subsection{The cost of forcing private prime signs}

A small distance between distributions need not be visible in
the positive cost of a particular conditional coupling.  Here
this distinction can be quantified within the arithmetic model.
We use $\kappa(n)=P^+(\sf(n))$, $P^+(1)=1$.

To state its precise logarithmic scale, let $H_X=\log X$ and
use the following definitions from \cite[Proposition~4.2 and
Appendix~F.3]{pouly2026c}.  For $v\ge e$, let $u(v)\ge1$ solve
$e^{u(v)}/u(v)=v$ and put
\[
 D(v)=vu(v)-\operatorname{Ei}(u(v)),\qquad
 V_X=D(H_X/V_X),\qquad \nu_X=H_X/V_X.
\]
The indicated positive solution is unique for large $X$.
At the precision used below,
\begin{equation}
 2V_X=\sqrt{2H_X(\log H_X+\log\log H_X-\log2-2)}
          +o(\nu_X),\qquad
 \log\log X=o(\nu_X).
 \label{p1:limits:pivot-scale}
\end{equation}

\begin{lemma}[Reciprocal odd-prime pivots]
\label{p1:limits:reciprocal}
As $X\to\infty$,
\begin{equation}
 \sum_{n\le X}\frac1{\kappa(n)}
       =X\exp\{-2V_X+o(\nu_X)\}.
 \label{p1:limits:reciprocal-sum}
\end{equation}
The sum over $N\le n<2N$ is
$N\exp\{-2V_N+o(\nu_N)\}$.
\end{lemma}
\begin{proof}
The asymptotic of Erd\H{o}s--Ivi\'c--Pomerance
\cite[Theorem~1 and equation~(1.3)]{erdosivicpomerance1986}
gives this logarithmic formula for
$S^+(X)=\sum_{n\le X}1/P^+(n)$.
Its first correction is the explicit square root in
\eqref{p1:limits:pivot-scale}; the later terms are
$O(\sqrt{H_X\log H_X}(\log\log H_X)^2/(\log H_X)^2)=o(\nu_X)$.
Since $\kappa(n)\le P^+(n)$, this proves the lower bound.

For the upper bound, write uniquely $n=a^2v$ with $v$
squarefree and use positivity to obtain
\[
 \sum_{n\le X}\frac1{\kappa(n)}
       \le\sum_{a\le\sqrt X}S^+(X/a^2).
\]
The terms $a>e^{3V_X}$ are $O(Xe^{-3V_X})$.
On the remaining range $\log(X/a^2)=H_X-O(V_X)$ and
$\nu_{X/a^2}\sim\nu_X$.  The explicit function
$G(H)=\sqrt{2H(\log H+\log\log H-\log2-2)}$ has
$G'(H)=o(1)$ uniformly there.  The stated classical remainder,
without differentiation, is uniformly $o(\nu_X)$.
Thus
\[
 S^+(X/a^2)\le
 X\exp\{-2V_X+o(\nu_X)\}a^{-2+o(1)}.
\]
These weights have bounded sum.  This also explains the
square-part reduction following Proposition~F.5 in
\cite{pouly2026c}.

For a dyadic lower bound, send each $n\le N$ with
$\kappa(n)>2$ to the unique $2^jn\in[N,2N)$.
The pivot is unchanged, and each image has $O(\log N)$
preimages.  The omitted integers have squarefree part $1$ or
$2$, so number $O(\sqrt N)$.  The factor $\log N$ is absorbed
in $\exp\{o(\nu_N)\}$.  The upper bound follows from the sum
up to $2N$, whose scales agree at this precision.
\end{proof}

Consider a specific forcing for the signed exact-run field.
Let
$E=\lceil3V_N/\log2\rceil$, $B_*=L+E+2$, and
$Y=2B_*$.  Retain starts whose maximal support
$\{x-1,\ldots,x+L+E\}$ has every pivot larger than $Y$.
Let $W_{x,e,\sigma}$ be the indicator of the exact word
\[
 (-\sigma,\underbrace{\sigma,\ldots,\sigma}_{L+e},-\sigma),
 \qquad 0\le e\le E,\quad \sigma\in\{-1,1\}.
\]
Its mean is $\lambda_{x,e,\sigma}=p2^{-e-2}$.
For a source $a=(x,e,\sigma)$, fix every nonpivot prime sign
and set each private pivot of its actual word to make that
word occur.  The resulting full field $W^a$ has exactly
$\Law(W\mid W_a=1)$: conditional on the nonpivots, its raw
word bits were independent and fair and its prescribed
probability was constant.  This is the private-pivot
conditional coupling of \cite[Proposition~F.7]{pouly2026c}.
Its nonnegative coordinate cost is
\[
 \mathcal E_N
       =\sum_a\lambda_a
               \E\|W-(W^a-\boldsymbol e_a)\|_1.
\]

\begin{proposition}[Exact logarithmic cost of this forcing]
\label{p1:limits:private-exact}
For this specified signed forcing,
\begin{equation}
 \mathcal E_N=\exp\{-2V_N+o(\nu_N)\}.
 \label{p1:limits:private-exact-eq}
\end{equation}
Consequently this positive coordinate-by-coordinate cost is
not $O(N^{-\alpha})$ for any fixed $\alpha>0$.
\end{proposition}
\begin{proof}
A target can change only if it contains a multiple of a
source pivot.  Summing multiples, then reciprocals, bounds
the number of directed source--target possibilities by
$O(NB_*^2\sum_{n\le3N}1/\kappa(n)+NB_*^2)$.
The finite Palm ledger
\cite[Lemma~F.2 and equation~(F.6)]{pouly2026c} therefore gives
\[
 \mathcal E_N\ll Np^2+
 p^2\left\{NB_*+
       NB_*^2\sum_{n\le3N}\frac1{\kappa(n)}
       +R^{\rm val}_2(N,B_*)\right\}.
\]
Here $R^{\rm val}_2$ is the sum of $2^{\rho^{\rm val}}-1$
over separated maximal supports.  The full-value estimate
\cite[Proposition~3.24]{pouly2026c} gives
$R^{\rm val}_2(N,B_*)\le N^{5/3+o(1)}$, since
$B_*=\log_2N+o(\log N)$.  The type weights sum to at most
one after the common factor $p$ has been removed.
Lemma~\ref{p1:limits:reciprocal} proves the upper bound.

For a lower bound use only the type $e=0,\sigma=1$, of
mean $p_0=p/4$, and restrict both maximal supports to pivots
larger than $Y_0=(\log N)^5$.  Rankin's inequality with
$s=9/10$ gives
\[
 \#\{n\le3N:\kappa(n)\le Y_0\}
 \le(3N)^{9/10}\zeta(9/5)
                    \prod_{\ell\le Y_0}(1+\ell^{-9/10})
       =N^{9/10+o(1)}.
\]
Thus only $N^{9/10+o(1)}$ starts are excluded.
For each source choose $q=\kappa(x-1)\le e^{3V_N}$
and count targets with $q\mid y-1$.
The dyadic reciprocal estimate shows that these pairs number
$N^2\exp\{-2V_N+o(\nu_N)\}$.

Deleting excluded sources or targets costs at most
$N^{19/10+o(1)}$ pairs; the tail $q>e^{3V_N}$ costs
$O(N^2e^{-3V_N})$, and nearby pairs cost $O(NB_*)$.
The condition $q^2\mid y-1$ removes at most a proportion
$Y_0^{-1}$ of the main count plus $O(N)$.
For any other source pivot $q'$ and target offset, the
Chinese remainder theorem counts at most $N/(qq')+1$
targets.  All these extra intersections cost at most
$O(B_*^2/Y_0)$ times the main count plus $O(NB_*^2)$.
Finally the full-value host estimate
\cite[Proposition~3.7]{pouly2026c} excludes at most
$N^{3/2+o(1)}$ pairs with a nonzero raw relation.
Every removed term is negligible.

For each remaining pair the two raw word vectors $U,V$
are independent and uniform, while the forcing changes
only the first target bit, by the fair bit $U_0+b_0$.
The prescribed target word and its first-bit translate
are disjoint, each of probability $p_0$.  Therefore
$\E|W_{y,0,+}-W^a_{y,0,+}|=p_0$ exactly.
Multiplication by the source weight $p_0$ and summation
over the surviving pairs prove the lower bound in
\eqref{p1:limits:private-exact-eq}.
\end{proof}

The two original words in each lower-bound pair are
independent.  Their positive contribution is caused by
the chosen forcing.  It is therefore not a lower bound
for their dependence, for total variation, or for the
scalar cost after summing all changed coordinates with
their signs.

There is also an obstruction that does not depend on the
choice of an arithmetic anchor.  Let the maximal support
have length $B_*=O_C(\log N)$, and retain starts with at
most one pivot at most $Y=2B_*$.  For $\iota\in\{0,1\}$,
use the unsigned word of length $B_\iota=L+1+\iota$:
the start event when $\iota=0$, and exact excess zero
when $\iota=1$.  Its mean on this retained set is
$q_\iota=p2^{-\iota}$.  Choose an anchor as any deterministic
function of the integer arguments and their factorizations;
if there is a pivot at most $Y$ in the actual word, it must
be the anchor.  Preserve its value and force the other
private pivots to impose the word up to a global sign.
All other prime signs are left fixed.  Let
\[
 \mathcal E_{\iota,N}
 =\sum_xq_\iota\sum_{y\ne x}
       \E|W_{y,\iota}-W^{x,\iota}_{y,\iota}|
\]
be the same-type, off-site part of its coordinate cost.

\paragraph{A sufficient ordinary smooth-window census.}
\label{p18:limits:smooth-census}
For fixed $a>0$ put $Y=\lfloor N^a\rfloor$ and
$h_N=\lfloor\log N/(4\log2)\rfloor$.
The precise internally proved statement is
\[
 \#\left\{j\in I_N:
   \#\{n\in[j,j+h_N)\cap\Z:P^+(n)\leq Y\}\leq1\right\}
   =o_a(N).
\]
Here $P^+(1)=1$.  Equivalently, for every $\tau>0$, one cutoff
precedes all $N$ for which the displayed exceptional cardinality
is at most $\tau N$.
Only this sufficient statement is needed for
Proposition~\ref{p1:limits:all-anchors}.

\begin{proof}
Write $\chi_Y$ for the ordinary smooth indicator and
$g_Y(d)=\mu(d)\ind_{\gcd(d,\prod_{p\leq Y}p)=1}$.
The exact identity $\chi_Y=1*g_Y$ gives, for
$S_Y(x)=\sum_{n\leq x}\chi_Y(n)$,
\[
 S_Y(x)=\sum_{d\leq x}g_Y(d)\lfloor x/d\rfloor .
\]
For a common cutoff $K\geq x,z>0$,
\[
 \left|\frac{S_Y(x)}x-\frac{S_Y(z)}z\right|
 \leq R_Y(K)\left(\frac1x+\frac1z\right),
 \quad
 R_Y(K)=\#\{d\leq K:\gcd(d,\prod_{p\leq Y}p)=1\}.
\]
A primorial Selberg upper-bound sieve with fixed exponent
$\min(a,1/8)$ gives $R_Y(4N)\ll_a N/\log N$.
Together with the positive dyadic ordinary smooth density, this proves
uniform positivity of averages on intervals of length
$H_N=\lceil N/\sqrt{\log N}\rceil$.
The explicit difference from the dyadic mean is
$O_a((\log N)^{-1/2})$.

For $b>0$ let
\[
 \gamma_b(v)=\frac1{\sqrt{4\pi b}}e^{-v^2/(4b)},\qquad
 G_b(u)=\sum_{N\leq n<4N}\frac{\chi_Y(n)}n
                  \gamma_b(u+\log n),
\]
and put
$F(t)=\sum_{N\leq n<4N}\chi_Y(n)n^{-1-it}$.
The exact finite Gaussian identity is
\[
 \int_{\mathbb R}|G_b-G_c|^2\,du
 =\frac1{2\pi}\int_{\mathbb R}|F(t)|^2
                 (e^{-bt^2}-e^{-ct^2})^2\,dt .
\]
The finite sums are integrable, so the Gaussian Fourier identity
can be applied term by term.
Choose $b=(\log N/(DN))^2$ for fixed $D>0$, and
$c=(H_N/N)^2$.
The off-central Dirichlet estimate, applied to the two dyadic
blocks, makes the frequency integral small away from a fixed central
band.  On that band the difference of the Gaussian weights tends to
zero.  The remaining tail is bounded by the unconditional Gaussian
mean-value estimate
$\sqrt{\pi/b}/N+4\pi$ and its exponential tail factor.
Consequently $\int|G_b-G_c|^2=o_D(1)$ for the ordinary smooth
indicator used here.

The long smooth averages give $G_c(-\log X)\geq\kappa_a>0$
for $N\leq X$ and $X+1\leq2N$.
If a window $[j,j+h_N)$ has at most one smooth integer, then for
$j+h_N/2\leq X\leq j+h_N/2+1$ its short Gaussian average is at most
$O_D(1/\log N)+O(e^{-c_0D^2})$ for a fixed $c_0>0$.
The first term bounds the possible single interior integer; the
second bounds all exterior integers using their actual logarithmic
distance.  First choose $D$ so the exterior term is at most $\kappa_a/4$;
then choose the common integer cutoff so the interior term is at most
$\kappa_a/4$.

The middle unit cells of distinct integer windows have disjoint
negative-logarithmic images, each of measure at least $1/(4N)$.
Thus the number of interior bad windows is at most
$16N\int|G_b-G_c|^2/\kappa_a^2=o_a(N)$.
The omitted endpoint windows number at most $h_N=o(N)$.
This proves the sufficient census.
\end{proof}
This argument proves only the two-candidate census at the displayed
logarithmic length. It does not claim a Dickman relative-density
asymptotic or a general sharp short-average theorem. The census is
used in the forcing obstruction below and is not an input to the
scalar reconstruction.

\begin{proposition}[An obstruction for every such anchor]
\label{p1:limits:all-anchors}
For every fixed $0<a<1/10$, uniformly over these arithmetic
anchor rules,
\begin{equation}
 \mathcal E_{\iota,N}\ge
       (1-o_{C,a}(1))\,2^{-2\iota}\lambda_N^2N^{-a}.
 \label{p1:limits:anchor-lower}
\end{equation}
The same assertion with an additional factor $1/2$ holds
if, instead, the orientation is any fair randomized
function of the source word alone and only its private
pivots are changed.
\end{proposition}
\begin{proof}
Use the sufficient logarithmic-window census proved just above.  For every fixed $a>0$, with
$h_N=\lfloor\log N/(4\log2)\rfloor$, all but $o_a(N)$ integers
$j\in I_N$ have at least two ordinary $N^a$-smooth integers in
$[j,j+h_N)$.  On the critical band, this interval is contained
in each actual word under consideration; translating the initial
boundary contributes at most one further source.
Thus the actual word contains the required two candidates for all
but $o_a(N)$ sources.  This argument uses only the two-candidate
census, which is proved internally by Gaussian comparison and
localization.  The stronger Dickman relative-density asymptotic
of Matom\"aki--Radziwi{\l}{\l}
\cite[Corollary~6]{matomakiradziwill2016} gives another route to
the same sufficient conclusion.

Restrict again to maximal supports whose every pivot
exceeds $Y_0=(\log N)^5$.  The preceding Rankin bound
loses only $N^{9/10+o(1)}$ starts.  At most one of the
two small pivots can be the anchor, so at each of the
remaining $(1-o_a(1))N$ sources choose a forced pivot
$Y_0<q_x\le N^a$.  There are at least
$(1-o_a(1))N^{2-a}$ pairs with $q_x\mid y-1$.
The clean-pair deletions from the preceding proof apply:
excluded targets cost $N^{19/10+o(1)}=o(N^{2-a})$;
nearby targets, valuations at least two, and other forced
pivots cost $o(N^{2-a})$ as well.  Full raw dependence
has at most $N^{3/2+o(1)}$ hosts.

For a surviving pair, $U,V$ are independent uniform
vectors in $\mathbb F_2^{B_\iota}$.  If the selected
source index is $i$ and the anchor is $j\ne i$, the flip
is $U_i+U_j+b_i+b_j$, a fair bit independent of $V$.
The allowed relative word is the two-point coset
$\{b,b+\boldsymbol1\}$, disjoint from its translate by
the first coordinate vector.  The expected target
indicator change is therefore exactly $q_\iota$.
Weighting by the source mean gives
$q_\iota^2(1-o(1))N^{2-a}$, proving
\eqref{p1:limits:anchor-lower}.

For the final assertion, write the prescribed source as
$b+g_x(U,\xi)\boldsymbol1$, where $\xi$ is independent
auxiliary randomness and $g_x$ is fair for uniform $U$.
This guarantees the exact conditional source law even
after all nonpivot bits are fixed.  For two distinct
small-pivot indices $i,j$, the flips satisfy
\[
 \delta_i+\delta_j=U_i+U_j+b_i+b_j.
\]
The right side is fair, so
$\P(\delta_i=1)+\P(\delta_j=1)\ge1/2$.
Choose, before drawing signs, an index with flip
probability at least $1/4$.  Repeat the clean-pair count,
now excluding every other source pivot from the target.
Since $V$ is independent of $(U,\xi)$, each target
change has expectation at least $2q_\iota/4$.
This gives the additional factor $1/2$.
\end{proof}

For any proposed power $\alpha>0$, choose a fixed
$a<\min(\alpha,1/10)$ in
\eqref{p1:limits:anchor-lower}; then
$N^\alpha\mathcal E_{\iota,N}\to\infty$.
The quantifier fixes $a$ before $N$ and invokes no
uniform short-interval smoothness estimate with a growing
parameter.  The result concerns only forcings that change
the indicated private pivots.  An orientation using the
outside environment, a different full conditional coupling,
or a signed scalar comparison is not covered.


\section{Offset geometry and residual relation spaces}
\label{p1:offset:section}

The group budget used in the main proof keeps the product of primes
sharing an offset. Two further geometric consequences explain what such
products can force. A line of offsets has a bounded total prime product;
a large product on three offsets gives a plane involving all three
windows. These statements also give a more detailed description of the
relation space after a covering set of coordinates has been removed.

Throughout this section $B=L+1$, $I_N=[N,2N)\cap\mathbb Z$, and
$H=C_0B$ for a fixed $C_0>1$. We consider three distinct regular windows
with starts $X=(x_1,x_2,x_3)\in I_N^3$ and pairwise direct relative
spaces. Their offsets lie in
$\mathcal O=\{-1,0,\ldots,L-1\}$, whose diameter is less than $B$.
Every prime $\ell>H$ having odd valuation in all three windows has a
unique offset $u_\ell\in\mathcal O^3$. Such primes are called shared
in the statements below. No independence of the three-window system
is assumed. Whenever a primitive mass is mentioned, it is the sum of
the original affine joint probabilities over ordered primitive triples.
Each of these probabilities is at most $p^2$.

\subsection{Whole prime products on lines and planes}

\begin{lemma}[The prime product on an offset line]
\label{p1:offset:line-product}
Let $\mathcal L$ be a line in offset space, and let
$\mathcal Q_{\mathcal L}$ be the shared primes whose offsets lie on it.
If their distinct offsets consist of one point, then
\[
 \prod_{\ell\in\mathcal Q_{\mathcal L}}\ell\le3N.
\]
If there are at least two distinct offsets, the line has a primitive
integer direction $d$ and
\begin{equation}
 \prod_{\ell\in\mathcal Q_{\mathcal L}}\ell
       \le6N\|d\|_\infty\le6NB.
 \label{p1:offset:line-product-bound}
\end{equation}
For $Y\ge H$, the number of distinct offsets on this line supplied by
primes greater than $Y$ is at most
\[
 M_Y=1+\frac{\log(6NB)}{\log Y}=O_{C_0,C}(B/\log B)=o(B).
\]
\end{lemma}
\begin{proof}
At a single offset the product of the distinct primes divides each
coordinate of $X+u$, all of which are at most $3N$ for large $N$.
Otherwise fix one occurring offset $u_0$. Every integer point on the
line has the form $u_0+td$, with $t\in\mathbb Z$. Each of its shared
primes divides the same integer vector
\[
 d\times(X+u_\ell)=d\times(X+u_0).
\]
This vector is nonzero. If it were zero, positivity of $X+u_0$ would
make all coordinates of $d$ nonzero. Two distinct offsets would give
two positive proportional value vectors, with two distinct values in
each block. Their projections on any two blocks would give equal
nonzero square classes of even selections. Injectivity and the direct
sum of those two spaces exclude this relation.

A nonzero coordinate of the displayed cross product has absolute value
at most $6N\|d\|_\infty$ and is divisible by every prime in the
product. Two offsets in $\mathcal O^3$ give $\|d\|_\infty<B$.
This proves \eqref{p1:offset:line-product-bound}. Choose one prime per
distinct offset and take logarithms to obtain the last assertion;
the extra one covers a line carrying only one offset.
\end{proof}

\begin{lemma}[An integer product tail]
\label{p1:offset:integer-tail}
For fixed $s>1$ and $T\ge1$,
\[
 \sum_{\substack{a,b\ge1\\ab>T}}(ab)^{-s}
       \ll_s T^{1-s}\log(2T).
\]
\end{lemma}
\begin{proof}
For $a\le T$, the inner sum is at most
$C_s(T/a)^{1-s}$, giving
$C_sT^{1-s}\sum_{a\le T}a^{-1}$ after multiplication by $a^{-s}$.
For $a>T$, use the full sum $\zeta(s)$ and
$\sum_{a>T}a^{-s}=O_s(T^{1-s})$.
\end{proof}

Let $\mathscr P_B$ be the family of distinct affine planes
\[
 A\cdot X=h,\qquad
 A=(u_2-u_1)\times(u_3-u_1)\ne0,
 \qquad h=-A\cdot u_1,
 \qquad u_1,u_2,u_3\in\mathcal O^3.
\]
There are at most $B^9$ such planes, and these representations satisfy
$\|A\|_\infty\le2B^2$ and $|h|\le6B^3$.

\begin{proposition}[A forced plane uses all three coordinates]
\label{p1:offset:forced-plane}
Two distinct shared primes satisfy $\ell\ell'\le6NB$.
If three distinct shared primes have product greater than $18NB^2$,
then $X$ belongs to a plane of $\mathscr P_B$ whose normal satisfies
$A_1A_2A_3\ne0$. Every further shared prime $\ell>12B^3$ has
\begin{equation}
 A\cdot u_\ell=-h.
 \label{p1:offset:subsequent-offset}
\end{equation}
The primitive mass of triples belonging to two distinct planes of
$\mathscr P_B$ is $O_C(N^{-1+o(1)})$.
\end{proposition}
\begin{proof}
For two primes, either they share a value in one of two selected blocks,
so their product is at most $3N$, or their occurrences are distinct in
both. In the latter case their product divides
$nm'-n'm$. Its vanishing gives the forbidden pair relation
$nn'mm'=(nm')^2$. Since $n,m\le3N$ and both within-block
differences have absolute value less than $B$, its nonzero absolute
value is at most $6NB$.

For three primes, their product divides
\[
 \det(X+u_1,X+u_2,X+u_3)
 =(X+u_1)\cdot((u_2-u_1)\times(u_3-u_1)),
\]
whose absolute value is at most $18NB^2$. Hence it vanishes.
The offsets cannot lie on a line, by
Lemma~\ref{p1:offset:line-product}, including the case of repeated
offsets. Thus $A\ne0$ and $A\cdot X=-A\cdot u_1$.

Suppose, for example, that $A_1=0$. The other two coefficients must
both be nonzero, since $A\cdot(X+u_i)=0$ cannot have just one
nonzero coefficient at a positive vector. The projections of all three
value vectors onto coordinates two and three are proportional. If two
projections differ, both their coordinates differ, giving a forbidden
pair relation. If none differ, the two offset differences are parallel
to the first axis, contradicting $A\ne0$. The other coordinates are
handled in the same way.

A further prime divides $A\cdot(X+u_\ell)=h+A\cdot u_\ell$,
an integer of absolute value at most $12B^3$. A larger prime forces
it to vanish. Finally, two distinct planes are disjoint or meet in a
line, which has at most $N$ integer points in $I_N^3$. There are at
most $B^{18}$ pairs of planes. Multiplication by the upper probability
$p^2$ gives the mass assertion.
\end{proof}

\begin{lemma}[Many planar supports force a small normal]
\label{p1:offset:normal-budget}
Let $A\cdot u=-h$ be a plane of offsets, with $A\in\mathbb Z^3$
primitive and nonzero, and $h\in\mathbb Z$. Let $m_{A,Y}$ count its
distinct offsets supplied by shared primes greater than $Y\ge H$.
Then
\begin{equation}
 m_{A,Y}\le M_Y+\frac{6B^2}{\|A\|_\infty}.
 \label{p1:offset:normal-count}
\end{equation}
If the hyperedges given by these offsets require at least $\delta L$
vertices to cover them, where $\delta>0$ is fixed, then for large $N$,
\[
 \|A\|_\infty\le24B/\delta,
 \qquad |h|\le72B^2/\delta.
\]
There are consequently $O_\delta(B^5)$ possible such integer planes.
\end{lemma}
\begin{proof}
The lattice $\Lambda=A^\perp\cap\mathbb Z^3$ has covolume
$\|A\|_2$. Indeed, a lattice basis $d,e$ can be completed by an
integer vector $w$ with $A\cdot w=1$ to a basis of $\mathbb Z^3$,
so $d\times e=\pm A$.
Choose a shortest nonzero lattice vector $d$, of length $\ell$.
It is primitive and extends to a lattice basis. Unless the cube has
at most one integer point on the plane, $\ell\le\sqrt3B$.
Parallel lattice lines in direction $d$ have spacing $\|A\|_2/\ell$.
At most $1+\sqrt3B\ell/\|A\|_2$ meet the cube, and each contains
at most $1+\sqrt3B/\ell$ integer points. By
Lemma~\ref{p1:offset:line-product}, each carries at most $M_Y$ supports.
It follows that
\[
 m_{A,Y}\le
 \left(1+\frac{\sqrt3B\ell}{\|A\|_2}\right)
       \min\left(M_Y,1+\frac{\sqrt3B}{\ell}\right)
 \le M_Y+\frac{6B^2}{\|A\|_\infty}.
\]
The omitted one-point case satisfies the same bound.
The covering number is no larger than the number of supports.
Since $M_Y=o(B)$ and $L=B-1$, a covering number at least $\delta L$
implies the asserted bound for $A$, with room in the constant.
The existence of an offset gives $|h|\le3B\|A\|_\infty$.
Counting the normals and the constants gives $O_\delta(B^5)$ planes.
\end{proof}

\subsection{Root separation and counting on a plane}

\begin{lemma}[At most one common root]
\label{p1:offset:root-separation}
Let $y=\alpha x+\beta$, with rational $\alpha\ne0$, pass through
a positive pair of regular windows with direct relative spaces.
Choose offsets distinct within each block. The two root lists
$-i$ and $-(\beta+j)/\alpha$ have at most one common element.
If these offsets are partitioned among equations with disjoint supports,
discarding at most two equations leaves squarefree polynomials with
pairwise disjoint root sets after substitution of the line.
\end{lemma}
\begin{proof}
A common root says $\beta+j=\alpha i$, and hence
$y+j=\alpha(x+i)$. Positivity makes $\alpha>0$.
Two common roots would give two distinct values on each side in the
same ratio. Their even square classes are nonzero by injectivity and
equal, contradicting directness. The sole possible common root belongs
to at most two equations. Removing those equations removes both
repeated roots within an equation and roots shared between equations.
\end{proof}

\begin{proposition}[Many transverse square products on prescribed planes]
\label{p1:offset:planar-components}
Fix $\eta,\varepsilon>0$ and $K>0$. Let a family of at most $B^K$
planes have integer equations
$A_1x+A_2y+A_3z=h$, with nonzero normal and all coefficients of
absolute value at most $B^K$.
The number of primitive triples on these planes admitting at least
$\eta B$ disjoint sets of values, each meeting all three windows and
having square product modulo primes at most $H$, is
$O_{C,C_0,\eta,K,\varepsilon}(N^{1+\varepsilon})$.
Their primitive mass is $O(N^{-1+\varepsilon})$.
\end{proposition}
\begin{proof}
A plane with only one nonzero coordinate coefficient cannot meet the
dyadic region for large $N$. Otherwise permute coordinates so that
$A_1A_2\ne0$. The budget of $3B$ values leaves at least $\eta B/2$
sets of size at most $D_0=\lceil6/\eta\rceil$.
Select $k+2$ of them, where $k$ is fixed below. Their equations are
\[
 w_\nu^2=d_\nu P_\nu(x)Q_\nu(y)R_\nu(z),
\]
with nonempty split products, disjoint root lists in each window,
total degree at most $D_0$, and squarefree $H$-smooth $d_\nu$.
Their offsets and coefficients have $N^{o(1)}$ possibilities for fixed
$k,D_0$, and $d_\nu\le\prod_{\ell\le H}\ell=N^{O_{C,C_0}(1)}$.

Fix $z$. The plane gives
$y=\alpha x+\beta$, where
$\alpha=-A_1/A_2$ and $\beta=(h-A_3z)/A_2$.
Lemma~\ref{p1:offset:root-separation} leaves $k$ equations whose
polynomials $P_\nu(x)Q_\nu(\alpha x+\beta)$ have simple disjoint
roots. Their square classes are geometrically independent, since
valuation at a root belonging to one equation detects that equation.
They define a geometrically integral curve of degree at least $2^k$
in the coordinates $(x,w_1,\ldots,w_k)$: its generic projection to
the $x$-line has that degree. Its degree is also bounded above in terms
of $k,D_0$ only.

Clearing powers of $A_2$ from the equations preserves the original
integer coordinates $w_\nu$. Their heights are at most $N^{q_0}$
for fixed $q_0=q_0(C,C_0,D_0,K)$ independent of $k,z$ and the motifs;
this follows from the original three-variable products.
Pila's uniform curve estimate \cite[Theorem~A]{pila1995} therefore
gives $O(N^{q_0/2^k+q_0\zeta})$ points for each motif and fixed $z$.
Choose $k$ sufficiently large, then $\zeta$ sufficiently small, to
make this exponent less than $\varepsilon/3$. Sum the $O(N)$ values
of $z$, the $N^{o(1)}$ motifs and the $B^K$ planes. Finally use
$q(\{x,y,z\})\le p^2$ to obtain the mass bound.
\end{proof}

\subsection{Deleting a cover of the three-block constraints}

For a separated primitive triple, let $\mathcal H_3$ have the $3B$
formal window positions as vertices and the three occurrences of each
shared prime as a hyperedge. Repeated supports may be identified.
Let $D$ be any vertex cover of these edges, with $t=|D|$.
For the relative relation space $R$, define
$R_D=\{v\in R:v|_D=0\}$.

\begin{lemma}[Covering cost and free components]
\label{p1:offset:cover-rank}
One has
\[
 \dim R\le t+\dim R_D.
\]
After coordinates in $D$ are set to zero, the equations from primes
greater than $H$ form a graph with possibly pinned vertices. If $c_3(D)$
counts its unpinned components meeting all three blocks, then
\[
 \dim R_D\le c_3(D)+O_{C,C_0}(B/\log B),
\]
uniformly in $D$. The product of the values in every free component
is a square modulo primes at most $H$.
\end{lemma}
\begin{proof}
The first statement is rank--nullity for the coordinate restriction
$R\to\mathbb F_2^D$. Every row previously of weight three meets
$D$, so its remaining weight is at most two. A row of weight two
equates its coordinates, and a row of weight one pins its coordinate
to zero; pinning propagates through the component.

A free isolated vertex was already $H$-defective. Otherwise a large
prime incident to it would remain either an edge or a pin. There are
$O(B/\log B)$ such vertices by Lemma~\ref{p1:arith:defects}.
A nonsingleton component cannot be contained in one block because
edges join distinct blocks. For a fixed pair of blocks, there are at
most $\pi(H)+2$ free components meeting exactly that pair: otherwise
their small-prime valuation vectors, augmented by the two block
parities, have a nonzero zero-sum combination. Disjointness of the
components would then give a forbidden relative pair relation.
These components also number $O(B/\log B)$.
All remaining free components meet the three blocks and each has
one binary coordinate before the additional equations are imposed.
This proves the dimension estimate. Every large-prime row meets a
free component zero or two times, proving its product assertion.
\end{proof}

\begin{proposition}[A six-vertex budget for the residual space]
\label{p1:offset:six-vertex}
For every fixed $\varepsilon>0$, there is a set of primitive triples
of total mass $O_{C,C_0,\varepsilon}(N^{-1+\varepsilon})$ outside
which every cover $D$ simultaneously satisfies
\begin{equation}
 \dim R_D\le\frac{3B-|D|}{6}
           +O_{C,C_0}(B/\log B)+O_{C,C_0,\varepsilon}(1).
 \label{p1:offset:six-vertex-bound}
\end{equation}
For fixed $\eta>0$, triples admitting some cover with
$\dim R_D\ge\eta L$ number $O_{C,C_0,\eta,\varepsilon}(N^{3/2+\varepsilon})$.
On a prescribed family of planes as in
Proposition~\ref{p1:offset:planar-components}, their number is
$O(N^{1+\varepsilon})$ and their primitive mass is
$O(N^{-1+\varepsilon})$.
\end{proposition}
\begin{proof}
For any fixed positive type $(a,b,c)$ with a singleton, sufficiently
many disjoint square products of that type have $O(N^{1+\varepsilon})$
primitive hosts. Fix the third centre and apply the degree-one part
of Lemma~\ref{p1:arith:geometric} to the other two, putting the
singleton in its linear variable. The specialized coefficients have
fixed polynomial height. The original smooth coefficients and offsets
have $N^{o(1)}$ choices; summing the third centre proves the count.
Multiplication by $p^2$ gives mass $O(N^{-1+\varepsilon})$.

Apply this exclusion to the finitely many permutations of
$(1,1,1)$, $(2,1,1)$, $(3,1,1)$ and $(2,2,1)$.
Outside their union, the number $s$ of free three-block components
of size at most five is bounded, for every cover $D$.
All other three-block components use at least six surviving vertices,
and the small ones use at least three. Therefore
\[
 6c_3(D)-3s\le3B-|D|.
\]
Lemma~\ref{p1:offset:cover-rank} proves
\eqref{p1:offset:six-vertex-bound}.

If $\dim R_D\ge\eta L$, the same lemma supplies a positive
proportion of disjoint three-block square products. Select a fixed
number of bounded-size ones. Fixing one centre and using the general
two-parameter estimate of Lemma~\ref{p1:arith:geometric} gives
$N^{3/2+\varepsilon}$ hosts after the same coefficient and offset
enumeration. On the prescribed planes use instead
Proposition~\ref{p1:offset:planar-components}.
The covers are never enumerated: their existence supplies the disjoint
sets, of which only a fixed number are selected for counting.
\end{proof}

\section{Domain dependence and polylogarithmic arithmetic cutoffs}
\label{p1:consequences:section}

Two features of the arithmetic estimates have uses beyond the dyadic
law theorem. Changing the interval of starting points changes the dominant
rational channel. Revealing more windows changes the set of available
prime coordinates, and the resulting loss must be measured by rank,
not just by the number of exceptional positions. We record these facts,
then explain how to cut off large ternary supports without assuming
uniformity of a fixed-order counting theorem.

\subsection{The aligned profile on a macroscopic interval}

For an interval of starts $A$, let $\sigma(x,y)$ be the dimension
of its rational even-unit code and put
\[
 R_{\rm rat}(A,B)=\sum_{\substack{x,y\in A\\|x-y|>L}}
             (2^{\sigma(x,y)}-1).
\]
The rational code is zero unless there are at least two exact units.
Let $R_{\rm al}(A,B)$ and $C_{\rm al}(A,L)$ be the sums over
aligned ordered separated pairs of $2^{\rho(x,y)}-1$ and
$\eta(x,y)2^{\rho(x,y)}-1$, respectively. Alignment has the
definition in Section~\ref{p1:arith:section}; $B=L+1$.

\begin{proposition}[Profiles depend on the interval]
\label{p1:consequences:macro}
In the dyadic critical band, writing $B=3k+r$, $0\le r<3$,
\begin{equation}
 R_{\rm rat}(I_N,B)\sim R_{\rm al}(I_N,B)
       \sim \frac{3+r}{18}NB2^k,
 \qquad
 R_{\rm al}(I_N,B)-R_{\rm rat}(I_N,B)
       =O_{C,\varepsilon}(N^{5/4+\varepsilon}).
 \label{p1:consequences:dyadic-positive}
\end{equation}
Fix $0<\delta<1$, let $A_{M,\delta}=[\lceil M^\delta\rceil,M)
\cap\mathbb Z$, and suppose $|L-\log_2M|\le C$.
For $B=2k+r$, $0\le r<2$,
\begin{equation}
 R_{\rm rat}(A_{M,\delta},B)\sim R_{\rm al}(A_{M,\delta},B)
       \sim\frac{2+r}{4}MB2^k,
 \label{p1:consequences:macro-positive}
\end{equation}
and their difference is $O_{C,\delta,\varepsilon}(M^{4/3+\varepsilon})$.
For the signed profile, write instead $L=2m+s$, $0\le s<2$.
Then
\begin{align}
 C_{\rm al}(A_{M,\delta},L)
   &=(M-2M^\delta)S_{1,2}(L)
           +O_{C,\delta,\varepsilon}(M^{4/3+\varepsilon}),
          \label{p1:consequences:macro-signed}\\
 S_{1,2}(L)
   &=2^{m-1}\big[(s+1)m+s^2+s+4\big]-3L+1
       \qquad(m\ge1),\label{p1:consequences:S12}\\
 C_{\rm al}(A_{M,\delta},L)
   &\sim\frac{s+1}{4}ML2^m.
          \label{p1:consequences:macro-equivalent}
\end{align}
\end{proposition}

\begin{proof}
In the channel $1:2$, each orientation permits
$M/2-M^\delta+O(B)$ translations of the unit parameter.
For $h=y-2x\in\{-1,\ldots,B-2\}$, the exact units have
the size of the congruence class of $h$ modulo two in this
interval. If these two sizes are $m_0,m_1$, their plateau
contribution is
\[
 \sum_{j=0}^1m_j(2^{m_j-1}-1)
       =2^{k-1}\big[(2+r)k+2r\big]-B.
\]
Outside the plateau, the unit count drops linearly with the
distance, so the sum of the weights is $O(2^{B/2})$.
Multiplying by both orientations proves the main term in
\eqref{p1:consequences:macro-positive}. The height-one
channel has no exact unit for separated windows; its residual
cost is $M^{1+o(1)}$. All other reduced heights are at least
three and contribute at most $M B^{O(1)}2^{L/3+o(L)}$.

At a fixed height, the full kernel differs from the rational
code only when the normalized windows contain a defective
integer at a threshold proportional to $B$. There are
$M^{1/2+o(1)}$ such hosts, by the square--smooth decomposition
in Lemma~\ref{p1:arith:defects}. At height two their weight
is $2^{L/2+o(L)}$, so the residual costs $M^{1+o(1)}$.
The proof of Lemma~\ref{p1:arith:alignedrank} applies uniformly
down to $x,y\ge M^\delta$: its Runge distance becomes
$\gg_\delta\log M/\log B$, while its redundancy remains
$O(B/\log B)$. Thus the other aligned heights have the
stated $M^{4/3+\varepsilon}$ total bound, including all
residual relations. This proves the comparison of the two
positive sums.

Retaining the predecessor sign replaces each placement factor
$L+1-D$ by $g_L(D)=L-1-D$. Hence the signed contribution
per translation is $S_{1,2}(L)$; the identity
\eqref{p1:arith:channelidentity} gives
\eqref{p1:consequences:S12}. The endpoints cost
$O(B^3 2^{L/2})$ and the other channels and residual kernels
have the same bounds in absolute value as above. This proves
\eqref{p1:consequences:macro-signed} and its equivalent.

For the dyadic assertion, the height-two channel is confined
to the endpoints. The channel $2:3$ has translation count
$N/3+O(B)$ over the two orientations, and three residue
classes of total size $B=3k+r$. Its plateau weight is
$2^{k-1}[(3+r)k+2r]-B$. The height-one residual costs
$N^{1+o(1)}$; the remaining heights are at least four and
cost $N^{5/4+o(1)}$. The fixed-height residual costs
$N^{5/6+o(1)}$. These estimates prove
\eqref{p1:consequences:dyadic-positive}.
\end{proof}

The residue classes of $B$ in the positive formulas and of
$L$ in the signed formulas are different. They cannot be
interchanged. With $\lambda_M=M2^{-L}$, the last formula gives
\[
 p^2C_{\rm al}(A_{M,\delta},L)
   \sim\frac{s+1}{4\,2^{s/2}}\lambda_M^{3/2}L M^{-1/2}.
\]
This is an aligned-pair statement. It does not assert a
macroscopic analogue of Theorem~\ref{p1:thm:main}. The actual
mean parameter for this interval would be $|A_{M,\delta}|p$,
whereas $\lambda_M$ expresses the ambient critical scale.

\subsection{How exceptional vertices bound the conditional rank}

Fix $j$ old windows and let $W$ be the square-class span of
all their raw values. Let $\mathcal P$ contain the primes at
most $B$ and every prime larger than $B$ dividing an old value.
Put
\[
 \mathcal S(\mathcal P)=\{n\le3N:\sf(n)\text{ is supported on }
                    \mathcal P\},\qquad
 t_x=|V_x\cap\mathcal S(\mathcal P)|.
\]

\begin{lemma}[Conditional rank loss]
\label{p1:consequences:conditional-rank}
Let $\rho_x$ be the dimension of the kernel of the map from
$E_x$ to square classes modulo $W$. Then
\[
 \rho_x\le(t_x-1)_+.
\]
Conditional on any realizable vector of the old raw signs,
the probability of a new start is either zero or $p2^{\rho_x}$.
It equals $p$ if $t_x\le1$.

Let $O$ be the set of starts at distance at most $L$ from an
old start, and put
$D_T=\#\{x\in I_N\setminus O:t_x>T\}$ for $1\le T\le L$.
Uniformly over the old signs,
\begin{equation}
 \E[Z_N\mid\mathcal F_{\rm old}]
   \le\lambda_N+|O|+p2^{T-1}B|\mathcal S(\mathcal P)|+D_T.
 \label{p1:consequences:conditional-budget}
\end{equation}
For fixed $j$ and $T=\lfloor(1-\delta)L\rfloor$, $0<\delta<1$,
the third term is $N^{1/2-\delta+o_{C,j}(1)}$.
\end{lemma}

\begin{proof}
A nonzero element of $E_x$ is an even nonempty selection of
values. If its square class lies in $W$, every prime larger
than $B$ with odd valuation in a selected value must divide
an old value: it cannot divide a second value of $V_x$.
Thus every selected value lies in $\mathcal S(\mathcal P)$.
The kernel injects into the even-subset space on these $t_x$
values, of dimension $(t_x-1)_+$. The conditional distribution
of the new relative signs is uniform on an affine image of
dimension $L-\rho_x$, proving the probability statement.

Use probability at most one on $O$ and on $\{t_x>T\}$.
On $\{t_x\le1\}$ the probability is $p$; on
$\{2\le t_x\le T\}$ it is at most $p2^{T-1}$.
Moreover $\#\{x:t_x\ge2\}\le\sum_xt_x
\le B|\mathcal S(\mathcal P)|$.
This proves \eqref{p1:consequences:conditional-budget}.
The Euler-product bound in Lemma~\ref{p1:limits:branching}
gives $|\mathcal S(\mathcal P)|=N^{1/2+o_{C,j}(1)}$ when
$j$ is fixed; substituting $T$ proves the final assertion.
\end{proof}

For example, if $Q(\mathcal P)$ counts pairs $n<m$ in
$\mathcal S(\mathcal P)$ with $m-n\le L$, double counting gives
\[
 D_T\le\frac{B Q(\mathcal P)}{\binom{T+1}{2}}.
\]
No independence of the two exceptional integers is assumed.
The rank factor explains why merely counting dependent windows
throws away useful conditional information.

\subsection{Cutting off large primes and large support groups}

Use the triples of blocks and the groups $G_u$ from
Proposition~\ref{p1:arith:ternary-budget}. The offsets and
width constants are fixed; a polynomial number in $B$ of
such patterns is allowed. Let $\mathcal Q$ be the set of
primes greater than the diameter threshold that occur with
odd valuation in all three blocks. A positive mass below
means a sum of nonnegative host weights, each at most
$N^{-\beta+o(1)}$, with $\beta$ fixed before $N$ and uniform
$o(1)$. For an original
primitive triple one can take $\beta=2$: two of its
window spaces are in direct sum.

\begin{lemma}[A Rankin bound for offset patterns]
\label{p1:consequences:rankin-patterns}
For fixed $c,D_0>0$, $d\in\{2,3\}$ and $X\ge1$, put
$s=1+(d-1)/c$. Then
\begin{equation}
 \sum_{\substack{Q\le X\ {\rm squarefree}\\P^-(Q)>B^c}}
       (D_0B)^{d\omega(Q)}
 \le X^s\exp\{O_{c,D_0}(B/\log B)\}.
 \label{p1:consequences:rankin-primes}
\end{equation}
If instead the factors are arbitrary distinct integers
$g>B^c$, then, for every fixed $0<\zeta<1$ and
$s_\zeta=1+(d-1+\zeta)/c$,
\begin{equation}
 \sum_{\substack{A\subset\mathbb N\ {\rm finite}\\
                g>B^c\ (g\in A),\ \prod_{g\in A}g\le X}}
       (D_0B)^{d|A|}
 \le X^{s_\zeta}\exp\{O_{c,d,\zeta,D_0}(B^{1-\zeta})\},
 \label{p1:consequences:rankin-groups}
\end{equation}
where $D_0$ is fixed.
\end{lemma}

\begin{proof}
Rankin's inequality bounds the first sum by
$X^s\prod_{q>B^c}(1+(D_0B)^dq^{-s})$.
Chebyshev's prime estimate and partial summation give
$\sum_{q>Y}q^{-s}\ll_s Y^{1-s}/\log Y$.
Since $d+c(1-s)=1$, taking the logarithm proves
\eqref{p1:consequences:rankin-primes}. For integer factors
the analogous product is
$\prod_{g>B^c}(1+(D_0B)^dg^{-s_\zeta})$.
Bound its logarithm by the corresponding sum over integers;
the exponent of $B$ is now $1-\zeta$.
\end{proof}

\begin{proposition}[Polylogarithmic confinement]
\label{p1:consequences:confinement}
For every fixed $c>0$ and $\varepsilon>0$, there is a fixed
integer $k$ such that, outside a mass
\begin{equation}
 O_{c,\beta,\varepsilon}
 \left(N^{-\min\{\beta-6/5,\,\beta-1-2/c\}+\varepsilon}\right),
 \label{p1:consequences:confinement-error}
\end{equation}
one has
\begin{equation}
 \#\{u:G_u>B^c\}
       \le k-1+\left\lfloor\frac{\log N}{c\log B}\right\rfloor=o(L).
 \label{p1:consequences:confinement-count}
\end{equation}
The same assertion holds for $\#\{q\in\mathcal Q:q>B^c\}$.
On a family of $B^{O(1)}$ fixed affine planes for the centers,
replace $2/c$ in the exceptional mass by $1/c$.
For original triples, $c>6$ is therefore sufficient for an
exceptional mass $o(LN^{-2/3})$; on these planes $c>3$
suffices. These are sufficient thresholds, not optimality claims.
\end{proposition}

\begin{proof}
First, the fixed-$k$ budget of
Proposition~\ref{p1:arith:ternary-budget} gives a mass
$O(N^{6/5-\beta+\varepsilon/3})$ for $\mathcal G_k>N$.
The identical budget holds for the product of the $k$ largest
individual primes. Indeed CRT fixes each center modulo their
product, exactly as in \eqref{p1:arith:groupcrt}. In the
three-prime tail, repeated offset supports cannot cause a
new case: if all supports coincide the product divides one
value; if exactly two coincide, or three are collinear, it
divides a nonzero coordinate of $d\times(X+u)$, of size
$O(NB)$. A zero cross product would give a forbidden pair
relation, as in the proof of the group budget. The large
product therefore forces a genuine plane. The CRT tail and
the induction on fixed $k$ apply unchanged.

Outside this first exception at most $k-1$ factors exceed
$N^{1/k}$. If the product of the remaining factors in
$(B^c,N^{1/k}]$ exceeds $N$, multiply them in a deterministic
order until their product first exceeds $N$. The resulting
$Q$ satisfies $N<Q\le X=N^{1+1/k}$.
For individual primes, choosing their three offsets costs
$(D_0B)^{3\omega(Q)}$ for a fixed width constant $D_0$.
CRT gives $O(1)$ choices for each center.
Equation~\eqref{p1:consequences:rankin-primes} bounds the
weighted mass by
\[
 N^{-\beta+(1+1/k)(1+2/c)+o(1)}.
\]
For groups, each factor has one support, costing at most
$(D_0B)^3$. Their pairwise coprimality is used in CRT first;
it may then be dropped in the positive sum over factor sets.
Equation~\eqref{p1:consequences:rankin-groups} gives instead
\[
 N^{-\beta+(1+1/k)(1+(2+\zeta)/c)+o(1)}.
\]
Choose $\zeta$ sufficiently small and then $k$ sufficiently
large, both before $N$, to absorb the extra exponent into
$\varepsilon$. If the centers lie on one of the specified
planes, enumerate only the two free coordinates and their
offsets; the third is fixed by the plane. This uses $d=2$
and gives $1/c$ in place of $2/c$. In the remaining
configurations the product is at most $N$, so the number
of factors greater than $B^c$ is at most
$\lfloor\log N/(c\log B)\rfloor$.
\end{proof}

The cutoff applies to whole support groups, even if a group
contains many primes. It can therefore be imposed before a
coordinate regularization. Reversing that order could merge
support labels and invalidate a bound stated only for the
original labels.

\section{Further uniform consequences of the arithmetic estimates}
\label{p1:uniform:appendix}

The following consequences record uniformities useful when the
counting argument is applied to selected subfamilies or to
conditional prime-sign systems. Their parameters and the
quantities being summed are kept explicit.

\subsection{Full-support relations at all cardinalities}

For \(r\ge2\), let \(\mathcal H_{r,N}\) be the ordered
lists of distinct, mutually separated regular windows for
which
\[
 s_i\in S_{x_i}\setminus\{0\},\qquad
                     s_1+\cdots+s_r=0
\]
has a solution. Minimality and affine compatibility are
not required.

\begin{proposition}[A weighted count of every full support]
\label{p1:uniform:full-support}
For fixed \(s\ge2\), \(u>0\), and \(\varepsilon>0\),
\[
 \sum_{r=s}^{|G|}\frac{u^r}{r!}
     \sum_{\boldsymbol x\in\mathcal H_{r,N}}p^r
       \ll_{s,u,C,\varepsilon}N^{-2s/(2s+1)+\varepsilon}.
\]
\end{proposition}

\begin{proof}
A full-support relation selects at least two true
values in every window, since selections are even
and nonzero. Consequently the moderate certificates
and the saturation graph of
Section~\ref{p14:cert:section} apply to these families,
without any primitivity assumption.

Use the parameters in
\eqref{p1:core:primitive-parameters}.
The total independent weight of moderate
certificates is \(N^{-2\beta+4/A_0+o(1)}
=N^{-91/50+o(1)}\), by
Lemma~\ref{p1:core:moderate}.
If a family has no moderate certificate, every
unmarked vertex has a neighbour in the saturation
graph. The tilted generating function of
Lemma~\ref{p1:core:generator}, now with no factor
for rank loss, bounds its weight beyond size \(s_0\)
by \(O_u(N^{-s_0\delta/2})\).
Choose a fixed \(s_0>\max(s,5000)\).
Both large-size contributions are smaller than
\(N^{-1+\varepsilon}\).

For each of the finitely many \(s\le r<s_0\),
Proposition~\ref{p1:arith:fixed-support} gives
\[
 |\mathcal H_{r,N}|
       \ll_{r,C,\varepsilon}
             N^{r-2r/(2r+1)+\varepsilon}.
 \]
Multiplying by \((up)^r/r!\), and using
\(2r/(2r+1)\ge2s/(2s+1)\), completes the proof.
The factor \(1/r!\) counts unordered sets; it is
not a bound for the growth of an arithmetic
constant in \(r\).
\end{proof}

\begin{lemma}[When no relation uses every block]
\label{p1:uniform:missing-block}
If a family \(I\) of \(r\) regular windows has no
relative relation meeting every block, then
\[
 q(I)\le p\sum_{x\in I}q(I\setminus\{x\}).
\]
Summed over any subfamily of ordered distinct
\(r\)-tuples, the resulting upper bound is at most
\(rNp\,\E(Z_G)_{r-1}\).
\end{lemma}

\begin{proof}
The relation space is covered by the kernels of
its \(r\) coordinate-block projections.
Hence \(2^{\rho(I)}\le\sum_x2^{\rho(I\setminus\{x\})}\).
If the full affine event is compatible, every
subevent is compatible and multiplication by
\(p^r\) proves the claim. If it is incompatible,
the left side is zero. Summing over the omitted
centre gives at most \(N\) choices, for each of
the \(r\) omitted positions.
\end{proof}

An elementary companion to such support bounds is the
following weighted triangle inequality. If
\((w_{xy})\) is a finite symmetric matrix of
nonnegative integers and \(W=\sum_{x,y}w_{xy}\), then
\[
 \sum_{x,y,z}\min(w_{xy},w_{xz},w_{yz})\le W^{3/2}.
\]
Indeed, for each integer \(h\ge1\), let \(A_h\) be the
symmetric zero-one matrix of entries at least \(h\),
and \(M_h\) its number of ones. Its triangle sum is
\(\operatorname{tr}(A_h^3)\le
(\operatorname{tr}(A_h^2))^{3/2}=M_h^{3/2}\).
Sum in \(h\) and use
\(\sum_hM_h^{3/2}\le(\sum_hM_h)^{3/2}\).

\subsection{Old square classes and a growing number of sources}

Let \(\mathcal P\) be a finite set of primes, and
\[
 A(X;\mathcal P)
   =\#\{n\le X:\sf(n)\text{ is supported on }\mathcal P\}.
\]

\begin{lemma}[Rankin under a logarithmic prime budget]
\label{p1:uniform:prime-budget}
For \(1/2<\sigma<1\),
\[
 A(X;\mathcal P)\le
 X^\sigma\zeta(2\sigma)
                  \prod_{\ell\in\mathcal P}(1+\ell^{-\sigma}).
\]
Uniformly for \(3/5\le\sigma\le3/4\), \(Y\ge3\), and
\(\sum_{\ell\in\mathcal P}\log\ell\le Y\),
\[
 \log A(X;\mathcal P)
       \le \sigma\log X+
              O\left(1+\frac{Y^{1-\sigma}}{\log Y}\right).
\]
\end{lemma}

\begin{proof}
Write \(n=du^2\), where \(d\) is squarefree and
supported on \(\mathcal P\), and sum
\((X/(du^2))^\sigma\).
For the second assertion, the prime-counting
upper bound and partial summation give
\(\sum_{\ell\le Y}\ell^{-\sigma}
\ll Y^{1-\sigma}/\log Y\) uniformly in the stated
compact interval. There are at most \(Y/\log Y\)
members of \(\mathcal P\) exceeding \(Y\);
their contribution satisfies the same bound.
Finally \(\zeta(2\sigma)\) is uniformly bounded.
\end{proof}

\begin{proposition}[Independent extensions at the critical order]
\label{p1:uniform:branching}
Put \(H=\log N\), \(h=\log H\), and fix \(K>0\).
For every \(j\le KH/h\) source windows, there is
an exceptional set of new centres of size at most
\(D_N^*\), where
\begin{equation}
 \log D_N^*
 \le\frac23H-\frac{H\log h}{9h}+O_{K,C}(H/h),
 \label{p1:uniform:branching-rate}
\end{equation}
such that outside it the \(L\) new relative bits
are exactly uniform and independent of all
previously exposed values.
\end{proposition}

\begin{proof}
Include in \(\mathcal P\) all primes at most \(B\)
and all primes dividing one of the \(jB\) old
values. Then
\[
 \sum_{\ell\in\mathcal P}\log\ell
       \ll_C B+jB\log(3N)\ll_{K,C}H^3/h.
\]
Take \(Y=C_{K,C}H^3/h\) and
\(\sigma=1-h/\log Y\). Eventually
\(\sigma\in[3/5,3/4]\) and \(Y^{1-\sigma}=H\).
Lemma~\ref{p1:uniform:prime-budget} gives
\[
 \log A(3N;\mathcal P)
 \le H-\frac{Hh}{\log Y}
                       +O_{K,C}(H/\log Y).
\]
Since \(\log Y=3h-\log h+O_{K,C}(1)\), the
right side equals the bound in
\eqref{p1:uniform:branching-rate}.

Exclude centres whose windows overlap an old window,
or contain a value whose square class is supported
on \(\mathcal P\). This costs at most
\(j(2L+1)+BA(3N;\mathcal P)\), with the additional
polylogarithmic factors absorbed in the error term.
Every value in a remaining window has an odd prime
outside \(\mathcal P\). Such a prime exceeds the
diameter and occurs nowhere else in the window.
These private coordinates form an identity minor
in the new valuation matrix and are independent
of all old values. This proves the exact
independence assertion.
\end{proof}

\begin{corollary}[A single pair at all critical orders]
\label{p1:uniform:one-pair}
Let \(A_N=2\sum_{e\in E_N}q_e\).
Construct ordered \(r\)-tuples by choosing the
two positions of a good pair and then adding
windows in position order, always outside the
exceptional set of
Proposition~\ref{p1:uniform:branching}.
Their total positive mass \(B_{r,N}\) satisfies,
uniformly for \(2\le r\le m\le KH/h\),
\[
 B_{r,N}=\binom r2\lambda_N^{r-2}A_N
                       \{1+O_{K,C}(rD_N^*/N)\}.
\]
Furthermore,
\[
 \sum_{r=2}^{m}\frac{2^{r-1}}{r!}
   \left|B_{r,N}-\binom r2\mu_N^{r-2}\alpha_N\right|
       =o_C(a_N).
\]
\end{corollary}

\begin{proof}
For a fixed pair each added window contributes
exactly \(p\). The number of extensions is
between \((N-D_N^*)^{r-2}\) and \(N^{r-2}\);
\(mD_N^*/N\to0\). The pair is the unique
dependent pair in the resulting tuple, so its
choice of positions is unique. This proves the
first formula.
The weighted branching error sums to
\[
 O_{K,C}\left(A_N\frac{D_N^*}{N}
       \sum_{r\ge2}\frac{2^{r-1}}{r!}
                r\binom r2\Lambda^{r-2}\right)=o(a_N),
\]
where \(\Lambda\) bounds the means.
The pair and mean estimates replace \(A_N\)
by \(\alpha_N\), and \(\lambda_N\) by \(\mu_N\),
at the same aggregate precision.
\end{proof}

This corollary controls a selected positive family.
Its independence guarantee is uniform at growing
orders, but it does not by itself estimate the full
factorial moments at those orders.

\subsection{Exact moments of the modified count}

Let \(E'_N\) be the event in Lemma~\ref{p1:core:exceptional},
let \(A=\{x\in G:J_x=1\}\), and put
\(\widetilde Z=Z_G\ind_{(E'_N)^c}\).
For \(r\ge1\), define
\[
 D_r(A)=\#\{I\subseteq A:|I|=r,\ I\text{ direct}\},
 \qquad L'_r=\E[D_r(A);E'_N].
\]
For \(r\ge2\), also put
\[
 P_r(A)=\sum_{e\in E_N}J_e
       \#\{U\subseteq A\setminus e:U\in\Gamma_e(r-2)\},
 \qquad H_r=\E[P_r(A);E'_N],
\]
and set \(H_1=0\).
Here \(\Gamma_e\), the independent deficits \(F_r,G_r\),
and \(\beta_N=\sum_eq_e\) are those used in the
stopped expansion.

\begin{proposition}[Moment decomposition and pair leakage]
\label{p1:uniform:modified-moments}
For every integer \(r\ge1\),
\begin{equation}
 \E\binom{\widetilde Z}{r}
   =\binom nrp^r-F_r-L'_r
     +\beta_N\binom{n-2}{r-2}p^{r-2}-G_r-H_r,
 \label{p1:uniform:modified-identity}
\end{equation}
where the last three terms are zero for \(r=1\).
All four deficits \(F_r,G_r,L'_r,H_r\) are
nonnegative. Moreover,
\[
 \sum_{r\ge2}2^rG_r
       \ll_\varepsilon N^{-1+\varepsilon},
 \qquad
 \sum_{r\ge2}2^rH_r=o(a_N).
\]
\end{proposition}

\begin{proof}
Outside \(E'_N\), an active family is either direct
or contains the unique active good pair and is
direct after its compression. Thus
\[
 \binom{Z_G}{r}=D_r(A)+P_r(A)
              \qquad\text{on }(E'_N)^c.
\]
The expectations without restriction to this event
are, respectively,
\[
 \E D_r(A)=\binom nrp^r-F_r,\qquad
 \E P_r(A)
     =\beta_N\binom{n-2}{r-2}p^{r-2}-G_r.
\]
Subtract the two contributions on \(E'_N\) to
obtain \eqref{p1:uniform:modified-identity}.
The bound for \(G_r\) was proved in
Lemma~\ref{p1:close:deficit-bounds}.

For each fixed \(r\), \(P_r(A)\le C_rZ_G^r\).
H\"older's inequality and the fixed moments give
\[
 H_r\le C_r(\E Z_G^{8r})^{1/8}
                  \P(E'_N)^{7/8}
       \ll_{r,\varepsilon}N^{-7/10+7\varepsilon/8}
       =o(a_N)
\]
when the margin is sufficiently small.
To sum all \(r\), use directness before taking
expectations:
\[
 \sum_{r\ge2}4^rH_r
 \le\sum_{r\ge2}4^r\E P_r(A)
 \le16\beta_N(1+4p)^{n-2}
       \ll_C a_N.
\]
For any fixed \(K\), the terms \(2\le r\le K\)
are \(o(a_N)\), and the remaining \(2^r\)-weighted
sum is \(O_C(2^{-K}a_N)\).
Take first \(N\to\infty\), then \(K\to\infty\).
This proves the asserted small \(o\), without
giving it an additional polynomial rate.
\end{proof}

The direct leakage \(L'_r\) and the deterministic
deficit \(F_r\) are different quantities.
The small probability of \(E'_N\) gives a bound
for \(L'_r\) at each fixed order, but does not
on its own control their exponentially weighted
sum. The stopped expansion avoids paying that
whole sum: it keeps the suffix inside a
positive kernel at the first excluded prefix.

\subsection{Replicas with common small-prime signs}

Fix a threshold \(Y\). Keep the signs at primes
\(\ell\le Y\) common, and give each member of a
partition \(\pi\) an independent copy of the signs
at primes \(\ell>Y\). Occurrences may have repeated
centres; their formal coordinates remain distinct.
For the resulting valuation map let \(K_\pi\)
be its kernel on all formal occurrence coordinates, before quotienting
by the internal relations of individual windows. Thus internal local
relations are included in \(K_\pi\). Let
\(\eta_\pi\in\{0,1\}\) denote affine compatibility.
If instead each local block is first replaced by its effective image,
the relative-kernel formula uses the product of the actual marginals
\(\prod_i p_{x_i}\) in place of the nominal factor \(p^r\).
Repeated centres retain distinct formal occurrence coordinates.

\begin{proposition}[The replica formula and positive propagation]
\label{p1:uniform:replicas}
For \(r\) prescribed occurrences,
\begin{equation}
 \E\prod_{A\in\pi}
     \P(J_{x_i}=1\ (i\in A)\mid\mathcal F_Y)
       =p^r\eta_\pi2^{\dim K_\pi}.
 \label{p1:uniform:replica-formula}
\end{equation}
After any fixed number of compatible colored
sources, the sum of the conditional probabilities
of targets is \(O_\varepsilon(N^\varepsilon)\),
uniformly in \(Y\) and the coloring.
Here conditioning means prescribing these source events
in the replicated space; the common small-prime signs
are still averaged. Thus a positive random sum
\[
 R=\sum_{\boldsymbol x\in\mathcal A}w_{\boldsymbol x}
       \prod_{A\in\pi}
          \P(J_{x_i}=1\ (i\in A)\mid\mathcal F_Y),
\]
with fixed tuple length, fixed partition \(\pi\), and
nonnegative deterministic weights bounded by a fixed
power of \(\log N\),
satisfies for every fixed integer \(q\ge1\)
\begin{equation}
 \E R^q\ll_\varepsilon N^\varepsilon\E R,
 \label{p1:uniform:replica-moment}
\end{equation}
when its index set is a subfamily of all ordered
tuples of the stated fixed length.
\end{proposition}

\begin{proof}
Conditionally on the common small-prime signs,
the independent large-prime copies realize the
product of conditional probabilities on the left
of \eqref{p1:uniform:replica-formula}.
Fourier inversion in this larger binary space
gives its right side, with its compatibility
factor retained.

Merging colors adds prime-parity equations
together, so a relation in the colored system
is also a relation in the merged one. The
projection of its kernel to the target block
is contained in the corresponding merged
projection. Consequently the target's colored
rank loss is no larger than its uncolored rank
loss. Proposition~\ref{p1:fixed:targets} and the
proof of \eqref{p1:fixed:conditional} apply.
Targets whose centre repeats a source account
for only a fixed number of choices and have
probability at most one.

Expand \(R^q\) positively, retain the first
replica product as a seed, and sum all further
occurrences successively. Each step has the
conditional bound just proved. Repeated centres
are treated as above. There are only a fixed
number of steps; splitting \(\varepsilon\)
among them and absorbing the fixed powers of
\(\log N\) proves \eqref{p1:uniform:replica-moment}.
\end{proof}

\paragraph{A local-kernel example.}
For a single start at \(x=18\) of length \(L=10\), the ten affine
rows have rank nine. At prime cutoff \(28\), the unique compatible
assignment makes prime \(17\) negative and all other prime signs
positive, so \(p_{18}=2^{-9}\). The formal relation kernel has
dimension one, giving \(2^{-10}2^1=2^{-9}\) in the nominal formula.
The effective relative nullity of a singleton is zero; its formula
therefore retains \(p_{18}\). The internal even square relation is
visible in \(18\cdot24\cdot25\cdot27=540^2\).
This is a finite algebraic illustration, not an assertion that these
parameters lie beyond an asymptotic threshold.

The same reasoning permits measurable gates in
positive estimates. For example, suppose
\(\E R\ll_\delta N^{-\eta+\delta}\) for every
\(\delta>0\), and let \(\mathcal A\in\mathcal F_Y\).
For an integer \(k\ge1\), interpolation gives
\[
 \E R^{1+1/k}\le(\E R)^{1-1/k}(\E R^2)^{1/k}
       \ll_\delta N^\delta\E R.
\]
H\"older's inequality with conjugate exponents
\(1+1/k\) and \(k+1\), followed by the fixed
moment of order \(d(k+1)\), gives
\(\E[RZ_G^d]\ll_\delta
 N^{-\eta k/(k+1)+O(\delta)}\).
Choose \(k\) after the desired error margin.
Since the gate is at most one, for fixed \(d\) this yields
\[
 \E[R\,Z_G^d\,\ind_{\mathcal A}]
       \ll_\varepsilon N^{-\eta+\varepsilon}.
\]
Here \eqref{p1:uniform:replica-moment} supplies the
second moment of \(R\). The estimate is an
unconditional weighted expectation; it makes
no division by \(\P(\mathcal A)\).

\section{Fixed moments do not give a polynomial rate}
\label{p1:moment-obstruction}

This example isolates the logical role of the growing-order
estimate. It is a family of probability laws, with no claimed
multiplicative realization.

\begin{proposition}[A perturbation invisible at each fixed order]
\label{p1:obstruction:fixed}
For \(0\le a\le1\) put
\[
 Q_a=\Pi_1+\frac a2\Pi_1*\Delta^{*2}.
\]
For every integer \(m\ge3\), define
\[
 \epsilon_{m,a}=\frac{(1-a/2)e^{-1}}{2m!},
 \qquad
 P_{m,a}=Q_a+\epsilon_{m,a}\Delta^{*m}.
\]
Then \(P_{m,a}\) is a probability measure satisfying
\(\frac12Q_a\le P_{m,a}\le\frac32Q_a\) pointwise.
Its factorial moments agree exactly with those of \(Q_a\)
at every order other than \(m\), and
\[
 \TV(P_{m,a},Q_a)
       =\frac{(1-a/2)2^{m-2}}{e\,m!}.
\]
For every fixed \(t\ge0\), its exponential moment
\(\E_{P_{m,a}}e^{tX}\) is bounded uniformly in \(m,a\).
\end{proposition}

\begin{proof}
The second Charlier polynomial at mean one is
\(k^2-3k+1\), whose minimum on \(\N\) is \(-1\).
Thus
\[
 Q_a\{k\}\ge(1-a/2)e^{-1}/k!.
\]
For \(0\le k\le m\),
\[
 \frac{\epsilon_{m,a}\binom mk}{Q_a\{k\}}
       \le\frac1{2(m-k)!}\le\frac12 .
\]
Outside this range the perturbation is zero. Positivity
and the two pointwise bounds follow. The perturbation
has total mass zero, so \(P_{m,a}\) has mass one.

The probability generating function of \(\Delta^{*m}\)
is \((z-1)^m\). Its derivatives at \(z=1\) vanish at
every order other than \(m\), where the derivative is
\(m!\). This proves the factorial-moment statement.
Its \(\ell^1\) norm is \(2^m\), yielding the exact
total variation distance. Finally,
\[
 \E_{P_{m,a}}e^{tX}
 \le\frac32
       \exp(e^t-1)\left(1+\frac12(e^t-1)^2\right),
\]
which is independent of \(m,a\).
\end{proof}

Take \(m_N=\lfloor\sqrt{\log N}\rfloor\) and
\(a_N^\circ=(\log N)N^{-2/3}\), for \(N\) sufficiently
large. At every fixed order \(r\), the factorial
moments of \(P_{m_N,a_N^\circ}\) eventually equal
\(1+\binom r2a_N^\circ\) exactly. Nevertheless,
Stirling's formula gives
\[
 \log\TV(P_{m_N,a_N^\circ},Q_{a_N^\circ})
   =-\tfrac12\sqrt{\log N}\log\log N
                   +O(\sqrt{\log N}).
\]
This distance exceeds \(N^{-\gamma}\) for every fixed
\(\gamma>0\), eventually. Fixed-order information,
even with the displayed exponential-moment control,
therefore cannot replace the uniform estimate on the
boundary of the stopped expansion.

\clearpage
\section*{Formalization and statement correspondence}
\addcontentsline{toc}{section}{Formalization and statement correspondence}
The repository \url{https://github.com/Vulkin-prog/paper-p1-lean}
contains the Lean development for both texts, checked with
Lean~4.36.0-rc2 and the pinned Mathlib revision.
It records exact statements, proof entries and validation sources;
the version~1.6 archive is preserved byte for byte.

Version~1.7 corrected two false original formulations:
the enclosing-interval radius in Lemma~\ref{p1:arith:asymmetric}
and the capped exponent in Lemma~\ref{p13:short:quadratic}.
Their counterexamples and corrected scopes are now discussed beside
those results in the companion. Version~1.8 selectively integrates
stronger estimates and shorter proofs; all 167 numbered-result labels
and their numbering are retained. Intermediate proof constructions
in the change ledger are not all separate scientific contributions.

Kernel checking and finite diagnostics do not replace checking
the correspondence between the formal and printed statements.
The audited proofs use only the standard kernel axioms recorded
in the repository. This editorial preview has no Palomar identifier.

\section*{Statement on AI assistance}
\addcontentsline{toc}{section}{Statement on AI assistance}

ChatGPT (GPT-6 Astra) was used to organize the research material,
develop and rewrite mathematical arguments, inspect references,
and prepare both manuscripts and their sources.
Parallel AI-assisted reviews and exact finite computations examined
selected identities, constants, and proof dependencies.
The development accompanying versions~1.7--1.8 also used AI assistance
for the Lean proofs and
the comparison of formal statements with the manuscripts.
AI-assisted review and machine checking are distinct from independent
human verification. The author remains responsible for the mathematical
claims, their correspondence with the formal statements, and any submission.

\ifdefined\PoneTechnicalCompanion\else
\paragraph{Sources and formalized proofs.}
The preserved archives \nolinkurl{P1_V1_6_Sources_2026-10-31.zip} and
\nolinkurl{P1_V1_6_Complete_2026-10-31.zip} document version~1.6,
including both documents' LaTeX sources and finite diagnostic checks.
Those scripts verify selected identities.
This version~1.8 working draft is accompanied by the Lean development,
statement correspondence and kernel-axiom audits described in the
preceding appendix and available in the associated repository.
\fi

\clearpage
\begin{thebibliography}{99}
\addcontentsline{toc}{section}{References}

\bibitem{agg1989}
R.~Arratia, L.~Goldstein, and L.~Gordon,
\emph{Two moments suffice for Poisson approximations:
the Chen--Stein method},
Annals of Probability \textbf{17} (1989), no.~1, 9--25.
\href{https://doi.org/10.1214/aop/1176991491}
{doi:10.1214/aop/1176991491}.

\bibitem{agw1990}
R.~Arratia, L.~Gordon, and M.~S.~Waterman,
\emph{The Erd\H{o}s--R\'enyi law in distribution, for coin tossing
and sequence matching},
Annals of Statistics \textbf{18} (1990), no.~2, 539--570.
\href{https://doi.org/10.1214/aos/1176347615}
{doi:10.1214/aos/1176347615}.
\href{https://dornsife.usc.edu/msw/wp-content/uploads/sites/236/2023/09/msw-093.pdf}
{Article text}.

\bibitem{barbour1987}
A.~D.~Barbour,
\emph{Asymptotic expansions in the Poisson limit theorem},
Annals of Probability \textbf{15} (1987), no.~2, 748--766.
\href{https://doi.org/10.1214/aop/1176992169}
{doi:10.1214/aop/1176992169}.
\href{https://user.math.uzh.ch/barbour/pub/Barbour/AsympExpanPois.pdf}
{Article text}.

\bibitem{barboureagleson1987}
A.~D.~Barbour and G.~K.~Eagleson,
\emph{An improved Poisson limit theorem for sums of dissociated
random variables},
Journal of Applied Probability \textbf{24} (1987), no.~3, 586--599.
\href{https://doi.org/10.2307/3214091}
{doi:10.2307/3214091}.
\href{https://user.math.uzh.ch/barbour/pub/Barbour/BEagleson_ImprovedPoisson.pdf}
{Article text}.

\bibitem{bkn2014}
A.~D.~Barbour, E.~Kowalski, and A.~Nikeghbali,
\emph{Mod-discrete expansions},
Probability Theory and Related Fields \textbf{158}
(2014), 859--893.
\href{https://doi.org/10.1007/s00440-013-0498-8}
{doi:10.1007/s00440-013-0498-8}.
Section numbers refer to the
\href{https://people.math.ethz.ch/~kowalski/mod-discrete.pdf}
{author version}, Sections~5 and~8.1.

\bibitem{basaknathzaharescu2024}
D.~Basak, K.~Nath, and A.~Zaharescu,
\emph{Poisson behavior for chains of small quadratic non-residues},
International Mathematics Research Notices \textbf{2024} (2024),
no.~4, 3356--3390.
\href{https://doi.org/10.1093/imrn/rnad168}
{doi:10.1093/imrn/rnad168}.
\href{https://drive.google.com/file/d/1-7VxBG_KspRvDj35EFb9s9p13sR4rRjM/view}
{Author version}, Theorem~1.

\bibitem{binyaminicluckersnovikov2024}
G.~Binyamini, R.~Cluckers, and D.~Novikov,
\emph{Bounds for rational points on algebraic curves, optimal
in the degree, and dimension growth},
International Mathematics Research Notices \textbf{2024} (2024),
no.~11, 9256--9265.
\href{https://doi.org/10.1093/imrn/rnae034}
{doi:10.1093/imrn/rnae034}.
The degree-uniform integer-point estimate is cited from
\href{https://arxiv.org/abs/2302.04209v2}
{arXiv:2302.04209v2} (20 September 2023),
Theorem~2, equation~(27).

\bibitem{bretwngxu2026}
R.~de~la~Bret\`eche, V.~Y.~Wang, and M.~W.~Xu,
\emph{Random multiplicative functions and making squares from
polynomial values},
\href{https://arxiv.org/abs/2607.06398v1}
{arXiv:2607.06398v1} (2026), Theorem~1.2 and
Lemmas~6.2--6.3.

\bibitem{dohmen1999}
K.~Dohmen,
\emph{Improved inclusion-exclusion identities and inequalities
based on a particular class of abstract tubes},
Electronic Journal of Probability \textbf{4} (1999), paper~5,
1--12.
\href{https://doi.org/10.1214/EJP.v4-42}
{doi:10.1214/EJP.v4-42}.
\href{https://www.maths.tcd.ie/EMIS/journals/EJP-ECP/article/download/42/42-83-1-PB.pdf}
{Sections~2--3}.

\bibitem{erdosivicpomerance1986}
P.~Erd\H{o}s, A.~Ivi\'c, and C.~Pomerance,
\emph{On sums involving reciprocals of the largest
prime factor of an integer},
Glasnik Matemati\v{c}ki \textbf{21}(41) (1986), 283--300.
\href{https://www.renyi.hu/~p_erdos/1986-10.pdf}
{Article text}.

\bibitem{erdosturk1984}
P.~Erd\H{o}s and J.~Turk,
\emph{Products of integers in short intervals},
Acta Arithmetica \textbf{44} (1984), no.~2, 147--174.
\href{https://doi.org/10.4064/aa-44-2-147-174}
{doi:10.4064/aa-44-2-147-174}.
\href{https://www.renyi.hu/~p_erdos/1984-17.pdf}
{Article text}.

\bibitem{godbole1991}
A.~P.~Godbole,
\emph{Poisson approximations for runs and patterns of rare events},
Advances in Applied Probability \textbf{23} (1991), no.~4,
851--865.
\href{https://doi.org/10.2307/1427680}
{doi:10.2307/1427680}.

\bibitem{heathbrown2012}
D.~R.~Heath-Brown,
\emph{Square-free values of \(n^2+1\)},
Acta Arithmetica \textbf{155} (2012), no.~1, 1--13.
\href{https://doi.org/10.4064/aa155-1-1}
{doi:10.4064/aa155-1-1}.
Preprint version:
\href{https://arxiv.org/abs/1010.6217v3}
{arXiv:1010.6217v3}, Lemma~2.

\bibitem{kempkemp1965}
C.~D.~Kemp and A.~W.~Kemp,
\emph{Some properties of the `Hermite' distribution},
Biometrika \textbf{52} (1965), nos.~3--4, 381--394.
\href{https://doi.org/10.1093/biomet/52.3-4.381}
{doi:10.1093/biomet/52.3-4.381}.

\bibitem{matomakiradziwill2016}
K.~Matom\"aki and M.~Radziwi{\l}{\l},
\emph{Multiplicative functions in short intervals},
Annals of Mathematics \textbf{183} (2016), 1015--1056.
\href{https://doi.org/10.4007/annals.2016.183.3.6}
{doi:10.4007/annals.2016.183.3.6}.

\bibitem{munsch2026}
M.~Munsch, I.~E.~Shparlinski, Y.-C.~Sun, and Y.~Xiao,
\emph{A large sieve inequality for sums of Legendre symbols
over short intervals},
\href{https://arxiv.org/abs/2604.23661v3}
{arXiv:2604.23661v3} (5 August 2026), Lemma~3.3 and its proof.

\bibitem{naimanwynn1997}
D.~Q.~Naiman and H.~P.~Wynn,
\emph{Abstract tubes, improved inclusion-exclusion identities
and inequalities and importance sampling},
Annals of Statistics \textbf{25} (1997), no.~5, 1954--1983.
\href{https://doi.org/10.1214/aos/1069362380}
{doi:10.1214/aos/1069362380}.

\bibitem{najnudel2020}
J.~Najnudel,
\emph{On consecutive values of random completely multiplicative
functions},
Electronic Journal of Probability \textbf{25} (2020), paper~59,
1--28.
\href{https://doi.org/10.1214/20-EJP456}
{doi:10.1214/20-EJP456}.
\href{https://arxiv.org/abs/1702.01470v2}
{arXiv:1702.01470v2}, Lemma~3.1 and Proposition~6.3.

\bibitem{pila1995}
J.~Pila,
\emph{Density of integral and rational points on varieties},
in \emph{Columbia University Number Theory Seminar
(New York, 1992)}, Ast\'erisque \textbf{228} (1995),
183--187.
\href{https://www.numdam.org/item/AST_1995__228__183_0/}
{Numdam: AST\_1995\_\_228\_\_183\_0}.

\bibitem{pouly2026c}
B.~Pouly,
\emph{Long runs and rare patterns of a random completely
multiplicative function}, with technical companion,
public preprint, version~3 (2026).
\href{https://doi.org/10.5281/zenodo.22872154}
{doi:10.5281/zenodo.22872154}.
\href{https://hexagonmath.org/2610.00223v1}
{hexagon:2610.00223v1}, Theorem~1.1 and Section~3.6, (3.22).

\bibitem{reuss2014}
T.~Reuss,
\emph{Pairs of \(k\)-free numbers, consecutive square-full
numbers},
\href{https://arxiv.org/abs/1212.3150v2}
{arXiv:1212.3150v2} (2014), Sections~2.1--2.3.

\bibitem{salberger2023}
P.~Salberger,
\emph{Counting integral points of affine hypersurfaces},
\href{https://arxiv.org/abs/2311.05561v1}
{arXiv:2311.05561v1} (2023), Lemma~2.

\bibitem{zacharovashwang2010}
V.~Zacharovas and H.-K.~Hwang,
\emph{A Charlier--Parseval approach to Poisson approximation
and its applications},
Lithuanian Mathematical Journal \textbf{50} (2010), no.~1,
88--119.
\href{https://doi.org/10.1007/s10986-010-9073-5}
{doi:10.1007/s10986-010-9073-5}.
Preprint version used for equation and remark numbers:
\href{https://arxiv.org/abs/0810.4756v1}
{arXiv:0810.4756v1}, (2.13) and Remark~4.4.

\bibitem{griesmer1960}
J.~H.~Griesmer,
\emph{A bound for error-correcting codes},
IBM Journal of Research and Development \textbf{4} (1960),
no.~5, 532--542.
\href{https://doi.org/10.1147/rd.45.0532}
{doi:10.1147/rd.45.0532}.
\href{https://bitsavers.trailing-edge.com/pdf/ibm/IBM_Journal_of_Research_and_Development/045/ibmrd0405M.pdf}
{Original article scan}, Section~2, equation~(2.10).

\bibitem{zeilberger1997}
D.~Zeilberger,
\emph{The abstract lace expansion},
Advances in Applied Mathematics \textbf{19} (1997),
no.~3, 355--359.
\href{https://doi.org/10.1006/aama.1997.0552}
{doi:10.1006/aama.1997.0552}.
\href{https://sites.math.rutgers.edu/~zeilberg/mamarim/mamarimPDF/lace.pdf}
{Author version}, formula~(*).

\bibitem{cilleruelogaraev2011}
J.~Cilleruelo and M.~Z.~Garaev,
\emph{Concentration of points on two and three dimensional
modular hyperbolas and applications},
Geometric and Functional Analysis \textbf{21} (2011),
no.~4, 892--904.
\href{https://doi.org/10.1007/s00039-011-0127-6}
{doi:10.1007/s00039-011-0127-6}.
Numbering refers to
\href{https://arxiv.org/abs/1007.1526v2}{arXiv:1007.1526v2},
Section~3, Lemma~4 and Proposition~1.

\bibitem{tao2026anatomy}
T.~Tao,
\emph{Products of consecutive integers with unusual anatomy},
\href{https://arxiv.org/abs/2603.27990v3}{arXiv:2603.27990v3}
(26 September 2026), Lemma~2.10.

\bibitem{cekanaviciusvellaisamy2015}
V.~\v{C}ekanavi\v{c}ius and P.~Vellaisamy,
\emph{Discrete approximations for sums of $m$-dependent random variables},
ALEA, Latin American Journal of Probability and Mathematical Statistics
\textbf{12} (2015), 765--792.
\href{https://alea.impa.br/articles/v12/12-29.pdf}{Journal text},
Theorem~3.1, equation~(3.5), and Theorem~3.2, equation~(3.6).
\href{https://arxiv.org/abs/1301.7196v3}{arXiv:1301.7196v3}.

\bibitem{chhaibietal2020}
R.~Chhaibi, F.~Delbaen, P.-L.~M\'eliot, and A.~Nikeghbali,
\emph{Mod-$\phi$ convergence: Approximation of discrete measures
and harmonic analysis on the torus},
Annales de l'Institut Fourier \textbf{70} (2020),
no.~3, 1115--1197.
\href{https://doi.org/10.5802/aif.3332}{doi:10.5802/aif.3332}.
\href{https://numdam.org/item/10.5802/aif.3332.pdf}{Journal text},
Theorem~3.11.

\end{thebibliography}

\end{document}
